\documentclass[12pt]{article}
% %%%%%%%%%%%%%%%%%%%%%%%%%%%%%%%%%%%%%%%%%%%%%%%%%%%%%%%%%%%%%%%%%%%%%%%%%%%%%%%%%%%%%%%%%%%%%%%%%%%%%%%%%%%%%%%%%%%%%%%%%%%%GGsection
%
\usepackage{amsmath,amssymb,amsthm,enumerate,graphicx}
\usepackage{ifthen,latexsym,syntonly}
\usepackage{setspace}
\usepackage{subcaption,caption}
\usepackage{mathabx}
%\usepackage[showrefs]{refcheck}  %use this to show equation and section labels
\usepackage{color}
\usepackage[round,comma,authoryear]{natbib}   % for natbib
%\usepackage{subfigure}
%\usepackage{graphics}
\usepackage{hyperref}
\usepackage{float}
\usepackage{epstopdf}
\usepackage{mathabx}
\bibliographystyle{mynat}
%\bibliographystyle{ormsv080}
\usepackage{psfrag}
\usepackage{tikz}
\usetikzlibrary{decorations.pathreplacing}


%\usepackage{mathtools}
\usepackage{multirow}
\usepackage{tabularx}
%\usepackage{soul}
\newcolumntype{C}{>{\centering\arraybackslash}X}

\onehalfspacing

\DeclareMathOperator*{\argmax}{arg\,max}
\DeclareMathOperator*{\argmin}{arg\,min}
\setlength{\textwidth}{6.5in}
\setlength{\textheight}{8.5in}
\usepackage{color}
\allowdisplaybreaks[3]
\usepackage{calc,xspace}
\definecolor{wacky}{rgb}{1,0,1}
\onehalfspacing

\newcommand{\wacky}[1]{{\large\textcolor{wacky}{#1}} }

\definecolor{ForestGreen}{RGB}{34, 139, 34}


\newtheorem{theorem}{Theorem}%[section]
\newtheorem{remark}[theorem]{Remark}
\newtheorem{assumption}{Assumption}
\newtheorem{case}[theorem]{Case}
\newtheorem{claim}[theorem]{Claim}
\newtheorem{corollary}{Corollary}
\newtheorem{condition}{Condition}
\newtheorem{criterion}[theorem]{Criterion}
\newtheorem{definition}{Definition}
\newtheorem{example}{Example}
\newtheorem{lemma}[theorem]{Lemma}
\newtheorem{problem}[theorem]{Problem}
\newtheorem{bound}[theorem]{Bound}
\newtheorem{proposition}{Proposition}
%\newtheorem{solution}[theorem]{Solution}
\newtheorem{thm}[theorem]{Theorem}
\newtheorem{restructure}{Debt Restructure}

\providecommand{\keywords}[1]{\textbf{\textit{Keywords---}} #1}


\setlength{\oddsidemargin}{.05in} \setlength{\topmargin}{-.45in}
\setlength{\textwidth}{6.4in} \setlength{\textheight}{8.5in}



\newcommand{\NW}[1]{{\textcolor{blue}{#1}}}
\newcommand{\WJ}[1]{{\textcolor{red}{#1}}}
\newcommand{\TS}[1]{{\textcolor{magenta}{#1}}}
%\newcommand{\JQ}[1]{{\textcolor{cyan}{#1}}}



\begin{document}
	
	\date{%December 2009, revised
		%TCIMACRO{\TeXButton{Today}{\today}}%
		%BeginExpansion
		\today%
		%EndExpansion
	} %n : A Discrete-time Formulation
	\title{Implementing an Optimal Tax Plan with Short Term Debt\thanks{We thank Andrew Atkeson, Fernando Alvarez, Marco Bassetto, Roberto Chang, Lars Peter Hansen, Zhiguo He, Robert Shimer,  Nancy Stokey, Pierre-Olivier Weill, Pierre Yared, and seminar participants at  Princeton University, the University of Chicago, 25th Macro Finance Society Workshop, and UCLA for helpful comments.}}  % Commitment Continuous-Time
	%} %Deterministic Model} % A 
	%\NW{scratch note: Let's discuss the paper's title. Power laws and Pareto distributions? Permanent Earnings Shocks }}
%\\ % and Liquidity Constraints
%\bigskip Preliminary and Incomplete \\ Do not quote
%} % \maketitle
\date{\today }
\author{	 Wei Jiang\thanks{%
Department of Industrial Engineering and Decision Analytics, Hong Kong University of Science and Technology. E-mail: weijiang@ust.hk } \and Thomas J. Sargent\thanks{New York University. E-mail: thomas.sargent@nyu.edu} \and Neng Wang\thanks{Cheung Kong Graduate School of Business. E-mail: {newang@gmail.com}} }

\maketitle



\begin{abstract}

  \citet*{debortoli2021optimal} showed that    when initial government debt is too high,  \citeauthor{LucasStokey1983}'s (1983) debt management policy fails to  implement a Ramsey plan for flat-rate taxes on labor.
       We show that  preventing continuation Ramsey planners from resetting current-period tax rates   and    using short-term debt to finance accumulated primary deficits implements all Ramsey plans, including  those with the high debt levels that trouble Lucas and Stokey's
 implementation. 
 Our  implementation  equalizes Lagrange multipliers on  distinct implementability constraints that  face the Ramsey planner and  the  continuation planners, an essential indication of a successful implementation. We provide examples of our implementation under  different initial debt structures and  government spending patterns.

\thispagestyle{empty}

\medskip

\noindent  \textbf{Keywords}: Ramsey planner, continuation Ramsey planner, implementability, short term debt  
\newline
\newline
JEL Classification: E3, E6, C7
\end{abstract}  





\clearpage

%\NW{Reset the page number.} 

	\pagenumbering{arabic}
%	
%\section*{Warehouse for additions}
%
%Eight examples from \citet[Sec.~3]{LucasStokey1983}. 
%
%Footnote 2 of \citet{LucasStokey1983}  about short-term term and \citet{turnovsky1980time}

\section{Introduction}



	%We  confront a difficulty   that \citet*{debortoli2021optimal} detected with 
	\citet{LucasStokey1983}    implemented an optimal plan for flat-rate taxes on labor   by managing the  
	term structure of government debt under a timing protocol that represents  an  ``intermediate situation''  of ``partial commitment''.\footnote{Because  they are not implemented under    sequential timing protocols that 
	actually characterize  government decision-making processes, \citet{kydland1977rules} and  \citet{calvo1978time} argued 
	% as it had \citet{kydland1977rules} and  \citet{calvo1978time},
	that optimal  plans are implausible.   }   In  their %\citeauthor{LucasStokey1983}'s 
	model, a government   finances exogenous  sequences  $ \{G_{i}\}_{i=0}^{\infty}$ of  government expenditures and initial debt service coupons $ \{b_{-1,i}\}_{i=0}^{\infty}$ with a sequence
	 of flat-rate taxes on labor, the only factor of  production. % for time $t \geq 0$. 
	{At time $0$,  a planner chooses a sequence  of flat-rate taxes $\{\tau_{0, i}\}_{i=0}^{\infty}$ and   restructures an initial  debt  coupon process  $ \{b_{-1,i}\}_{i=0}^{\infty}$ to become $\{{b}_{0, i}\}_{i=0}^{\infty}$. In each  period $ i >  0$, a continuation  planner must  honor  the  debt service sequence  that it inherits from a period $i-1$ planner,}\footnote{This is the ``commitment'' part of Lucas and Stokey's ``partial commitment.''} but is  free to  redesign the continuation of  the  flat-rate tax sequence and  to reschedule government debt service from period $i $ onward.\footnote{This is the ``partial'' part of Lucas and Stokey's ``partial commitment.''}  Lucas and Stokey said that  an optimal  plan for setting tax rates  is ``time consistent''\footnote{We say that it is ``implemented.''} if all continuation planners choose to continue it.  {  \citet[Sec.~3]{LucasStokey1983}  
	provided eight examples in which restructuring government debts and issuing consols  
	implements a Ramsey plan for flat-rate taxes on labor.
	% in which	 a  Ramsey planner and continuation Ramsey planners  reschedule government debt so that  the Ramsey plan is implemented. %NW{Tom: the preceding sentence is a bit not clear and maybe somewhat repetitive?} % ß\NW{Tom: is the preceding sentence clear?}  
	%In that way, they neatly linked   a prescription for managing the term  structure of government debt with a timing protocol that  implements an optimal plan. 
	}
	
	\citet{debortoli2021optimal} constructed examples in which initial government debts are so high that in \citeauthor{LucasStokey1983}'s way restructuring government debt  fails to motivate  continuation planners   to continue the Ramsey tax plan; coincidentally, % in those examples, %\citeauthor{debortoli2021optimal} noted that  
	 the Ramsey plan sets a tax rate above the peak of the Laffer curve; the continuation Ramsey planner instead  sets that tax rate below the  peak of the Laffer curve.
	%an outcome that was always coincident with failure of    \citeauthor{LucasStokey1983}'s debt management policy to implement the Ramsey plan in their examples.  % confirm the Ramsey plan.  5
	%call for this tax policy.} %it is ex ante optimal to do so.} 
	
%	
%	\WJ{Also discuss a bit how minor change our timing protocol is to LS's? In some sense, our timing protocol may be even more natural? you mentioned  some reference at some point -- cite?} 
%	

%To enable a debt management policy to  implement a Ramsey plan, 
We modify Lucas and Stokey's timing protocol  by
%maintain all aspects of  \citeauthor{LucasStokey1983}'s timing protocol for decision making except one that we alter in an arguably plausible way {by 
withdrawing from each continuation planner the authority to reset today's tax rate and adding the authority to set tomorrow's tax rate. 
% to the protocol for setting tax rates.
Under our timing protocol, there exist  debt management policies,  %used by the Ramsey planner and by a sequence of continuation Ramsey planners 
different  from \citeauthor{LucasStokey1983}'s, that implement a Ramsey plan whenever one exists.  
%Under our timing protocol and debt management policy, a Ramsey plan is always implemented.
%Our debt management policy ``takes a short route'' in the sense that, 
%Unlike  \citeauthor{LucasStokey1983}'s, 
We focus on one such   policy that each period involves minimally  rescheduling the initial debt sequence.  This policy  uses short-term debt to finance  primary government deficits accumulated under the Ramsey tax plan and payments due each period under the   initial debt structure 	 $ \{b_{-1,i}\}_{i=0}^{\infty}$.\footnote{Subsection \ref{sec:otherpolicies} describes other possible policies.}
	
	To appreciate  this  debt management policy, it is useful to recall how
	\citet*{aguiar2019take}   contrasted  \citeauthor{LucasStokey1983}'s (1983) model with theirs: 
	{\it 
\begin{quotation}
	\citet{LucasStokey1983} studied optimal fiscal policy with complete markets and discussed
	at length how maturity choice is a useful tool to provide incentives to a government 
	that lacks commitment to taxes and debt issuance, but cannot default. The government
	has an incentive to manipulate the risk-free real interest rate, by changing taxes which affects
	investors' marginal utility, to alter the value of outstanding long-term bonds, something
	ruled out by our small open-economy framework with risk-neutral investors. Their
	main result is that the maturity of debt should be spread out, resembling the issuance of
	consols. Our model instead emphasizes default risk, something absent from their work.
	Our main result is also the reverse, providing a force for the exclusive use of short-term
	debt.
\end{quotation}
}


Short-term debt takes center stage  in our model too, but   for different reasons  because our economic environment differs from \citeauthor{aguiar2019take}'s. 
 Unlike \citet{aguiar2019take},
we retain almost all parts  of   \citeauthor{LucasStokey1983}'s and \citeauthor{debortoli2021optimal}'s structure including complete markets,
a closed economy with endogenous interest rates, the presence of 
incentives that  continuation Ramsey planners  have  to use  tax rates to manipulate interest rates, 
 obligations of  governments to pay their debts, and \citeauthor{LucasStokey1983}'s  ``partial commitment'' assumption:

{\it
\begin{quotation}
  We focus on a situation in which there are no binding commitments
on future taxes but in which debt commitments are fully binding.
 $\ldots$ (this) seems to come closest to the institutional arrangement we
observe in any stable, democratically governed countries. $\ldots$
 Our main finding, in this intermediate situation, is that being unable
to make commitments about future tax rates is not a constraint. In
the absence of any ability to bind choices about tax rates directly,
each government restructures the debt in a way that  induces its
successors to continue with the optimal tax policy.
  For this to be possible, a rich enough mix of debt instruments must
be available, where a rich enough means, roughly, one security for
each dated, state-contingent good being traded.

\quad \quad \citet[p.~69]{LucasStokey1983}

\end{quotation}
}

 \citeauthor{LucasStokey1983} concluded that %under their timing protocol 
  commitment to servicing debt suffices to induce successor governments to  implement the Ramsey tax plan,  but \citeauthor{debortoli2021optimal}'s counterexample shows that actually it may not   be enough. Provided that we also substantially modify their
  debt management policy,  by marginally altering    
 \citeauthor{LucasStokey1983}'s timing protocol we can make it be enough. 


\subsection*{Related literature}

%\TS{Tom and Neng and Wei: Maybe add that one of our examples in section XXXX recover a `tax smoothing'' outcome like Barro's}

\citet{barro1979determination},  \citet{LucasStokey1983}, and \citet{debortoli2021optimal} are the most pertinent references for us, but  many other papers are related to  aspects of this paper.
\citet{angeletos2002fiscal}
 shows that by choosing the maturity structure of non-state-contingent bonds, a  government can smooth tax rates in response to  shocks to its expenditures  as effectively  as  if it had access
 to complete markets in state-contingent government debt. He provides an example in which the government does that by trading a risk-free one-period bond and a consol. 
 \citet{buera2004optimal} showed quantitatively that such a government portfolio policy   entails large trades, hundreds of times GDP), that would be very  sensitive to a model's parameter values. \citet{aiyagari2002optimal} 
study  optimal taxation in a setting in which  a government can issue only one-period risk-free debt.    They show that an  optimal policy   combines features  of \citeauthor{barro1979determination}'s (1979)  model (a near-unit root component in taxes and debt)  and  \citeauthor{LucasStokey1983}'s (1983) (strong dependence of taxes and deficits on current shocks). 
Like \citet{angeletos2002fiscal} and \citet{buera2004optimal}, 
\cite{shin2007managing} shows  that a government can implement complete-market Ramsey allocations by  actively managing  just a few bonds. %Using data from 18th century British government expenditure, 
\citeauthor{shin2007managing} uses his model to explain why  in the 18th century Britain issued long-term bonds during peacetime and short-term bonds during wars.
\citet{bassetto2014optimal} studies  how heterogeneity among economic agents affects optimal government  responses to war-time surges in expenditures. If  a government  favors  taxpayers, it  runs larger deficits during wars, while if a government favors rentiers who earn most of their income from assets, it raises taxes immediately. He uses his model to understand  differences between British and French fiscal policies during the 17th and 18th centuries.


\citet{bhandari2017fiscal} provide conditions under which government debt under a Ramsey plan converges to an invariant distribution. They
 construct an approximation to that asymptotically invariant distribution as well as  an approximation to the  rate of convergence to it.
They show how binding implementability constraints push government debt in a direction that eventually  lets the government use fluctuations in equilibrium interest rates to insure against shocks to government expenditures,  rather   than using fluctuations in the par value of debt.
\citet{aparisiLannoy2024}
study  optimal taxation and government debts of  different  maturities, specifying  preferences and shock process to match macro and asset pricing facts, something that previous papers had not done. They provide formulas that express an optimal government  portfolio policy  as a function of statistics formed from financial and macroeconomic data. For US data, they find that  it is optimal for the government to  issue bonds in amounts that decline approximately exponentially with increases in their times to maturity.




%\citet{barro1979determination}

%
%\citet{buera2004optima}
% investigates whether a government can achieve the same welfare benefits as state-contingent debt by strategically managing the maturity structure of non-contingent debt. While the authors prove theoretically that this is possible, their numerical simulations reveal that implementing such a strategy would require enormous financial transactions (hundreds of times GDP) and positions that are extremely sensitive to small parameter changes, suggesting the approach may not be practically feasible.



\citet{zhu1992optimal}
constructs  a  complete markets model  with   capital. He finds that while optimal capital income tax rates are indeterminate,  their one-period ahead  conditional expectation equals zero. 
%The paper also demonstrates that the Ramsey equilibrium with wealth constraints is time consistent. 
\citet{chari1994optimal} construct a  quantitative version of a similar model. They  
  find that in a stationary equilibrium the optimal {\em ex ante} tax rate on capital income is approximately zero, that optimal  labor tax rates exhibit the same persistence as exogenous economic shocks and have low volatility,  and that  most welfare gains come from setting the  initial tax rate at an arbitrary exogenous upper bound. % rather than long-term tax structure changes.
%
%
%\citet{chari1990sustainable}
%This paper develops a framework for analyzing time-consistent policy in dynamic economies with competitive private agents and governments lacking commitment technologies. The authors propose a concept called "sustainable equilibrium," characterize the conditions under which arbitrary policies can be sustained, and show that with sufficient patience, even Ramsey allocations can be sustained in an infinite-horizon tax model.
%
%\citet{chari1994optimal}
% analyzes a business cycle model to determine optimal fiscal policy, finding that in a stationary equilibrium the optimal ex ante tax rate on capital income is approximately zero, labor tax rates should fluctuate very little and inherit the persistence properties of economic shocks, and most welfare gains come from high initial capital taxation rather than long-term tax structure changes.
\citet{chari2020optimal}
 also study optimal capital taxation. They find  that when a government has access to enough   instruments, capital should not be taxed. %if households have standard macroeconomic preferences. 
 %The authors show that the widely-held presumption in the literature that capital should be taxed for some length of time arises only when the tax system is restricted to fewer instruments.
\citet{barro2024taxation} show  that zero taxation of capital income is optimal when a government can't impose a capital levy. 
%The authors show that with appropriate constraints on initial household wealth (measured in utility units), and when extended to a multi-period setting, optimal policy features zero tax rates on asset income, constant tax rates on consumption, and achieves time-consistency from a "timeless perspective."
%
%\citet{aparisiLannoy2024}
% develops a framework for determining optimal government debt portfolios across various maturities, deriving formulas that express these portfolios in terms of observable statistics from financial and macroeconomic data, and shows that for US data, an optimal portfolio has bond shares that decline approximately exponentially with maturity.

%WJ{Broadly speaking, our paper is also related to the durable good monopoly literature (see, e.g., Coase, Stokey). absent commitment,  In that literature, the monopoly has incentives to overproduce.  current government has incentives to manipulate .} %the government's incentive } 

%\NW{Tom: shall we add a cite to DeMarzo-He (in corp/sovereign debt) and relatedly durable goods monopolist literature, e.g., Stokey? How likely is Zhiguo chosen to be referee?} 

%In corporate finance, to

%To avoid the challenges posed by  %an arbitrary
{An essential aspect  of \cite{LucasStokey1983}  is that the  initial debt structure $\{ b_{-1,i} \}_{i=0}^\infty$ can be any admissible function of calendar time, which shapes and substantially complicates their analysis and our analysis here. By way of contrast,  some  related  papers about debt dynamics have for  tractability   %e.g., corporate finance and some international finance settings, 
 assumed  exponentially maturing term debt with a constant amortization rate.\footnote{See  \cite{arellano2012default} and \cite{chatterjee2012maturity} for examples in sovereign debt settings.} 
% we assume that debt takes the form of  %a constant average maturity  where  debt balance decays exponentially (absent new issuances.) 
% where the level of perpetual (risky) debt is  chosen  at time 0 and remains fixed until the firm defaults, %by a firm's shareholders protected by limited liability, 
Using such exponential  term  debt structures,  \cite{demarzo2021leverage}  extended the classic capital-structure model  of \citep{leland1994corporate} in a corporate finance context, while  \cite*{demarzo2023sovereign}  studied  a sovereign finance environment. In their settings,  a borrower continuously ratchets  its term debt upward  until it eventually defaults.  A borrower's incentive to dilute the value of existing  debt claims shape dynamics in these papers.\footnote{This debt ratcheting mechanism resembles a mechanism that drives the price of a durable good to zero in a durable goods monopoly setting without commitment studied by \citet{coase1972durability,stokey1981rational}.  }  %\cite{
} % that generates debt ratcheting. }
%Like    \citet*{aguiar2019take} ,
  \cite*{demarzo2023sovereign} specified an exogenous %and  the corporate finance models mentioned above, 
discount factor process in the tradition of ``small open economy'' models. In contrast, in \cite{LucasStokey1983}, \citet{aiyagari2002optimal},  and \citet{debortoli2021optimal} and our  paper,   the government manipulates an  equilibrium discount factor process within a closed economy. 




\subsection*{Organization}


Section \ref{lssetting}  describes a setting with   an infinite horizon and a discrete  time increment $\Delta >0$. A representative household and benevolent government   participate in competitive markets with distorting flat-rate taxes on labor. The household supplies labor to produce goods that are  allocated between consumption and government expenditures. The government finances an exogenous sequence of government expenditures and an exogenous  initial sequence of debt payments by collecting a flat-rate tax on labor.  A manifold of competitive equilibria with  distorting taxes is indexed by sequences of budget-feasible flat-rate taxes on labor.
%The government presides over the best competitive equilibrium with distorting taxes, an allocation and price system that solves a Ramsey problem.

{%Section \ref{sec:RamseycontRamsey} 
Subsection \ref{subsec31}
 describes how a Ramsey planner chooses tax rates to maximize household utility subject to resource and implementability constraints implied by the household's first-order conditions.     Subsection \ref{subsec32} introduces continuation Ramsey planners who must honor debt structures inherited from previous planners when they decide whether to confirm  the Ramsey planner's intertemporal plan for setting  tax rates. Understanding that, the Ramsey planner chooses how to  restructure the exogenous initial government debt structure. This subsection also  introduces a different timing protocol that 
 lets continuation planners set taxes one period in advance and prevents them  from resetting current-period tax rates.} %and induces % Consequently,  
% continuation planners to voluntarily  confirm  the Ramsey plan, as we show later. %The timing protocol  \NW{To analyze ``time consistency,'' this section
%} %The Ramsey planner also chooses how to  restructure government debt.  
%
%\WJ{Tom: below is your original version of the paragraph for section 3. Let's explain what we have in mind.}  
%Section \ref{sec:RamseycontRamsey} 
%describes how a Ramsey planner chooses tax rates to maximize household utility subject to resource and implementability constraints implied by the household's first order conditons. The Ramsey planner also chooses how to  restructure government debt.  To analyze ``time consistency'', this section introduces  continuation Ramsey planners who must honor debt structures inherited from previous planners when deciding whether to confirm  the Ramsey planner's intertemporal plan for setting  tax rates. \TS{Section \ref{sec:RamseycontRamsey}  also  derives optimal decisions of the Ramsey and continuation planners and  conditions for implementing a  Ramsey plan. }
%%

%Section \ref{sec:RamseycontRamsey}  derives optimal decisions of the Ramsey and 

Section \ref{newimple} presents a  debt management policy that under the new timing protocol implements the optimal  plan.  Our debt management policy  ``takes a short route'' by leaving longer-term debts unchanged and using short-term debt to finance  the resulting accumulated unpaid government liabilities under the optimal plan, defined as 
{%Conditional on the government not having rescheduled  the initial longer-term debt process  $\{b_{-1, i}\}_{i=0}^{\infty}$ before period $j-1$, accumulated unpaid government liabilities from period $0$ to $j-1$ is 
$\Pi_j^\ast= \sum_0^{j-1} \frac{q_{0,i}^\ast}{q_{0,j}^\ast} \big(b_{-1,i}- S_{0,i}^\ast\big) $, where $\{q_{0,i}^\ast\}_{i= 0}^\infty$ and $\{S_{0,i}^\ast\}_{i= 0}^\infty$ are  bond price and  primary surplus processes associated with the Ramsey plan. Note how a  period $i$ flow liability $\big(b_{-1,i}- S_{0,i}^\ast\big) $ times     $\frac{q_{0,i}^\ast}{q_{0,j}^\ast} $  contributes to  the period $j$ value of government liabilities. 
The  $\{\Pi_j^\ast\}_{j= 0}^\infty$ process consolidates  all    equilibrium price and   quantity   information about the  exogenous initial term debt structure  that concerns the Ramsey planner. 
} %conditional on the government not having rescheduled  the initial  term debt process  $\{b_{-1, i}\}_{i=0}^{\infty}$ before period $j-1$.}   conjecture and then 
We use  forward induction to prove that this debt management  policy implements the optimal plan under our  timing protocol.
%which we describe below. 

%\WJ{ 
%A bond price process $\{q_{0,i}^\ast; i\geq 0\}$ and a  primary surplus process $\{S_{0,i}^\ast; i\geq 0\}$ are affiliated with \citeauthor{LucasStokey1983}'s Ramsey plan.  Together with the initial term debt structure $\{b_{-1,i}; i\geq 0\}$, these two objects uniquely determine a process: $\Pi_j^\ast= \sum_0^{j-1} \frac{q_{0,i}^\ast}{q_{0,j}^\ast} \big(b_{-1,i}- S_{0,i}^\ast\big) $ that   accumulates the   government's unpaid  liabilities from period $0$ to $j-1$, conditional on the government not having rescheduled  the initial  term debt process  $\{b_{-1, i}\}_{i=0}^{\infty}$ before period $j-1$.  Our ``take a short route" debt management policy refrains from  restructuring longer-term debts in an initial  debt term structure and instead  instructs continuation governments to issues only  short-term debt. % (i.e., zero in the limit). 
%%to   track  the government's cumulative liability over time. 
%}


{Subsection \ref{sec:im_pro}  connects conditions for implementing the optimal  plan with the invariance and equality  of   Lagrange multipliers on the  implementability constraints that confront the  planner and the continuation planners.}   Subsection \ref{sec:mechanics} conducts a reverse engineering exercise that reveals  mechanics that underlie our implementation. %\NW{Tom: do we need the following details on how we reversed engineered the strategy? W} 
{We work backwards from decision rules that express the tax rates chosen by the  Ramsey planner and continuation Ramsey planners as functions of period-$i$ government expenditures and debt payments as well as the  Lagrange multipliers for  their respective implementability constraints. While the functional forms of these decision rules are identical, for the Ramsey planner and for continuation Ramsey planners, different arguments  enter these functions. For a Ramsey plan to be implemented, these arguments must be  aligned appropriately.  This logic puts restrictions on debt management policies and Lagrange multipliers that imply the  ``short route'' debt management policy together with  the invariance of the respective Lagrange multipliers. }
%that emerges as the tell-tale signature of implementability.}
%\TS{Neng and Wei: our reverse engineering exercise uses {\em our} implementation, not the others that in subsection \ref{sec:otherpolicies} we show also sometimes work. What should we say about this and where should we say it?}
%\NW{Tom: how about adding the following in the next paragraph? please see the magenta part.}



%
%\WJ{Tom: below is your original version of the paragraph for section 3. Let's explain what we have in mind. }  
%
%Section \ref{newimple} describes the timing protocol together with our debt management policy and how it implements the Ramsey plan.   We describe a debt management policy that ``takes a short route'' by leaving longer-term debts unchanged while using short-term debt to finance  accumulated unpaid government liabilities. We use  forward induction to prove that this debt management  policy implements the Ramsey plan under our  timing protocol.
%\TS{Subsection \ref{sec:mechanics}  connects conditions for implementation with the invariance and equality  of   Lagrange multipliers on  implementability constraints facing the Ramsey planner and continuation Ramsey planners.}   Subsection \ref{sec:mechanics} also conducts a reverse engineering exercise that exposes the mechanics that underlie our implementation. 
%\TS{We work backwards from decision rules that express the tax-rate decisions of the  Ramsey planner and continuation Ramsey planners as functions of time-$t$ government expenditures and debt payments as well as the  Lagrange multipliers for  their respective implementability constraints.} Functional forms for these decision rules are all alike, but arguments differ for the Ramsey planner and for continuation Ramsey planners. For a Ramsey plan to be implemented, those arguments must align properly.  \TS{This logic puts restrictions on debt management policies and Lagrange multipliers that induce us to work backwards and recover exactly our ``short route'' debt management policy.  In addition, this analysis identifies the invariance of the respective Lagrange multipliers that emerges as the tell-tale signature of implementability.}
%
%on their respective decision rules

Section \ref{sec:DNYexamples} % describe DNY's examples and how our debt management policy and timing protocol always secure implementability. 
studies  \citeauthor{debortoli2021optimal}'s {counterexamples  in which Lucas and Stokey's debt management policy  fails to implement Ramsey plans when initial debts are too high.} 
We show  how our  modified timing protocol and debt management policy   implement a  Ramsey plan regardless of whether Ramsey tax rates fall above or below the peak of a Laffer curve.      {It is  invariance of Lagrange multipliers for Ramsey and continuation plans that is essential for implementability  (see subsection \ref{subLM}). Subsection \ref{sec:otherpolicies} shows that for some
initial debt structures, but not for others, policies that issue both 
short and  longer term debts  also implement a Ramsey plan. Using both short- and long-term debt can  fail to implement a Ramsey plan, for example, when  initial debts are too  high.}  %If a Ramsey plan exists, our short-term debt management strategy implements it. F

%\TS{Tom: rewrite these two paragraphs after you have edited the examples sections.}
Section \ref{sec:moreexamples} describes additional examples that alter initial debt structures and government expenditure patterns. We  designed these examples to illustrate covariation of outcomes under different exogenous government expenditure and initial term debt structures. {By slightly enriching initial debt structures in \citeauthor{debortoli2021optimal}'s counterexamples, subsection \ref{subsec6.1} shows that  
whether  Lucas-Stokey's implementation  works does not depend on  whether the Ramsey planner sets a tax rate on the `bad' side of the Laffer curve. %is irrelevant for the implementability of a Ramsey plan. 
Subsection \ref{subsecdecliningb} studies Ramsey tax plans and debt management policies that implement them when
initial debt payments decline exponentially toward zero, while subsection \ref{subsecperiodicb} 
 periodic initial debt sequences, and subsection \ref{subsecperiodicG} studies periodic  government spending sequences.  Associated optimal plans display   different patterns of government  accumulations of one-period liabilities  and  consequent issues of one-period debt that illustrate how our implementation works.} 

Section \ref{sec:conclusion} discusses possible extensions. A technical appendix ends our paper. 


\section{Setting} \label{lssetting} 


%\subsection{Environment}

There is no uncertainty.\footnote{\citet[ftnt.~1]{LucasStokey1983}  said that many, if not most, of the issues analyzed in their paper could be analyzed in a context without uncertainty,
like the setting analyzed by \citet{turnovsky1980time}.  \citeauthor{LucasStokey1983}'s footnote 2 about  the consequences of  \citeauthor{turnovsky1980time}'s restricting continuation governments to issue only short term debt is especially interesting in view of the substantial difference between Lucas and Stokey's debt management policy and ours.} Time  $t_i$ is discrete with increment $\Delta >0$, so $\{ t_i=i\Delta, i=0, 1, \cdots\}$.  Period  $i$  means time $t_i.$ A representative household and a benevolent government participate in a complete set of perfectly competitive markets. 
%The government must finance    exogenous  sequences of  expenditures and  distorting flat tax rates. There is no uncertainty. In period $0$, a Ramsey planner selects the best competitive equilibrium. 
The representative household supplies a labor supply sequence $\{N_{0, i}{\Delta}\}_{i=0}^\infty$ that produces a sequence  of a single nonstorable good that must be  allocated between  a consumption sequence  $\{C_{0, i}{\Delta}\}_{i=0}^\infty$ and an exogenous  government expenditure sequence $\{G_{i}{\Delta}\}_{i=0}^\infty$:
\begin{equation}\label{resource_constraint}
	C_{0, i}\Delta + G_{i}\Delta = N_{0, i}\Delta\,, \text{for all}\,  i\geq 0\,,
\end{equation}
{where a first subscript indicates a  time  that a variable is chosen and a second subscript indicates a time  that a variable is realized.} %\NW{Our discrete-time setting follows most of the specifications of \cite{debortoli2021optimal} but uses a different time grid with an interval of $\Delta$.}
%At time $0$, 
The representative household orders consumption and labor supply streams $\{C_{0, i}\}_{i=0}^\infty$ and $\{N_{0, i}\}_{i=0}^\infty$ according to
\begin{equation} \label{utility}
	\sum_{i=0}^\infty \, e^{-\rho i\Delta} U(C_{0, i}, N_{0, i})\Delta\,,
\end{equation}
where $\rho>0$ and $U(\cdot,\cdot)$ is strictly increasing in consumption $C$, strictly decreasing in labor supply $N$, globally concave, and continuously differentiable. Along with \cite{debortoli2021optimal}, we assume that labor supply $\{N_{0, i}\}$ has no upper bound.


{In  period $0$, the  government inherits an initial term-debt structure {$\{b_{-1, i}{\Delta}\}_{i=0}^\infty$} that it must service: %\WJ{we require $b_{-1, i}\Delta\geq 0$ for all $i\geq 0$ and holds strict inequality for some $i\geq 0$.} %where $-1$ in the subscript refers to the  the debt was purchased at time $t_{-1}=-\Delta$. 
	%In period $i \geq 0$, 
	the government must pay its creditors   an amount  $b_{-1, i}\Delta$  during time interval $[i\Delta, {(i+1)}\Delta)$}. %with a coupon rate $b_{-1, t_i}\geq 0$
Let $\tau_{0, i}$ denote a period $i$ flat tax rate that the government sets at period $i=0$ and  let $q_{0, i}$ be the period $0$ value of a zero-coupon bond with a unit payoff at time  $i$. 
The representative household faces a single intertemporal budget constraint:
%\begin{equation*} \label{hh_constraint}
%\sum_{i=0}^\infty C_{0, i}q_{0, i}\Delta \leq \sum_{i=0}^\infty b_{-1, i} q_{0, i}\Delta + \sum_{i=0}^\infty (1-\tau_{0, i})N_{0, i}q_{0, i}\Delta\,.
%\end{equation*}
\begin{equation} \label{hh_constraint}
	\sum_{i=0}^\infty \big( C_{0, i}\Delta \big) \cdot q_{0, i} \leq \sum_{i=0}^\infty \big( b_{-1, i} \Delta \big) \cdot q_{0, i}  + \sum_{i=0}^\infty\big( (1-\tau_{0, t_i})N_{0, i}  \Delta \big) \cdot q_{0, i}  \,.
\end{equation}
The left side of \eqref{hh_constraint} is the present value of the household's consumption stream $\{C_{0, i}\}_{i=0}^\infty$ and the right side is the sum of the household's financial wealth and its human wealth, the present value of after-tax labor income.

Given  tax rate sequence $\{\tau_{0, i}\}_{i=0}^\infty$ and bond price sequence $\{q_{0, i}\}_{i=0}^\infty$, the household chooses $\{C_{0, i}\Delta, N_{0, i}\Delta\}_{i=0}^\infty$ to maximize (\ref{utility}) subject to constraint (\ref{hh_constraint}). First-order necessary conditions for $N_{0,i}$ and $C_{0,i}$ are
\begin{eqnarray}
	1-\tau_{0, i} & = &  - \frac{U_{N, i}}{U_{C, i}}  \label{hhfoc1}  \\
	q_{0, i} & = & e^{-\rho i \Delta} \frac{U_{C, i}}{U_{C, 0}} \,.  \label{hhfoc2}
\end{eqnarray}

%
%The period $0$ government finances the  exogenous spending and debt service sequences that it faces by imposing a sequence $\{\tau_{0,i}\}_{i=0}^\infty$ of flat-rate taxes on labor income.
%\TS{Tom -- tidy up where tax sequence first introduced.}
Let 
\begin{equation} \label{S0i} 
	S_{0, i}\Delta=\left(\tau_{0, i}N_{0, i}-G_{i}\right)\Delta
\end{equation} 
denote the primary government budget surplus over time interval $[i \Delta, (i+1)\Delta )$. Constraint \eqref{hh_constraint} becomes\footnote{Competitive equilibrium and Walras' Law underlies
	this assertion.} 
\begin{equation} \label{gbc} %\left( \tau_{0, t_i}N_{0, i}-G_{i} \right)
	\sum_{i=0}^\infty \, \big(S_{0, i}  \Delta \big) \cdot
	q_{0, i} \geq \sum_{i=0}^\infty \big( b_{-1, i}\Delta \big) \cdot q_{0, i}\,.
\end{equation}
The left side is the time $0$ value of primary government budget surplus $\{{S_{0, i}\Delta}\}_{i=0}^\infty$ and the right side is the  time $0$ value of the government debt  payout sequence  $\{b_{-1, i}\Delta\}_{i=0}^\infty$. 
%\WJ{In addition, the period $0$ government chooses a debt structure $\{b_{0, i}\}_{i=1}^\infty$ for the period $1$ government.} \TS{tom move Wei's sentence.}

\begin{definition}
	Given the government's initial  term-debt structure  $\{b_{-1, i}\Delta\}_{i=0}^\infty$ and  expenditure sequence $\{G_{i}\Delta\}_{i=0}^\infty$, a competitive equilibrium is a feasible  allocation $\{C_{0, i}\Delta, N_{0, i}\Delta\}_{i=0}^\infty$, % satisfying (\ref{resourceconstraint}),
	a flat-rate tax sequence $\{\tau_{0, i} \}_{i=0}^\infty$,  and a bond price sequence $\{q_{0, i}\}_{i=0}^\infty$ for which	\begin{itemize}
		\item  the government's budget constraint \eqref{resource_constraint} is satisfied, and 
		\item  given  government spending,  tax, and bond price sequences, the allocation solves the household's optimization problem
	\end{itemize}
\end{definition}


We follow \cite{LucasStokey1983} and use first-order conditions (\ref{hhfoc1}) and (\ref{hhfoc2}) together with feasibility constraint (\ref{resource_constraint}) to eliminate tax rates and bond prices from the government budget constraint (\ref{gbc}) and thereby to deduce the following implementability constraint on competitive equilibrium allocations $\{C_{0, i}\Delta, N_{0, i}\Delta\}_{i=0}^\infty$:
\begin{equation} \label{impletconstraint}
	\sum_{i=0}^\infty \, e^{-\rho i\Delta} \left[ C_{0, i}U_{C, i}+(C_{0, i}+G_{i})U_{N, i} - b_{-1, i} U_{C, i} \right]\cdot\Delta \geq 0 \,.
\end{equation}

%\begin{equation} \label{impletconstraint}
%	\sum_{i=0}^\infty \, e^{-\rho i\Delta} \left[ C_{0, i}U_{C, i}+(C_{0, i}+G_{i})U_{N, i} \right]\cdot\Delta \geq \sum_{i=0}^\infty\, e^{-\rho i\Delta} \big( b_{-1, i} U_{C, i}\big) \cdot \Delta\,.
%\end{equation}

\begin{proposition}\label{propo::1}
	Given an initial  term-debt  structure  $\{b_{-1, i}\Delta\}_{i=0}^\infty$ and government spending sequence $\{G_{i}\Delta\}_{i=0}^\infty$, an allocation $\{C_{0, i}\Delta, N_{0, i}\Delta\}_{i=0}^\infty$ is a competitive equilibrium allocation if and only if it satisfies  (\ref{impletconstraint}) at $i=0$ and  (\ref{resource_constraint}) for all $i\geq 0$. % and (\ref{impletconstraint}) at $i=0$. 
\end{proposition}

%\paragraph{Initial Debt Structure.}
%For the east of notation, denote $t_i=i\Delta, i=0, 1, \cdots$. The time$-0$ government confronts a term-debt structure $\{b_{0, t_i}; i=0, 1, \cdots\}$. We follow most of the specifications of DNY (2021) but use a different time grid with an interval of $\Delta$. At time $t_i$, an amount of $b_{0, t_i}\Delta$ will be due where $b_{0, t_i}\geq 0$ for all $i=0, 1, \cdots$.

\section{Ramsey  and Continuation Ramsey Plans\label{sec:RamseycontRamsey}}
%\paragraph{Ramsey Problem.}

A Ramsey planner and a sequence of continuation Ramsey planners are linked together by a  structure of `partial commitments'  like those described by \citet[p.~69]{LucasStokey1983} in which debts incurred by previous planners must be paid or rescheduled at equilibrium prices, but sequences of tax rates can be redesigned each period. 
\subsection{Ramsey Plan} \label{subsec31}
A Ramsey planner chooses the best competitive equilibrium. It can accomplish this by  choosing a sequence   ${\{C_{0, i}\Delta\}}_{i=0}^\infty$  that maximizes  (\ref{utility}) 
%%\eqref{utility}.
%\begin{align} \label{obj}
%	&  \sum_{0}^{\infty}  e^{-\rho t} U(C_{0, t}, C_{0, t}+G_t) dt  %\\
%	%s.t. &  %\\
%\end{align} %+B_{0}U_{C}
%\begin{equation} \label{utility2}
%\sum_{i=0}^\infty \, e^{-\rho i\Delta} U(C_{0, i}, C_{0, i}+G_i)\Delta
%\end{equation}
subject to resource constraint (\ref{resource_constraint}) and 
%the smoothness condition (\ref{smooth}) and the 
implementability constraint \eqref{impletconstraint}.  Given that sequence, we can use  first-order necessary conditions
\eqref{hhfoc1} and \eqref{hhfoc2} to 
compute   associated equilibrium tax rate  and bond price sequences. 

%\NW{Can we eliminate (\ref{utility2}) and refer to (\ref{utility})?} 

We denote competitive equilibrium objects associated with a Ramsey plan as 
$$ \{C_{0, i}^\ast\Delta, N_{0, i}^\ast\Delta, \tau_{0, i}^\ast, q_{0, i}^\ast\}_{i= 0}^\infty. 
$$
%Notice that among these objects  we did not reschedule debt service sequence $\{b_{0,i}^\ast\}_{i=1}^\infty$. Actually, 
An    equivalence class $\mathcal{D}$ of continuation debt structures $\{b_{0,i}\}_{i=1}^\infty$   supports a Ramsey tax plan.\footnote{In  complete market structures like ours, this is  a typical portfolio irrelevance  outcome.} %\NW{Not sure if this is a perfect analogy... Maybe Arrow-Debreu if you prefer?} 

%
%
%\paragraph{Short-term Debt as a Forward Contract.}
%\WJ{\bf Discuss!!!} Next, we discuss short-term debt in this discrete-time environment. In fact, the short-term debt is a forward contract subject to cash flow below
%\[ (- 1, \beta_{0,t_{i+1}}/ \beta_{0,t_i} )  = (-1, e^{r_{0,t_i}\Delta}) \] for $(t_i, t_{i+1})$, where 
%\[ \beta_{0,t_{i+1}}  = e^{r_{0,t_i}\Delta} \beta_{0,t_i}  \] 
%and $\beta_{0, t_i}$ is the value of short-term debt at period $t_i$ for one unit a government borrowed at time $0$, and $r_{0, t_i}$ is the short-rate over the period $[t_i, t_{i+1}]$\,.
%%In our setting, it is key to have short-term debt, by which we mean 
%%\[  \beta_{0,t_{i+1}} =e^{r_{0,t_i}\Delta} \beta_{0,t_i}  \] 
%This actually is a forward contract agreed at $t_0=0$ with a forward date $r_{0, t_i}$ at $t_i$, and the contract matures in one period ahead. 

%\WJ{What about figures and earlier part in our cont time version?} 
\paragraph{Characterizing Ramsey plan.}  
%\subsection{Ramsey Plan} 
%\WJ{Different greek letters for the multiplier?} 
Attach a  nonnegative multiplier  $\Phi_0$  to implementability constraint  (\ref{impletconstraint}) and form  a Lagrangian: 
\begin{align}\label{L0}
	{\mathcal L}^{}_0 & =  \sum_{i=0}^{\infty}  e^{-\rho i\Delta}\left[ U(C_{0, i}, C_{0, i}+G_{i})+ \Phi_0^{}\left( C_{0, i}U_{C, i}+(C_{0, i}+G_{i})U_{N, i}-b_{-1, i}^{} U_{C, i}\right)  \right] \Delta \, . %\nonumber \\%- \Psi_0(b_{0, t}^{} U_C- b_{0-, t} U_C)	\right]dt \\
	%&  \quad \quad \;  \, .
\end{align} %  to  minimize   \eqref{L0}
The Ramsey planner  maximizes the right side of  \eqref{L0} over consumption plans $\{C_{0, i}\}_{i=0}^\infty$   and minimizes over multiplier $\Phi_0$.  Call the extremizing values: $\{C_{0, i}^*\}_{i=0}^\infty$ and $\Phi_0^*$.
% Maximize under the integral sign on the right side of \eqref{L0} to get a continuum of static problems: 
%% We characterize the optimal policy for time $0+$ onward as follows.  For a given time interval $(t, t+dt), t>0$, the planner chooses consumption to solve
%\begin{equation} \label{optp-L}
%	\max_{C_{0, t}} \;\; \left[ U(C_{0, t}, C_{0, t}+G_t)+ \Lambda_0^{}\left( C_{0, t}U_C+(C_{0, t}+G_t)U_N-b_{0, t} U_C \right) \right]dt\,, \;\; t \geq 0
%\end{equation} %If there are no constraints on $C_{0, t}$, then 
For a given $\Phi_0$, the optimal consumption plan satisfies%\footnote{We assume that the optimal consumption plan for the problem defined in   \eqref{Csoln}   is unique for a  given $\Phi_0$ under certain regularity condition.} 
\begin{equation}\label{Csoln}
	C(\Phi_0; b_{-1, i}, G_{i}):=\argmax_{C_{0, i}}  \left[ U(C_{0, i}, C_{0, i}+G_{i})+ \Phi_0\left( C_{0, i}U_{C, i}+(C_{0, i}+G_{i})U_{N, i}-b_{-1, i} U_{C, i} \right) \right] .
\end{equation}
%solve problem (\ref{optp-L}) 
%, and $\tau(\Lambda_0^{}, b_{0, t}, G_t)$, respectively.
Substituting  \eqref{Csoln} into  implementability condition (\ref{impletconstraint}), we obtain
\begin{align}\label{eqPhi}
	& \sum_{i=0}^{\infty} e^{-\rho i\Delta} \big[ C(\Phi_0; b_{-1, i}, G_{i})U_{C, i}+(C(\Phi_0; b_{-1, i}, G_{i})+G_i)U_{N, i}-b_{-1, i} U_{C, i} \big]\Delta \geq  0\,.
	%&&  +B_{0}U_{C}( C(\Lambda_0^{}, b_{0, 0}, G_{0}) ,C(\Lambda_0^{}, b_{0, 0}, G_{0})+G_{0}) .
\end{align}
%\footnote{Contradiction arises when the Lagrangian multiplier $\Phi_0=0$. That is, when $\Phi_0=0$, $C(0; b_{-1, i}, G_i)=\argmax_{C}\, U(C, C+G_i)$, which implies $U_{C}+U_{N}=0$. Substituting into the implementability condition (\ref{impletconstraint}), one obtains $\sum_{i=0}^{\infty} e^{-\rho i\Delta} \big[-G_i-b_{-1, i}\big]U_{C, i}\Delta\geq 0  $, which is a contradiction.    } 
%\begin{eqnarray}\label{eqPhi}
%	&& \int_{0}^{\infty} e^{-\rho t} ( C(\Lambda_0^{*}, b_{0, t}, G_t)U_C+(C(\Lambda_0^{*}, b_{0, t}, G_t)+G_t)U_N )  dt  =  \int_{0}^{\infty} e^{-\rho t}b_{0, t} U_C dt\,.
%	%&&  +B_{0}U_{C}( C(\Lambda_0^{}, b_{0, 0}, G_{0}) ,C(\Lambda_0^{}, b_{0, 0}, G_{0})+G_{0}) .
%\end{eqnarray}
%where we impose the smooth consumption condition
%\begin{equation}
%	b_{0-, t} \text{is cadlag sequence.}
%\end{equation}
%\NW{Wei: maybe this is left continuous with right limit? cadlag sequence? } 
%\WJ{Need to justify this condition.  }
Substituting  $C(\Phi_0; b_{-1, i}, G_{i})$ into the right side of   (\ref{L0}) lets us  write  $\mathcal L_0$ as a function %of $\Lambda_0^*$:
$\mathcal{L}_0(\Phi_0)$. 
An optimal  $\Phi_0^{*}$ that maximizes $\mathcal{L}_0(\Phi_0)$ must satisfy (\ref{eqPhi}). 
The Ramsey allocation can be represented with two functions, namely,  $C_{0, i}^*=C(\Phi_0^{*}; b_{-1, i}, G_{i})$ and $N_{0, i}^* = C_{0, i}^*+G_{i} = N(\Phi_0^{*}; b_{-1, i}, G_{i})$. 
%\WJ{As we will show below, under regularity conditions,} the existence of a Ramsey plan requires   % Lagrange multiplier 
%$\Phi_0^*>0$, which implies (\ref{eqPhi}) holds with equality. 
%\NW{	$C(\Lambda_0, b_{0, t}, G_t)  =  \argmax_{C_{0, t}}  \left[ U(C_{0, t}, C_{0, t}+G_t)+ \Lambda_0^{}\left( C_{0, t}U_C+(C_{0, t}+G_t)U_N-b_{0, t} U_C \right) \right]$.} 

%In summary, we conclude with

%\NW{make this prettier and maybe add an eqn as you noted.} 




%_{\{\Lambda_0 =  \textrm{a positive root of}  (\ref{eqPhi})\}}
\begin{proposition} \label{proplag} 
	%When %Assuming  , 
	%	Under plausible regularity conditions,\footnote{We provide two regularity conditions to guarantee that  a Ramsey plan  exists. First,  for a given initial debt structure, the admissible allocation set is  not empty, so  there exists allocations $\{C_{0, i}\Delta\}_{i=0}^\infty$ that satisfy the implementability condition (\ref{impletconstraint}). Second, %for the given utility function $U(C, C+G)$, 
		%		$\left(CU_C + (C+G) U_N - bU_C \right)$ is concave in $C$ when $G>0, b\geq 0$. 
		%	For  economies analyzed  in \citet{debortoli2021optimal} these two regularity conditions are satisfied when initial debt is not too high.} 
	%	When a Ramsey plan exists, 
	{Under  regularity conditions \ref{condition1} and \ref{condition2} described in Appendix \ref{appendix:A}}, 
	(\ref{eqPhi}) holds with equality and there  exists  a positive root  $\Phi_0^*$ of  (\ref{eqPhi}) that solves
	\begin{equation*}
		\Phi_0^*  = \argmin_{\Phi_0}  \,\,  \mathcal{L}_0(\Phi_0)\, , %(\{C_{0,t}^*\}_{t=0}^\infty,  
	\end{equation*}	
	where the    Ramsey allocation,  tax rate, and bond price  sequence satisfy
	\begin{eqnarray}
		&	 {C}^\ast_{0, i} =     C(\Phi_0^{*}; b_{-1, i}, G_{i}); \, \quad  
		{N}^*_{0, i}  =    {C}_{0, i}^\ast +G_{i} ;  \label{RamseyC} \\
		&	  {\tau}^*_{0, i} =  1+\frac{U_{N, i}({C}^*_{0, i}, {N}^*_{0, i})}{U_{C, i}({C}^*_{0, i}, {N}^*_{0, i})};\,   \quad
		{q}^*_{0, i}  = e^{-\rho i\Delta}\frac{U_{C, i}({C}^*_{0, i}, {N}^*_{0, i})}{U_{C, 0}({C}^*_{0, 0}, {N}^*_{0, 0})} \label{Ramseytaxprice} \, .
	\end{eqnarray}
	%	        \begin{eqnarray*}
		%	        	 C(\Phi_0^{\mathcal L}, b_{0-, 0}, G_0)= \lim _{s\downarrow 0} C(\Phi_0^{\mathcal L}, b_{0-, s}, G_s)
		%	        \end{eqnarray*}	
	%	where 
	%and the marginal utilities $U_C$ and $U_N$ are evaluated at $\widetilde{C}_{0, t},\widetilde{N}_{0, t}$.
	The value ${\mathcal L}^{*}_0 $ of a Ramsey plan is
	%	\begin{equation*}
		$	 \sum_{i=0}^{\infty}  e^{-\rho i\Delta} U(	{C}^*_{0, i}, {C}^*_{0, i}+G_{
			i})  \Delta.
		$	%\end{equation*}
	\end{proposition}
	See Appendix \ref{appendix:A} for a proof and supporting technical details. 
	
	
	\subsection{Continuation Ramsey Plans} \label{subsec32}
	%The Ramsey planner is a Stackelberg leader and the  representative household is   Stackelberg follower.
	%\TS{Tom -- here bring in material on slides about portfolio indeterminacy.}
	
	
	There is a sequence of continuation Ramsey planners indexed by $j\in \{1, 2, \cdots \}$, each of whom % are Stackelberg followers. 
	%A continuation Ramsey planner 
	manipulates   a  continuation representative household. %another Stackelberg follower.   %which is a subservient Stackelberg follower. 
	A period $j$ continuation planner  confronts a debt structure {$\{ b_{j-1,i}\}_{i=j}^\infty$} chosen by the preceding  planner and the tail government spending rate process {$\{ G_{i}\}_{i=j}^\infty$}. 
	The period $j$ continuation planner  faces an intertemporal budget constraint:
	% This implies the following time$-0$ government budget constraint:
	\begin{equation}\label{ic} %\scriptsize{
			\sum_{i=j}^\infty \big( S_{j,i} \Delta  \big) \cdot q_{j, i} \geq \sum_{i=j}^\infty \big({ b_{j-1, i} }\Delta \big) \cdot q_{j, i}  \,, \,  %}
	\end{equation}
	where %\begin{equation}
	$S_{j, i}\Delta=\left(\tau_{j, i}N_{j, i}-G_{i}\right)\Delta
	$ %\end{equation} 
	is the primary  surplus for period $j$ government over time interval $[i \Delta, (i+1)\Delta )$.  
	The period $j$ continuation planner  seeks  continuation sequences    $\{ C_{j, i}\, , N_{j, i}\}_{i=j}^\infty$ that  maximize
	\begin{equation}\label{eqn:preferences}%\scriptsize{
			\sum_{i=j}^\infty \, e^{-\rho (i-j)\Delta} U(C_{j, i}, N_{j, i}) \Delta  % \NW{\Delta}  %}
	\end{equation}
	subject to continuation implementability constraint (CIC): 
	\begin{align}\label{con_IC0}
		& \sum_{i=j}^{\infty} e^{-\rho(i-j) \Delta} \left(  C_{j, i}U_{C, i} + (C_{j, i}+G_{i})U_{N, i} -b_{j-1, i} U_{C, i} \right) \Delta \geq 0\, .
	\end{align}
	% in a minor but substantial way.  
	Our period $j$ continuation planner confronts the tail government spending rate process of {$\{ G_{i}\}_{i=j}^\infty$}  and  a debt structure {$\{ b_{j-1,i}\}_{i=j}^\infty$} inherited from the preceding  planner.
	The continuation planner cannot reset  {$\tau_{j-1, j}$}. The period $j$ planner  sets  a flat-rate tax {$\tau_{j, j+1}$} that the  period $j+1$ continuation planner {cannot} reset and sets (actually, just ``recommends'') a continuation tax plan  {$   \{\tau_{j, i}\}_{i=j+2}^{\infty} $} that  a  period $j+1$ continuation planner is free to reset.
	Knowing the period $j$ continuation government's tax plan, a period $j$  continuation household  orders $\{ C_{j, i}\, , N_{j, i}\}_{i=j}^\infty$ according to the continuation utility functional
	\eqref{eqn:preferences}.
	
	\medskip
	
	
	
	
	%In period $0$, the Ramsey planner confronts an initial debt \WJ{$\{ b_{-1,j}\}_{j=0}^\infty$} and government spending rate processes \WJ{ $\{ G_{j}\}_{j=0}^\infty$}, pays \WJ{$b_{-1, 0}\Delta $} and $G_0\Delta$, and sets \TS{$\tau_{0, 0}$}. She then sets \TS{$(b_{0, 1}, \tau_{0, 1})$}  that time $t_1 = \Delta$ continuation planner \TS{cannot}  reset. Finally, she recommends  \NW{$  \big( \{ {b}_{0, j}\}_{j=2}^{\infty} \;, \{\tau_{0, j}\}_{j=2}^{\infty}\big) $}   that time $t_1 = \Delta$ continuation planner \NW{can} reset. 
	%
	%Then for a period $1$ continuation planner that confronts a debt structure \WJ{$\{ b_{0,j}\}_{j=1}^\infty$} inherited from precedent government and the tail government spending rate processes of \WJ{$\{ G_{j}\}_{j=1}^\infty$}, she pays \WJ{$b_{0, 1}\Delta$} and $G_0\Delta$. Unlike the protocol in Lucas-Stokey, the continuation planner in our timing protocol cannot reset \WJ{$\tau_{0, 1}$}. Instead, she sets \TS{$(b_{1, 2}, \tau_{1, 2})$}  that time $t_2 = 2 \Delta$ continuation planner \TS{cannot} reset and recommends \NW{$  \big( \{ {b}_{1, j}\}_{j=3}^{\infty} \;, \{\tau_{1, j}\}_{j=3}^{\infty}\big) $} that $t_2 = 2 \Delta$ continuation planner \NW{can} reset. 
	
	
	%Then for a period $i$ continuation planner that confronts a debt structure \WJ{$\{ b_{i-1,j}\}_{j=i}^\infty$} inherited from precedent government and the tail government spending rate processes of \WJ{$\{ G_{j}\}_{j=i}^\infty$}, she pays $\WJ{b_{i-1, i}}$ and $G_{i}\Delta$. Unlike the protocol in Lucas-Stokey, the continuation planner cannot reset  \WJ{$\tau_{i-1, i}$}. Instead, she sets sets  \TS{$(b_{i, i+1}, \tau_{i, i+1})$} that period $i+1$ continuation planner \TS{cannot} reset and recommends \NW{$  \big( \{ {b}_{i, j}\}_{j=i+2}^{\infty} \;, \{\tau_{i, j}\}_{j=i+2}^{\infty}\big) $} that  period $i+1$ continuation planner \NW{can} reset.
	
	
	
	
	
	
	
	%\subsection{\normalsize Continuation {Debt  Structure}} % in the spirit of \citeauthor{LucasStokey1983}}
%\paragraph{Continuation debt term structure}

%At   Ramsey  allocation $\{C_{0, j}^\ast\Delta, N_{0, j}^\ast\Delta\}_{j= 0}^\infty$,  flat tax rate sequence $\{\tau_{0, j}^\ast\}_{j=0}^\infty$, and  price sequence $\{q_{0, j}^\ast\}_{j= 0}^\infty$, many  {debt   structures} $\{ b_{{i-1},j}\Delta \}_{j=i}^\infty$  satisfy  period $i\geq 1$ continuation governments' budget constraints: 
%\begin{eqnarray}%\scriptsize{
%	%	  \sum_{j=0}^\infty \left( \hat b_{{i-1}, j} \Delta\right) 	q_{0, j}^*  &\leq & 	\sum_{j=i}^{\infty}  \left(\tau_{0, j}^\ast N_{0, j}^*-G_{j}\right)\Delta	q_{0, j}^*  \, \label{eqhatb2}\,,  %}
%	%		\end{equation*} 	\begin{equation*}%\scriptsize{
%		  	\sum_{j=i}^{\infty}  \left(\tau_{0, j}^\ast N_{0, j}^*-G_{j}\right)\Delta	\cdot q_{0, j}^* &\geq & \sum_{j=i}^\infty \big(  b_{{i-1}, j} \Delta\big) \cdot	q_{0, j}^*    \, \label{eqhatb2}   \, . %}
%\end{eqnarray}
%\NW{Here we use $\widehat B_{t_i}$ to denote the balance/stock of short-term debt purchased at time $t_{i-1}$ and due at time $t_i$. This notation differentiates from \cite{LucasStokey1983} that uses a small amount of one-period debt, together with other term debt, denoted as $\{\hat b_{0,t_i}\Delta \}_{i=1}^\infty$. }

\begin{definition} \label{defprotocol}
A continuation of a Ramsey plan in period $j$ is the {tail  
	of the Ramsey plan for $i \geq j$:  $\{C_{0, i}^\ast\Delta, N_{0, i}^\ast\Delta, \tau_{0, i}^\ast, q_{0, i}^\ast\}_{i= j}^\infty$}.
A continuation Ramsey planner in period $j$ must  administer tax rate $\tau_{j-1, j}$ for period $j$ and finance an inherited continuation {debt  structure}  $\{ b_{{j-1}, i}\Delta \}_{i=j}^\infty$  but is free to  choose a   continuation tax rate  plan    $\{\tau_{j, i}\}_{i=j+1}^\infty$. %    that need  not equal the continuation $\{\tau_{0, j}^\ast\}_{j=i+1}^\infty$ of the original Ramsey tax-rate plan.
It also passes   a  rescheduled continuation  {debt  structure} $\{ b_{j, i}\Delta \}_{i=j+1}^\infty$ on   to a  period ${i+1}$ continuation planner.
A Ramsey plan  is said to be \textbf{implemented} if each continuation Ramsey planner confirms the  Ramsey planner's tax rate plan {$\{\tau_{0,i}^\ast\}_{i=0}^\infty$}
\end{definition}

%\WJ{Neng: Tom suggests that we consider removing Figure 1, as it appears to be redundant with the accompanying text. }
%
%Figure \ref{fig:timing}  illustrates the   timing protocol from period $0$ to period $1$. In period $0$, the Ramsey planner confronts an initial debt {$\{ b_{-1,i}\}_{i=0}^\infty$} and government spending rate processes i{ $\{ G_{i}\}_{i=0}^\infty$}, pays {$b_{-1, 0}\Delta $} and $G_0\Delta$, and sets {$\tau_{0, 0}$}. She then sets {$(b_{0, 1}, \tau_{0, 1})$}  that time $t_1 = \Delta$ continuation planner {cannot}  reset. Finally, she recommends  {$  \big( \{ {b}_{0, i}\}_{i=2}^{\infty} \;, \{\tau_{0, i}\}_{i=2}^{\infty}\big) $}   that time $t_1 = \Delta$ continuation planner {can} reset. Knowing these tax rate  and debt structure   sequences,  at  time $0$ the representative agent chooses
%$\{C_{0, i}\Delta, N_{0, i}\Delta\}_{i=0}^\infty$.\footnote{In  \citeauthor{debortoli2021optimal}'s laboratory, a single continuation Ramsey planner in period $1$ confronts    $\{G_{i}\Delta\}_{i=1}^\infty$ and $\{ b_{0, i}\Delta\}_{i=1}^\infty$ that it must finance
%with a period $1$ continuation tax plan $\{\tau_{1, i}\}_{i=1}^\infty$.}
%


%
%%	\NW{Remover the hat in the figure?} 
%
%\begin{figure}[]
%\caption{Timing protocol for Ramsey planner and time $t_1 = \Delta$ continuation planner.}
%\centering
%\begin{tikzpicture}[scale=0.95]
%
%\draw [thick,->] (-1.2,0) -- (9.7,0);
%%\draw [thick,-] (-0.19,0.1) -- (-0.19,0);
%
%%\node[align=center] at (4.1,1.3) {\bf  \ZL{ \scriptsize Step 2}};
%\node[align=center] at (4.3, 1.3) {\bf  {\tiny Household}};
%\node[align=center] at (5.2, 0.9) {\bf  {\tiny chooses $\{C_{0, i}, N_{0, i}\}_{i=0}^{\infty}$}};
%
%\draw [thick,->] (-1.2,0.1) -- (-1.2,0.5);
%\node[align=center] at (-1.2, 2.7) {\bf \WJ{\tiny $\{ b_{-1, i}\}_{i=0}^{\infty},\; \{G_{i}\}_{i=0}^{\infty}$  }};
%\node[align=center] at (-1.2, 2.1) {\bf \TS{\tiny Ramsey planner,}};
%\node[align=center] at (-1.2, 1.7) {\bf {\tiny pays $\WJ{b_{-1, 0}}$ and sets $\TS{\tau_{0, 0}}$,}};
%\node[align=center] at (-1.2, 1.3) {\bf {\tiny sets \TS{$(b_{0,1}, \tau_{0,1})$}, }};
%\node[align=center] at (-1.2, 0.9) {\bf {\tiny sets $  \NW{\left( \{ {b}_{0, i}\}_{i=2}^{\infty} \;, \{\tau_{0, i}\}_{i=2}^{\infty}\right) } $	
	%}};
	%
	%%\node[align=center] at (-1, 1.3) {\bf  \NW{ \scriptsize  Step 3}};
	%
	%%\draw [thick,-] (5.20,0.23) -- (5.20,0);
	%\draw [thick,decorate,decoration={brace,amplitude=12pt,raise=4pt}] (-1.15,0) -- (9.5,0);
	%
	%
	%%\node[align=center] at (8.9, 1.3) {\bf   \YC{\scriptsize  Step 1}};
	%%	\node[align=center] at (8.9, 0.9) {\bf   \TS{\scriptsize Agent pays $M_{\tau}$  }};
	%%	\draw [thick,-] (9.4,0.23) -- (9.4,0);
	%%\draw [thick,decorate,decoration={brace,amplitude=12pt,raise=4pt}] (5.25,0) -- (9.35,0);
	%
	%
	%%\draw [thick,decorate,decoration={brace,amplitude=12pt,raise=4pt}] (9.45,0) -- (10.75,0);
	%
	%\draw[fill,color=blue] (-1.2,0.005) circle (.05);	
	%%\draw[fill,color=black] (-0.2,0.005) circle (.05);
	%
	%%\node[align=center] at (5.2,-0.3) {\footnotesize $ \tau^d$};
	%%\node[align=center] at (4.2,-0.3) {$\tau_F$};
	%\node[align=center] at (-1.2,-0.3) {\footnotesize $0 $};
	%\node[align=center] at (10.1,0) {\;\footnotesize  \; time};
	%\node[align=center] at (9.6,-0.3) {\footnotesize $\infty $};
	%
	%\end{tikzpicture}
	%
	%\bigskip
	%
	%\begin{tikzpicture}[scale=1]
	%
	%
	%\draw [thick,->] (1.5,0) -- (9.7,0);
	%%\draw [thick,--] (-1.2,0) -- (9.7,0);
	%%\draw [thick,-] (-0.19,0.1) -- (-0.19,0);
	%
	%%\node[align=center] at (4.1,1.3) {\bf  \ZL{ \scriptsize Step 2}};
	%\node[align=center] at (5.7, 1.3) {\bf  {\tiny Household}};
	%\node[align=center] at (6.5, 0.9) {\bf  {\tiny chooses $\{C_{1, i}, N_{1, i}\}_{i=1}^{\infty}$}};
	%
	%\draw [thick,->] (1.5,0.1) -- (1.5,0.5);
	%\node[align=center] at (0.6, 2.5) {\bf \WJ{\tiny $\{ {b}_{0, i}\}_{i=1}^{\infty},\; \{G_{i}\}_{i=1}^{\infty}$  }};
	%\node[align=center] at (-0.8, 0.9) {\bf \color{white}{\tiny Ramsey planner commits }};
	%\node[align=center] at (0.6, 2.1) {\bf \TS{\tiny Continuation Ramsey planner}}; 
	%\node[align=center] at (0.6, 1.7) {\bf {\tiny confronts \WJ{$\left(b_{0, 1}, \tau_{0, 1}\right)$},}};
	%\node[align=center] at (0.6, 1.3) {\bf {\tiny sets \TS{$(b_{1,2}, \tau_{1,2})$}, }};
	%\node[align=center] at (0.6, 0.9) {\bf {\tiny  sets  $ \NW{ \left( \{ {b}_{1, i}\}_{i=3}^{\infty} \;, \{\tau_{1, i}\}_{i=3}^{\infty}\right) } $}};
	%
	%%\node[align=center] at (-1, 1.3) {\bf  \NW{ \scriptsize  Step 3}};
	%
	%%\draw [thick,-] (5.20,0.23) -- (5.20,0);
	%\draw [thick,decorate,decoration={brace,amplitude=12pt,raise=4pt}] (1.6,0) -- (9.5,0);
	%
	%
	%%\node[align=center] at (8.9, 1.3) {\bf   \YC{\scriptsize  Step 1}};
	%%	\node[align=center] at (8.9, 0.9) {\bf   \TS{\scriptsize Agent pays $M_{\tau}$  }};
	%%	\draw [thick,-] (9.4,0.23) -- (9.4,0);
	%%\draw [thick,decorate,decoration={brace,amplitude=12pt,raise=4pt}] (5.25,0) -- (9.35,0);
	%
	%
	%%\draw [thick,decorate,decoration={brace,amplitude=12pt,raise=4pt}] (9.45,0) -- (10.75,0);
	%
	%%\draw[fill,color=white] (-1.2,0.005) circle (.05);	
	%\draw[fill,color=blue] (1.5,0.005) circle (.05);
	%%\draw[fill,color=black] (-0.2,0.005) circle (.05);
	%
	%%\node[align=center] at (5.2,-0.3) {\footnotesize $ \tau^d$};
	%%\node[align=center] at (4.2,-0.3) {$\tau_F$};
	%%\node[align=center] at (-1.2,-0.3) {\footnotesize $0 $};
	%\node[align=center] at (1.5,-0.3) {\footnotesize $\Delta $};
	%\node[align=center] at (10.1,0) {\;\footnotesize \; time};
	%\node[align=center] at (9.6,-0.3) {\footnotesize $\infty $};
	%
	%\end{tikzpicture}
	%
	%\label{fig:timing}	
	%\end{figure}
	
	%A continuation Ramsey planner in period $i$ must honor an inherited continuation {debt  structure}  $\{ b_{{i-1}, j}\Delta \}_{j=i}^\infty$   but is free to  choose a   continuation tax rate  plan    $\{\tau_{i, j}\}_{j=i}^\infty$    that need  not equal the continuation $\{\tau_{0, j}^\ast\}_{j=i}^\infty$ of the original Ramsey tax-rate plan. It also chooses   a  continuation new {debt  structure} $\{ b_{i, j}\Delta \}_{j=i+1}^\infty$ to hand off  to a  period ${i+1}$ continuation planner.
	%A Ramsey plan is said to be	 implemented 	 when  continuation Ramsey planners  at  all period $i>0$ choose to confirm it.
	
	
	%  \WJ{Tom: time $0$ Ramsey  
% \noindent Thus, a  continuation Ramsey planner at $s>0$ that must honor 
%  debt term structure $\{\hat b_{0,t}\}_{t=s}^\infty$ constructs a %that was chosen by Ramsey a 
% continuation  plan.


%	\TS{\tiny Do not mention  Ricardian equivalence and MM here as they refer to the first-best world allocation where taxes are not distortionary?} 

%	\TS{Neng and Wei: in the above limits, I want $0$ in the lower limit in all 3 places}

%\TS{I changed the lower bound to 1 from 0 for the integrals} 




\paragraph{Constructing continuation Ramsey plan.} 


Attach a  nonnegative multiplier  $\Phi_j$  to CIC (\ref{con_IC0}) and form  a Lagrangian: 
\begin{align}\label{Lj}
	{\mathcal L}^{}_j & =  \sum_{i=j}^{\infty}  e^{-\rho (i-j)\Delta}\left[ U(C_{j, i}, C_{j, i}+G_{i})+ \Phi_j^{}\left( C_{j, i}U_{C, i}+(C_{j, i}+G_{i})U_{N, i}-b_{j-1, i}^{} U_{C, i}\right)  \right] \Delta \, . %\nonumber \\%- \Psi_0(b_{0, t}^{} U_C- b_{0-, t} U_C)	\right]dt \\
	%&  \quad \quad \;  \, .
\end{align} %  to  minimize   \eqref{L0}
Call the extremizing values: $\{C_{j, i}^*\}_{i=j}^\infty$ and $\Phi_j^*$. Note that under our timing protocol, continuation households' period $j$ consumption satisfies $C_{j, j}^*=C_{j-1, j}^*$ because of the FOCs.
%The Ramsey planner  maximizes the right side of  \eqref{L0} over consumption plans $\{C_{0, i}\}_{i=0}^\infty$   and minimizes over multiplier $\Phi_0$.  
% Maximize under the integral sign on the right side of \eqref{L0} to get a continuum of static problems: 
%% We characterize the optimal policy for time $0+$ onward as follows.  For a given time interval $(t, t+dt), t>0$, the planner chooses consumption to solve
%\begin{equation} \label{optp-L}
%	\max_{C_{0, t}} \;\; \left[ U(C_{0, t}, C_{0, t}+G_t)+ \Lambda_0^{}\left( C_{0, t}U_C+(C_{0, t}+G_t)U_N-b_{0, t} U_C \right) \right]dt\,, \;\; t \geq 0
%\end{equation} %If there are no constraints on $C_{0, t}$, then 
For a given $\Phi_j$, the optimal consumption plan for $i\geq j+1$ satisfies%\footnote{We assume that the optimal consumption plan for the problem defined in   \eqref{Csoln}   is unique for a  given $\Phi_0$ under certain regularity condition.} 
\begin{equation}\label{Csolnj}
	C(\Phi_j; b_{j-1, i}, G_{i}):=\argmax_{C_{j, i}}  \left[ U(C_{j, i}, C_{j, i}+G_{i})+ \Phi_j\left( C_{j, i}U_{C, i}+(C_{j, i}+G_{i})U_{N, i}-b_{j-1, i} U_{C, i} \right) \right] .
\end{equation}
%solve problem (\ref{optp-L}) 
%, and $\tau(\Lambda_0^{}, b_{0, t}, G_t)$, respectively.
Substituting  \eqref{Csoln} into  the CIC (\ref{con_IC0}), we obtain
\begin{align}\label{eqPhij}
	& \sum_{i=j+1}^{\infty} e^{-\rho (i-j)\Delta} \big[ C(\Phi_j; b_{j-1, i}, G_{i})U_{C, i}+(C(\Phi_j; b_{j-1, i}, G_{i})+G_i)U_{N, i}-b_{j-1, i} U_{C, i} \big]\Delta \nonumber\\
	&\quad \quad + \big[ C_{j, j}^*U_{C, j}+(C_{j, j}^*+G_j)U_{N, j}-b_{j-1, j} U_{C, j} \big]\Delta \geq  0\,.
	%&&  +B_{0}U_{C}( C(\Lambda_0^{}, b_{0, 0}, G_{0}) ,C(\Lambda_0^{}, b_{0, 0}, G_{0})+G_{0}) .
\end{align}
%\footnote{Contradiction arises when the Lagrangian multiplier $\Phi_0=0$. That is, when $\Phi_0=0$, $C(0; b_{-1, i}, G_i)=\argmax_{C}\, U(C, C+G_i)$, which implies $U_{C}+U_{N}=0$. Substituting into the implementability condition (\ref{impletconstraint}), one obtains $\sum_{i=0}^{\infty} e^{-\rho i\Delta} \big[-G_i-b_{-1, i}\big]U_{C, i}\Delta\geq 0  $, which is a contradiction.    } 
%\begin{eqnarray}\label{eqPhi}
%	&& \int_{0}^{\infty} e^{-\rho t} ( C(\Lambda_0^{*}, b_{0, t}, G_t)U_C+(C(\Lambda_0^{*}, b_{0, t}, G_t)+G_t)U_N )  dt  =  \int_{0}^{\infty} e^{-\rho t}b_{0, t} U_C dt\,.
%	%&&  +B_{0}U_{C}( C(\Lambda_0^{}, b_{0, 0}, G_{0}) ,C(\Lambda_0^{}, b_{0, 0}, G_{0})+G_{0}) .
%\end{eqnarray}
%where we impose the smooth consumption condition
%\begin{equation}
%	b_{0-, t} \text{is cadlag sequence.}
%\end{equation}
%\NW{Wei: maybe this is left continuous with right limit? cadlag sequence? } 
%\WJ{Need to justify this condition.  }
%Substituting  $C(\Phi_0; b_{-1, i}, G_{i})$ into the right side of   (\ref{L0}) lets us  write  $\mathcal L_0$ as a function %of $\Lambda_0^*$:
%$\mathcal{L}_0(\Phi_0)$. 
%{An optimal  $\Phi_0^{*}$ that maximizes $\mathcal{L}_0(\Phi_0)$ must satisfy (\ref{eqPhi}). }
%The Ramsey allocation can be represented with two functions, namely,  $C_{0, i}^*=C(\Phi_0^{*}; b_{-1, i}, G_{i})$ and $N_{0, i}^* = C_{0, i}^*+G_{i} = N(\Phi_0^{*}; b_{-1, i}, G_{i})$. 
%\WJ{As we will show below, under regularity conditions,} the existence of a Ramsey plan requires   % Lagrange multiplier 
%$\Phi_0^*>0$, which implies (\ref{eqPhi}) holds with equality. 
%\NW{	$C(\Lambda_0, b_{0, t}, G_t)  =  \argmax_{C_{0, t}}  \left[ U(C_{0, t}, C_{0, t}+G_t)+ \Lambda_0^{}\left( C_{0, t}U_C+(C_{0, t}+G_t)U_N-b_{0, t} U_C \right) \right]$.} 

%In summary, we conclude with

%\NW{make this prettier and maybe add an eqn as you noted.} 


Similarly, we obtain the continuation Ramsey plan:

%_{\{\Lambda_0 =  \textrm{a positive root of}  (\ref{eqPhi})\}}
\begin{proposition} \label{proplagj} 
	%When %Assuming  , 
	%	Under plausible regularity conditions,\footnote{We provide two regularity conditions to guarantee that  a Ramsey plan  exists. First,  for a given initial debt structure, the admissible allocation set is  not empty, so  there exists allocations $\{C_{0, i}\Delta\}_{i=0}^\infty$ that satisfy the implementability condition (\ref{impletconstraint}). Second, %for the given utility function $U(C, C+G)$, 
		%		$\left(CU_C + (C+G) U_N - bU_C \right)$ is concave in $C$ when $G>0, b\geq 0$. 
		%	For  economies analyzed  in \citet{debortoli2021optimal} these two regularity conditions are satisfied when initial debt is not too high.} 
	%	When a Ramsey plan exists, 
	{Under  regularity conditions \ref{condition1} and \ref{condition2j} described in Appendix \ref{appendix:A}} {and timing protocol given in Definition \ref{defprotocol}}, 
	(\ref{eqPhij}) holds with equality and there  exists  a positive root  $\Phi_j^*$ of  (\ref{eqPhij}) that solves
	\begin{equation*}
		\Phi_j^*  = \argmin_{\Phi_j}  \,\,  \mathcal{L}_j(\Phi_j)\, , %(\{C_{0,t}^*\}_{t=0}^\infty,  
	\end{equation*}	
	where the continuation  Ramsey allocation and tax rate  sequence for $i\geq j+1$ satisfy
	\begin{eqnarray} \label{RamseyCj}
		{C}^\ast_{j, i} =     C(\Phi_j^{*}; b_{j-1, i}, G_{i}); \, \quad  
		{N}^*_{j, i}  =    {C}_{j, i}^\ast +G_{i} ; \,\quad
		{\tau}^*_{j, i} =  1+\frac{U_{N, i}({C}^*_{j, i}, {N}^*_{j, i})}{U_{C, i}({C}^*_{j, i}, {N}^*_{j, i})}; \, .
	\end{eqnarray}
	%	        \begin{eqnarray*}
		%	        	 C(\Phi_0^{\mathcal L}, b_{0-, 0}, G_0)= \lim _{s\downarrow 0} C(\Phi_0^{\mathcal L}, b_{0-, s}, G_s)
		%	        \end{eqnarray*}	
	%	where 
	%and the marginal utilities $U_C$ and $U_N$ are evaluated at $\widetilde{C}_{0, t},\widetilde{N}_{0, t}$.
	The value ${\mathcal L}^{*}_j $ of a continuation Ramsey plan is
	%	\begin{equation*}
		$	 \sum_{i=j}^{\infty}  e^{-\rho (i-j)\Delta} U(	{C}^*_{j, i}, {C}^*_{j, i}+G_{i})  \Delta.
		$	%\end{equation*}
	\end{proposition}
	
	
	\subsection{Lucas and Stokey's Timing Protocol}
	%To overcome a difficulty detected by \citet*{debortoli2021optimal}, we have  perturbed \citeauthor{LucasStokey1983}'s timing protocol. 
	{In \citeauthor{LucasStokey1983}'s (1983) timing protocol, a  period $j$ planner  sets  {$\tau_{j, j}$} as well as 
		%  a flat-rate tax \TS{$\tau_{i, i+1}$} that the  period $i+1$ continuation planner \underline{\TS{can}} reset and 
		a continuation tax plan  {$   \{\tau_{j, i}\}_{i=j+ 1}^{\infty} $} that a  period $j+1$ continuation planner {can} reset.  %The only difference between \citeauthor{LucasStokey1983}'s timing 
		%protocol and ours is that in ours the  period $i$ planner sets a  a flat-rate tax \TS{$\tau_{i, i+1}$} that the  period $i+1$ continuation planner \TS{{cannot}} reset, while 
		%in theirs, the period $i+1$ continuation planner {\TS{can}} reset it. 
		For instance, \citet{LucasStokey1983}  constructed examples of government expenditure processes and restructured debt  processes $\{ b_{0, i}\Delta \}_{i=1}^\infty$ that induce period $1$ continuation planner to adhere to continuations of  a Ramsey tax plan $\{\tau_{0, i}^\ast\}_{i= 1}^\infty$ and its associated allocation and price system $\{C_{0, i}^\ast\Delta, N_{0, i}^\ast\Delta, q_{0, i}^\ast \Delta\}_  {i= 1}^\infty$. 
		\citet*{debortoli2021optimal} constructed examples  in which   continuation term debt   structures  $\{ b_{0, i}\Delta \}_{i=1}^\infty$ like those   recommended by     \citet{LucasStokey1983} fail to  induce  continuation Ramsey planners to  continue  a Ramsey  plan.} In order to implement  an optimal  plan, we shall modify  \citeauthor{LucasStokey1983}'s
	timing protocol and construct a different  debt-management strategy. %\TS{Tom -- this paragraph still needs completion and rewriting.} \WJ{Tom: How about we defer this paragraph and introduce our strategy first and return to discuss how ours differs from theirs?} 
	
	
	%\NW{Next, we introduce  our  implementation.} % based on active management of one-period debt.} % without touching debts with periods to  maturity (weakly) larger than two.} 
	
	%\section{A New Implementation}
	
	%We first present a new timing protocol, and then state our implementation strategy.
	
	
	
	
	%\medskip
	
	
	
	\section{Implementation} %based on A}
\label{newimple} 

%\NW{Tom: we moved the subsection 4.1 title forward to the beginning of this section.} 
\subsection{Debt Management} \label{sec:activedebt}  %that relies on active short-term debt}

% for period $i \geq 1$ 
{A sequence of primary surplus $S_{0, i}^*\Delta$   for period $i$  given in \eqref{S0i} is an important  component of our implementation of  a Ramsey plan.  
Another key component is a scalar   $\Pi_j^*$ that  
% strategy,  then verify how it implements the Ramsey plan under our timing protocol. 
%show
% (described in subsubsection \ref{sec:activedebt}.)%: In , we propose a feasible 
%debt management strategy  Second, we introduce a covenant condition that   
%Together, they make implement the Ramsey plan. 
tracks   %to first define 
unpaid government liabilities %  $\Pi_j^*$
accumulated under  a Ramsey plan  from period $0$ to $j-1$, conditional on no  initial term debts  before period $j-1$, $\{ b_{-1,i}\}_{i\leq j-1}$, having been rescheduled. It is defined by the following equation:
\begin{equation} \label{Pistar}
\Pi_j^* = \sum_{i=0}^{j-1} \frac{q_{0, i}^*}{q_{0, j}^*} \left(b_{-1, i}-S_{0, i}^*\right)\Delta \quad \text{with} \quad \Pi_0^*= 0\,.
\end{equation} % unpaid government liabilities 
Note that  $ \left(b_{-1, i}-S_{0, i}^*\right)$ describes the part of the  initial term debt in period $i$, $b_{-1, i}$, that is not paid by the contemporaneous period's primary surplus $S_{0, i}^*$.} Evaluating it at  bond price ratio $ \frac{q_{0, i}^*}{q_{0, j}^*}$, we obtain its period-$j$ value. {The stock $\Pi_j^*$ sums   period-$j$  values of these unpaid liabilities from periods $0$ through $j-1$. 
The time $j-1$ continuation planner must  finance $\Pi_j^*$. %these liabilities.
Let $r_{0, j}$ denote the equilibrium interest rate in period $j$. Along the Ramsey plan, $r_{0, j}^*=\frac{1}{\Delta}\log(q_{0, j}^*/q_{0, j+1}^*)$ and $\Pi_j^*$ follows a recursion: 
	\begin{equation}\label{dynamicPi}
		\Pi_{j+1}^* = \left(\Pi_j^* +(b_{-1, j}-S_{0, j}^*)\Delta \right)e^{r_{0, j}^*\Delta}\,.
	\end{equation}
}

%which accumulates the government's  unpaid liabilities from period $0$ to $j-1$, conditional on the government not having rescheduled the initial term debt process before period $j-1$. 
%= \tau_{0, j}^*N_{0, j}^* - G_i^*$\,.


%\WJ{Tom: $\Pi_j^*$ is accumulative liability up to period $j-1$ valued at period $j$. As a result, $j-1$ government has to issue some debt to finance this $\Pi_j^*$. }

We  describe  a debt management policy that implements
the Ramsey plan under our timing protocol. First,  
this policy  leaves  term debts with periods to maturity weakly larger than two unchanged, so that  %As a result, the period $i-1$ government chooses 
\begin{equation}\label{btermuntouched} 
b_{j-1, i}\Delta =b_{-1, i}\Delta, \quad j\geq 1, i\geq j+1\,.
\end{equation}
%The strategy sets terms debts with period to maturity equal to $1$ 
%
% accumulated  from period $0$ to period $i-1$   calculated  in the following way. 
%A  period $j \geq 0 $ government   pays all  debts due in period $j$, so under  the strategy (\ref{btermuntouched}) 
%the accumulated liability, $\Pi_i^*$ is  %the following expression for  %is given by  
%
%
%
%%Let  denote the period $i$ value of We then have Second, because the initial term debt due gradually, the government faces an accumulated liabilities from past. 
Second, to finance the sum of (i) the accumulated liability $\Pi_j^*$ given in (\ref{Pistar}) and (ii) the initial term debt $b_{-1, j}\Delta$ due in period $j$, the debt management strategy tells the   period $j-1$ government to  issue just enough  one-period (short-term) debt that will fall due  in period $j$. Consequently, one-period debt issuance satisfies
\begin{equation} \label{bPi}
b_{j-1, j}\Delta = \Pi_j^*+b_{-1, j}\Delta\,, \quad j\geq 1\, . 
\end{equation}
Equation \eqref{bPi}  evidently ``takes the short route.''
%As we discuss in detail later, this debt management strategy is very different from \citeauthor{LucasStokey1983}'s.
%Note that this three-step procedure is feasible as it satisfies the government budget constraints in all periods. An induction 
%We then state the debt management strategy in our implementation satisfies the government dynamic budget constraint under the Ramsey plan.
%\begin{lemma}
%	At any period $i-1$ for $i\geq 1$, the following continuation debt structure satisfies the period $i-1$ government's dynamic budget constraint under the Ramsey plan:
%	\begin{eqnarray} 
%		b_{i-1, i}\Delta &=& \Pi_{i}^*+b_{-1, i}\Delta; \label{finance_d1} \\ 
%		\, b_{i-1, j}\Delta &=& b_{-1, j}\Delta, \quad j\geq i+1\,. \label{finance_d2}
%	\end{eqnarray}
%\end{lemma}

%Next, we introduce a covenant condition that is necessary for a government to issue short-term debt. 





%state a theorem
%verify our debt management strategy with short-term debt only with a covenant condition \eqref{covenant} 
%Next, we prove that  the  implementation proposed above  works. %as long as a Ramsey plan  exists. 
\subsection{Implementation Proposition} \label{sec:im_pro}

%\NW{Mention L method and show recursion....} 
%that bound to honor short-term debt $b_{i-1, i}\Delta$ with a covenant condition and term debt sequence $\{b_{-1, j}\Delta; j\geq i\}$,
%Recall that
We use forward induction starting with the first continuation Ramsey planner. 
Assume that  the period $j-1$ planner confirmed the Ramsey tax plan {by setting $\{\tau_{j-1, i}^{}=\tau_{0, i}^*\}_{i= j}^\infty$}. Consider the  period $j\geq 1$ continuation problem.
The period $j$ government inherits longer term debt $\{b_{j-1, i}\Delta=b_{-1, i}\Delta\}_{i= j+1}^\infty$ {and one-period debt $b_{j-1, j}\Delta$}. %, of which it  must pay $b_{j-1, j}\Delta$ in period $j$. % = \Pi_i^*+b_{-1, i}\Delta$,  the period $i$ continuation government  
Because the period $j$ continuation planner cannot reset the period $j$ tax rate $\tau_{j-1, j}$,   the period $j$'s government's tax rate equals %$\tau_{i,i} = 
$\tau_{j-1, j} =\tau_{0, j}^*$ and consequently $ C_{j, {j}}=C_{j-1, {j}}=C_{0, j}^*$.  The period $j$ continuation planner is free to  choose a new continuation  consumption sequence $\{C_{j, i}\Delta\}_{i=j+1}^\infty$ % whose first component is  $C_{i, i+1}\Delta$ 
that need not equal  $\{C_{0, i}^\ast \Delta\}_{i=j+1}^\infty$.  %to maximize the household's utility.  
Substituting  the debt management policy described in equations \eqref{btermuntouched} and \eqref{bPi}  into  % continuation implementability constraint 
{CIC} (\ref{con_IC0}), we obtain: 
%\begin{equation}\label{con_IC}
%	\sum_{j=i}^{\infty} e^{-\rho(j-i) \Delta} \left(  C_{i, j}U_{C, j} + (C_{i, j}+G_{j})U_{N, j} -b_{-1, j} U_{C, j} \right) \Delta- \Pi_{{i}}^*U_{C, i}\geq 0\,.
%\end{equation}
\begin{align}\label{con_IC}
& \sum_{i=j+1}^{\infty} e^{-\rho(i-j) \Delta} \left(  C_{j, i}U_{C, i} + (C_{j, i}+G_{i})U_{N, i} -b_{-1, i} U_{C, i} \right) \Delta \nonumber\\
& \quad  + \left (C_{0, j}^*U_{C, j} + (C_{0, j}^*+G_{j})U_{N, j} - b_{-1,j} U_{C, j} \right)\Delta - \Pi_{{j}}^*U_{C, j}\geq 0\, .
\end{align} 
Attach a  nonnegative  Lagrange multiplier $\Phi_j $ to (\ref{con_IC}), take the period $j$ continuation government's objective  function (\ref{eqn:preferences}),
and  form  a continuation Lagrangian $\mathcal L_j$ : %is then given by
% \begin{equation}\label{Lagran_t}
% 	\sum_{j=i}^{\infty} \, e^{-\rho(j-i) \Delta} \left[ U(C_{i, j}, N_{i, j}) + \Phi_i  \left (C_{i, j}U_{C, j} + (C_{i, j}+G_{j})U_{N, j} -b_{-1, j} U_{C, j} \right) \right]\Delta - \Phi_i \Pi_{{i}}^*U_{C, i} ,
% \end{equation}
%Facing $\{b_{-1, j}\Delta; j\geq i\}$ and $ \Pi_i^*$, %b_{i-1, i}\Delta = \Pi_i^*+b_{-1, i}\Delta$, 
%the period $i$ continuation government  chooses a sequence $\{C_{i, j}\Delta\}_{j=i+1}^\infty$ to maximize the household's utility, subject to honoring its debt obligation and the covenant condition (\ref{covenant}). 
%The continuation Lagrangian is then given by
\begin{align}%\label{Lagran_t0}
%	& \sum_{j=i}^{\infty} \, e^{-\rho(j-i) \Delta} \left[ U(C_{i, j}, N_{i, j}) + \Phi_i  \left (C_{i, j}U_{C, j} + (C_{i, j}+G_{j})U_{N, j} -b_{i-1, j} U_{C, j} \right) \right]\Delta  \notag \\= 
{\mathcal L}_j\; & =\;
\sum_{i=j+1}^{\infty} \, e^{-\rho(i-j) \Delta} \left[ U(C_{j, i}, N_{j, i}) + \Phi_j  \left (C_{j, i}U_{C, i} + (C_{j, i}+G_{i})U_{N, i} -b_{-1, i} U_{C, i} \right) \right]\Delta \nonumber\\
& \quad\quad\,  + \left[ U(C_{0, j}^*, N_{0, j}^*) + \Phi_j  \left (C_{0, j}^*U_{C, j} + (C_{0, j}^*+G_{j})U_{N, j} - b_{-1,j} U_{C, j} \right) \right]\Delta -\Phi_j \Pi_{{j}}^*U_{C, j} \, .\label{Lagran_t}
\end{align}
%\NW{Under our  timing protocol, for a given $\Phi_j\geq 0$, the period $j$ Ramsey planner takes   its predecessor's  decisions for period $j$ as given and  makes her decisions for all future periods  $i\geq j+1$, as described by  the right side of (\ref{Lagran_t}).} 
The continuation planner maximizes \eqref{Lagran_t}
with respect to $\{C_{j, i}\}_{i=j+1}^\infty$ and minimizes it with respect to $\Phi_j$.
For $i\geq j+1$, the continuation planner sets $C_{j,i}$ to solve %   maximizer of  (\ref{Lagran_t}) is  %continuation consumption plan for $j\geq i+1$ satisfies
\begin{equation}\label{consum_t}
%C_{i, j} = \argmax
\max_{C_{j,i}} \left\{ U(C_{j, i}, C_{j, i} + G_i)  +\Phi_j \left (C_{j, i}U_{C, i} + (C_{j, i}+G_{i})U_{N, i} -b_{-1, i}U_{C, i}\right)  \right\} \, , 
\end{equation} %as for the Ramsey planner, 
so the maximizer takes the form
% %we  can rewrite (\ref{consum_t}) as 
$C_{j, i} = C(\Phi_j; b_{-1, i}, G_{i}) $ for $i\geq j+1$, where $C(\cdot; \cdot, \cdot )$ is given in \eqref{Csoln}. 
%%When the {CIC} (\ref{con_IC}) binds, 
%%Call the extremizing values: $\{C_{0, i}^*\}_{i=0}^\infty$ and $\Phi_0^*$.
%% Maximize under the integral sign on the right side of \eqref{L0} to get a continuum of static problems: 
%%% We characterize the optimal policy for time $0+$ onward as follows.  For a given time interval $(t, t+dt), t>0$, the planner chooses consumption to solve
%%\begin{equation} \label{optp-L}
%%	\max_{C_{0, t}} \;\; \left[ U(C_{0, t}, C_{0, t}+G_t)+ \Lambda_0^{}\left( C_{0, t}U_C+(C_{0, t}+G_t)U_N-b_{0, t} U_C \right) \right]dt\,, \;\; t \geq 0
%%\end{equation} %If there are no constraints on $C_{0, t}$, then 
%for a given $\Phi_i$, the optimal consumption plan $C(\Phi_i; \cdot, \cdot )$  satisfies%\footnote{We assume that the optimal consumption plan for the problem defined in   \eqref{Csoln}   is unique for a  given $\Phi_0$ under certain regularity condition.} 
%\begin{equation}\label{CsolnPhii}
%	C(\Phi_i; b_{-1, i}, G_{i}):=\argmax_{C_{j, i}}  \left[ U(C_{j, i}, C_{j, i}+G_{i})+ \Phi_j\left( C_{j, i}U_{C, i}+(C_{0, i}+G_{i})U_{N, i}-b_{-1, i} U_{C, i} \right) \right] .
%\end{equation}


%

\medskip

%\TS{Neng and Wei: I edited the following proposition}


\begin{proposition} \label{proplag2} 
Under  regularity conditions \ref{condition1} and \ref{condition3} described in Appendix \ref{appendix:A}, %Under regularity conditions provided in Proposition \ref{proplag}, 
(\ref{con_IC}) holds with equality and   $I_j(\Phi_j^*)=0$ 
where
\begin{align}
I_j(\Phi_j) 
%	& = & \sum_{j=i+1}^{\infty} e^{-\rho(j-i) \Delta} \left(  C_{i, j}U_{C, j} + (C_{i, j}+G_{j})U_{N, j} -b_{-1, j} U_{C, j} \right) \Delta \label{IC01}\\
%	& & \quad \quad + \left[\left(  C_{i, i}U_{C, i} + (C_{i, i}+G_{j})U_{N, i}  \right) \Delta- \left(b_{-1, i}\Delta +\Pi_{{i}}^* \right)U_{C, i}\right] \label{IC02}\\
=&   \sum_{i=j+1}^{\infty} e^{-\rho(i-j) \Delta} \left(  C(\Phi_j; b_{-1, i}, G_{i})U_{C, i} + (C(\Phi_j; b_{-1, i}, G_{i})+G_{i})U_{N, i} -b_{-1, i} U_{C, i} \right) \Delta \nonumber \\
& \quad\quad + \left (C_{0, j}^*U_{C, j} + (C_{0, j}^*+G_{j})U_{N, j} - b_{-1,j} U_{C, j} \right)\Delta -\Pi_{{j}}^*U_{C, j} \,.\label{con_IC2}
\end{align}
%Where $\Phi_j^* >0 $ solves $I_j(\Phi_j^*) =0$, % (\ref{con_IC2}) at which  
The continuation Ramsey allocation for period $j$ satisfies 
\begin{eqnarray} \label{ContinueC}
&	 {C}^\ast_{j, i} =     C(\Phi_j^{*}; b_{-1, i}, G_{i}); \, \, 
{N}^*_{j, i}  =    {C}_{j, i}^\ast +G_{i} ; \, i\geq j+1 
%	&	  {\tau}^*_{0, i} =  1+\frac{U_{N, i}({C}^*_{0, i}, {N}^*_{0, i})}{U_{C, i}({C}^*_{0, i}, {N}^*_{0, i})};\,   \quad
%	{q}^*_{0, i}  = e^{-\rho i\Delta}\frac{U_{C, i}({C}^*_{0, i}, {N}^*_{0, i})}{U_{C, 0}({C}^*_{0, 0}, {N}^*_{0, 0})}\, .
\end{eqnarray}
%	        \begin{eqnarray*}
%	        	 C(\Phi_0^{\mathcal L}, b_{0-, 0}, G_0)= \lim _{s\downarrow 0} C(\Phi_0^{\mathcal L}, b_{0-, s}, G_s)
%	        \end{eqnarray*}	
%	where 
%and the marginal utilities $U_C$ and $U_N$ are evaluated at $\widetilde{C}_{0, t},\widetilde{N}_{0, t}$.
%	The value ${\mathcal L}^{*}_0 $ of a Ramsey plan is
%	\begin{equation*}
%	$	 \sum_{i=0}^{\infty}  e^{-\rho i\Delta} U(	{C}^*_{0, i}, {C}^*_{0, i}+G_{
%		i})  \Delta.
%	$	%\end{equation*}

\end{proposition}
See Appendix \ref{appendix:A} for a proof.

\begin{lemma} \label{lemma1}
The Lagrange multiplier $\Phi_0^\ast$ associated with the Ramsey plan satisfies  $I_j(\Phi_0^\ast)=0$. 
%Ramsey $\Phi_0^\ast$ is a root of $I_i(\,\cdot\,)=0$. %If $I_i(\,\cdot\,)=0 $ has unique root, then the Ramsey plan is implemented. 
\end{lemma} 
See Appendix \ref{appendix:A} for a proof.

%\NW{TO continue.} 

%\begin{lemma} 
%If the root of (\ref{}) is unique, Ramsey $\Phi_0^\ast$ is a root of XX. 
%\end{lemma} 

%\NW{Comment the next three paragraphs after editing Lemma and proof in Appendix.} 
%First note that  if $\Phi_i$ equals  the  Lagrangian multiplier for the Ramsey plan, $\Phi_0^*,$ in that $I_i(\Phi_0^*)=0$,  then $C_{i, j}=C(\Phi_i; b_{-1, j}, G_{j})=C_{0, j}^*$ and the Ramsey plan is implemented. %As the  Lagrange multiplier for the Ramsey planner, $\Phi_0^*,$ is a root of $I(\Phi)=0 $, 
%But there may exist multiple roots for $I_i(\,\cdot\,)=0$. %This is why it is difficult to  implementing a Ramsey plan is that 
%
%Therefore, if $I_i(\,\cdot\,)=0 $ has a unique root,   the optimal continuation Lagrange multiplier in period $i$,  $\Phi_i^*$, must equal $\Phi_0^*$. This  means that the period  $i$ continuation Ramsey planner indeed chooses to confirm the continuation of  the  Ramsey plan. 
%




%	\begin{figure}[h!]
%	\caption{{\bf Cash flows from $t_{i-1}$ to $t_i$.} Black terms refer to cash flows when time $t_{i-1}$ government issues an amount of short-term debt $\hat b_{i-1, i}\Delta$ that pays off at time $t_{i}$. Blue terms are cash flows for consumption of households. \label{fig3}}
%	\bigskip
%	\centering
%	\begin{tikzpicture}[scale=1]
%		
%		\draw [thick,->] (-2.5,0) -- (9.7,0);
%		%\draw [thick,-] (-0.19,0.1) -- (-0.19,0);
%		
%		%\node[align=center] at (4.1,1.3) {\bf  \ZL{ \scriptsize Step 2}};
%		%		\node[align=center] at (2.3, 1) {\bf  \NW{ $C_{0, i-1}^*$}};
%		%		\node[align=center] at (5.4, 1) {\bf  \NW{$C_{i, i}$}};
%		
%		\draw [thick,->] (1,-0.5) -- (1,-1.4); 
%		\draw [thick,->] (1, 0.1) -- (1, 1.2);
%		\draw [thick,->] (4.5,0.1) -- (4.5,1.2 );
%		%	\draw [thick,->] (6,0.1) -- (6,0.5);
%		
%		%	\node[align=center] at (1, 1.4) {\bf  \NW{\footnotesize $\widehat{b}_{i-1, i}\Delta =\Pi_{{i}}^* + \WJ{b_{-1, i}\Delta}$}};
%		%	\node[align=center] at (-0.1, 1.9) {\bf  \WJ{\footnotesize ${b}_{-1, i+1}\Delta$}};
%		\node[align=center] at (0.9, -1.8) {\bf  {\footnotesize $\hat{b}_{i-1, i}\Delta e^{-r_{0, {i-1}}^*\Delta}$}}; 
%		\node[align=center] at (0.9, 2) {\bf  \NW{\footnotesize $U_{C, i-1}(C_{0, i-1}^*, N_{0, i-1}^*)$}};
%		
%		%	\node[align=center] at (5.83, 0.8) {\bf  \NW{\footnotesize $ = \Pi_{i+1}^* + \WJ{{b}_{-1, i+3}\Delta}$}};
%		%	\node[align=center] at (6.9, 1.4) {\bf  \NW{\footnotesize $\widehat{b}_{i, i+1} \Delta = \left(\widehat{b}_{i-1, i}\Delta  -S_{i-1, i}\Delta\right)e^{-r_{0, i}^*\Delta} + \WJ{{b}_{-1, i+1}\Delta}$}};
%		%	\node[align=center] at (3.8, 1.9) {\bf  \WJ{\footnotesize ${b}_{-1, i+2}\Delta$}};
%		\node[align=center] at (4.4, 1.6) {\bf  {\footnotesize $\hat{b}_{i-1, i }\Delta$}};
%		\node[align=center] at (4.4, 2) {\bf  \NW{\footnotesize $U_{C, i}(C_{i, i }, N_{i, i})$}};		
%	%	\node[align=center] at (4.9, 1.4) {\bf  \NW{\footnotesize $U_{i, i }\Delta$}};		
%		
%		%		\node[align=center] at (-1.2, 1.9) {\bf \WJ{\scriptsize $\{ b_{-1, j}\Delta\}_{j=0}^{\infty},\; \{G_{j}\Delta\}_{j=0}^{\infty}$  }};
%		%		\node[align=center] at (-1.2, 1.3) {\bf \TS{\scriptsize Ramsey planner  }};
%		%		\node[align=center] at (-1.2, 0.9) {\bf \TS{\scriptsize chooses $\{\tau_{0, j}\Delta\}_{j=0}^{\infty},\; \{\hat{b}_{0, j}\Delta\}_{j=1}^{\infty}$  }};
%		%\node[align=center] at (-1, 1.3) {\bf  \NW{ \scriptsize  Step 3}};
%		
%		%\draw [thick,-] (5.20,0.23) -- (5.20,0);
%		%	\draw [thick,decorate,decoration={brace,amplitude=12pt,raise=4pt}] (1.05,0) -- (3.9,0);
%
%		%   \draw [thick,decorate,decoration={brace,amplitude=12pt,raise=4pt}] (4.1,0) -- (6.9,0);		
%		
%		%\node[align=center] at (8.9, 1.3) {\bf   \YC{\scriptsize  Step 1}};
%		%	\node[align=center] at (8.9, 0.9) {\bf   \TS{\scriptsize Agent pays $M_{\tau}$  }};
%		%	\draw [thick,-] (9.4,0.23) -- (9.4,0);
%		%\draw [thick,decorate,decoration={brace,amplitude=12pt,raise=4pt}] (5.25,0) -- (9.35,0);
%		
%		
%		%\draw [thick,decorate,decoration={brace,amplitude=12pt,raise=4pt}] (9.45,0) -- (10.75,0);
%		
%		\draw[fill,color=blue] (-2.5,0.005) circle (.05);	
%		\draw[fill,color=blue] (1,0.005) circle (.05);	
%		\draw[fill,color=blue] (4.5, 0.005) circle (.05);	
%	%	\draw[fill,color=blue] (7, 0.005) circle (.05);	
%		%\draw[fill,color=black] (-0.2,0.005) circle (.05);
%		
%		%\node[align=center] at (5.2,-0.3) {\footnotesize $ \tau^d$};
%		%\node[align=center] at (4.2,-0.3) {$\tau_F$};
%		\node[align=center] at (-2.5,-0.3) {\footnotesize $0 $};
%		\node[align=center] at (1,-0.3) {\footnotesize $t_{i-1} $};
%		\node[align=center] at (4.5,-0.3) {\footnotesize $t_{i} $};
%	%	\node[align=center] at (7,-0.3) {\footnotesize $t_{i+1} $};
%		\node[align=center] at (10.1,0) {\;\footnotesize  \; time};
%		\node[align=center] at (9.7,-0.3) {\footnotesize $\infty $};
%		
%	\end{tikzpicture}
%\end{figure}


%For the utility function $U(C, N)=\log C-\eta \frac{N^\gamma}{\gamma}$ with $\eta> 0, \gamma\geq 1$ used in \cite{debortoli2021optimal}, we can show that the associated function $I_i(\,\cdot\,)=0$ has only one root. As a result,  our strategy described in subsection \ref{newimple} implements the Ramsey plan. 


%. Problems (\ref{}) for the Ramsey planner, the two problems are


\medskip


{Problem   (\ref{consum_t})  that confronts a period $j$ continuation planner who is  choosing $C_{j,i} $ for $i \geq j+1$ is  static, as  is problem  (\ref{Csoln}) that  had confronted the Ramsey planner who chose $C_{0,i}$ for $i \geq 0$.   %the Ramsey planner's %for all $\{G_i, b_{-1,j}\}$ 
Consequently, $C_{j,i} = C_{0, i} $ for $i \geq j+1$ if and only if $\Phi_j$ in   (\ref{consum_t}) equals $\Phi_0^\ast$ in (\ref{Csoln}).} 

%\TS{Wei -- please check how I wrote the above in color.}

%\TS{Tom paused here, mid afternoon March 23.}
%

\setcounter{theorem}{0}
\begin{theorem}\label{thm1}
Suppose that %Consider a Lucas-Stokey setting defined in Section \ref{lssetting} where 
\begin{equation} \label{DNYutil}
	U(C, N)=\log C-\eta \frac{N^\gamma}{\gamma} \quad \text{with} \quad \eta>0  \quad \text{and} \quad \gamma\geq 1\, . 
\end{equation}
A Ramsey plan exists	under  conditions \ref{condition1} and \ref{condition3}  described in Appendix \ref{appendix:A}.  Under the timing protocol described in Definition \ref{defprotocol},
the  plan can be implemented by using the debt management strategy provided in \eqref{btermuntouched} and \eqref{bPi}. {Lagrange multipliers of continuation   planners equal the Lagrange multiplier of the  Ramsey planner:  $\Phi_j^\ast = \Phi_0^\ast$ for $j\geq 1$}.
\end{theorem} 


\begin{proof}
We use forward induction. %First, 
We suppose that the Ramsey plan has been implemented up to period ${j-1}$ and then %(the initial problem  automatically satisfies.) 
verify that it can also be implemented in period $j$. 
We begin by verifying that (\ref{DNYutil}) satisfies  Condition \ref{condition1}. %required in Lemma \ref{proplag2}. 
Together with Condition \ref{condition3}, Proposition \ref{proplag2} implies that the continuation planner's optimal Lagrange multiplier  $\Phi_j^*>0$ is positive and  that the CIC (\ref{con_IC}) must hold with equality so that  $I_j(\Phi_j^*)=0$. 
%the consumption $\{C_{i, j}\}_{j\geq i+1}$ are functions of $\Phi_i$: 
Substituting  (\ref{DNYutil}) into \eqref{consum_t}, we obtain $C_{j, i}= C(\Phi_j; b_{-1, i}, G_{i})$ for $i\geq j+1$, where $C(\Phi; b, G)$ satisfies the following first-order necessary condition for the  problem defined on the right side of equation  \eqref{consum_t}:
\begin{equation} \label{FF} 
	F(C, \Phi_j):=\frac{1}{C} - \eta(C+G)^{\gamma-1} +\Phi_j \left(\frac{b}{C^2}- \eta\gamma (C+G)^{\gamma-1}  \right) =0\,.
\end{equation}
%Because of the resource constraint, labor supply $\{N_{i, j}\}_{j\geq i+1}$ are given by  
%Recall that under our implementation,
The period $j$ continuation planer confronts   the tail of the initial term debt service sequence  $\{b_{-1, i}\Delta; i\geq j+1\}$ as well as a  short-term debt balance of $ b_{j-1, j}\Delta = \Pi_j^* + b_{-1,j} \Delta $. %First note that the covenant condition requires $C_{i, i}=C_{0, i}^*$. In addition, 
Using $N(\Phi_j; b_{-1, i}, G_{i})= C(\Phi_j; b_{-1, i}, G_{i}) + G_{i}$ to  
simplify the CIC (\ref{con_IC2}), we obtain: 
%Then the continuation implementability condition $I_i(\Phi_i)=0$ follows
\begin{equation}\label{IC_p3}
	I_j(\Phi_j) =\left( 1-\eta(C_{j, j}+G_{j})^\gamma-\frac{b_{-1, j}}{C_{j, j}} \right)\Delta -\frac{\Pi_{j}^*}{C_{j, j}}+  \sum_{i=j+1}^\infty e^{-\rho(i-j)\Delta}\left( 1-\eta(C_{j, i}+G_{i})^\gamma-\frac{b_{-1, i}}{C_{j, i}} \right)\Delta \, , 
\end{equation}
{where $C_{j, j}=C_{j-1, j}=C_{0, j}^*$ under our timing protocol.} 
%\TS{Neng and Wei -- I want to discuss the above couple of sentences. Haven't really used  $C_{i, i}=C_{i-1, i}=C_{0, i}^*$ and $\ldots$}



%We then prove the term on 
To show that $I_j(\Phi_j)  =0$  has a  unique positive root where $I_j(\Phi_j)  =0$ is defined by  CIC (\ref{IC_p3}), it is necessary and sufficient to show that $\sum_{i=j+1}^\infty e^{-\rho(i-j)\Delta} \left( 1-\eta(C_{j, i}+G_{i})^\gamma-\frac{b_{-1, i}}{C_{j, i}} \right)$, the last term in \eqref{IC_p3}, 
is strictly monotone in $\Phi_j$. Denote $P(\Phi_j; b, G) = 1-\eta\left(C(\Phi_j; b, G)+G\right)^\gamma-b/C(\Phi_j; b, G)$. It is then sufficient to prove that $P(\Phi_j; b, G)$ is a strictly monotone function of $\Phi_j$. Recognizing that $C$ is a function of $\Phi_j$ and differentiating $F(C, \Phi_j)=0$ given in (\ref{FF}), we obtain  % with respect to $\Phi_i$
\begin{equation}\label{Fdiff}
	\frac{\partial F}{\partial C}\frac{\partial C}{\partial \Phi_j} + \frac{\partial F}{\partial \Phi_j} =0\,,
\end{equation}
where 
$
\frac{\partial F}{\partial C} = -{1}/{C^2} -\eta(\gamma-1)(C+G)^{\gamma-2}-\Phi_j\eta\gamma(\gamma-1) (C+G)^{\gamma-2}-{2\Phi_j b}/{C^3}< 0\,$ and %\label{FpartialC}\\
$\frac{\partial F}{\partial \Phi_j}  =  \left(b/C^2- \eta\gamma (C+G)^{\gamma-1}  \right)\,.$ % \label{FpartialPhi}
%\end{eqnarray}
Differentiate $P(\Phi_j; b, G)$ and use   \eqref{Fdiff} and $\frac{\partial F}{\partial C}<0$ to obtain
\begin{eqnarray}
	\frac{\partial P}{\partial \Phi_j}  =  \bigg[  b/C^2- \eta\gamma (C+G)^{\gamma-1}   \bigg] \frac{\partial C}{\partial \Phi_j} = - \frac{\partial F}{\partial C}\left( \frac{\partial C}{\partial \Phi_j} \right)^2 >0\, . 
\end{eqnarray}
We conclude  that $P(\Phi_j; b, G)$ and therefore $\sum_{i=j+1}^\infty P(\Phi_j; b_{-1, i}, G_{i})$  are both strictly monotone in $\Phi_j$. %As a result,  also monotone in $\Phi_i$.  
Since $I_j(\Phi_j)$  is strictly monotone in $\Phi_j$, $I_j(\Phi_j)=0$ has a unique root. Applying  Lemma \ref{lemma1} lets us conclude that  the Ramsey plan is implemented. 

%Note when we set $\Phi_i=\Phi_0^*$, $C_{i, j}=C_{0, j}^*$ and the continuation implementability condition \eqref{con_IC} holds equality. Therefore, $\Phi_i=\Phi_0^*$ is the only root for $I_i(\,\cdot\,)=0$, and the continuation Ramsey planner optimally choose $\Phi_i^* = \Phi_0^*$ to confirm the Ramsey plan.
\end{proof}

%\NW{Edit.} 



\subsection{Mechanics }\label{sec:mechanics}

%\NW{Shall we change this section to  subsection 4.3?} 

%\NW{Possibly make this a subsection -- not a big issue -- to be discussed later.} 

%\subsection{Reverse Engineering Our Implementation}
%In this section, we %inspect the mechanism that underpins our implementation. 
%What's the intuition behind our implementation? An reverse engineering approach helps us to design the method.  . That is, the period $j$ continuation Ramsey planner  (for any $j\geq 1$) 
Our implementation  induces a  period $j\geq 1$ 
 planner to % continue  the Ramsey plan by 
set continuation  allocation  $\{ C_{j, i}^*\}_{i\geq j}$ to the tail of the Ramsey allocation $\{C_{0, i}^*\}_{i\geq j}$,  so 
%Because that $\{ C_{j, i}^*=C_{0, i}^*\}_{i\geq j}$, 
continuation tax rates, asset prices, and debt management plans align to support the Ramsey tax  plan and associated outcomes.  
%(Recall that the Ramsey allocation is $\{C_{0, i}^*\}_{i\geq 0}$  and )% the tail of Ramsey plan, 
To understand mechanics  underlying these outcomes, first recall from Propositions \ref{proplag} and \ref{proplag2} %under the Ramsey plan and the period $j$'s continuation Ramsey plan,
that for {$i\geq j+1$}
\begin{eqnarray}
C_{0, i}^*  =  C(\Phi_0^*; b_{-1, i}, G_i)   \quad \text{and} \quad  
%	\label{sec4C0i}\\
C_{j, i}^* =  C(\Phi_j^*; b_{j-1, i}, G_i)  . \label{sec4Cji}
\end{eqnarray}
% (\ref{sec4C0i}) and
In  (\ref{sec4Cji}), different {\em arguments}    in the same {\em function} determine   $C_{0, i}^* $ and $C_{j, i}^*$. A good way to understand our implementation is to  find   arguments of  \eqref{sec4Cji} that guarantee that 
$C_{0, i}^*=C_{j, i}^* $ for all $j \geq 1$.  % it is sufficient to make the arguments be the same. % for the  two functions in  (\ref{sec4Cji}) the same. %if feasible.  %Below we proceed in two steps.  (\ref{sec4C0i})  and 
To  do  that, we reverse engineer %second step. Comparing the two functional forms  is to 
($1$) continuation  debt-management policies that satisfy     $b_{j-1, i}=b_{-1, i}$  for  $i\geq j+1$, and ($2$) assure that %determinants underlying $\Phi_0^*$ and $\{\Phi_j^*\}_{j=1}^\infty$ that will 
$\Phi_j^*=\Phi_0^*$ for all $j \geq 1$.
% Doing that  leaves us with figuring out how      equate $C_{j, j}^*$ with $C_{0, j}^*;$  our timing protocol for setting continuation tax plans accomplishes that.  %, as we discuss below.  

If we impose reverse engineering specification ($1$) and also
% requires that  all debts with times to maturity that weakly exceed two  inherited by the period $j$ planner, $b_{j-1, i}$, are the same the initial debt structure $b_{-1, i}$. %A simple way to achieve this debt management outcome is for all preceding continuation Ramsey planners leave all debts with times to maturity (weakly) exceeding two untouched. %But then 
%Without touching its debts with  times to maturity (weakly) larger than two, how does the period $j-1$ planner balance its budget? 
assume  that the Ramsey plan has been followed up to period $j-1$, {at the beginning of  period $j-1$ the government has accumulated} 
$\Pi_j^*$ given by (\ref{Pistar}), evaluated at time $j$ prices. 
%\WJ{Tom: the previous sentence may mislead the reader that decision maker is period $j$ government. Actually, it's not. It is period $j-1$ planner's decision to management debt to induce period $j$ planner to follow Ramsey plan.} 
The period $j-1$ planner must finance both $\Pi_j^*$ and the initial term  debt due  $b_{-1, j}\Delta$.  To do that, the period $j-1$ planner    issues  one-period debt equal to $b_{j-1, j}\Delta = \Pi_j^*+b_{-1, j}\Delta$. This is  the short-term-debt-only policy given in (\ref{bPi}).
% under the assumption that the Ramsey plan has been followed up to period $j-1$. 
% To verify that the Ramsey plan has been followed up to period $j-1$, it is sufficient that reverse engineering condition ($2$) hold.
%\TS{Wei and Neng: I suspect that Imisstated some $j$ and $j-1$'s in the preceding sentences.  Please repair!}


To establish that  $\Phi_j^*=\Phi_0^*$, recall that period $j$'s continuation planner's Lagrange multiplier $\Phi_j^*$ satisfies 
\begin{align}
&\sum_{i=j+1}^{\infty} e^{-\rho(i-j) \Delta} \left(  C(\Phi_j^*; b_{-1, i}, G_i)U_{C, i} + (C(\Phi_j^*; b_{-1, i}, G_i)+G_{i})U_{N, i} -b_{-1, i} U_{C, i} \right) \Delta \nonumber \\
& \quad\quad + \left (C_{j, j}^*U_{C, j} + (C_{j, j}^*+G_{j})U_{N, j} - b_{-1,j} U_{C, j} \right)\Delta -\Pi_{{j}}^*U_{C, j} =0 \, \tag{CIC} \label{sec4CIC}
\end{align}
and that the Ramsey planner's Lagrange multiplier $\Phi_0^*$ satisfies 
\begin{align}
&\sum_{i=j+1}^{\infty} e^{-i \rho \Delta} \left(  C(\Phi_0^*; b_{-1, i}, G_i)U_{C, i} + (C(\Phi_0^*; b_{-1, i}, G_i)+G_{i})U_{N, i} -b_{-1, i} U_{C, i} \right) \Delta \nonumber\\
& \quad \quad +  \left (C_{0, j}^*U_{C, j} + (C_{0, j}^*+G_{j})U_{N, j} - b_{-1,j} U_{C, j} \right)\Delta -\Pi_{{j}}^*U_{C, j} =0
\,,\tag{IC} \label{sec4IC}	
\end{align}
Setting $C_{j, j}^*=C_{0, j}^*$ makes the second term in   \eqref{sec4CIC} for period $j$'s continuation planner  identical to the second term in \eqref{sec4IC} for the Ramsey planner. 
%\NW{Edit.} 
Furthermore, if    $\Phi_j^*=\Phi_0^*$,
the  first term in   \eqref{sec4CIC}  for period $j$'s continuation planner  equals the first term in \eqref{sec4IC} for the Ramsey planner. Thus,  both terms agree, so  
both equations hold.  
%
%\NW{Edit.} 
%
%Notice that   if $C_{j, j}^*=C_{0, j}^*$ holds, then the second term in the CIC for period $j$'s continuation planner, \eqref{sec4CIC}, is identical to the second term in the  IC for the Ramsey planner \eqref{sec4IC}. As a result, 
%%\NW{Edit.} 
%equation \eqref{sec4CIC} of $\Phi_j^*$ is the same as equation \eqref{sec4IC} of $\Phi_j^*$, suggesting  that   $\Phi_j^*=\Phi_0^*$ may hold. 
%
%
% the first term in the CIC for period $j$'s continuation planner,  ,   to the first term in the  IC for the Ramsey planner

%\NW{Edit.} 
%perfectly line up with each other,. 
%where $\Pi_j^*$ is that under the Ramsey allocation  $\{C_{0, i}^*\}_{i\geq 0}$. 
%To further ensure the second the third terms of \eqref{sec4CIC} perfectly line up with those of \eqref{sec4IC}, one needs to let $C_{j, j}^*=C_{0, j}^*$, which effectively prevent the continuation planner from diluting the value of one-period deb. 

%and show that under this  protocol
%The last step in our reverse engineering project is to
We can now verify that   $\Phi_j^*=\Phi_0^*$ and $C_{j, j}^*=C_{0, j}^*$ for all $j \geq 1$ under our timing protocol for setting taxes, which requires that 
%On the one hand, using the households' FOCs, we can show that the new timing protocol induces t
the period $j$ continuation planner must administer tax rate $\tau_{j-1, j}$ for period $j$ which is set by the period $j-1$ planner. This induces continuation households in period $j$ to  choose $C_{j, j}^*=C_{j-1, j}^*=C_{0, j}^*$, which follows from the household's first-order necessary condition $1-\tau_{j-1, j}=-U_{N, j}(C_{j, j}, C_{j, j}+G_j)/U_{C, j}(C_{j, j}, C_{j, j}+G_j)$.\footnote{Note that for $j=1$, our timing protocol implies $C_{1, 1}^*=C_{0, 1}^*$, which ensures the forward induction.}  Thus, taken together, our  timing protocol and debt management strategy
imply $\Phi_j^*=\Phi_0^*$.\footnote{This  follows from Lemma \ref{lemma1}, which states that  Ramsey planner's Lagrange multiplier $\Phi_0^*$ is a root of  \eqref{sec4CIC}, and from Theorem \ref{thm1}, which states that \eqref{sec4CIC} has a unique root for the utility function given in \eqref{DNYutil}.} 
We have thus reverse engineered a ``take the short route''  debt-management policy that implements the Ramsey plan.


%\section{Applying our Implementation to \citeauthor{debortoli2021optimal}'s Setting} % Implementation} %Additional  Examples} 
%Ramsey Plan
\section{%\citeauthor{debortoli2021optimal}'s 
\citeauthor{debortoli2021optimal}'s Examples \label{sec:DNYexamples}} % Implementation} %Additional  Examples} 

\subsection{Analysis}
%In this section, we show how our Implementation works in \citeauthor{debortoli2021optimal}'s  counterexamples. 
%Next, we  apply our implementation to 
\citeauthor{debortoli2021optimal} constructed  examples  in which  \citeauthor{LucasStokey1983}'s debt management policy does not implement the Ramsey plan under their timing protocol.
These  examples feature an initial debt structure that takes the form 
\begin{equation}\label{DNYdebt}
	b_{-1, 0}=b>0   \;\;  \text{and} \;\; b_{-1, i}=0\;\; \text{for all} \;\; i \geq 1 \,  
\end{equation} 
a  high value of $b$, 
 a constant government spending process $G_{i}=G$ for all $i\geq 0$,  and the utility function  (\ref{DNYutil}). %A key feature is that the initial debt is high enough. 
After first reviewing why Lucas and Stokey's  debt management policy  fails to implement the Ramsey plan in  \citeauthor{debortoli2021optimal}'s examples under their timing protocol,
 %why   \cite{LucasStokey1983}'s implementation does not work 
we explain how  our  debt management policy succeeds in  implementing the Ramsey plan under our timing protocol.
%  our implementation works in  \citeauthor{debortoli2021optimal}'s counterexamples. %by making a minor modification to the timing protocol in   \cite{LucasStokey1983}. 

		\begin{figure}[h!] 
	\centering
	\includegraphics[width=0.6\linewidth]{figure_202403/laffer_curve3}
	\caption{{\bf Laffer curve.} Parameter values: $\gamma=1$, $\eta=1,$  and  $G=0.2 $  }
\label{laffercurve} \end{figure}


To begin, let $\{C_{0, 0}^*=C_0^*, C_{0, i}^*=C_1^*; i\geq 1\}$ denote the Ramsey plan and let $C^{\text{Laffer}}$ denote  consumption at the peak of the  Laffer curve (see Figure \ref{laffercurve}.) 
\citeauthor{debortoli2021optimal}
show (i) that  there exists $b^* \in (0, \bar b)$ such that the Ramsey planner sets  $C_1^*> C^{\text{Laffer}}$ if $b < b^\ast$ and  $C_1^*< C^{\text{Laffer}}$ if $b > b^\ast$ (their Proposition 1);\footnote{ \citeauthor{debortoli2021optimal} show that a Ramsey plan exists if and only if $b<	\bar{b}$,  
where
\begin{equation*}\label{bbar}
	\bar{b}  \; \; = \; \; \max_{\widetilde{C}} \quad \widetilde{C}\left[ [1-\eta(\widetilde{C}+G)^\gamma] + \frac{e^{-\rho\Delta}}{1-e^{-\rho\Delta}}\left(1-\eta G^\gamma \right)   \right]\, . 
\end{equation*}} and (ii)  that
if $b < b^\ast$, then Lucas and Stokey's debt management policy  implements the Ramsey plan, but that  if $b > b^\ast$, then Lucas and Stokey's debt management policy does
not implement the Ramsey plan (their Proposition 2.)  
 \citeauthor{debortoli2021optimal} highlight a coincidence that occurs  between $C_1^*$ being tied to a tax rate that is  above the peak of the Laffer curve and the failure of Lucas and Stokey's debt management 
 policy to implement the Ramsey plan.  

	

{Does  that   coincidence persist  for  other initial debt structures under Lucas and Stokey's timing protocol?  Subsection \ref{subsec6.1} shows that the answer is no. Does that   coincidence persist  under  our timing protocol and debt management policy?
 The following proposition shows that the answer to this question is also no.} 
% }} \WJ{Neng: The above sentences are not clear.}% 
%  is that  the Ramsey plan cannot be implemented when  the Ramsey planner sets its tax rate  above the peak of a Laffer curve. In contrast, under our timing protocol, the short-term debt based management  strategy described in Theorem \ref{thm1} (and also Proposition \ref{corollary1} below) always implements  the Ramsey plan, making which side of  the  Laffer curve the Ramsey allocation is on irrelevant for implementability. 
%which side of the Laffer curve the Ramsey planner sets its tax rate does not matter for the implementability of  

%Our implementation   sustains a Ramsey plan even    a situation in which 


%For simplicity in this section we focus on the \citet*{debortoli2021optimal}'s 
%\subsection{Applying our Implementation to \citeauthor{debortoli2021optimal}'s Setting} % Implementation} %Additional  Examples} 

%First, we summarize \citeauthor{debortoli2021optimal}'s  main results. 

  % economy.

%The second part of our argument 

%Debortoli et al. (2021) further establish that the Ramsey plan cannot implemented when a Ramsey tax rate is above the peak of a Laffer curve.

%
% \NW{Edit -- 
%define  \eqref{eqPhi} can hold for a feasible allocation:  
%A value  $\bar{b}$ is sustained with $C_{0, i}=0$ for all $i\geq 1$ and $C_{0, 0} = \tilde C$ set to maximize  the right side of equation  \eqref{bbar}.
%}
% 
 
%\cite{debortoli2021optimal} further show that  the Ramsey plan is implemented if and only if  the Ramsey  plan sets the tax rate above the peak of the Laffer curve (see their Proposition 1 and our Appendix \ref{appendix:B} for details of \cite{debortoli2021optimal}'s Laffer curve argument.) 
%
%In contrast, whether %our implementation overcomes this Laffer curve issue.(see 
%

%Our debt management strategy adjusts only one-period debt, so  %and leaves all term debt with maturity longer than one period untouched, i.e.,
%\subsection{Our Implementation for \citeauthor{debortoli2021optimal}'s  Setting} %Parametric Example and  counterexample.} \citeauthor{debortoli2021optimal} 
%\subsection{Irrelevance of Laffer Curve in Our Implementation} %Parametric Example and  counterexample.} \citeauthor{debortoli2021optimal}   \citeauthor{debortoli2021optimal}'s  

\begin{proposition}\label{corollary1}
	%In an economy that households follow utility function $U(C, N)=\log C-\eta \frac{N^\gamma}{\gamma}$ with $\gamma\geq 1$ and government confronts an initial debt $b_{-1, 0}=b>0 $ and $b_{0, i}=0, \forall i\geq 1$ and constant spendings $G_{i}=G, \forall i\geq 0$, 
	 For  all $b\in (0, \bar b)$, %where $\bar b$ is given in (\ref{bbar}), 
	 the following debt management policy  implements the  Ramsey plan under our  Definition \ref{defprotocol} timing protocol: % in period $1$:  the Ramsey plan is unique and 
	\begin{equation}\label{restructure_debt_dny}
	{b}_{0, 1}\Delta= \Pi_{{1}}^* \, \;\; \text{and} \;\;  b_{0, i}\Delta = b_{-1, i}\Delta=0 ; \; i\geq 2\,,
\end{equation}
where $\Pi_{1}^* = \frac{e^{\rho \Delta}C_1^*}{C_0^*}\left(b-C_0^*\left( 1-\eta(C_0^*+G)^\gamma\right) \right)\Delta$. {The Lagrange multiplier for the continuation   planner's problem equals the Lagrange multiplier for the  Ramsey problem:  $\Phi_1^\ast = \Phi_0^\ast$}.
% \eqref{restructure_debt_dny}. 
\end{proposition} 

%which we further discuss   in detail. %We return to discuss this point later in Section XX. 

%Moreover, we show which side of the Laffer curve the Ramsey planner sets its tax rate does not determine whether the Ramsey plan is implementable. \TS{tom edit following} 

%
%Second, we show that %, for utility function $U(C, N)=\log C-\eta \frac{N^\gamma}{\gamma}$ with $\eta>0$ and $\gamma\geq 1$, 
%$\Phi_0^*>0$ is the unique positive root for the CIC (\ref{}), which implies  $\Phi_i^*=\Phi_0^*$ (see Lemma \ref{}). %equation $I_i(\Phi_i)=0$. 
%%Given $\Phi_i$, according to 
%

%\begin{proof}
%we only need to verify that under utility $U(C, N)=\log C-\eta \frac{N^\gamma}{\gamma}$, the function $I(\Phi)$ is monotone in $\Phi$. Note that $C(\Phi; b, G)$ solve the following equation
%\[ \frac{1}{C} - \eta(C+G)^{\gamma-1} -\Phi \left( \eta\gamma (C+G)^{\gamma-1} + b/C^2 \right) =0\,. \]
%Let $K(\Phi; b, G) = 1-\eta\left(C(\Phi; b, G)+G\right)-b/C(\Phi; b, G)$. We show that $\partial K/\partial \Phi>0$ using the implicit differentiation.
%\end{proof}

%To illustrate, for the case $U(C, N)=\log C-\eta \frac{N^\gamma}{\gamma}$ with $\gamma=1$, the function $I(\cdot)$ is given by
%\begin{equation}
%I(\Phi) = \sum_{j=i}^\infty \, e^{-t_j\Delta} \left( 1-\eta(C(\Phi; b_{0, t_j}, G_{t_j})+G_{t_j})-\frac{b_{0, t_j}}{C(\Phi; b, G_{t_j})}  \right)
%\end{equation}
%where $C(\Phi; b, G) = \frac{2b\Phi}{\sqrt{1+4\eta b\Phi(1+\Phi)}-1} $. Because $\left( 1-\eta(C(\Phi; b, G_{t_j})+G_{t_j})-\frac{b_{0, t_j}}{C(\Phi; b, G_{t_j})}  \right)$ is an increasing function of $\Phi$, $I(\Phi)$ is a monotone function of $\Phi$ such that $I(\Phi)=0$ has only one root. As a result, we can show that our implementation strategy implements the Ramsey plan for a DNY economy.
% 



%Why our implementation overcomes 

%Irrelevance of \cite{debortoli2021optimal}'s 


%\subsection{Irrelevance of the Laffer curve in our implementation} 
%\paragraph{Why our implementation overcomes Laffer curve issue.} 
%Next, we illustrate the irrelevance of \cite{debortoli2021optimal}'s Laffer curve argument in our implementation. %e use \cite{debortoli2021optimal}'s example given in Section \ref{newimple}.  the period 1 government cannot choose 

%{This proposition can be directly derived by applying Theorem \ref{thm1}. To highlight the irrelevance of the `™wrong' side of the Laffer curve for implementability, we provide an alternative proof below.}

\noindent{We can establish this proposition directly  by applying Theorem \ref{thm1}, but to highlight the irrelevance for implementability of being on  the {bad} side of the Laffer curve, %irrelevance of the  
 we instead provide the following alternative proof.
} 

%\noindent{\bf Proof:}
\begin{proof}
Under our timing protocol, the Ramsey planner sets $C_{1, 1}=C_{1}^*$. 
Confronting  debt structure  \eqref{restructure_debt_dny},  
the period 1 continuation planner chooses  $\{ C_{1, i}, i\geq 2\}$ to maximize %the household's utility from period 1 onward: %(\ref{contplannerDNYexample}) 
%\begin{equation}\label{contplannerDNYexample} 
$$	\sum_{i=1}^{\infty} e^{-\rho (i-1)\Delta} \left( \log C_{1, i}- \eta \frac{N_{1, i}^\gamma}{\gamma} \right)\Delta \, $$ 
%\end{equation}
%$\sum_{i=1}^{\infty} e^{-\rho (i-1)\Delta} \left( \log C_{1, i}- N_{1, i} \right)\Delta$,
subject to the CIC: %continuation implementability condition: 
\begin{equation}
	\left(1-\eta(C_{1, 1}+G)^\gamma-\frac{b_{0, 1}}{C_{1, 1}}\right)\Delta+\sum_{i=2}^\infty e^{-\rho (i-1)\Delta}\left( 1-\eta(C_{1, i}+G)^\gamma-\frac{b_{0, i}}{C_{1, i}} \right)\Delta =0 \, . \label{IC_dny2}
\end{equation}
%Unlike in Lucas-Stokey's implementation, both the first term in the objective  (\ref{contplannerDNYexample}) and the first term in the CIC (\ref{IC_dny2}) are beyond the control of period 1 continuation planner.  
The associated Lagrangian is 
extremized by setting  $C_{1, i}=C_{1, 2}$ for all $i\geq 3$. Consequently, the %continuation implementability condition 
CIC \eqref{IC_dny2} and the debt management strategy \eqref{restructure_debt_dny} %can be written as a monotone function of $C_{1, 2}$:
jointly imply a quadratic equation for $C_{1, 2}$: 
%	\begin{equation} \label{Squad} 
	%	\frac{1}{C_{1, 1}}\left(S(C_{1, 1}) -\hat b\right) =  \frac{1}{C_{1, 1}} \left( -C_{1, 1}^2-(G+1)C_{1, 1} -\hat b \right)=0 \,. 
	%\end{equation}
	%\begin{align} \label{Squad2}
	%\frac{e^{-\rho\Delta}}{1-e^{-\rho\Delta}}\frac{1}{C_{1, 2}}\bigg(S(C_{1, 2})-b_{-1, 2}\bigg)  & =   \frac{\Pi_1^*/\Delta}{C_1^*}-\left(1-(C_1^*+G)-\frac{b_{-1, 1}}{C_1^*}\right)  \nonumber\\
	%& =  e^{\rho \Delta } \frac{1}{C_0^*}\bigg(b_{-1, 0} -S(C_0^*)\bigg) - \frac{1}{C_1^*}\bigg(S(C_1^*)-b_{-1, 1} \bigg) \,.
	%\end{align} % with the Ramsey planner's
	%\small
		\begin{align} \label{Squad2}
			\frac{1}{1-e^{-\rho\Delta}}\frac{1}{C_{1, 2}}\big(S(C_{1, 2})-b_{-1, 2}\big)%  & =   \frac{\Pi_1^*/\Delta}{C_1^*}-\left(1-(C_1^*+G)-\frac{b_{-1, 1}}{C_1^*}\right)  \nonumber\\
			=- \bigg[ \frac{e^{2\rho\Delta}}{C_{0, 0}^*}\big(S(C_{0, 0}^*) - b_{-1, 0} \big) + \frac{ e^{\rho \Delta } }{C_{0, 1}^*}\big(S(C_{0, 1}^*)-b_{-1, 1} \big) \bigg]\,,
	\end{align} 
	where $S(C) = C\left(1-\eta(C+G)^\gamma \right)$ is the primary surplus rate.   
	%, we are able to keep $b_{-1,2}$ and indeed all 
	Equation \eqref{Ramseytaxprice} for setting taxes confirms  that a quadratic equation $S(C)$ implies  a Laffer curve. 
	%\citet*{debortoli2021optimal}' Laffer curve observation follows from the fact
   
    Now rewrite the Ramsey plan's  IC as  % can be written as
    		\begin{align} \label{Squad2IC}
    	\frac{1}{1-e^{-\rho\Delta}}\frac{1}{C_{0, 2}^*}\big(S(C_{0, 2}^*)-b_{-1, 2}\big)%  & =   \frac{\Pi_1^*/\Delta}{C_1^*}-\left(1-(C_1^*+G)-\frac{b_{-1, 1}}{C_1^*}\right)  \nonumber\\
    	=- \bigg[ \frac{e^{2\rho\Delta}}{C_{0, 0}^*}\big(S(C_{0, 0}^*) - b_{-1, 0} \big) + \frac{ e^{\rho \Delta } }{C_{0, 1}^*}\big(S(C_{0, 1}^*)-b_{-1, 1} \big) \bigg]\, . 
    \end{align} 
Since  the right sides of (\ref{Squad2}) and (\ref{Squad2IC}) are identical,  their left sides must be identical too. 
Consequently, $C_{1, 2}=C_1^*$ is a root of the quadratic equation (\ref{Squad2IC}). It must be the larger root because if it were  not, there would exist another root of \eqref{Squad2}, $\widehat{C}_1$, larger than $C_1^*$. But if that were true, 
		%$C_{1, 2}=\widehat{C}_1>C_1^*$, we then obtain an allocation $\{ C_{0, 0}=C_{0}^*, C_{0, 1}=C_1^*, C_{0, i}=\widehat{C}_1; i\geq 2 \}$ for 
		the Ramsey planner would have chosen $\{ C_{0, 0}=C_{0}^*, C_{0, 1}=C_1^*, C_{0, i}=\widehat{C}_1; i\geq 2 \}$, contradicting the hypothesis that $\{ C_{0, 0}=C_{0}^*, C_{0, 1}=C_1^*, C_{0, i}={C}_1^\ast; i\geq 2 \}$ is the Ramsey plan and thus verifying Proposition \ref{corollary1}.%whether the Ramsey consumption plan is above or below the peak of%is a side show.} %This allocation satisfies the IC and achieves higher utility than the Ramsey allocation, giving rise of the contradiction.} 
%\WJ{The invariance of Lagrange multiplier is naturally obtained by applying Theorem \ref{thm1} here.}
\end{proof}

%Technically, setting $C_{1, 2}=C_1^*$ makes the CIC \eqref{Squad2} identical to the IC (\ref{Squad2IC}) under the Ramsey plan. %As a result, \WJ{$C_{1, 2}=C_1^*$ as one of its too roots of \eqref{Squad2}. %Furthermore, 

%, which 



% so that the associated equilibrium discount factor (marginal utility) , which equals $\Pi_2^\ast/C_{0,2}^\ast$ and is beyond the period 1 planner's control
{The validity of Proposition  \ref{corollary1} does not depend on  whether the Ramsey planner sets the tax rate on the good or the bad side  of  the Laffer curve. 
Even though a high period 1 tax rate and associated low period 1 consumption are associated with large accumulated government  liabilities,  under   our timing protocol the period 1 continuation planner cannot reset   period 1 consumption.} This makes the right side of the CIC \eqref{Squad2} equal to the right side of the IC (\ref{Squad2IC}). By managing only short-term debt and leaving 
	inherited debts  maturities weakly larger than two periods  {unchanged}, the period 1 continuation  planner  continues the Ramsey tax plan,  aligning  the left side of its CIC \eqref{Squad2} with the IC  (\ref{Squad2IC}) under the Ramsey plan. %}%future consumption  % % with the Ramsey planner's
	%In summary, the new timing protocol and debt management strategy together make our implementation work.}


\subsection{{Three} Lagrange Multipliers} 
\label{subLM} 
%In the preceding two subsections, we show 

 Lagrange multipliers of the Ramsey planner and continuation Ramsey planners behave differently under our timing protocol  and debt management policy, on the one hand,  and under    Lucas and Stokey's, on the other hand.    Here we focus on the   $\eta=\gamma=1$ case, in which case 
	the  Lagrange multiplier for the Ramsey plan satisfies %$\Phi_0^*>0$, 
		\begin{equation} \label{PhiRamsey} 
		{\Phi}_0^\ast= \frac{C_1^*-(C_1^*)^2}{(C_1^*)^2-b_{-1, i}}\,.
	\end{equation}
%$\frac{1}{C_{0, i}^*} -1 +\Phi_0^*\left( \frac{b_{-1, i}}{(C_{0, i}^*)^2} -1 \right) =0, i\geq 2\,.$
	
%\citeauthor{debortoli2021optimal}  also use  a continuation planner's Lagrange multiplier to explain why  \cite{LucasStokey1983}'s implementation can fail. % their  key results. %is to  %we turn to {Lucas-Stokey's implementation.}  
%To ease exposition we focus on 
%to summarize their analysis based on the continuation planner's Lagrange multiplier. % in the remainder of this subsection. 
%To implement the Ramsey plan,
%the Ramsey planner pays $b_{-1,0}$ and restructures the remaining debt $\{b_{-1, i}\}_{i\geq 1}$ to $\{b_{0, i}\}_{i\geq 1}$. Thus, 
For  both Lucas and Stokey's setup and our setup, the   Lagrangian for a  period-$1$ continuation planner is %is then given by
\begin{align}\label{L1}
	{\mathcal L}^{}_1 & =  \sum_{i=1}^{\infty}  e^{-\rho (i-1)\Delta}\left[ \log C_{1, i}-  (C_{1, i}+G_{i})+ \Phi_1^{}\left( 1-(C_{1, i}+G_{i})-\frac{b_{0, i}^{}}{C_{1, i}}\right)  \right] \Delta \, , %\nonumber \\%- \Psi_0(b_{0, t}^{} U_C- b_{0-, t} U_C)	\right]dt \\
	%&  \quad \quad \;  \, .
\end{align}
where  $\{b_{0, i}\}_{i\geq 1}$ is the   debt structure that the Ramsey planner passes on to the period 1  continuation planner. % after %pays $b_{-1,0}$. 
The period 1  continuation planner's first-order condition for $C_{1,i}$ is %Ramsey consumption $C_{1, i}$ for $i\geq 1$ satisfies 
\begin{align}\label{FOC53}
	\frac{1}{C_{1, i}} -1 +\Phi_1\left( \frac{b_{0, i}}{C_{1, i}^2} -1 \right) =0.
\end{align}
%the following Lucas-Stokey's debt management strategy 
%}

We first study outcomes under Lucas and Stokey's timing protocol and debt management policy.  
\citeauthor{debortoli2021optimal}  showed that the Ramsey planner  sets 
\begin{equation}\label{b0iLS} 
	b_{0, i} = \frac{(C_{1}^*)^2-C_1^*}{\Phi_1}+(C_{1}^*)^2\,   \quad    \text{for all}\;\; i\geq 1,  
\end{equation} %we obtain the CIC \eqref{IC_dny}, 
an  equation that can be  obtained by substituting $C_{1, i}=C_{0, i}^*=C_1^*$, which has to hold for   the Ramsey plan to be implemented, %$C_{1, i}=C_{0, i}^*=C_1^*$ has to hold for all $i\geq 1$. 
into the first-order necessary condition (\ref{FOC53}). 
Substituting (\ref{b0iLS}) into the CIC \eqref{IC_dny2}, we can  obtain  what  \citeauthor{debortoli2021optimal} call a  `constructed' Lagrange multiplier  associated with  the  period $1$ continuation plan,
namely,  %Lagrange multiplier: % $\Phi_1$, denoted by $\widehat\Phi_1$: 
\begin{equation}  \label{hatPhi1LS} 
	\widehat\Phi_1 =\frac{1}{2} \frac{1-C_1^*}{C_1^*-C^{\text{Laffer}}}\, , 
\end{equation}
where $C^{\text{Laffer}}=(1-G)/2$. %which we refer to  as  $\widehat \Phi_1$. 
{The Lagrange multiplier ${\Phi}_0^\ast$  associated with the Ramsey plan given by \eqref{PhiRamsey} does not equal the ``constructed'' Lagrange multiplier 
	$\widehat\Phi_1$ associated with the period 1 continuation Ramsey plan   given by \eqref{hatPhi1LS}.}
	
%\TS{Note that the optimal Lagrange multiplier $\Phi_1^*$ associated with the continuation Ramsey plan may not equal the constructed Lagrange multiplier 
%	$\widehat\Phi_1$. } 

Equation (\ref{hatPhi1LS}) implies  that  $\widehat\Phi_1<0$ if and only if  $C_1^*<C^{\text{Laffer}}$, which by virtue of \eqref{Ramseytaxprice} means that the Ramsey plan puts $\tau_{0,1}^*$  on the %(conventionally perceived)
 `wrong' side of the Laffer curve. % $C_1^*<C^{\text{Laffer}}$. 
When   $\widehat\Phi_1<0$, it is not optimal for the continuation planner to  continue the Ramsey plan. %Later we show that in richer settings, whether  
%(as $C_1^*<1$ always holds.)
In this way, 
\cite{debortoli2021optimal}   establish that the period 1 continuation planner would never  choose a  tax rate $\tau_{1,1}^*$  above the peak of the Laffer curve.
This verifies that  if the initial debt is too high, then under Lucas and Stokey's timing protocol and debt management strategy, the Ramsey plan is not implemented. 
{When $\widehat{\Phi}_1<0$, the  Lagrange multiplier ${\Phi}_1^\ast$  associated with the continuation Ramsey plan does not equal the constructed Lagrange multiplier 
 $\widehat\Phi_1$ associated with the  period 1 continuation Ramsey plan given by \eqref{hatPhi1LS}, an 
inequality that indicates that the   Ramsey plan will  not be implemented under this  
  timing protocol and debt management policy.} 
%\TS{Tom: Say much more, pick up rents.}

%\NW{Tom: they did not really explicitly show (\ref{}) and make this explicit connection, but they more or less know this.} 
%We show that the policy is not time consistent if the initial debt is high enough. The first part of our argument establishes that if initial debt is high enough, optimal policy under commitment implies that the date 1 government chooses high tax rates above the peak of the Laffer curve.The second part of our argument 



%\WJ{We then turn to analyze the Lagrange multiplier of our implementation.
%In our implementation,
	%The Ramsey planner pays $b_{-1,0}$ and restructures the remaining debt to $\{b_{0, i}\}_{i\geq 1}$ that follows \eqref{restructure_debt_dny}. Thus, the   Lagrangian for  period-$1$ continuation planner is %is then given by
%	\begin{align}\label{L1_o}
%		{\mathcal L}^{}_1 & =  \sum_{i=1}^{\infty}  e^{-\rho (i-1)\Delta}\left[ \log C_{1, i}-  (C_{1, i}+G_{i})+ \Phi_1^{}\left( 1-(C_{1, i}+G_{i})-\frac{b_{0, i}^{}}{C_{1, i}}\right)  \right] \Delta \, . %\nonumber \\%- \Psi_0(b_{0, t}^{} U_C- b_{0-, t} U_C)	\right]dt \\
%		%&  \quad \quad \;  \, .
%	\end{align}
	Turning now to  our  timing protocol given in Definition \ref{defprotocol}, the period 1 continuation Ramsey planner sets  $C_{1, 1}=C_{0,1}^*$ and  does not adjust term debts with times to maturity (weakly) larger than two, so $b_{0,i}=b_{-1, i}$ for all $i\geq 2$. 
%	The FOC for period-1   planner is: %Ramsey consumption $C_{1, i}$ for $i\geq 1$ satisfies 
%	$
%	\frac{1}{C_{1, i}} -1 +\Phi_1\left( \frac{b_{0, i}}{C_{1, i}^2} -1 \right) =0, i\geq 2.
%	$ %the following Lucas-Stokey's debt management strategy 
%	To implement the Ramsey plan, $C_{1, i}=C_{0, i}^*=C_1^*$ has to hold for all $i\geq 2$. 
	Substituting these two conditions into (\ref{FOC53}) %$C_{1, i}=C_1^*$ and $b_{0,i}=b_{-1, i}$ for $i\geq 2$ into the preceding FOC, 
	delivers the following `constructed' Lagrange multiplier for the period 1 planner: %our implementation: 
	\begin{equation} \label{hatPhi1_our}
		\widehat{\Phi}_1 = \frac{C_1^*-(C_1^*)^2}{(C_1^*)^2-b_{-1, i}}\,, \; i\geq 2. 
	\end{equation}
Evidently,   $\widehat{\Phi}_1$ is positive and equals the Ramsey planner's  Lagrange multiplier $\Phi_0^*$ given in \eqref{PhiRamsey}, so $\widehat{\Phi}_1=\Phi_0^*>0$. {Consequently,  the  Lagrange multiplier ${\Phi}_1^\ast$  associated with the continuation Ramsey plan equals the constructed Lagrange multiplier 
	$\widehat\Phi_1$ associated with the  period 1 continuation Ramsey plan given by \eqref{hatPhi1_our}.}

{Taking stock,   under Lucas and Stokey's timing protocol,     when $b$ is too high, {$\widehat{\Phi}_1<0$} and  the Ramsey plan is not implemented. 
However, under our timing protocol,   $\widehat{\Phi}_1=\Phi_0^*>0$ and our ``short route'' debt management policy  implements the Ramsey plan. 
} %\WJ{Tom: We changed the statement about the LM to a sign statement, as this is why LS implementation fails.} 

%	The period 1 continuation planner  sets  $\Phi_1^*=\widehat{\Phi}_1$ because $\widehat{\Phi}_1>0$. 
%	In contrast, in LS implementation, when $C_1^*<C^{\text{Laffer}}$, $\widehat{\Phi}_1<0$ and therefore the period-1 continuation planner would never  choose to confirm the continuation of the Ramsey plan. 
	
	%Next, we use the numerical counterexample  in \cite{debortoli2021optimal} to highlight the differences between the two implementations. 
	%In summary, our implementation implies invariant Lagrange multipliers: $\Phi_1^*=\widehat{\Phi}_1=\Phi_0^*$.



%	\begin{equation}\label{b0iLS} 
%		b_{0, i} = \frac{(C_{1}^*)^2-C_1^*}{\Phi_1}-(C_{1}^*)^2\,  , \quad    \text{for all}\;\; i\geq 1. 
%	\end{equation} %we obtain the CIC \eqref{IC_dny}, 
%	Substituting (\ref{b0iLS}) into the CIC \eqref{con_IC0}, we obtain  the following expression for a  `constructed' Lagrange multiplier  for   period-$1$ continuation planner: %Lagrange multiplier: % $\Phi_1$, denoted by $\widehat\Phi_1$: 
%	\begin{equation}  \label{hatPhi1LS} 
%		\widehat\Phi_1 = -\frac{1-C_1^*}{C_1^*-C^{\text{Laffer}}}\, , 
%	\end{equation}
%	where $C^{\text{Laffer}}=(1-G)/2$. %which we refer to  as  $\widehat \Phi_1$.  

%that can also implement the Ramsey plan. In this subsection, we show that using 

\subsection{Other Debt-Management Policies Can Also Work}\label{sec:otherpolicies}
If a Ramsey plan exists, our  short-term debt management strategy  implements  it. %a Ramsey plan for  all initial government expenditure and initial debt sequences  for which  which an optimal  plan exists. 
 For some initial debt structures,  but not for others,   debt management policies that issue a mixture of  short and some  longer term debts can also implement a Ramsey plan. 
%However, other initial debt structures,   the initial  debt structure in order for  an  implementation that uses  longer-term debt to work. In this sense,  using term debt makes implementation more difficult.} %In other words, provided that a Ramsey plan exists, using short-term debt to implement a Ramsey plan always works, but using a combination of short- and long-term debts does not necessarily work. 
To understand this claim,  suppose that a Ramsey planner  finances a fraction $\alpha\in [0, 1]$ of accumulated  liabilities  up to period $1$, $\Pi_1^*$, by issuing one-period debt %$b_{0, 1}\Delta $ 
\begin{equation} \label{b01_n}
	b_{0, 1}\Delta = \alpha \Pi_1^*+b_{-1, 1}\Delta\, , 
\end{equation}
  and  finances   fraction $(1-\alpha)$ of $\Pi_1^*$ by issuing longer-term debts, $\{b_{0, i}\}_{i\geq 2}$. Substituting \eqref{b01_n} into  CIC \eqref{IC_dny2} gives 
  \begin{equation} \label{b02_n}
	b_{0, i}\Delta = (1-\alpha)(e^{\rho\Delta}-1)\Pi_1^*+b_{-1, i}\Delta\,,\,  i\geq 2\,.
\end{equation} 
Use  the continuation planner's optimality condition \eqref{FOC53} to form the  `constructed' Lagrange multiplier: 
\begin{equation} \label{hatPhi1_n}
	\widehat{\Phi}_1 = \frac{1-C_1^*}{((1-\alpha)e^{\rho\Delta}+1)C_1^*-(1-\alpha)e^{\rho\Delta}(1-G)} \,.
\end{equation}
%Consider two special cases. 
If $\alpha = 1$, then the debt management   strategy corresponds to our  short-term debt only strategy, %and $\widehat{\Phi}_1$ given in \eqref{hatPhi1_n}; 
 since $b_{-1,i} = 0$ for $i \geq 2$,  \eqref{hatPhi1_n} simplifies  to \eqref{hatPhi1_our}. 
If  $\alpha = 1 - e^{-\rho \Delta}$, then the debt management   strategy corresponds to  Lucas and Stokey's consol-based implementation strategy in which $b_{0,1} = b_{0,i}$ for $i \geq 2$ and $\widehat{\Phi}_1$ given in \eqref{hatPhi1_n} simplifies to  \eqref{hatPhi1LS}. 


%Note that when $\alpha=1$, it corresponds to our implementation and $\widehat{\Phi}_1$ in \eqref{hatPhi1_n} boils down to \eqref{hatPhi1_our} (because $b_{-1, i}=0$ for $i\geq 2$.) When $\alpha=1-e^{-\rho\Delta}$, we have $b_{0, 1}=b_{0, i}, i\geq 2$, and the $\widehat{\Phi}_1$ in \eqref{hatPhi1_n} boils down to that in Lucas-Stokey implementation given in \eqref{hatPhi1LS}.


%To ensure $\widehat{\Phi}_1>0$, a condition that is necessary for the implementation method works, we require that
%$((1-\alpha)e^{\rho\Delta}+1)C_1^*-(1-\alpha)e^{\rho\Delta}(1-G)>0$, that is
%\begin{equation*}
%	C_1^*>\frac{(1-\alpha)e^{\rho\Delta}(1-G)}{(1-\alpha)e^{\rho\Delta}+1} \,.
%\end{equation*}
%According to \cite{debortoli2021optimal}, $C_1^*$ is an decreasing equation of $b$. As a result, for each $\alpha \leq 1$, there exist a $b^*(\alpha)$, which is an increasing equation of $\alpha$, such that when $b=b^*(\alpha)$, $C_1^*=\frac{(1-\alpha)e^{\rho\Delta}(1-G)}{(1-\alpha)e^{\rho\Delta}+1}$. Moreover, for given $\alpha$, the implementation method only works for $b<b^*(\alpha)$ where $\widehat{\Phi}_1>0$.



Sometimes using  both short- and long-term debt fails to implement a Ramsey plan.  
A necessary condition for a successful implementation is 
  $\widehat{\Phi}_1 > 0$. This  condition requires that 
$((1 - \alpha)e^{\rho\Delta} + 1)C_1^* - (1 - \alpha)e^{\rho\Delta}(1 - G) > 0$ or equivalently 
\begin{equation} \label{C1sub53} 
	C_1^* > \frac{(1 - \alpha)e^{\rho\Delta}(1 - G)}{(1 - \alpha)e^{\rho\Delta} + 1} \,.
\end{equation}
Note that $C_1^*$ is strictly decreasing in $b$ (see \citeauthor{debortoli2021optimal}) and  that the right side of (\ref{C1sub53}) is increasing in $\alpha$. Consequently, there exists a  threshold $\tilde{b}(\alpha)$ that  is an increasing function of $\alpha$  for all $\alpha \leq 1$ such that 
inequality (\ref{C1sub53}) prevails if $b < \tilde{b}(\alpha)$. %we have $C_1^* = \frac{(1 - \alpha)e^{\rho\Delta}(1 - G)}{(1 - \alpha)e^{\rho\Delta} + 1}$.
Thus, for a given $\alpha$, $\widehat{\Phi}_1 > 0$ and  a mixed short- and long-term debt issuance policy  (\ref{b01_n})-(\ref{b02_n}) implements a Ramsey plan only 
when $b < \tilde{b}(\alpha)$.\footnote{For the two polar special cases discussed in the text, $\tilde{b}(1)=\bar b$ and $\tilde{b}(1-e^{-\rho\Delta})=b^*$.}







%'s
\subsection{Quantitative Example}\label{sec:DNYquantexample}
%\subsection{%Illustration via 
We turn briefly to
\citeauthor{debortoli2021optimal}'s numerical example of a situation in which the Ramsey plan exists but is not implemented. 
  \citet{debortoli2021optimal} set an initial debt  structure: $b_{-1, 0}=0.6$ and $b_{-1, i}=0$ for $i\geq 1$ with $G=0.2$ and $\eta=\gamma=1.$\footnote{See their section 3 and their figure 3.} %and figure XX.%We explain why our implementation works in this counterexample where \citeauthor{LucasStokey1983}'s implementation does not work. 
% planner sets the period $0$ consumption on the right side of Laffer curve,
The Ramsey plan is  $C_{0, 0}^* =0.8172, \tau_{0, 0}^* =0.1828$  and $C_{0, i}^*=0.3125, \tau_{0, 0}^* =0.6875$ for all $i\geq 1$. At the peak of the Laffer curve, $C^{\text{Laffer}}=0.4$ and $\tau^{\text{Laffer}}=0.6$, so the period 0   Ramsey tax rate  is  on the `right' side of the Laffer curve and the period 1 Ramsey  tax rate  is  on the `wrong' side of the Laffer curve. 
The Ramsey plan's  Lagrange multiplier  is $\Phi_0^\ast=2.19$. %lower than $C_{\text{Max}}^{\text{Laffer}}=0.4$. 
% $C_{\text{Max}}^{\text{Laffer}}=0.4$ but
%\TS{Tom: Shall we drop $Max$ subscript as DNY to line up with DNY notation and also to avoid unintended potential Max/Min confusion?} 


%and the consumption from period $1$ onward on the wrong side of Laffe curve, $C_{0, i}^*=0.3125<C_{\text{Max}}^{\text{Laffer}}, i\geq 1$,  where $C_{\text{Max}}^{\text{Laffer}}$ denotes consumption at the peak of the Laffer curve.
%the  tax rate  above the peak of Laffer curve in  period $0$ and  below the peak of Laffer curve from period $1$ onward. 
%As a result, Lucas-Stokey recommends a restructured term debt 
%Using  formulas in this section, we have 
Lucas and Stokey's  Ramsey planner leaves the time 1 continuation planner  with $b_{0, i}=0.1523$  for all $i\geq 1$, which differs from the tail  $b_{-1, i}=0$ for all $i\geq 1$ of the initial debt structure. Consequently, the continuation planner's   `constructed' Lagrange multiplier  $\widehat{\Phi}_1=-3.93 < 0$. Facing debt structure $\{b_{0, i}=0.1523\}_{i=1}^\infty$,   the continuation planner sets  $C_{1, i}=0.4874$  instead of continuing the Ramsey plan $C_{0, i}^\ast=0.3125$  for all $i\geq 1$. {The associated Lagrange multiplier for   the continuation planner is ${\Phi}_1^\ast=2.93>0$, which 
differs  from  both the `constructed' Lagrange multiplier 
$\widehat{\Phi}_1=-3.93$ and the Ramsey planner's Lagrange multiplier $\Phi_0^\ast=2.19$.} 
 %${\Phi}_1^\ast=2.93>0$. 

%C_{\text{Max}}^{\text{Laffer}}$ for $i\geq 1$. 	that is on the wrong side of Laffer curve 

Turning to our timing protocol and debt management policy, the Ramsey planner sets   $b_{0, 1}= 0.3871$ and  $b_{0, i}=b_{-1, i}=0$  for all $i\geq 2$, leaving term debts with time to maturities (weakly) than two unchanged. The continuation planner sets $C_{1, i}=C_{0, i}^*=0.3125$ for $i\geq 1$, confirming the tail of the Ramsey plan. 
{The associated `constructed' Lagrange multiplier for the continuation planner now equals the resulting Lagrange multiplier for   the continuation planner, ${\Phi}_1^\ast=2.19>0$, which equals  the Lagrange multiplier for Ramsey planner: $\widehat{\Phi}_1=\Phi_1^*=\Phi_0^*=2.19$.}  %Importantly,   $\Phi_1^*=\widehat{\Phi}_1=\Phi_0^*=2.98$. 
%% and  Note that the restructured debt plan 
%	
%%By doing this, it effectively lowers the price of period $0$ debt,  and meanwhile, increases the value of surplus for period $1$ onward.  
%Note that the constructed Lagrange multiplier $\widehat{\Phi}_1=-3.93$ is negative. This is why  \citeauthor{LucasStokey1983}'s implementation does not work in this case. 
Thus, as we would expect, \citeauthor{debortoli2021optimal}'s quantitative example confirms  equality of    Lagrange multipliers for the Ramsey plan and the continuation Ramsey plan
as a tell-tale sign of implementability. % the same for the Ramsey and continuation planners. 

%\paragraph{Summary.}

%^Knowing that the Ramsey plan and continuation plan satisfy \eqref{sec4Cji}, we finally obtain $C_{j, i}^*=C_{0, i}^*$ for $i\geq j+1$.
%(i.e., , i.e., $\Phi_j^*=\Phi_0^*$
%Moreover, we verify that combining 
%perfectly line up with each other,. 
%where $\Pi_j^*$ is that under the Ramsey allocation  $\{C_{0, i}^*\}_{i\geq 0}$. 
%To further ensure the second the third terms of \eqref{sec4CIC} perfectly line up with those of \eqref{sec4IC}, one needs to let $C_{j, j}^*=C_{0, j}^*$, which effectively prevent the continuation planner from diluting the value of one-period deb. 
%
%To achieve this goal, we then introduce a new timing protocol that the period $j$ continuation planner must administer tax rate $\tau_{j-1, j}$ for period $j$ which is set by the former planner. On the one hand, using the households' FOCs, we can show that the new timing protocol induces the continuation households choose $C_{j, j}^*=C_{j-1, j}^*=C_{0, j}^*$.\footnote{Note that for $j=1$, our timing protocol implies $C_{1, 1}^*=C_{0, 1}^*$, which ensures the forward induction.} Moreover, we verify that combining the new timing protocol and our debt management strategy, the invariant of Lagrange multiplier can be obtained by the fact that the optimal Ramsey Langrane multiplier $\Phi_0^*$ is the root of the \eqref{sec4CIC} (i.e., Lemma \ref{lemma1}) and \eqref{sec4CIC} has unique root under utility function \eqref{DNYutil} (i.e., Theorem \ref{thm1}), i.e., $\Phi_j^*=\Phi_0^*$. Knowing that the Ramsey plan and continuation plan satisfy \eqref{sec4Cji}, we finally obtain $C_{j, i}^*=C_{0, i}^*$ for $i\geq j+1$.

%by dynamically managing its one-period debt by following (\ref{}), 


%How to implement a Ramsey plan 






%also 



%for any $i\geq j+1$. %and 2.) 


%$j\geq 1$ and 

%For , all $b_{-1, i}, G_i$, 

%Then, according to , we know that Ramsey consumption and continuation Ramsey consumption both are functions of Lagrange multipliers, initial debt, and government spendings, i.e.,



%is implemented up to period $j-1$ for $j\geq 1$. Given continuation Ramsey allocation  $\{C_{j, i}^*\}_{i\geq j}$ for 



%Keeping the objective in mind, if we let $b_{j-1, i}=b_{-1, i}$ for all $i\geq j+1$ and assume $\Phi_j^*=\Phi_0^*$, we obtain $C_{j, i}^*=C_{0, i}^*$ for $i\geq j+1$. 


%Moreover, our debt management strategy that leaving period to maturity weakly larger than two periods untouched effectively line up  the first term of Ramsey planner's IC and that of the continuation planner's CIC, i.e.,  

%\clearpage


	
%	\section*{Comparing our Implementation  with Lucas-Stokey's} % Implementation} %Additional  Examples} 
\section{More Examples\label{sec:moreexamples}} % Settings} % Implementation} %Additional  Examples} 
%
%\NW{In this section, we analyze more examples going beyond \citeauthor{debortoli2021optimal} to show that our short-term debt based management strategy always works and  whether the Ramsey tax plan is above or below the peak of the Laffer curve is not the key driving force for the implementability of the Ramsey plan.} 
%
%\NW{Tom: do we need some type of motivation like the above for this section?}  
This section adds  quantitative examples to those presented in subsection \ref{sec:DNYquantexample}.
Subsection \ref{subsec6.1} studies the consequences of modestly  perturbing \citeauthor{debortoli2021optimal}'s initial debt structure (\ref{DNYdebt}) while retaining their assumption that government expenditures are constant.  Our experiments in this setting show  that  an optimal tax rate on the wrong side of the Laffer curve is  not a tell-tale sign that Lucas and Stokey's debt management policy fails to implement the Ramsey plan, but that  {variations in pertinent Lagrange multipliers on the Ramsey and continuation problems are a tell-tale sign.} %\citeauthor{debortoli2021optimal}'s
Remaining subsections turn to  looking  under the hood of our implementation, in the same  spirit of  the suite of examples provided in \citet[Sec.~3]{LucasStokey1983}.
Subsection \ref{subsecdecliningb}  studies implementable debt and tax sequences that emerge when government expenditures are constant but  initial debt payments $\{b_{-1,i}\}$ decline exponentially. 
Subsection \ref{subsecperiodicb} studies implementable debt and tax sequences that emerge when government expenditures are constant but  initial debt payments $\{b_{-1,i}\}$ fluctuate periodically.  
Subsection \ref{subsecperiodicG} studies implementable debt and tax sequences that emerge when initial debt payments $\{b_{-1,i}\}$   are constant but government expenditures  fluctuate periodically.  {Like example 4 of  \citet[sec.~3]{LucasStokey1983}, this example  presents  an optimal  tax sequence and an  equilibrium interest rate sequence that matches prescriptions and other aspects 
of  \citet{barro1979determination}.} 
%
% \citet{barro1979determination}, where the   %in which  the intertemporal properties of the optimal  tax rate are  disconnected from the intertemporal properties of that exogenous government expenditure process that is to  be financed. 
%the interest rate is an exogenous constant, and example 4 of  \citet[sec.~3]{LucasStokey1983} ,





\subsection{Perturbing DNY's Example} \label{subsec6.1}
%highlight why Lucas-Stokey's Implementation use 
We now parameterize  a class of initial debt structures in the following way:
\begin{equation} \label{richerb} 
b_{-1, 0}=b_0,\; b_{-1, 1}=b_1, \; b_{-1, i}=b_2,\,  i\geq 2  .
\end{equation} 
%and retain all other assumptions  in the preceding section. 
Comparing debt structure  (\ref{richerb})  with  \citeauthor{debortoli2021optimal}'s initial debt structure (\ref{DNYdebt}) that we  studied in  section \ref{sec:DNYexamples}, notice  that now $b_{-1, i}=b_2\neq b_1$ for periods $i\geq 2$. {Let $\{C_{0, 0}^*=C_0^*, C_{0, 1}^*=C_1^*, C_{0, i}^*=C_2^*; i\geq 2\}$ denote the Ramsey plan.}  %This allows us to generate new insights absent in the preceding section. % is a bit richer and the  \ref{}.  %with three segments 
%For the ease of exposition, we consider an economy with constant government spendings $G_{i}=G, \forall i\geq 0$ and period utility function $U(C, N)=\log C-\eta \frac{N^\gamma}{\gamma}$ with $\eta=\gamma= 1$.

%Following Proposition \ref{proplag}, the Ramsey plan is given by $C_{0, i}^* = \frac{2 b_{-1, i}\Phi_0^*}{\sqrt{1+4 b_{-1, i}\Phi_0^*(1+\Phi_0^*)}-1}$ for $i\geq 0$. Given the initial debt structure, we denote $C_{0, 0}^*=C_0^*, C_{0, 1}^*=C_1^*, C_{0, i}^*=C_2^*, i\geq 2$ for simplicity.
Under Lucas and Stokey's timing protocol and debt management policy, a  Ramsey planner  that confronts initial debt structure  (\ref{richerb})   restructures debt according to %management in 
\begin{eqnarray}
b_{0, 1} = \frac{(C_{1}^*)^2-C_{ 1}^*}{\Phi_1}+(C_{1}^*)^2 \quad \text{and} \quad 
b_{0, i} %=   b_{0, 2} 
= \frac{(C_{ 2}^*)^2-C_{ 2}^*}{\Phi_1}+(C_{2}^*)^2,\, i\geq 2 \, .  \label{sect6bPhi} 
\end{eqnarray}
%following the same procedure as in subsection \ref{subLM}, 
We have reverse engineered this restructuring plan from %by working backwards from 
 choices that the period 1  continuation planner must make to continue the Ramsey plan.  We must verify that Lagrangian \eqref{L1} is maximized  when  the consumption sequence $\{C_{1, i}\}_{i\geq 1}$ satisfies  $C_{0,1}^\ast=C_1^\ast$ and $C_{0,i}^\ast=C_2^\ast$ for $i\geq 2$.\footnote{Subsection \ref{subLM} describes an analogous reverse-engineering 
 exercise for  \citeauthor{debortoli2021optimal}'s initial debt structure (\ref{DNYdebt}).} %and the associated FOCs \eqref{FOC53},   assumption that  allocations are 
Substituting (\ref{sect6bPhi}) into the CIC \eqref{IC_dny2}, we obtain the following `constructed' Lagrange multiplier for the period 1 continuation planner: 
\begin{eqnarray}\label{hatphi2}
{\widehat\Phi_1}=\frac{1}{2}\frac{\left(1-C_{ 1}^*\right)+\left(1-C_{ 2}^*\right)\frac{1}{1-e^{-\rho\Delta}}}{\left(C_{ 1}^*-C^{\text{Laffer}}\right)+\left(C_{ 2}^*-C^{\text{Laffer}}\right)\frac{1}{1-e^{-\rho\Delta}}}  ,
\end{eqnarray}
where $C^{\text{Laffer}}=(1-G)/2$. Compare (\ref{hatphi2}) under  \citeauthor{debortoli2021optimal}'s initial debt structure (\ref{richerb}) with (\ref{hatPhi1LS}) under the debt structure \eqref{DNYdebt} and notice  that the sign of the `constructed' Lagrange multiplier ${\widehat\Phi_1}$ shapes  Laffer curves  for all periods $i\geq 1$. {Evidently, $\frac{1}{1-e^{-\rho\Delta}}$ represents a  discounted value of  $\left(C_{ 2}^*-C^{\text{Laffer}}\right)$  from period 2 onward.} 

If $\widehat \Phi_1>0$, the period 1 continuation planner chooses to continue the  Ramsey tax plan and sets $\Phi_1^*=\widehat \Phi_1$. However, if $\widehat \Phi_1<0$, the period 1 continuation planner chooses not to continue the Ramsey plan but instead chooses a new  plan that extremizes its Lagrangian \eqref{L1} and provides a positive  period 1 continuation plan Lagrange multiplier  %$\Phi_1^*$ be  positive: 
$\Phi_1^*>0>\widehat \Phi_1$. %reneging

% Expression \eqref{hatphi2} implies that 
%The constructed Lagrange multiplier  $\widehat \Phi_1$  can be either positive or negative. 
If the Ramsey planner  sets all  tax rates on the `wrong' side of the Laffer curve, so that $C_1^*<C^{\text{Laffer}}$ and $C_2^*<C^{\text{Laffer}}$, then $\widehat\Phi_1 <0$. If  $C_i^*>C^{\text{Laffer}}$ for all $i\geq 1$,  $\widehat\Phi_1 >0$. %In summary, even f simple debt structure, the Laffer curve argument does not work. %This result indicates that  may not necessary be negative and therefore the continuation planner may confirm the Ramsey plan. In contrast to \cite{debortoli2021optimal}, we show a government in the future may follow the Ramsey plan choose to tax above the peak of the Laffer curve for a period.

%either  period 1 or  future periods ($i\geq 2$),   
Depending on  values of the initial debt structure parameters $b_0, b_1, b_2$ in specification \eqref{richerb},  a Ramsey planner may choose to  set some tax rates on the `wrong' side of the Laffer curve and other tax rates  on the `good' side, while  $\widehat\Phi_1$ can  be  either positive or negative. Whether a  Ramsey plan's tax rates are  on the `wrong' side of 
a Laffer curve  does not indicate  the plan can be implemented. %as ${\widehat\Phi_1}$ can be either positive or negative. %or not.} %which in turn determines whether the Ramsey plan is i
{Table \ref{table2}} illustrates this.  Case I sets $b_{-1, 0}=0.3$, $b_{-1, 1}=0,$ and $b_{-1, i}=0.15$ for $i\geq 2$, while 
case II  increases $b_{-1, 0}$ to $0.4$ and keeps all future debts at their case I values. %b_{-1, 1}=0, b_{-1, i}=0.15, i\geq 2$.% to demonstrate why  Lucas-Stokey's implementation sometimes does not work.  
Under  Lucas and  Stokey's timing protocol and debt management policy, the Ramsey plan is implemented  in case I  but not  in case II. With our timing protocol, our debt management
policy implements a Ramsey plan  in both cases. 



\begin{table}[h!]
%\centering
\caption{Lucas-Stokey's and implementation and ours} %and continuation Ramsey plan.}
\caption*{	Parameter values: $\eta=\gamma=1, G=0.2, \rho=0.5, \Delta=1$: {$C^{\text{Laffer}}_{\text{Max}}=0.4$}. \label{table2}
}
\centering
\renewcommand{\arraystretch}{1.4} 
%\resizebox{0.8\textwidth}{!}{
%\begin{tabular}{c|c|c|c|c|c|c}
\begin{tabularx}{\textwidth}{C | C | C | C | C | C | C }
\hline
{	Case} &  \multicolumn{3}{c|}{I} & \multicolumn{3}{c}{ II}    \\ 		
\hline 

Period $i$ &    $0$ & $1$ & $\geq 2$ &  $0$ & $1$ & $\geq 2$  \\
\hline
$b_{-1, i}$ &  	  0.3  & 0 & 0.15  &   0.4  & 0 & 0.15  		 	  \\
\hline 
\multicolumn{7}{c}{\bf A. Ramsey plan}\\ 
\hline
$C_{0, i}^*$ &   0.6297 &  {0.2930}  &  0.5035  &  0.6459 & { 0.0702} &  0.4102  \\
\hline
LM $\Phi_0^*$ &   \multicolumn{3}{c|}{2.41}  &   \multicolumn{3}{c}{13.24}    \\
\hline  
\multicolumn{7}{c}{ \bf B. LS implementation and continuation Ramsey plan}\\ \hline
$b_{0, i}$ &   n.a.   & 0.0710 & 0.2357  &  n.a.   & 0.0272 & 0.2509     \\
\hline 
%Constr' & \multicolumn{3}{c|}{}  &   \multicolumn{3}{c|}{}  &  \multicolumn{3}{c}{} \\
LM $\widehat\Phi_1$  &   \multicolumn{3}{c|}{{13.96}}  &   \multicolumn{3}{c}{{-2.92}}   \\    \hline
$C_{1, i}^*$  &    n.a.  &  {0.2930}  &  0.5035  &  n.a. &  {0.3877} &  0.6116\\
\hline
LM $\Phi_1^*$  &   \multicolumn{3}{c|}{{13.96}}  &   \multicolumn{3}{c}{1.92}   \\  
\hline 	
Works?  &  \multicolumn{3}{c|}{{Yes: $\widehat\Phi_1=\Phi_1^*$}}  &   \multicolumn{3}{c}{{No: $\widehat\Phi_1\neq \Phi_1^*$}}   \\   
\hline  

\multicolumn{7}{c}{ \bf C. Our implementation}\\ \hline
{$\Pi_i^*$} &  0  & {0.1478} & {-0.0018}   & 0& {0.0538} & {0.0251}  \\\hline
$b_{0, i}$ &   n.a.   & {0.1478} & 0.15  &  n.a.   & {0.0538} & 0.15    \\ \hline
LMs &   \multicolumn{3}{c|}{{ ${\widehat \Phi}_1 =   \Phi_1^\ast = \Phi_0^\ast  = 2.41$}}  &   \multicolumn{3}{c}{{${\widehat \Phi}_1 =   \Phi_1^\ast = \Phi_0^\ast  =13.24$}}    \\  \hline 
%		$C_{1, i}^*$ &   0.6297 &  {0.2930}  &  0.5035  &  0.6459 & { 0.0702} &  0.4102  \\		\hline
%   LM $\Phi_1^*$  &    \multicolumn{3}{c|}{{2.41}}  &   \multicolumn{3}{c}{{13.24}}   \\   \hline				
Works?&  \multicolumn{3}{c|}{{Yes}}  &   \multicolumn{3}{c}{{{Yes}}}    \\  \hline    
%\end{tabular}
\end{tabularx}
%	}
\end{table}

%Recall that  \citeauthor{LucasStokey1983}'s timing protocol and debt management strategy (see definitions in Appendix \ref{appendix:B}) 

% 	\item  Using Case I, we reproduce \citet{debortoli2021optimal}'s counterexample that shows why   \citeauthor{LucasStokey1983}'s implementation may not work. 
%
%		Note that the constructed Lagrange multiplier $\widehat{\Phi}_1=-3.93$ is negative. This is why  \citeauthor{LucasStokey1983}'s implementation does not work in this case. 
%			
%Facing initial debt $b_{-1, 0}=0.6$ and no other debt, the Ramsey planner sets the period $0$ consumption on the right side of Laffer curve, $C_{0, 0}^*=0.8172$, and the consumption from period $1$ onward on the wrong side of Laffe curve, $C_{0, i}^*=0.3125<C_{\text{Max}}^{\text{Laffer}}, i\geq 1$,  where $C_{\text{Max}}^{\text{Laffer}}$ denotes consumption at the peak of the Laffer curve.
%%the  tax rate  above the peak of Laffer curve in  period $0$ and  below the peak of Laffer curve from period $1$ onward. 
%By doing this, it effectively lowers the price of period $0$ debt,  and meanwhile, increases the value of surplus for period $1$ onward. To implement the Ramsey plan, Lucas-Stokey recommends a restructured term debt $\{b_{0, i}=0.1523\}_{i=1}^\infty$ that does not equal to the tail of initial debt $\{b_{-1, i}=0\}_{i=1}^\infty$, which suggests a negative  $\widehat{\Phi}=-3.93$, where $\widehat{\Phi}$ denotes the constructed Lagrange multiplier if $\{C_{1, i}=C_{0, i}^*\}_{i\geq 1}$ is imposed to the continuation Lagrangian (i.e., Appendix \ref{appendix:B}). Confronts the restructured debt  $\{b_{0, i}=0.1523\}_{i=1}^\infty$, the continuation Ramsey planner optimally choose $\{C_{1, i}=0.4874\}_{i=1}^\infty$, reneging the Ramsey plan that is on the wrong side of Laffer curve $C_{0, i}=0.3125<C_{\text{Max}}^{\text{Laffer}}$ for $i\geq 1$. 

%\item In case I, we analyze a numerical counterexample with an initial debt  structure: $b_{-1, 0}=0.3, b_{-1, 1}=0, b_{-1, i}=0.15, i\geq 2$ with $G=0.2$ and $\eta=\gamma=1$ analyzed in \citet{debortoli2021optimal}'s section 3 (and figure 3). %and figure XX.%We explain why our implementation works in this counterexample where \citeauthor{LucasStokey1983}'s implementation does not work. 

% planner sets the period $0$ consumption on the right side of Laffer curve,
%The Ramsey plan is $C_{0, 0}^* =0.6297$ (so that $\tau_{0, 0}^* =0.3703$), $C_{0, 1}^*=0.2930$ (so that $\tau_{0, 1}^* =0.7070$) and $C_{0, i}^*=0.5035$ (so that $\tau_{0, 0}^* =0.4965$) for all $i\geq 1$. At the peak of the Laffer curve, $C_{}^{\text{Laffer}}=0.4$ so that $\tau_{}^{\text{Laffer}}=0.6$. The period-0   Ramsey allocation  is  thus on the `right' side of the Laffer curve, the period-1 Ramsey  allocation is  on the `wrong' side of the curve, and period-2 onward Ramsey allocations are on the `right' side of the curve. The Lagrange multiplier for the Ramsey planner is $\Phi_0^\ast=2.41$. %lower than $C_{\text{Max}}^{\text{Laffer}}=0.4$. 
% $C_{\text{Max}}^{\text{Laffer}}=0.4$ but



%and the consumption from period $1$ onward on the wrong side of Laffe curve, $C_{0, i}^*=0.3125<C_{\text{Max}}^{\text{Laffer}}, i\geq 1$,  where $C_{\text{Max}}^{\text{Laffer}}$ denotes consumption at the peak of the Laffer curve.
%the  tax rate  above the peak of Laffer curve in  period $0$ and  below the peak of Laffer curve from period $1$ onward. 
%As a result, Lucas-Stokey recommends a restructured term debt 
%Using  formulas in this section, we have 
%In LS implementation, the Ramsey planner leaves its successor with $b_{0, 1}=0.0710$ and $b_{0, i}=0.2357$  for all $i\geq 2$, which differs from the tail of the initial debt structure $b_{-1, 1}=0$ and $b_{0, i}=0.15$ for all $i\geq 2$. Importantly, the resulting `constructed' Lagrange multiplier for its successor is negative: $\widehat{\Phi}_1=13.96$. Facing the inherited  debt structure, the continuation planner optimally chooses $C_{1, 1}=0.2930$ and $C_{1, i}=0.5035$ for $i\geq 2$ to confirm the continuation of the Ramsey plan. The resulting Lagrange multiplier for   the continuation planner is ${\Phi}_1^\ast=13.96$, which is same with  the `constructed' Lagrange multiplier 
%$\widehat{\Phi}_1=13.96$.  
%%${\Phi}_1^\ast=2.93>0$. 
%
%%C_{\text{Max}}^{\text{Laffer}}$ for $i\geq 1$. 	that is on the wrong side of Laffer curve 
%
%In our implementation, the Ramsey planner leaves its successor with short-term debt amount of $b_{0, 1}= 0.1478$, leaving term debts with time to maturities (weakly) than two unchanged: $b_{0, i}=b_{-1, i}=0.15$  for all $i\geq 2$. The resulting `constructed' Lagrange multiplier for the continuation planner equals  the Lagrange multiplier for Ramsey planner: $\widehat{\Phi}_1=\Phi_0^*=2.41$. Facing the inherited  debt structure, the continuation planner sets $C_{1, 1}=0.2930$ and $C_{1, i}=0.5035$ for $i\geq 2$, confirming the tail of the Ramsey plan. Importantly,   $\Phi_1^*=\widehat{\Phi}_1=\Phi_0^*=2.41$.
%

%\item In case II, we analyze a numerical counterexample with an initial debt  structure: $b_{-1, 0}=0.4, b_{-1, 1}=0, b_{-1, i}=0.15, i\geq 2$ with $G=0.2$ and $\eta=\gamma=1$ analyzed in \citet{debortoli2021optimal}'s section 3 (and figure 3). %and figure XX.%We explain why our implementation works in this counterexample where \citeauthor{LucasStokey1983}'s implementation does not work. 

% planner sets the period $0$ consumption on the right side of Laffer curve,

%The Ramsey plan is $C_{0, 0}^* =0.6459$ (so that $\tau_{0, 0}^* =0.3541$), $C_{0, 1}^*=0.0702$ (so that $\tau_{0, 1}^* =0.9298$) and $C_{0, i}^*=0.4102$ (so that $\tau_{0, 0}^* =0.5898$) for all $i\geq 1$. 
 Panel A in Table \ref{table2} reports  Ramsey allocations and  associated Ramsey Lagrange multipliers $\Phi_0^\ast>0$ for both cases. %. In both cases
The peak of the Laffer curve $C_{}^{\text{Laffer}}=0.4$ so that $\tau_{}^{\text{Laffer}}=0.6$ and  the  Ramsey allocation in period 1 is on the `wrong' side of the Laffer curve, while the Ramsey plan  consumptions for all future periods ($i\geq 2$) are on the `right' side of the Laffer curve. 
This opens  the possibility that the `constructed' Lagrange multiplier  $\widehat{\Phi}_1$ for the two cases under Lucas and Stokey's implementation can be either positive or negative, %suggesting that Lucas-Stokey's implementation  
{as  Panel B of Table \ref{table2} confirms.} %compare  Lucas-Stokey implementations for the two cases. 
%The period-0  Ramsey allocation  is  thus on the `right' side of the Laffer curve, the period-1 Ramsey  allocation is  on the `wrong' side of the curve, and period-2 onward Ramsey allocations are on the `right' side of the curve. The Lagrange multiplier for the Ramsey planner is $\Phi_0^\ast=13.24$. %lower than $C_{\text{Max}}^{\text{Laffer}}=0.4$. 
% $C_{\text{Max}}^{\text{Laffer}}=0.4$ but
%$b_{-1,i}, i\geq 1$ works in case I but does not work in case II 
In case I, the Ramsey planner restructures  debt to $b_{0, 1}=0.0710$ and $b_{0, i}=0.2357$  for all $i\geq 2$. Facing this debt structure, the continuation planner  chooses to continue  the Ramsey plan and $\Phi_1^*=\widehat{\Phi}_1=13.96$. Although the  period 1 Ramsey allocation is on the `wrong' side of the Laffer curve, the continuation planner  chooses to continue the Ramsey plan because it recognizes that the  benefits of devaluing its  future debts, $b_{0,i}$ for all  periods $i\geq 2$, outweighs the costs of being on  the `wrong' side of the Laffer curve in period 1. The Lagrange multipliers  $\Phi_1^*=\widehat{\Phi}_1>0$   express how the continuation Ramsey planner accepts this tradeoff and continues the Ramsey plan.
Consequently,  Lucas and Stokey's debt management policy  implements the Ramsey plan  in case I. %Because the `constructed 
%This is an example  where LS implementation works 

%It is worth noting that 
In case II, the Ramsey planner restructures  initial debts $\{b_{-1,i}\}_{i\geq 1}$ to $b_{0, 1}=0.0272$ and $b_{0, i}=0.2509$  for all $i\geq 2$. 
The planner backloads its term debt more than that in case I: $b_{0,1}$ decreases from $0.0710$ in case I to  $0.0272$ and  $b_{0,i}$ for all $i\geq 2$ increases from $0.2357$ in case I to $0.2509$.
The period 1 continuation planner's  `constructed' Lagrange multiplier is negative in case II: $\widehat{\Phi}_1=-2.92$. Facing the restructured term debt, instead of continuing  the Ramsey plan the continuation planner sets  $C_{1, 1}=0.3877$ and $C_{1, i}=0.6116$ for $i\geq 2$.  The resulting  Lagrange multiplier for   the continuation planner is ${\Phi}_1^\ast=1.92$. Since ${\Phi}_1^\ast>0>\widehat{\Phi}_1$, Lucas and Stokey's debt management policy does not  implement the Ramsey plan in case II.  
Notice that although it doesn't continue  the Ramsey plan,
the continuation  planner  chooses to stay on the `wrong' side of the Laffer curve in period 1:  $C_{1, 1}=0.3877 <C^{\text{Laffer}}$. 
Relative  to case I, in case II future (restructured) debts $b_{0, i}=0.2509$ for all $i\geq 2$ are   so high that keeping  $C_{1, 1}$ on the `wrong' side of the Laffer curve is optimal because  
that   reduces the value of those debts. %so much that it is optimal for the continuation planer  to do so. %, the period-1 planner chooses 
%In summary, in case II, %has taught us 

%In summary, but the corresponding `constructed'  Lagrange multiplier is negative: $\widehat{\Phi}_1=-2.92$.  LS implementation thus fails in this case. 

Panel C Table \ref{table2} confirm that under our timing protocol, our debt management policy implements the Ramsey plan in both cases. 
By keeping term debts with times to maturity (weakly) than two unchanged, i.e.,  $b_{0, i}=b_{-1, i}=0.15$  for all $i\geq 2$, while  using short-term (one-period) debt to track accumulated liabilities so that $b_{0, 1}= \Pi_1^\ast = 0.1478$ in case I and $b_{0, 1}= \Pi_1^\ast = 0.0538$ in case II,  the period 1  continuation planner's Lagrange multiplier  equals  its  `constructed' Lagrange multiplier, which also equals Ramsey planner's Lagrange multiplier:  $\Phi_1^*=\widehat{\Phi}_1=\Phi_0^*$. % This is why our implementation always works  under our implementation. 

In summary, the quantitative examples in Table \ref{table2} % using  two examples slightly richer than \cite{debortoli2021optimal}'s counterexamples, we 
show first, 
%using a Laffer curve for the primary surplus to determine whether a Ramsey plan can be implemented or not under  Lucas-Stokey's implementation is unreliable. %A time-invariant  Lagrange multiplier for continuation planners is a reliable criterion. 
that tax rates on the `wrong' side of the Laffer curve are not universally coincident  symptoms of failure of Lucas and Stokey's debt management policy  to implement a  Ramsey plan  under
their timing protocol; and second, that under our timing protocol and debt management policy, a Ramsey plan can always  be implemented, even if  the Ramsey planner sets some tax rates above the peak of the Laffer curve. 
%In contrast, our implementation always works. 

%the Ramsey allocations and the associated Ramsey Lagrange multiplier $\Phi_0^\ast>0$ for the two cases. For both cases, 
%at the peak of the Laffer curve, $C_{}^{\text{Laffer}}=0.4$ so that $\tau_{}^{\text{Laffer}}=0.6$. In both cases, the period-1 Ramsey allocation is on the `wrong' side of the Laffer curve and Ramsey allocations for all future periods ($i\geq 2$) are on the `right' side of the Laffer curve. 



%and the consumption from period $1$ onward on the wrong side of Laffe curve, $C_{0, i}^*=0.3125<C_{\text{Max}}^{\text{Laffer}}, i\geq 1$,  where $C_{\text{Max}}^{\text{Laffer}}$ denotes consumption at the peak of the Laffer curve.
%the  tax rate  above the peak of Laffer curve in  period $0$ and  below the peak of Laffer curve from period $1$ onward. 
%As a result, Lucas-Stokey recommends a restructured term debt 
%Using  formulas in this section, we have 

%In LS implementation, the Ramsey planner which differs from the tail of the initial debt structure $b_{-1, 1}=0$ and $b_{0, i}=0.15$ for all $i\geq 2$. 

%Importantly, %${\Phi}_1^\ast=2.93>0$. 

%C_{\text{Max}}^{\text{Laffer}}$ for $i\geq 1$. 	that is on the wrong side of Laffer curve 
%, For example, in case II, the Ramsey planner leaves its successor with short-term debt: 

% and 


%This is because of our short-term debt management policy leaves long-term debt untouched, minimizing the term-debt adjustment and preserving the shadow price of debt.

% Facing the inherited  debt structure, the continuation planner confirms the tail of the Ramsey plan. 
% `constructed' Lagrange multiplier for

%	\item Using  case II, we show  that  the Ramsey plan is implemented even when the Ramsey plan consumption is on the wrong side of Laffer curve. This example shows that whether it is above the Laffer curve is not relevant .... 
%	
%	Note that the constructed Lagrange multiplier $\widehat{\Phi}_1=13.96$ and continuation Ramsey Lagrange multiplier $\Phi_1^*=13.96$ are the same. This is necessary for the Ramsey plan to be implemented. 
%	
%	The case II Ramsey planner faces an initial debt sequence that includes a high payment that is constant from  period 2 onward. The Ramsey plan sets period 1 consumption to be on the wrong side of Laffer curve: $C_{0, 1}^*=0.2930<C_{\text{Max}}^{\text{Laffer}}$. \citeauthor{LucasStokey1983}'s debt management strategy prescribes a restructured term debt $\{ b_{0,1}=0.0710, b_{0, i}=0.2357, i\geq 2\}$ that will confront a  period 1 continuation planner. The continuation Ramsey plan confirms the tail of the Ramsey plan and period 1 consumption remains on the wrong side of Laffer curve  $C_{1, 1}^*=C_{0, 1}^*=0.2930$. The continuation planner wants to stay on the wrong side of Laffer curve because that prospective  debt payments induce  the government to choose a low level of the period 1 consumption together with  high levels of consumption from  period 2 onward. That  increases period   $1$   value of the  period 1 surplus and lowers  value of the debt payment stream from period 2 onward.
%	
%	\item Using Case III, 
%	the debt due in period $0$ increases to 0.4 from its value of  0.3 in case II. The Ramsey plan puts the time $1$  tax rate on the wrong side of Laffer curve.
%The continuation planner chooses  not to follow the Ramsey plan, but  nevertheless  sets a time $1$ tax rate that is still  the wrong side of Laffer curve \TS{ and sustains a high continuation debt sequence  for period 2 onward. } \TS{Tom -- edit and refine last part of the previous sentence.}







%\TS{Neng and Wei -- please check how I edited previous paragraph.}




%
%\subsection{Illustrating Mechanics of our Implementation} 
%
%
%Table \ref{table3}  describes outcomes under our timing protocol and debt management strategy.
%In case I,
%Moreover, the continuation planner  implements the Ramsey plan in  all three cases. Lagrange multipliers on implementability constraints are now invariant in each case. 

%\renewcommand\arraystretch{1.65}
%\begin{table}[h!]
%	\caption{Our implementation and continuation Ramsey plan. $\eta=\gamma=1, G=0.2, \rho=0.5, \Delta=1$: {$C^{\text{Laffer}}_{\text{Max}}=0.4$}. \label{table3}}
%	\centering
%	\resizebox{0.58\textwidth}{!}{
%		\begin{tabular}{l|c|c|c|c|c|c|c|c|c}
%			\hline
%			{	Case} & \multicolumn{3}{c|}{I} 	& \multicolumn{3}{c|}{II} & \multicolumn{3}{c}{ III }  \\ 		
%			\hline 
%			
%			Period $i$ &  $0$ & $1$ & $\geq 2$   &  $0$ & $1$ & $\geq 2$ &  $0$ & $1$ & $\geq 2$  \\
%			\hline
%%			$b_{-1, i}$ &  0.6  & 0 & 0   &  0.3  & 0 & 0.15  &   0.4  & 0 & 0.15   \\
%%			\hline
%			\multicolumn{10}{c}{  \bf A. Ramsey plan}\\ 
%			
%			\hline 
%%			$C_{0, i}^*$ &  0.8172  & {0.3125} & {0.3125} &  0.6297 &  {0.2930}  &  0.5035  &  0.6459 &  {0.0702} &  0.4102 \\
%%			\hline
%			LM $\Phi_0^*$ &  \multicolumn{3}{c|}{{2.98}}  &  \multicolumn{3}{c|}{{2.41}}  &   \multicolumn{3}{c}{{13.24}}   \\
%			\hline  
%			\multicolumn{10}{c}{  \bf B. Our implementation}\\ \hline
%%			$b_{0, i}$ &  n.a.  & 0.3871 & 0  &  n.a.   & 0.1478 & 0.15  &  n.a.   & 0.0538 & 0.15     \\
%%			\hline 
%			%	Constr' & \multicolumn{3}{c|}{}  &   \multicolumn{3}{c|}{}  &  \multicolumn{3}{c}{} \\
%			LM $\widehat \Phi_1$&  \multicolumn{3}{c|}{{2.98}}  &  \multicolumn{3}{c|}{{2.41}}  &   \multicolumn{3}{c}{{13.24}}    \\  \hline
%			Works?&  \multicolumn{3}{c|}{{{Yes}}} & \multicolumn{3}{c|}{{Yes}}  &   \multicolumn{3}{c}{{{Yes}}}    \\ 
%			\hline    
%			\multicolumn{10}{c}{  \bf C. Cont Ramsey plan}\\ \hline
%%			
%%			$C_{1, i}^*$ &  n.a. &  {0.3125} & {0.3125} &  n.a.  & { 0.2930 } &  0.5035  &  n.a. &  { 0.0702} &  0.4102  \\
%%			\hline
%			LM $\Phi_1^*$  &  \multicolumn{3}{c|}{{2.98}} &  \multicolumn{3}{c|}{{2.41}}  &   \multicolumn{3}{c}{{13.24}}   \\   
%			\hline
%			
%		\end{tabular}
%}
%\end{table} %Our implementation

%\WJ{In case I, which is the example featured in 
%\citet{debortoli2021optimal}, for initial debt $b_{-1, 0}=0.6$, the Ramsey planner sets the period $0$ consumption on the right side of Laffer curve, $C_{0, 0}^*=0.8172$, and the consumption from period $1$ onward on the wrong side of Laffe curve, $C_{0, i}^*=0.3125<C_{\text{Max}}^{\text{Laffer}}, i\geq 1$,  where $C_{\text{Max}}^{\text{Laffer}}$ denote consumption at the peak of the Laffer curve.
%%the  tax rate  above the peak of Laffer curve in  period $0$ and  below the peak of Laffer curve from period $1$ onward. 
%By doing this, it effectively lowers the price of period $0$ debt,  and meanwhile, increases the value of surplus for period $1$ onward. To implement the Ramsey plan, our implementation leaves an one-period debt $b_{0, 1}=0.3871$ and untouched term debt $\{b_{0, i}=b_{-1, i}=0; i\geq 2\}$. Confronting this initial debt, the continuation Ramsey planner optimally choose $\{C_{1, i}=0.3125\}_{i=1}^\infty$, confirming the Ramsey plan, even when $C_{0, i}^*$ is on the wrong side of Laffer curve. Moreover, the continuation Lagrange multiplier equals the Ramsey Lagrange multiplier, $\Phi_1^*=\Phi_0^*$. 
%In case II and III, where Ramsey consumption in period $1$  $C_{0, 1}^*$ is also on the wrong side of Laffer curve, we show that our implementation also works and holds invariant Lagrange mulitipliers.
%
%}








%\subsection{Debt Management Strategies in the Two Implementations} % in Mechanics of our Implementation} 


\subsection{Declining $\{b_{-1, i}\}$ bring increasing $\{\tau_{0,i}\}$ }\label{subsecdecliningb} 
%exponential Consider
Here government spending unfolds at a constant rate $G_i=G=0.2$ for all $i\geq 0$ and the initial   debt structure $\{b_{-1, i} \}_{i\geq 0}$  starts with $b_{-1, 0} = b_0=0.2$ in period 0 and decays exponentially at a rate of $\theta=0.5$ each period. 
%where $b_{-1, i}=b_0e^{-\theta i\Delta}$ with $b_{-1,0}=b_0$. We set % in that the government's debt maturity distribution satisfies an exponential function.\footnote{Empirical studies document that the US government's debt  service stream is a decreasing function of maturity. } 
Figure \ref{figexpb} reports the Ramsey  outcome  and our implementation. %how a Ramsey planner  finances an initial debt term structure $b_{0, t}=0.2\times e^{-0.6 t}$ and constant government spending    and  labor supply $N_{0, i}^\ast=C_{0, i}^\ast+G_i$
As   $b_{-1, i} $ declines monotonically  as time period $i$ advances (panel A), the tax rate $\tau_{0, i}^\ast$  increases  (panel B.) %and    consumption $C_{0, i}^\ast=1-\tau_{0, i}^\ast$  decreases over time. 
%The intuition for this result is as follows. 
 Since $C _{0,i}^\ast=1-\tau_{0,i}^\ast$, by setting a higher tax rates  when debt $b_{-1, i}$ is higher, the Ramsey planner makes the marginal utility $U_C(C _{0,i}^\ast)$ higher. The Ramsey planner thereby minimizes the time 0 discounted value of $\{b_{-1, i} \}_{i\geq 1}$. 
Consumption growth rate $g_{0,i}^\ast= \log({C}_{0,i+1}^\ast/{C}_{0,i}^\ast) $ is   negative,\footnote{In this example, $C_{0, i}^* = 2 b_0e^{-\theta i\Delta}\Phi_0^*/ (\sqrt{1+4b_0e^{-\theta i\Delta}\Phi_0^*(1+\Phi_0^*)}-1)$. Therefore, the consumption growth rate  is negative: $g^*_{0, t} = \log({C}_{0,i+1}^\ast/{C}_{0,i}^\ast)<0$   because $d C_{0, i}^* /di = -b_0\Phi_0^*\theta\Delta e^{-\theta i \Delta}/\sqrt{1+4b_0e^{-\theta i\Delta}\Phi_0^*(1+\Phi_0^*)}<0$. } %lower than the (minus) household's discount rate ($-\rho$) at $t=0$, 
but becomes less negative over time and  approaches zero. 
Coincidentally, the  equilibrium interest rate $r_{0,i}^\ast = \rho + g_{0,i}^\ast$   increases over  time and   approaches the discount rate  $\rho=0.5$  (see panel C). 
By keeping the interest rate low in early periods when debt is high, the planner lowers the market value of its liabilities. %liability$b_{-1,i}$ 
%high debt level 
% As  %By letting consumption    consumption $C_{0, i}^\ast=1-\tau_{0, i}^\ast$  decreases over time. 
%, consumption decreases as $b_{-1, i}$ decreases over time, 
% 

\begin{figure}[h!]
\caption{\label{figexpb} { Ramsey plan and our debt management for exponentially decaying debt structure: $b_{-1, i}=b_{0}e^{-\theta  i\Delta}$.}
Parameter values: $b_{0}=b^H=0.2, \, \theta=0.5, \eta=\gamma=1, \rho=0.5, G=0.2,$ and $ \Delta=1$. %The higher the value of $\kappa$, the lower the debt capacity, the higher the tax rate, and the lower the drift of debt.
%	Other than $\zeta-r$,	all other parameter value are reported in Table \ref{tabelparam}. 
}	


\begin{center}
\includegraphics[height=0.75\textheight,width=\textwidth]{figure_202405/beyondDNY_expb}
\end{center}
%The higher the value of $\zeta$, the more impatient the government, the higher the jumpy debt issuance boundary $\underline{b}$, the lower the tax, and the higher the drift of the debt.  The equilibrium debt capacity is independent of the value of $\zeta$. %$\kappa=7> \frac{\rho-g+\xi}{c(\gamma)} \left(\frac{c((r-g)\bar{b}+\gamma)}{\rho-g} +\frac{\xi v(0)}{\rho-g+\xi} \right) = 4.7078$.
\end{figure}


%the government's debt stream down. %a low marginal utility $U_C(C _{0,i}^\ast)$ to 
%As consumption starts high and decreases 


The Ramsey planner  front loads  consumption so much that  increases in   output $N_{0,i}^\ast=C_{0,i}^\ast +G$  overwhelm  decreases in the tax rate $\tau _{0,i}^\ast=1-C_{0,i}^\ast$, making the primary surplus $S_{0,i}^\ast= \tau _{0,i}^\ast N_{0,i}^\ast - G= (1-C_{0,i}^\ast) (C_{0,i}^\ast +G_{})- G$  increase over time.
Since $b_{-1,i}$ decreases over time,  $S_{0,i}^\ast - b_{-1,i}$ is  negative at $i=0$, turns positive at $i=2$, and  converges to $15\% $ (see panel D).\footnote{As $i\to \infty$, $b_{-1, i}\to 0$. We can show that $C_{0, i}^*\to \frac{1}{1+\Phi_0^*}$ and 
$\lim_{i\to\infty}S_{0,i}^\ast - b_{-1,i} =\frac{1}{1+\Phi_0^*}\left(1-\frac{1}{1+\Phi_0^*}-G_{} \right)=15\%$, using $G_i=0.2$ and $\Phi_0^*=0.99$.} %

%Recall that  %multiplied by $
%Next, we turn to our implementation. 
 Government liabilities accumulate according to  %%By inducing an increasing equilibrium $r_{0,i}^\ast$ path (panel C), the planner  lowers the time-0 value of its wedge $S_{0,i}^\ast - b_{-1,i}$ shown in panel D.  Panel E shows that  the cumulative liability 
$	\Pi_{i+1}^* = \left(\Pi_i^* +(b_{-1, i}-S_{0, i}^*)\Delta \right)e^{r_{0, i}^*\Delta}\,, 
$ 
%$\Pi_i^\ast=\sum_0^{i-1} e^{\sum_0^{u-1} r_{0,v}^\ast } (b_{-1,u}-S_{0,u}^\ast)  $ 
starting from $\Pi_0^\ast=0$. Panel E shows that   $\Pi_i^\ast$ increases over time and converges to $0.38$ as $i\rightarrow \infty$.\footnote{{By letting $i\to\infty$ in $\Pi_{i}^* = \left(\Pi_{i-1}^* +(b_{-1, i-1}-S_{0, i-1}^*)\Delta \right)e^{r_{0, i-1}^*\Delta}$, we obtain $\lim_{i\to\infty}\Pi_{i}^* = \left(\lim_{i\to\infty}\Pi_i^* +\lim_{i\to\infty}(b_{-1, i-1}-S_{0, i-1}^*)\Delta \right)e^{\rho\Delta}$ and thus $\lim_{i\to\infty}\Pi_{i}^*=\frac{\Delta e^{\rho\Delta}}{e^{\rho\Delta}-1}\lim_{i\to\infty}\left(S_{0,i-1}^*-b_{-1, i-1}\right)= \frac{\Delta e^{\rho\Delta}}{e^{\rho\Delta}-1}\frac{1}{1+\Phi_0^*}\left(1-\frac{1}{1+\Phi_0^*}-G_{} \right)$ .} } %\Pi_i^\ast$
The government honors its liabilities by issuing one-period debt $b_{i-1,i} \Delta $  at time $i-1$ that  equals the sum of cumulative liability $\Pi_i^\ast$ (in panel E) and the initial debt due in period $i$: $b_{-1,i} \Delta$. %converging to 0.38.  
As $i\rightarrow \infty$, $b_{-1,i} \Delta \rightarrow 0$,  $\lim_{i\to \infty} \; b_{i-1,i} \Delta  = \lim_{i\to \infty} \Pi_i^* = 0.38.$  %\NW{add a similar footnote as we did for the other case?} 
% the state variable $X_t^\ast = \Pi_t^\ast U_{C,t}^\ast$ also increases over time (panel F). This is because both $\Pi_t^\ast$ and $U_{C,t}^\ast$ increase over time. 


%Next, we turn to an example where the initial debt structure $\{b_{0, t}\}$ is cyclical. 

\begin{figure}[h!]
\caption{\label{figcyclicb} { Ramsey plan and  debt management in an example where $b_{-1, i}$ is high ($b^H$) in even periods and is low ($b^L$)  in odd periods.}
Parameter values: $b^H=0.2, b^L=0, \eta=\gamma=1, \rho=0.5, G=0.2,$ and $ \Delta=1$. %The higher the value of $\kappa$, the lower the debt capacity, the higher the tax rate, and the lower the drift of debt.
%	Other than $\zeta-r$,	all other parameter value are reported in Table \ref{tabelparam}. 
}	

\begin{center}
\includegraphics[height=0.75\textheight,width=\textwidth]{figure_202405/beyondDNY_cyclicb}
\end{center}
%The higher the value of $\zeta$, the more impatient the government, the higher the jumpy debt issuance boundary $\underline{b}$, the lower the tax, and the higher the drift of the debt.  The equilibrium debt capacity is independent of the value of $\zeta$. %$\kappa=7> \frac{\rho-g+\xi}{c(\gamma)} \left(\frac{c((r-g)\bar{b}+\gamma)}{\rho-g} +\frac{\xi v(0)}{\rho-g+\xi} \right) = 4.7078$.
\end{figure}


\subsection{Periodic $\{b_{-1,i}\}$ brings periodic $\{\tau_{0,i}\}$}\label{subsecperiodicb}  %: $b_{-1, i} = b^H$ if $i$ is even and $b_{-1, i} = b^L$ if $i$ is odd} 
%, $b^H$ and $b^L<b^H$ over time:   
Here initial debts   $b_{-1, i}$ oscillate between two levels: 
$b_{-1, i} = b^H$ in even  periods and $b_{-1, i} = b^L<b^H$  in odd periods. We set $b^H=0.2$ and $b^L=0$, start with $b_{-1,0}=b^H$.  Figure \ref{figcyclicb} plots outcomes under the  Ramsey plan and our timing protocol and debt management policy.

% that starts with b0 “ 0.2 in period 0 and decays exponentially at a rate of ξ“0.6 each period (and fix government spending at a constant level: Gi “ 0.2 for all i ě 0.) 
%This subsection  takes on a high value, $b^H$, $b^L<b^H$, i $b_{-1, i}$ is high in odd periods and low in even periods, so that

%but in a lock step with  $
%As the debt structure $b_{-1, i}$ oscillates,   between two levels:
The tax rate $\tau_{0, i}^\ast$   takes  %ing structure of $b_{-1, i}$, but takes 
a high value when  $b_{-1, i}$ is low.  Panel B of Figure \ref{figcyclicb} indicates that $\tau_{0, i}^\ast = \tau^H =0.47$  in odd periods and $\tau_{0, i}^\ast = \tau^L =0.33$ in   even periods.  % As ,   so that  $\tau_{0, i}^\ast = \tau^H $ when $b_{-1, i} =b^L$ and  $\tau_{0, i}^\ast = \tau^L$ when $b_{-1, i} =b^H$ %it reaches a peak whenever  the $\{b_{-1, i}\}$ reaches a trough 
Since $C _{0,i}^\ast=1-\tau_{0,i}^\ast$, the Ramsey planner  sets    marginal utility $U_C(C _{0,i}^\ast)$ to be low in  high debt periods  and   high %marginal utility $U_C(C _{0,i}^\ast)$ 
in  low debts period. This policy makes  the time 0 value of the government's debt payments be   low.


%chooses a high level of consumption   in periods when $b_{-1, i} = b^H$ to 

%down. %is optimal because the Ramsey planner 
%the equilibrium consumption process has the same peaks and troughs as the $\{b_{0, i}\}$ process. %\footnote{In this example, $C_{0, i}^* = \left(\sqrt{1+4(b_0+\theta_b \sin\left(\frac{\pi}{2} t\right)\Phi_0^*(1+\Phi_0^*)}-1\right) / (2(1+\Phi_0^*))$, which is cyclical  with the same period as $b_{0, t}$. %We can verify the cyclical properties The Lagrangian multiplier is $\Phi_0^*=0.7256$. } 

%\end{document}   Therefore,  in a high $b_{-1, i}$ state, 
%What determines  $r_{0,i}^\ast $? %is thus when $b_{-1,i}$ is low, as consumption is expected to decline 
%Recall that when debt $b_{-1, i}$ is high, consumption $C _{0,i}^\ast=1-\tau_{0,i}^\ast$ is  high, resulting in a negative consumption growth rate: $g_{0,i}^\ast= \log({C}_{0,i+1}^\ast/{C}_{0,i}^\ast)<0$. As 
% g_{0,i}^\ast$ where  is  
The  equilibrium interest rate  $r_{0,i}^\ast$ equals the sum of  $\rho$ and the consumption growth rate $g_{0,i}^\ast= \log({C}_{0,i+1}^\ast/{C}_{0,i}^\ast)$.   This sum   moves in lock step with   $ g_{0,i}^\ast$. Evidently,  $ g_{0,i}^\ast$ moves inversely with  $b_{-1, i} $  because 
$g_{0,i}^\ast<0$  when  consumption is high and consumption is high when initial  debt due  is high. %which is also debt (and also a high) 
%$g_{0,i}^\ast= \log({C}_{0,i+1}^\ast/{C}_{0,i}^\ast)<0 
Thus,   $r_{0, i}^\ast$ is high  in odd periods: $r^H =0.74$ when $b_{-1, i} = b^L$, but  is low ($r^L=0.26$) in even periods  (see panel C). 
% $r_{0, i} = 
%,  resulting an equilibrium interest rate  of $r_{0,i}^\ast =0.26$ lower than $\rho=50\%$
%when  is high. 
%so that marginal utility $U_C()$ is low, which  is optimal because the Ramsey planner lowers the value of its debt obligations. 
%which is also cyclical with the same period as but different peaks/troughs  from those for $b_{-1,i}$. %We mark the peaks/troughs for $r_{0,t}^\ast $ with triangles to indicate that they are different from the peaks/troughs in panels A, B, and D for $b_{0,t}$, $\tau_{0,t}^\ast$, and  the wedge between primary surplus and $b_{0,t}$: $S_{0,t}^\ast-b_{0,t}$.\footnote{Note that this wedge shares the same peaks and troughs as the tax rate $\{\tau_{0, t}^\ast\}$ process.} 


Panel D  shows that $S_{0,i}^\ast - b_{-1,i}$  is  
$S_{0,i}^\ast - b_{-1,i}=-0.11 < 0$ in even periods when debt is high ($b^H=0.2$) and $S_{0,i}^\ast - b_{-1,i}=0.14 >0$ in odd periods when $b_{-1,i}=b^L=0$. 
Panel E reports  cumulative liabilities %%By inducing an increasing equilibrium $r_{0,i}^\ast$ path (panel C), the planner  lowers the time-0 value of its wedge $S_{0,i}^\ast - b_{-1,i}$ shown in panel D.  Panel E shows that  the cumulative liability 
$	\Pi_{i}^* = \left(\Pi_{i-1}^* +(b_{-1, i-1}-S_{0, i-1}^*)\Delta \right)e^{r_{0, i-1}^*\Delta}\,, 
$ 
%$\Pi_i^\ast=\sum_0^{i-1} e^{\sum_0^{u-1} r_{0,v}^\ast } ( b_{-1,u}-S_{0,u}^\ast )  $, 
starting from $\Pi_0^\ast=0$. It shows that   $\Pi_{i}^\ast$ is 0.14 in an odd periods  and zero  in an even periods.  A high level of $\Pi_{i}^\ast$ accompanies high debt ($b_{-1,i-1}=b^H=0.2$) and low surplus ($S_{0,i-1}^\ast=0.09$) in  period $i-1$. A zero level of $\Pi_{i}^\ast$  is an outcome of zero debt ($b_{-1,i-1}=b^L=0$) and  a high surplus ($S_{0,i-1}^\ast=0.14$) in  period $i-1$.  Here the government policy uses the surplus  to finance  all of its accumulated  liabilities $\Pi_{i-1}^*=0.14$. % $\Pi_{i+1}^\ast=0.$  %Similarly, when debt is zero  ($b_{-1,i}=b^L=0$) and surplus is high  ($S_{0,i}^\ast=0.14$), 

%as the government pays down its debt 
%debt is high in the preceding period $i-1$
%increases  and converges to  $XX$  as $i\rightarrow \infty$. %\Pi_i^\ast$
Finally, in our implementation, the government honors its liabilities by issuing one-period debt $b_{i-1,i} \Delta $ that  equals the sum of its  cumulative liabilities $\Pi_i^\ast$ (in panel E) and initial debt term debt due in period $i$: $b_{-1,i} \Delta$. %converging to 0.38.  
In an even period $i$, the government simply issues one-period debt to pay for its initial debt $b_{-1,i}$, since $\Pi_{i}^\ast=0$. After that,  in  odd periods $i$, the government  issues one-period debt to pay for its cumulative liability $\Pi_{i}^\ast=0.14$, since $b_{-1,i}=0$. Consequently,  $b_{i-1,i}$ oscillates between 0.2 in  even periods and 0.14 in  odd periods (panel F.) 


%As $i\rightarrow \infty$, $b_{-1,i} \Delta \rightarrow 0$, we thus have $\lim_{i\to \infty} \; b_{i-1,i} \Delta  = \lim_{i\to \infty} \Pi_i^* = 0.38.$  %\NW{Explain how we got this number.} 
% the state variable $X_t^\ast = \Pi_t^\ast U_{C,t}^\ast$ also increases over time (panel F). This is because both $\Pi_t^\ast$ and $U_{C,t}^\ast$ increase over time. 

%
%
%that the  cumulative liability $\Pi_i^\ast$ under our debt management and the one-period debt $b_{i-1,i} \Delta $ are also cyclical. Following our debt management strategy, $b_{i-1, i}\Delta=\Pi_i^*+b_{-1, i}\Delta$ is cyclical because both $\Pi_i^*$ and $b_{-1, i}$ are cyclical.
%
%Finally, panels E and F show that the  cumulative liability $\Pi_i^\ast$ under our debt management and the one-period debt $b_{i-1,i} \Delta $ are also cyclical. Following our debt management strategy, $b_{i-1, i}\Delta=\Pi_i^*+b_{-1, i}\Delta$ is cyclical because both $\Pi_i^*$ and $b_{-1, i}$ are cyclical.

\subsection{Periodic $\{G_i\}$  and constant $\{b_{-1,i}\}$ brings constant $\{\tau_{0,i}\}$}\label{subsecperiodicG} %  \citep{barro1979determination}} 
% ($G^H$)  is $G_{i}=G^L$
%(high in even periods and low  in odd periods) 
This example assumes  a cyclical government spending process  $G_{i}$ that oscillates between two levels: $G_{i} = G^H$ in even  periods and $G_{i} = G^L<G^H$  in odd periods. We set 
 $G_{0}=G^H$ and $G^H=0.3$ and $G^L=0.1$, fix $b_{-1,i}=b=0.1$ for all $i$, and plot Ramsey outcomes and continuation debt structures under our implementation in Figure \ref{figcycyclicg}. 

%Figure \ref{figcycyclicg}  presents an example in which the government  finances a cyclical government spending process (high in even periods and low  in odd periods) and  a constant $b_{-1,i}$ process. 
The Ramsey planner  sets a constant  flat tax rate  $\tau _{0,i}^\ast =61\%$ for all $i$, since $\tau _{0,i}^\ast = 1 - C_{0,i}^\ast$ and $C_{0,i}^\ast=0.39$ (panel B). Outcomes resemble the tax smoothing   in  \citet{barro1979determination}, where the   %in which  the intertemporal properties of the optimal  tax rate are  disconnected from the intertemporal properties of that exogenous government expenditure process that is to  be financed. 
the interest rate is an exogenous constant, and example 4 of  \citet[sec.~3]{LucasStokey1983} , where the interest rate is   endogenous. % unlike  \citet{barro1979determination} where it  is .  
In our example  the  interest rate is endogenous but, because    consumption $C_{0,i}^\ast$ is constant over time,  turns out to be constant: $r_{0,i}^\ast=\rho + \log({C}_{0,i+1}^\ast/{C}_{0,i}^\ast) = \rho$  (see panel C). 

% Panel D shows that the primary surplus process 
Since  $\tau_{0,i}^\ast$ and $N_{0,i}^\ast = G_i + C_{0,i}^\ast$  are constant, the primary surplus $S_{0,i}^\ast =\tau_{0,i}^\ast N_{0,i}^\ast - G_i  = C_{0, i}^*\left(1-C_{0, i}^*-G_i \right)$ moves inversely to  $G_i $. Consequently,  $S_{0,i}^\ast - b_{-1, i}$ is also cyclical, attaining a high value when spending  is low ($G_i=G^L$) and a low value when spending  is high ($G_i=G^H$). %Solid dots highlight the troughs/peaks in panels A and D.  Triangles in panel D highlights points where $S_{0,i}^\ast=b_{0,i}=0.1$.
Panel E reports the  government's accumulated  liabilities %%By inducing an increasing equilibrium $r_{0,i}^\ast$ path (panel C), the planner  lowers the time-0 value of its wedge $S_{0,i}^\ast - b_{-1,i}$ shown in panel D.  Panel E shows that  the cumulative liability 
$\Pi_i^\ast$  (with $\Pi_0^\ast=0$), which  It shows that   $\Pi_i^\ast$ 
oscillate between  0.08 in an odd periods  and zero  in even periods.  A high $\Pi_i^\ast$  accompanies high spending ($G^H$) and a negative $S_{0,i-1}^\ast - b_{-1, i-1}$.
%in the preceding period $i-1$.
 Similarly, if government  spending   $G_{i-1}$ is low, the government manages its surplus in period $i-1$ to cover all of  its accumulated liabilities $\Pi_{i-1}^\ast$, making  $\Pi_i^\ast=0$. % \NW{Check the subscripts -- maybe a bit more explanation about the mechanism..} 
% and surplus is high  as 

%Finally, in our implementation, the government honors its liability by issuing one-period debt: $b_{i-1,i} \Delta $ that  equals the sum of cumulative liability $\Pi_i^\ast$ (in panel E) and the initial debt due in period $i$: $b_{-1,i} \Delta$. %converging to 0.38.  

Turning  to the government's debt management strategy under our implementation, since $b_{-1,i} = b=0.1 $ for all $i \geq 0$, in  even periods $i$, the government  issues one-period debt to finance $b=0.1$ and $\Pi_{i}^\ast=0$, but in odd periods $i$, the government  issues one-period debt to finance the sum of its accumulated liabilities $\Pi_{i}^\ast = 0.08$ and its time $i$ initial debt obligation $b=0.1$. Consequently,  $b_{i-1,i}$ oscillates between 0.1 in  even periods and 0.18 in  odd periods (panel F) .


%
%
%%While we have focused on settings where $\{b_{0,t}, G_t\}$ are continuous in $t$, 
%Our technical machinery also applies   when $\{b_{0,t}, G_t\}$ are discontinuous. 
%In the next two  sections, we describe implementations of  Ramsey plans in settings with C\`adl\`ag  $\{b_{0,t}\}$ processes.
%Section \ref{sec:instantaneous}  shows how to use instantaneous debt to finance the cumulative liability $\Pi_t^\ast$. Section \ref{sec:implement}  connects our implementation   for C\`adl\`ag  $\{b_{0,t}, G_t \}$ processes to  
%%In the next two sections, we will utilize our  Ramsey 
%appropriate limits of  continuous $ \{b_{0,t}, G_t\}$ processes. We show how settings with continuous $ \{b_{0,t}, G_t\}$ processes in Section \ref{sec:quantitative}   can approximate Ramsey plans for discontinuous $ \{b_{0,t}, G_t\}$ processes. 

\begin{figure}[h!]
\caption{\label{figcycyclicg} Ramsey plan and  debt management in an example where government spending $G_{i}$ is high ($G^H$) in even periods and is low ($G^L$)  in odd periods, and $b_{-1, i}=b$ for all ${i\geq 0}$. 
%		Parameter values: $b^H=0.2, b^L=0, \theta=0.5, \eta=\gamma=1, \rho=0.5, G=0.2,$ and $ \Delta=1$. %The higher the value of $\kappa$, the lower the debt capacity, the higher the tax rate, and the lower the drift of debt.
%	Other than $\zeta-r$,	all other parameter value are reported in Table \ref{tabelparam}. 
%	}	
%	\caption{\label{figcycyclicg} { Ramsey plan and our debt management for cyclical $G_{i}=g_0$ when $i$ is even and $G_{i}=g_1$ when $i$ is odd.}
Parameter values: $G^H=0.3, G^L=0.1, b=0.1, \eta=\gamma=1, \rho=0.5,$ and $\Delta=1$. %The higher the value of $\kappa$, the lower the debt capacity, the higher the tax rate, and the lower the drift of debt.
%	Other than $\zeta-r$,	all other parameter value are reported in Table \ref{tabelparam}. 
}	
\begin{center}
\includegraphics[height=0.75\textheight,width=\textwidth]{figure_202405/beyondDNY_cyclicg}
\end{center}
%The higher the value of $\zeta$, the more impatient the government, the higher the jumpy debt issuance boundary $\underline{b}$, the lower the tax, and the higher the drift of the debt.  The equilibrium debt capacity is independent of the value of $\zeta$. %$\kappa=7> \frac{\rho-g+\xi}{c(\gamma)} \left(\frac{c((r-g)\bar{b}+\gamma)}{\rho-g} +\frac{\xi v(0)}{\rho-g+\xi} \right) = 4.7078$.
\end{figure}


%
%\begin{figure}[h!]
%	\caption{\label{figcycyclicg2} { Ramsey plan and our debt management for cyclical $G_{i}=g_0$ when $i$ is even and $G_{i}=g_1$ when $i$ is odd.}
%		Parameter values: $g_0=0.3, g_1=0.1, \{b_{0, i}=0.1\}_{i\geq 0}, \theta=0.5, \eta=\gamma=1, \rho=0.5, \Delta=1/2$. %The higher the value of $\kappa$, the lower the debt capacity, the higher the tax rate, and the lower the drift of debt.
%		%	Other than $\zeta-r$,	all other parameter value are reported in Table \ref{tabelparam}. 
%	}	
%	
%	\begin{center}
%		\includegraphics[height=0.75\textheight,width=\textwidth]{figure_202405/beyondDNY_cyclicg_delta}
%	\end{center}
%	%The higher the value of $\zeta$, the more impatient the government, the higher the jumpy debt issuance boundary $\underline{b}$, the lower the tax, and the higher the drift of the debt.  The equilibrium debt capacity is independent of the value of $\zeta$. %$\kappa=7> \frac{\rho-g+\xi}{c(\gamma)} \left(\frac{c((r-g)\bar{b}+\gamma)}{\rho-g} +\frac{\xi v(0)}{\rho-g+\xi} \right) = 4.7078$.
%\end{figure}



\section{Concluding Remarks\label{sec:conclusion}}

 {\cite{debortoli2021optimal} showed that Lucas and Stokey's (1983) way of implementing a Ramsey plan for flat-rate taxes on labor   fails for a simple initial government debt structure in which debt is too big. Our paper shows how to implement that Ramsey plan by slightly modifying   Lucas and Stokey's  timing protocol for setting tax rates and by using a ``short route" debt management strategy instead of Lucas and Stokey's debt management plan that uses consols.
 Our implementation  always works, regardless of initial debt levels or whether  the Ramsey plan sets  tax rates above or below  peaks of  Laffer curves.} 


Our  innovations involve two complementary modifications to Lucas and Stokey's analysis. First, we expand    continuation Ramsey planners' authority by allowing them to pre-set a one period-ahead tax rate, while also shrinking their authority by  preventing them  from resetting current-period tax rates. Second, we propose a debt management strategy that leaves longer-term debts unchanged while using only short-term debt  to finance accumulated unpaid government liabilities. This ``short route" debt management policy implements all Ramsey plans and equalizes  Lagrange multipliers that appear in  the constrained optimum problems faced by the  Ramsey planner and all continuation planners. We have explained how equality  of those Lagrange multipliers is decisive evidence for implementability. % \WJ{\bf Tom and Neng: We could include a comparison of timing protocols here, as discussed today.}


%The examples presented in this paper, including scenarios with exponentially decaying debt, periodic high-low debt structures, and cyclical government spending, demonstrate the robustness of our implementation across various initial conditions and government expenditure patterns. In each case, the equivalence of Lagrange multipliers serves as a reliable indicator of successful implementation.

%simplifying the debt management, 
The leading role of  short-term debt in our implementation opens a way to formulate the Ramsey problem  recursively. {By constructing a ``state'' variable based on accumulated  government liabilities $\Pi_i^\ast$  and establishing the invariance of Lagrange multipliers across time, our approach makes it possible to formulate a tractable Bellman equation for continuation planners. Such a recursive representation   facilitates quantitative versions of richer  variants of our model, potentially bridging  gaps between theoretical fiscal policy design and practical implementation.\footnote{\citet{kydland1980dynamic}   described a concise  way of representing  some Ramsey plans  recursively.
That paper is silent about implementation and did nothing to damped their skepticism about the plausibility of Ramsey plans that they expressed in \citet{kydland1977rules}.}   }
%\WJ{This may be unclear to readers who have not read our continuous time paper.} 

Our analysis opens other research possibilities. A useful  extension would be   to incorporate uncertainty.\footnote{Again, see \citet[ftnt.~1]{LucasStokey1983}.}  This would let us  study how stochastic shocks to government spending, productivity, or other economic fundamentals affect optimal tax policies and debt management strategies under limited commitment. The invariance of Lagrange multipliers and associated recursive formulations promise to provide insights into implementation conditions under uncertainty as well.
Further, by computing   limits as the time increment $\Delta$ approaches zero, we can develop continuous-time formulations of our model that  should bring insights and convenient mathematical expressions.\footnote{\citet{JiangSargentWang2024} analyzed  a continuous time analysis that interprets the timing protocol that restricts the current tax rate as a covenant on government debt. We prefer the interpretation in the present  paper. \citeauthor{JiangSargentWang2024} presented  a recursive formulation of the Ramsey problem and an associated HJB equation.}
 

%We may move the discussion about recursive framework here to the last paragraph about future work.}
%  and  facilitate connections with the continuous-time macro-finance literature on optimal government debt management. 
%This approach promises further insights about the role of debt maturity in designing a government fiscal policy.
%
%The central message of our analysis is that short-term debt, previoulsy overlooked in favor of more complex maturity structures, can play a crucial role in implementing optimal fiscal policy. This finding contributes to the ongoing debate about optimal debt management and highlights the importance of carefully designed institutional constraints in achieving desirable policy outcomes.
%
%\NW{Tom: Conclusion is great. Just wonder if this is too long and may invite referees to ask for more? Maybe they'll ask us to develop the recursive structure? We're happy to do these but I just want to raise this, as it is related to our `grand' strategy regarding the LS agenda. What we do for the followup papers and etc... You make the call and we march on.} 

\newpage 



%What's the role of short-term debt in our strategy? %Because the issuance of short-term debt captures the cumulative deficits, it effectively serves as a stock variable, rather than the flow variables for long term debt.
%Note that the issuance of short-term debt satisfies
%
% As a result, the balance of short-term debt, ${b}_{{i-1}, i}\Delta$, is more like a stock variable, rather than other flow variables. 
%
%This ``stock-flow" relationship is exacerbated when the time invertal shrinks. To see this, for a fixed length of time $[0, T]$, letting $\Delta$ decrease to zero, the stock variable $\Pi^*$ at time $T$ must converge to a finite value. Note $b_{-1, i}\Delta$ shrinks to zero because $b_{-1, i}$ is finite and exogenous given. We obtain ${b}_{{i-1}, i}\to \infty$, which suggests that the one-period debt balance is effectively a stock variable. In this regard, our short-term debt only strategy differs from Lucas-Stokey's term debt strategy in which no stock variable is used. 
%
%To illustrate, we take \cite{debortoli2021optimal}'s example and let the number of steps $n$ increases such that $n\times\Delta=1$. We assume $b_{-1, i}=0.5$ for $0\leq i\leq n-1$ and $b_{-1, i}=0$ for $i\geq n$. This assumption of initial debt implies that over the time interval $[0, 1)$ the government faces a total debt payout of $0.5$, without considering discounting. We plot the short-term debt flow that mature at time $t=1$, $b_{n-1, n}$, in Figure \ref{fig:shorttermdebtvsdelta}. It shows the short-term debt flow increases to infinity if time step size shrinks to zero.
%
%\begin{figure}
%	\centering
%	\includegraphics[width=0.5\linewidth]{figure_202405/shorttermdebt_vs_delta}
%	\caption{{\bf Short-term debt flow $b_{n-1, n}$ increases to infinity} where $n=1/\Delta$.}
%	\label{fig:shorttermdebtvsdelta}
%\end{figure}

%\newpage 
%
%\section*{To be put elsewhere when discussing  \cite{LucasStokey1983}.} 
%\NW{Do we need this paragraph here? I can see the value of keeping it here. Just want to raise this issue as skipping this paragraph allows us to move more quickly to our analysis.} % Another option is to have a section 

%After the Ramsey planner chooses   a tax  plan and associated price system,  
%the representative agent's best  responses confirm the associated  competitive equilibrium allocation.

%\NW{a question  similar to the one we raised earlier: Do we need this paragraph here? I can see the value of keeping it here. Just want to raise this issue as skipping this paragraph allows us to move more quickly to our analysis.} % Another option is to have a section 

%Next we formulate a Ramsey plan, then describe the debt management strategy that implements it. 

%\NW{Maybe merge subsection 1.2 with subsection 2.1; and similarly merge subsection 1.3 with subsection 2.2; ?} 



%\newpage 

\newpage
\bigskip
%\footnotesize
%\begin{center}
%\textbf{References}
%\end{center}
%	\singlespacing
%\setlength{\bibsep}{0.3ex}
%\bibliographystyle{chicago}
\bibliography{LSreference}

\pagebreak
\section*{Appendix}
%\appendix{\bf \centering \LARGE \noindent Appendices}
\appendix{}
%\setcounter{page}{1} 
%\setcounter{Theorem}{0} 
%\renewcommand{\thepage}{I-\arabic{page}}
%\setcounter{theorem}{0}
%\renewcommand{\thetheorem}{\Alph{section}.\arabic{theorem}}
\setcounter{equation}{0}
\renewcommand{\theequation}{A-\arabic{equation}}
\setcounter{section}{0}
\renewcommand{\thesection}{\Alph{section}}
\renewcommand{\thesubsection}{\Alph{section}.\arabic{subsection}}
\setcounter{figure}{0}
\renewcommand{\thefigure}{A-\arabic{figure}}
\setcounter{footnote}{0}
\renewcommand{\thefootnote}{A-\arabic{footnote}}
\bigskip%


\section{Technical Details} \label{appendix:A}


We  state four regularity conditions that guarantee that  a Ramsey plan and a continuation Ramsey plan exist. 

\begin{condition} \label{condition1}
For a given utility function $U(C, C+G)$ that is concave and differentiable, 
the function $\left(CU_C + (C+G) U_N - bU_C \right)$ is also concave and differentiable in $C$ when $G>0, b\geq 0$.
\end{condition}


\begin{condition} \label{condition2}
For a given initial debt structure $\{b_{-1, i}\}_{i=0}^\infty$, there exists allocations $\{C_{0, i}\Delta\}_{i=0}^\infty$ that strictly satisfy the IC (\ref{impletconstraint}). 
\end{condition}

\begin{condition} \label{condition2j}
For period $j$ $(j\geq 1)$ continuation planner that inherits a debt structure $\{b_{j-1, i}\}_{i=j}^\infty$ and is subject to new timing protocol given in Definition \ref{defprotocol}, there exists allocations $\{C_{j, i}\Delta\}_{i=j}^\infty$ that strictly satisfy the CIC (\ref{con_IC0}). 
\end{condition}

\begin{condition} \label{condition3}
For a period $j$ $(j\geq 1)$ continuation planner who inherits a debt structure \eqref{btermuntouched}-\eqref{bPi} and is subject to new timing protocol given in Definition \ref{defprotocol}, there exists allocations $\{C_{j, i}\Delta\}_{j=i}^\infty$ that strictly satisfy the CIC (\ref{con_IC}). 
\end{condition}
Condition \ref{condition1} ensures both Ramsey problem and continuation Ramsey problem being concave problems. Conditions \ref{condition2}, \ref{condition2j}, and \ref{condition3} are Slater's conditions for the Ramsey problem and continuation Ramsey problem, respectively.
For example, the economy studied in \citet{debortoli2021optimal} satisfies these four regularity conditions when initial debt is not too high. 


\begin{proof}[Proof of Proposition \ref{proplag}.]


Note that the first condition states the Ramsey problem is a concave problem, and the second condition is a Slater's condition that ensures the strong duality for this concave problem (i.e., \citealp{boyd2004convex}).
Then, for the Ramsey plan $\{C_{0, i}^*\}_{i=0}^\infty$, there must exist a $\Phi_0^*$ such that the following
Karush-Kuhn-Tucker condition holds:
% Karush-€“Kuhn€-Tucker condition holds:
\begin{align*}
& C_{0, i}^* = \argmax_{C_{0, i}}  \left[ U(C_{0, i}, C_{0, i}+G_{i})+ \Phi_0^*\left( C_{0, i}U_{C, i}+(C_{0, i}+G_{i})U_{N, i}-b_{-1, i} U_{C, i} \right) \right] \\
& 	\Phi_0^* \cdot \left( \sum_{i=0}^\infty \, e^{-\rho i\Delta} \left[ C_{0, i}^*U_{C, i}+(C_{0, i}^*+G_{i})U_{N, i} - b_{-1, i} U_{C, i} \right]\cdot\Delta \right) =0 \\
& 	\Phi_0^*\geq 0\\
&	\sum_{i=0}^\infty \, e^{-\rho i\Delta} \left[ C_{0, i}^*U_{C, i}+(C_{0, i}^*+G_{i})U_{N, i} - b_{-1, i} U_{C, i} \right]\cdot\Delta \geq 0
\end{align*}
Hence,  optimal consumption satisfies $C_{0, i}^*=C(\Phi_0^*; b_{-1, i}, G_i)$, where $C(\cdot; \cdot, \cdot)$ is defined in \eqref{Csoln}.

Lastly, we show that the optimal Lagrange multiplier must be positive, $\Phi_0^*>0$,  and thus
the implementability condition \eqref{eqPhi} must hold with equality. We prove it by the method of contradiction. Assume $\Phi_0^*=0$. Then optimality requires the optimal consumption $C_{0, i}^*$ must satisfy
\begin{equation*}
U_{C}(C_{0, i}^*, C_{0, i}^*+G_i) + U_{N}(C_{0, i}^*, C_{0, i}^*+G_i) = 0\,
\end{equation*}
and thus the utility $U(C_{0, i}^*, C_{0, i}^*+G_i)$ achieves its global maximum. Contradiction then arises because the implementability condition \eqref{impletconstraint} does not hold because
\begin{eqnarray*}
\sum_{i=0}^\infty \, e^{-\rho i\Delta} \left[ C_{0, i}^*U_{C, i}+(C_{0, i}^*+G_{i})U_{N, i} - b_{-1, i} U_{C, i} \right]\cdot\Delta = - \sum_{i=0}^\infty \, \left(G_i+b_{-1, i} \right)U_{C, i}\Delta <0  \,.
\end{eqnarray*}

\end{proof}



\begin{proof}[Proof of Proposition \ref{proplag2}]

Given the inherited short-term balance $ b_{j-1, j}\Delta = \Pi_j^* + b_{-1,j} \Delta $ and the tail of the initial term debt $\{b_{-1, i}\Delta; i\geq j+1\}$, the period ${i}$ continuation Ramsey planner's Lagrangian  is given by \eqref{Lagran_t}, and the optimal consumptions for period $i\geq {j+1}$, given a Lagrange multiplier $\Phi_j$, satisfy \eqref{consum_t}. 



We then show that optimality requires $\Phi_j^*>0$ and the continuation implementability condition binds, $I_j(\Phi_j^*)=0$, where $I_j(\, \cdot \,)$ is given by \eqref{con_IC2}. We prove this by the method of contradiction. If $\Phi_j^*=0$, the optimal consumption $C_{j, i}^*$ defined in \eqref{consum_t} must satisfy 
\begin{equation}
U_{C}(C_{j, i}^*, C_{j, i}^*+G_i)+U_N(C_{j, i}^*, C_{j, i}^*+G_i)=0, \quad i\geq j+1\,,
\end{equation} 
and $U(C_{j, i}^*, C_{j, i}^*+G_i)$ achieves the global maximum. On the other hand, using \eqref{Pistar}, the continuation implementatibity condition \eqref{con_IC}, if holds, implies that
\begin{align*}
& \sum_{i=j+1}^{\infty} e^{-\rho(i-j) \Delta} \left(  C_{j, i}^*U_{C, i} + (C_{j, i}^*+G_{i})U_{N, i} -b_{-1, i} U_{C, i} \right) \Delta \nonumber\\
& \quad \quad + \left (C_{0, j}^*U_{C, j} + (C_{0, j}^*+G_{j})U_{N, j} - b_{-1,j} U_{C, j} \right)\Delta - \Pi_{{j}}^*U_{C, j}  \\
& =  \sum_{i=j+1}^{\infty} e^{-\rho(i-j) \Delta} \left(  C_{j, i}^*U_{C, i} + (C_{j, i}^*+G_{i})U_{N, i} -b_{-1, i} U_{C, i} \right) \Delta \nonumber\\
& \quad \quad +	 \sum_{i=0}^{j} e^{-\rho(i-j) \Delta} \left(  C_{0, i}^*U_{C, i} + (C_{0, i}^*+G_{i})U_{N, i} -b_{-1, i} U_{C, i} \right) \Delta \geq 0  \,.
\end{align*}
We then obtain an allocation $\{C_{0, 0}^*, \cdots, C_{0, j}^*, C_{j, j+1}^*, C_{j, j+2}^*, \cdots,\}$ for the Ramsey planner than satisfies the implementability condition \eqref{impletconstraint}. Note that $U(C_{j, i}^*, C_{j, i}^*+G_i)>U(C_{0, i}^*, C_{0, i}^*+G_i); i\geq j+1$ because $U(C_{j, i}^*, C_{j, i}^*+G_i)$ achieves the global maximum. As a result, the households utility \eqref{utility} under this constructed allocation is larger than that given in Proposition \ref{proplag}. Contradiction arises. 

\end{proof}

\begin{proof}[Proof of Lemma \ref{lemma1}]
Let  $\Phi_j = \Phi_0^*$,  then $C_{j, i}=C(\Phi_j; b_{-1, i}, G_{i})=C_{0, i}^*$ for $i\geq j+1$. Substituting into \eqref{con_IC2} and using \eqref{Pistar} , we obtain
\begin{align*}
I_j(\Phi_0^*) 
%	& = & \sum_{j=i+1}^{\infty} e^{-\rho(j-i) \Delta} \left(  C_{i, j}U_{C, j} + (C_{i, j}+G_{j})U_{N, j} -b_{-1, j} U_{C, j} \right) \Delta \label{IC01}\\
%	& & \quad \quad + \left[\left(  C_{i, i}U_{C, i} + (C_{i, i}+G_{j})U_{N, i}  \right) \Delta- \left(b_{-1, i}\Delta +\Pi_{{i}}^* \right)U_{C, i}\right] \label{IC02}\\
=&   \sum_{i=j+1}^{\infty} e^{-\rho(i-j) \Delta} \left(  C_{0, i}^*U_{C, i} + (C_{0, i}^*+G_{i})U_{N, i} -b_{-1, i} U_{C, i} \right) \Delta \nonumber \\
& \quad\quad + \left (C_{0, j}^*U_{C, j} + (C_{0, j}^*+G_{j})U_{N, j} - b_{-1,j} U_{C, j} \right)\Delta -\Pi_{{j}}^*U_{C, j} \,.\\
= & \sum_{i=j+1}^{\infty} e^{-\rho(i-j) \Delta} \left(  C_{0, i}^*U_{C, i} + (C_{0, i}^*+G_{i})U_{N, i} -b_{-1, i} U_{C, i} \right) \Delta \nonumber \\
& \quad\quad + \sum_{i=0}^{j} e^{-\rho(i-j)\Delta}\left(  C_{0, i}^*U_{C, i} + (C_{0, i}^*+G_{i})U_{N, i} -b_{-1, i} U_{C, i} \right) \Delta \\
= & e^{j\Delta } I_0(\Phi_0^*) =0\,.
\end{align*}
When $I_j(\, \cdot, \,)=0$ has a unique root, the period $j$ planner  chooses $\Phi_j^*=\Phi_0^*$ and thereby confirms the Ramsey plan.
\end{proof}



\subsection{Other Debt-Management Policies For General Initial Debt Structure}\label{appsec:otherpolicies}

\WJ{We focus on period 0 to period 1 with two assumptions:
\begin{itemize}
\item DNY general utility function: $\log C - \eta \frac{N^\gamma}{\gamma}$
\item General initial debt: $\{ b_{-1, i}\Delta \}_{i=0}^\infty$
\end{itemize}
}

%\begin{equation} 
%\Pi_j^* = \sum_{i=0}^{j-1} \frac{q_{0, i}^*}{q_{0, j}^*} \left(b_{-1, i}-S_{0, i}^*\right)\Delta \quad \text{with} \quad \Pi_0^*= 0\,.
%\end{equation} 
%In particular, we have
%\begin{equation} \label{Pistar1}
%\Pi_1^* = \frac{1}{q_{0, 1}^*} \left(b_{-1, 0}-S_{0, 0}^*\right)\Delta  = \frac{e^{\rho \Delta}C_{0, 1}^*}{C_{0, 0}^*}\left(b_{-1, 0}-C_{0, 0}^*\left( 1-\eta(C_{0, 0}^*+G_1)^\gamma\right) \right)\Delta  \,.
%\end{equation} 



If a Ramsey plan exists, our  short-term debt management strategy  implements  it. %a Ramsey plan for  all initial government expenditure and initial debt sequences  for which  which an optimal  plan exists. 
\WJ{ For some initial debt structures,  but not for others,   debt management policies that issue a mixture of  short and some  longer term debts can also implement a Ramsey plan. }  
%However, other initial debt structures,   the initial  debt structure in order for  an  implementation that uses  longer-term debt to work. In this sense,  using term debt makes implementation more difficult.} %In other words, provided that a Ramsey plan exists, using short-term debt to implement a Ramsey plan always works, but using a combination of short- and long-term debts does not necessarily work. 
To understand this claim,  suppose that a Ramsey planner  finances a fraction $\alpha\in [0, 1]$ of accumulated  liabilities  up to period $1$, $\Pi_1^*$, by issuing one-period debt %$b_{0, 1}\Delta $ 
\begin{equation} \label{b01_n1}
	b_{0, 1}\Delta = \alpha \Pi_1^*+b_{-1, 1}\Delta\, , 
\end{equation}
  and  finances   fraction $(1-\alpha)$ of $\Pi_1^*$ by issuing longer-term debts, $\{b_{0, i}\}_{i\geq 2}$, where
\begin{equation} \label{Pistar1}
\Pi_1^* = \frac{1}{q_{0, 1}^*} \left(b_{-1, 0}-S_{0, 0}^*\right)\Delta  = \frac{e^{\rho \Delta}C_{0, 1}^*}{C_{0, 0}^*}\left(b_{-1, 0}-C_{0, 0}^*\left( 1-\eta(C_{0, 0}^*+G_1)^\gamma\right) \right)\Delta  \,.
\end{equation}    
According to FOC for the continuation planner, the continuation Lagrange multiplier $\Phi_1$ must satisfy
\begin{align*}
	\frac{1}{C_{1, i}} -\eta (C_{1, i}+G_i)^{\gamma-1} +\Phi_1\left( \frac{b_{0, i}}{C_{1, i}^2} -\eta \gamma(C_{1, i}+G_i)^{\gamma-1} \right) =0, i\geq 2
\end{align*}
Confirming the Ramsey plan requires $C_{1, i}=C_{0, i}^*$ for $i\geq 2$. As a result, $b_{0, i}$ must satisfy
\begin{equation}\label{b0i_n1}
b_{0, i}\Delta = \frac{C_{0, i}^*}{\Phi_1} \left[ -1+\eta \left( C_{0, i}^* + G_i  \right)^{\gamma-1}C_{0, i}^* + \Phi_1 \eta \gamma \left( C_{0, i}^* + G_i  \right)^{\gamma-1}C_{0, i}^*  \right]\Delta
\end{equation}

%\begin{equation}\label{b0i_n1}
%\frac{b_{0, i} } {C_{0, i}^*}= \frac{1}{\Phi_1} \left[ -1+\eta \left( C_{0, i}^* + G_i  \right)^{\gamma-1}C_{0, i}^*\right] +  \eta \gamma \left( C_{0, i}^* + G_i  \right)^{\gamma-1}C_{0, i}^* \,.
%\end{equation}


\begin{proposition}
%Suppose that the conditions of Theorem \ref{thm1} hold, and assume that $\Pi_1^*\geq 0$.
Suppose that the conditions of Theorem \ref{thm1} hold, and assume that $\Pi_1^*\geq 0$.
For $1\geq \alpha_1>\alpha_2\geq 0$, if a Ramsey plan is implemented at period $1$ for debt strategy \eqref{b01_n1}-\eqref{b0i_n1} with $\alpha=\alpha_2$, it is also implemented for debt strategy \eqref{b01_n1}-\eqref{b0i_n1} with $\alpha=\alpha_1$.
\end{proposition}

\begin{proof}
	Substituting \eqref{b0i_n1} into  CIC \eqref{IC_dny2} and solving for $\Phi_1$, we obtain the constructed Lagrangian $\widehat{\Phi}_1$:
	\begin{align}\label{hatPhi1}
	\widehat{\Phi}_1 = \frac{  - \sum_{i=2}^\infty e^{-\rho (i-1)\Delta}\left[ 1-\eta \left( C_{0, i}^* + G_i  \right)^{\gamma-1}C_{0, i}^*\right]      }{ \left(1-\eta(C_{0, 1}^*+G)^\gamma-\frac{b_{0, 1}}{C_{0, 1}^*}\right)+\sum_{i=2}^\infty e^{-\rho (i-1)\Delta}\left( 1-\eta(C_{0, i}^*+G)^\gamma  -\eta \gamma \left( C_{0, i}^* + G_i  \right)^{\gamma-1}C_{0, i}^*  \right)   }
	\end{align}
As pointed out by \cite{debortoli2021optimal}, a Ramsey plan is implemented if the constructed Lagrangian  $\widehat{\Phi}_1$ is positive. 
Note that, when $\Pi_1^*\geq 0$, the numerator of \eqref{hatPhi1} is always positive. This is because that
\begin{eqnarray*}
& & - \sum_{i=2}^\infty e^{-\rho (i-1)\Delta}\left[ 1-\eta \left( C_{0, i}^* + G_i  \right)^{\gamma-1}C_{0, i}^*\right] \\
& <  &  - \sum_{i=2}^\infty e^{-\rho (i-1)\Delta}\left[ 1-\eta \left( C_{0, i}^* + G_i  \right)^{\gamma}-\frac{b_{-1, i}}{C_{0, i}^*}\right] =- \frac{\Pi_1^*}{C_{0, 1}^*}\leq 0 \,,
\end{eqnarray*}
where the equality holds due to the IC \eqref{impletconstraint}. On the other hand, using \eqref{b01_n1}, the denominator of\eqref{hatPhi1} as a function of $\alpha$, $D(\alpha)$, satisfies
	\begin{eqnarray*}
D(\alpha)  & = &   \left(1-\eta(C_{0, 1}^*+G)^\gamma-\frac{\alpha \Pi_1^*+b_{-1, 1}\Delta}{C_{0, 1}^*}\right) \\
	& & +\sum_{i=2}^\infty e^{-\rho (i-1)\Delta}\left( 1-\eta(C_{0, i}^*+G)^\gamma  -\eta \gamma \left( C_{0, i}^* + G_i  \right)^{\gamma-1}C_{0, i}^*  \right)   \,,
\end{eqnarray*}
and $D(\alpha)$ is decreasing in $\alpha$. 

For debt strategy \eqref{b01_n1}-\eqref{b0i_n1} with $\alpha=\alpha_2$, if Ramsey plan is implemented, the associated constructed Lagrangian $\widehat\Phi_1$ is positive and we have $D(\alpha_2)<0$. Because $D(\alpha_1)<D(\alpha_2)<0$ for $\alpha_1>\alpha_2$, the constructed Lagrangian $\widehat \Phi_1$ for $\alpha=\alpha_1$ is positive and the Ramsey plan for $\alpha_1$ is also implemented. 
\end{proof}


\begin{proposition}
Suppose that the conditions of Theorem \ref{thm1} hold, and assume that $\Pi_1^*> 0$. There exists $\alpha^*\in[0, 1)$, such that a Ramsey plan is implemented at period $1$ for debt strategy \eqref{b01_n1}-\eqref{b0i_n1} with any $\alpha\in (\alpha^*, 1]$.
\end{proposition}

\begin{proof}

Substituting \eqref{b0i_n1} into  CIC \eqref{IC_dny2} and solving for $\Phi_1$, we obtain the constructed Lagrangian $\widehat{\Phi}_1$:
	\begin{align}\label{hatPhi11}
	\widehat{\Phi}_1 = \frac{  - \sum_{i=2}^\infty e^{-\rho (i-1)\Delta}\left[ 1-\eta \left( C_{0, i}^* + G_i  \right)^{\gamma-1}C_{0, i}^*\right]      }{ \left(1-\eta(C_{0, 1}^*+G)^\gamma-\frac{b_{0, 1}}{C_{0, 1}^*}\right)+\sum_{i=2}^\infty e^{-\rho (i-1)\Delta}\left( 1-\eta(C_{0, i}^*+G)^\gamma  -\eta \gamma \left( C_{0, i}^* + G_i  \right)^{\gamma-1}C_{0, i}^*  \right)   }
	\end{align}
	As pointed out by \cite{debortoli2021optimal}, a Ramsey plan is implemented if the constructed Lagrangian  $\widehat{\Phi}_1$ is positive. We then verify there exist $\alpha^*\in[0, 1)$ such that $\widehat \Phi_1>0$ for $\alpha^*\in (\alpha^*, 1] $.
	
	Note that, when $\Pi_1^*\geq 0$, the numerator of \eqref{hatPhi1} is always positive. This is because that
	\begin{eqnarray*}
		& & \sum_{i=2}^\infty e^{-\rho (i-1)\Delta}\left[ 1-\eta \left( C_{0, i}^* + G_i  \right)^{\gamma-1}C_{0, i}^*\right] \\
		& >  &   \sum_{i=2}^\infty e^{-\rho (i-1)\Delta}\left[ 1-\eta \left( C_{0, i}^* + G_i  \right)^{\gamma}-\frac{b_{-1, i}}{C_{0, i}^*}\right] = \frac{\Pi_1^*}{C_{0, 1}^*}\geq 0 \,,
	\end{eqnarray*}
	where the equality holds due to the IC \eqref{impletconstraint}. On the other hand, using \eqref{b01_n1}, the denominator of\eqref{hatPhi1} as a function of $\alpha$, $D(\alpha)$, satisfies
	\begin{eqnarray*}
		D(\alpha)  & = &   \left(1-\eta(C_{0, 1}^*+G)^\gamma-\frac{\alpha \Pi_1^*+b_{-1, 1}\Delta}{C_{0, 1}^*}\right) \\
		& & +\sum_{i=2}^\infty e^{-\rho (i-1)\Delta}\left( 1-\eta(C_{0, i}^*+G)^\gamma  -\eta \gamma \left( C_{0, i}^* + G_i  \right)^{\gamma-1}C_{0, i}^*  \right)   \,,
	\end{eqnarray*}
	and $D(\alpha)$ is strictly decreasing in $\alpha$. 	According to Theorem \ref{thm1}, when $\alpha = 1$, $D(1)<0$ and $\widehat \Phi_1 =\Phi_0^*>0$. Consider two cases: 1) $D(0)\geq 0$, and 2) $D(0)<0$. For case 1), there exists $1>\alpha^*\geq 0$ such that $D(\alpha^*)=0$ and therefore, for any $\alpha> \alpha^*$, $D(\alpha)<0$  and resulted $\widehat \Phi_1>0$. For case 2), $\alpha^* =0$, and, for any $\alpha> \alpha^*=0$, $D(\alpha)<0$  and $\widehat \Phi_1>0$.
\end{proof}




%
%\begin{equation}
%	\left(1-\eta(C_{0, 1}^*+G)^\gamma-\frac{b_{0, 1}}{C_{0, 1}^*}\right)\Delta+\sum_{i=2}^\infty e^{-\rho (i-1)\Delta}\left( 1-\eta(C_{0, i}^*+G)^\gamma-\frac{b_{0, i}}{C_{0, i}^*} \right)\Delta =0 \, . \label{IC_dny21}
%\end{equation}
%
%\begin{eqnarray*}
%&&	\left(1-\eta(C_{0, 1}^*+G)^\gamma-\frac{b_{0, 1}}{C_{0, 1}^*}\right)\Delta+\sum_{i=2}^\infty e^{-\rho (i-1)\Delta}\left( 1-\eta(C_{0, i}^*+G)^\gamma \right)\Delta \\
%& = & \sum_{i=2}^\infty e^{-\rho (i-1)\Delta}\left( \frac{b_{0, i}}{C_{0, i}^*} \right)\Delta  \\
%&	 = &  \sum_{i=2}^\infty e^{-\rho (i-1)\Delta}\left(   \frac{1}{\Phi_1} \left[ -1+\eta \left( C_{0, i}^* + G_i  \right)^{\gamma-1}C_{0, i}^*\right] +  \eta \gamma \left( C_{0, i}^* + G_i  \right)^{\gamma-1}C_{0, i}^*          \right)\Delta 
%\end{eqnarray*}

%\begin{equation}
%\widehat{\Phi}_1 = \frac{   \sum_{i=2}^\infty e^{-\rho (i-1)\Delta}\left[ -1+\eta \left( C_{0, i}^* + G_i  \right)^{\gamma-1}C_{0, i}^*\right]     \Delta }{ \left(1-\eta(C_{0, 1}^*+G)^\gamma-\frac{b_{0, 1}}{C_{0, 1}^*}\right)\Delta+\sum_{i=2}^\infty e^{-\rho (i-1)\Delta}\left( 1-\eta(C_{0, i}^*+G)^\gamma  -\eta \gamma \left( C_{0, i}^* + G_i  \right)^{\gamma-1}C_{0, i}^*  \right)\Delta   }
%\end{equation}




\WJ{Below is the discussion when initial debt is the special one in DNY.}


  \begin{equation} \label{b02_n1}
	b_{0, i}\Delta = (1-\alpha)(e^{\rho\Delta}-1)\Pi_1^*+b_{-1, i}\Delta\,,\,  i\geq 2\,.
\end{equation} 


Use  the continuation planner's optimality condition \eqref{FOC53} to form the  `constructed' Lagrange multiplier: 
\begin{equation} \label{hatPhi1_n1}
	\widehat{\Phi}_1 = \frac{1-C_1^*}{((1-\alpha)e^{\rho\Delta}+1)C_1^*-(1-\alpha)e^{\rho\Delta}(1-G)} \,.
\end{equation}
%Consider two special cases. 
If $\alpha = 1$, then the debt management   strategy corresponds to our  short-term debt only strategy, %and $\widehat{\Phi}_1$ given in \eqref{hatPhi1_n}; 
 since $b_{-1,i} = 0$ for $i \geq 2$,  \eqref{hatPhi1_n} simplifies  to \eqref{hatPhi1_our}. 
If  $\alpha = 1 - e^{-\rho \Delta}$, then the debt management   strategy corresponds to  Lucas and Stokey's consol-based implementation strategy in which $b_{0,1} = b_{0,i}$ for $i \geq 2$ and $\widehat{\Phi}_1$ given in \eqref{hatPhi1_n} simplifies to  \eqref{hatPhi1LS}. 


%Note that when $\alpha=1$, it corresponds to our implementation and $\widehat{\Phi}_1$ in \eqref{hatPhi1_n} boils down to \eqref{hatPhi1_our} (because $b_{-1, i}=0$ for $i\geq 2$.) When $\alpha=1-e^{-\rho\Delta}$, we have $b_{0, 1}=b_{0, i}, i\geq 2$, and the $\widehat{\Phi}_1$ in \eqref{hatPhi1_n} boils down to that in Lucas-Stokey implementation given in \eqref{hatPhi1LS}.


%To ensure $\widehat{\Phi}_1>0$, a condition that is necessary for the implementation method works, we require that
%$((1-\alpha)e^{\rho\Delta}+1)C_1^*-(1-\alpha)e^{\rho\Delta}(1-G)>0$, that is
%\begin{equation*}
%	C_1^*>\frac{(1-\alpha)e^{\rho\Delta}(1-G)}{(1-\alpha)e^{\rho\Delta}+1} \,.
%\end{equation*}
%According to \cite{debortoli2021optimal}, $C_1^*$ is an decreasing equation of $b$. As a result, for each $\alpha \leq 1$, there exist a $b^*(\alpha)$, which is an increasing equation of $\alpha$, such that when $b=b^*(\alpha)$, $C_1^*=\frac{(1-\alpha)e^{\rho\Delta}(1-G)}{(1-\alpha)e^{\rho\Delta}+1}$. Moreover, for given $\alpha$, the implementation method only works for $b<b^*(\alpha)$ where $\widehat{\Phi}_1>0$.



Sometimes using  both short- and long-term debt fails to implement a Ramsey plan.  
A necessary condition for a successful implementation is 
  $\widehat{\Phi}_1 > 0$. This  condition requires that 
$((1 - \alpha)e^{\rho\Delta} + 1)C_1^* - (1 - \alpha)e^{\rho\Delta}(1 - G) > 0$ or equivalently 
\begin{equation} \label{C1sub531} 
	C_1^* > \frac{(1 - \alpha)e^{\rho\Delta}(1 - G)}{(1 - \alpha)e^{\rho\Delta} + 1} \,.
\end{equation}
Note that $C_1^*$ is strictly decreasing in $b$ (see \citeauthor{debortoli2021optimal}) and  that the right side of (\ref{C1sub53}) is increasing in $\alpha$. Consequently, there exists a  threshold $\tilde{b}(\alpha)$ that  is an increasing function of $\alpha$  for all $\alpha \leq 1$ such that 
inequality (\ref{C1sub53}) prevails if $b < \tilde{b}(\alpha)$. %we have $C_1^* = \frac{(1 - \alpha)e^{\rho\Delta}(1 - G)}{(1 - \alpha)e^{\rho\Delta} + 1}$.
Thus, for a given $\alpha$, $\widehat{\Phi}_1 > 0$ and  a mixed short- and long-term debt issuance policy  (\ref{b01_n})-(\ref{b02_n}) implements a Ramsey plan only 
when $b < \tilde{b}(\alpha)$.\footnote{For the two polar special cases discussed in the text, $\tilde{b}(1)=\bar b$ and $\tilde{b}(1-e^{-\rho\Delta})=b^*$.}


For  both Lucas and Stokey's setup and our setup, the   Lagrangian for a  period-$1$ continuation planner is %is then given by
\begin{align}\label{L11}
	{\mathcal L}^{}_1 & =  \sum_{i=1}^{\infty}  e^{-\rho (i-1)\Delta}\left[ \log C_{1, i}-  (C_{1, i}+G_{i})+ \Phi_1^{}\left( 1-(C_{1, i}+G_{i})-\frac{b_{0, i}^{}}{C_{1, i}}\right)  \right] \Delta \, , %\nonumber \\%- \Psi_0(b_{0, t}^{} U_C- b_{0-, t} U_C)	\right]dt \\
	%&  \quad \quad \;  \, .
\end{align}
where  $\{b_{0, i}\}_{i\geq 1}$ is the   debt structure that the Ramsey planner passes on to the period 1  continuation planner. % after %pays $b_{-1,0}$. 
The period 1  continuation planner's first-order condition for $C_{1,i}$ is %Ramsey consumption $C_{1, i}$ for $i\geq 1$ satisfies 
\begin{align}\label{FOC531}
	\frac{1}{C_{1, i}} -1 +\Phi_1\left( \frac{b_{0, i}}{C_{1, i}^2} -1 \right) =0.
\end{align}


Confronting  debt structure  \eqref{restructure_debt_dny},  
the period 1 continuation planner chooses  $\{ C_{1, i}, i\geq 2\}$ to maximize %the household's utility from period 1 onward: %(\ref{contplannerDNYexample}) 
%\begin{equation}\label{contplannerDNYexample} 
$$	\sum_{i=1}^{\infty} e^{-\rho (i-1)\Delta} \left( \log C_{1, i}- \eta \frac{N_{1, i}^\gamma}{\gamma} \right)\Delta \, $$ 
%\end{equation}
%$\sum_{i=1}^{\infty} e^{-\rho (i-1)\Delta} \left( \log C_{1, i}- N_{1, i} \right)\Delta$,
subject to the CIC: %continuation implementability condition: 
\begin{equation}
	\left(1-\eta(C_{1, 1}+G)^\gamma-\frac{b_{0, 1}}{C_{1, 1}}\right)\Delta+\sum_{i=2}^\infty e^{-\rho (i-1)\Delta}\left( 1-\eta(C_{1, i}+G)^\gamma-\frac{b_{0, i}}{C_{1, i}} \right)\Delta =0 \, . \label{IC_dny21}
\end{equation}


\end{document}

\section{Additional Discussions}



%\WJ{$i$-the gov't is allowed to choose anything beyond $i$; it is sufficient to let the gov't to make one-period-ahead policy. That solves the problem. in other words, let each gov't choose its debt issuance and also the supporting fiscal policy. rolling over then works.} 
%
%
%%	\begin{figure}[h!]
%	%	\caption{Illustration of Short-term Debt Only Strategy.}
%	%	\bigskip
%	%	\centering
%	%	\begin{tikzpicture}[scale=1]
%		%		
%		%		\draw [thick,->] (-2.5,0) -- (9.7,0);
%		%		%\draw [thick,-] (-0.19,0.1) -- (-0.19,0);
%		%		
%		%		%\node[align=center] at (4.1,1.3) {\bf  \ZL{ \scriptsize Step 2}};
%		%%		\node[align=center] at (4.3, 1.3) {\bf  \NW{\scriptsize Household}};
%		%%		\node[align=center] at (4.3, 0.9) {\bf  \NW{\scriptsize chooses $\{C_{0, j}\Delta, N_{0, j}\Delta\}_{j=0}^{\infty}$}};
%		%		
%		%		\draw [thick,->] (1,0.1) -- (1,0.5);
%		%		\draw [thick,->] (4.5,0.1) -- (4.5,0.5);
%		%	%	\draw [thick,->] (6,0.1) -- (6,0.5);
%		%		
%		%		\node[align=center] at (1, 1.4) {\bf  \NW{\footnotesize $\widehat{b}_{i-1, i}\Delta =\Pi_{{i}}^* + \WJ{b_{-1, i}\Delta}$}};
%		%		\node[align=center] at (-0.1, 1.9) {\bf  \WJ{\footnotesize ${b}_{-1, i+1}\Delta$}};
%		%		\node[align=center] at (0.5, 2.4) {\bf  {\footnotesize ${b}_{-1, j}\Delta; j\geq i+2$}};
%		%		
%		%		\node[align=center] at (5.83, 0.8) {\bf  \NW{\footnotesize $ = \Pi_{i+1}^* + \WJ{{b}_{-1, i+1}\Delta}$}};
%		%		\node[align=center] at (6.9, 1.4) {\bf  \NW{\footnotesize $\widehat{b}_{i, i+1} \Delta = \left(\widehat{b}_{i-1, i}\Delta  -S_{i-1, i}\Delta\right)e^{-r_{0, i}^*\Delta} + \WJ{{b}_{-1, i+1}\Delta}$}};
%		%		\node[align=center] at (3.8, 1.9) {\bf  \WJ{\footnotesize ${b}_{-1, i+2}\Delta$}};
%		%		\node[align=center] at (4.4, 2.5) {\bf  {\footnotesize ${b}_{-1, j}\Delta; j\geq i+1$}};
%		%		
%		%		
%		%%		\node[align=center] at (-1.2, 1.9) {\bf \WJ{\scriptsize $\{ b_{-1, j}\Delta\}_{j=0}^{\infty},\; \{G_{j}\Delta\}_{j=0}^{\infty}$  }};
%		%%		\node[align=center] at (-1.2, 1.3) {\bf \TS{\scriptsize Ramsey planner  }};
%		%%		\node[align=center] at (-1.2, 0.9) {\bf \TS{\scriptsize chooses $\{\tau_{0, j}\Delta\}_{j=0}^{\infty},\; \{\hat{b}_{0, j}\Delta\}_{j=1}^{\infty}$  }};
%		%		%\node[align=center] at (-1, 1.3) {\bf  \NW{ \scriptsize  Step 3}};
%		%		
%		%		%\draw [thick,-] (5.20,0.23) -- (5.20,0);
%		%	%	\draw [thick,decorate,decoration={brace,amplitude=12pt,raise=4pt}] (-1.15,0) -- (9.5,0);
%		%		
%		%		
%		%		%\node[align=center] at (8.9, 1.3) {\bf   \YC{\scriptsize  Step 1}};
%		%		%	\node[align=center] at (8.9, 0.9) {\bf   \TS{\scriptsize Agent pays $M_{\tau}$  }};
%		%		%	\draw [thick,-] (9.4,0.23) -- (9.4,0);
%		%		%\draw [thick,decorate,decoration={brace,amplitude=12pt,raise=4pt}] (5.25,0) -- (9.35,0);
%		%		
%		%		
%		%		%\draw [thick,decorate,decoration={brace,amplitude=12pt,raise=4pt}] (9.45,0) -- (10.75,0);
%		%		
%		%		\draw[fill,color=blue] (-2.5,0.005) circle (.05);	
%		%		\draw[fill,color=blue] (1,0.005) circle (.05);	
%		%		\draw[fill,color=blue] (4.5, 0.005) circle (.05);	
%		%	%	\draw[fill,color=blue] (6, 0.005) circle (.05);	
%		%		%\draw[fill,color=black] (-0.2,0.005) circle (.05);
%		%		
%		%		%\node[align=center] at (5.2,-0.3) {\footnotesize $ \tau^d$};
%		%		%\node[align=center] at (4.2,-0.3) {$\tau_F$};
%		%		\node[align=center] at (-2.5,-0.3) {\footnotesize $0 $};
%		%		\node[align=center] at (1,-0.3) {\footnotesize $t_{i} $};
%		%		\node[align=center] at (4.5,-0.3) {\footnotesize $t_{i+1} $};
%		%	%	\node[align=center] at (6,-0.3) {\footnotesize $t_{i+1} $};
%		%		\node[align=center] at (10.1,0) {\;\footnotesize  \; time};
%		%		\node[align=center] at (9.7,-0.3) {\footnotesize $\infty $};
%		%		
%		%	\end{tikzpicture}
%	%\end{figure}
%	
%	
%	%Meanwhile, the time $t_i$ government also inherits the tail of the initial term debt structure $\{b_{0, t_j}; j=i, \cdots\}$. 
%	
%	
%	\NW{Time $t_{i-1}$ gov't commits time $t_i$ gov't to choose tax rate $\tau_{i, i}=\tau_{0,i}^\ast$.} 
%	
%	
%	\NW{Using short-term debt, which is a stock variable, the incentive to renege is large. this is why we need local commitment condition. it is in the interest of the gov't to commit in order to issue debt and implement Ramsey plan. In contrast, using term debt, the benefit of reneging locally is relatively speaking not large. therefore, it is possible that the government may get away (for certain settings where debt is not too high) with term debt restructuring... } 
%	
%	
%	\NW{Add a discussion and switch $b_{i-1,i} \Delta$ to $B_i$. 
%	} 
%	
%	\NW{In discrete time, one-period debt is ``stock'' and multi-period debt is ``flow.'' This prepares us for why continuous time formulation is more natural and clearer. In discrete time, the ``stock'' nature of one-period debt is hidden in the original LS formulation. } %This is why we can combine } 
%
%\NW{In discrete time, when we shrink $\Delta$, $b_{i-1,i}$ has to go to infinity in order for $b_{i-1,i} \Delta$  to be finite.} 
%
%\NW{To make sure the gov't has no incentives to deviate from the Ramsey plan, it is sufficient for the gov't not to renege locally (over one period) on its debt due. So long as locally (one period) the gov't honors Ramsey, globally (over the entire sequence), the gov't is able to implement Ramsey.} 
%
%\NW{This is in essence how to resolve debt overhang.} 
%
%\NW{Delta function and Dirac delta} 
%





\paragraph{Necessity of the covenant condition}

We show that it is the combination of short-term debt-based management strategy and the one-period-ahead tax covenant that implements the Ramsey plan.
The reasoning is as follows. First, without covenant condition our short-term debt only strategy  cannot implement a Ramsey plan for what ever initial debt. This is because, as pointed by Lucas-Stokey, the debt restructure method \eqref{LS_im} is unique if one wants to implement a Ramsey plan. Second, applying covenant condition to Lucas-Stokey's debt management strategy cannot implement a Ramsey plan neither when initial debt is too high. To illustrate the second point, we use \cite{debortoli2021optimal} economy with $\eta=\gamma=1$ for example. In Lucas-Stokey method, period $0$ government leaves period $1$ government a restructured debt \eqref{LS_im}. If covenant condition is imposed, period $1$ government must commit $C_{1, 1}=C_1^*$ but free to choose $C_{1, i}=C_2$, for $i\geq 2$. 
As a result, the continuation implementability conditoin follows $\hat I_1(\Phi_1) = 0$ where
\begin{eqnarray}\label{IC_1_dny}
\widehat{I}_1({\Phi}_1) %& = & \sum_{j=1}^\infty e^{-\rho (j-1)\Delta} \left( 1-\eta \left( C_{1, j}+G\right)^\gamma -\frac{\hat b}{C_{1, j}} \right)\Delta \notag\\
& = &  
\frac{e^{-\rho \Delta}}{1-e^{-\rho \Delta}} \left( 1-  \left( C_2+G\right) -{\hat b}/{C_2} \right)\Delta \,.
\end{eqnarray}
Then the optimal consumption $C_2$ is a function of $\Phi_1$, $C_2=\frac{2\hat b\Phi_1}{\sqrt{1+4\hat b\Phi_1(1+\Phi_1)}-1} $,
%defined implicitly by
%\begin{equation}\label{consum_10_dny}
%	\frac{1}{C_2} -\eta(C_2+G)^{\gamma-1} + \Phi_1 \left(\frac{\hat b}{C_2^2}-\eta \gamma (C_2+G)^{\gamma-1} \right)=0\,,
%\end{equation}
%which is a first-order condition from solving
which solves the following optimization problem,
\begin{equation}\label{consum_1_dny}
\max_{C_2} \left\{ \log C_2 - (C_2+G)  +\Phi_1 \left( 1- \left( C_2+G\right) -{\hat b}/{C_2} \right)  \right\}\,.
\end{equation}
When $b>b^*$ such that $C_1^*<C^{\text{Laffer}}=(1-G)/2$, requesting $C_2=C_1^*$ implies a negative constructed Lagrange multiplier:
\begin{equation}
\Phi_1 = -\frac{1}{2}\frac{1-C_1^*}{C^{\text{Laffer}}-C_1^*} <0 \,.
\end{equation}
As a result, adding covenant condition to Lucas-Stokey method cannot implement Ramsey plan when $b>b^*$. 



\section*{Warehoused stuff for intro}


A notable contribution of   \cite{LucasStokey1983}  was to introduce continuation Ramsey planners and continuation representative households
whose best responses  to a particular  $\{b_{0,i}\}_{i=1}^\infty \in {\mathcal D}$  induce them  to set 
$\{\tau_{1,i}\}_{i=1}^\infty = \{\tau_{0,i}^\ast\}_{i=1}^\infty $, thus implementing the Ramsey plan. 
The role of continuation planners and households in \citeauthor{LucasStokey1983}'s analysis is to determine a unique debt structure $\{b_{0,j}\}_{j=0}^\infty$. 
\citeauthor{LucasStokey1983}'s timing protocol requires a period $1$ continuation planner to finance   debt structure {$\{ b_{0,i}\}_{i=1}^\infty$}  and the tail government spending rate process {$\{ G_{i}\}_{i=1}^\infty$}.

\section*{old abstract}

\cite{LucasStokey1983} 
motivated  governments to confirm  an optimal tax  plan by rescheduling  government debt appropriately.
\cite{debortoli2021optimal} showed  that does not work when an   initial government debt  is so  high that a Ramsey planner  sets the tax rate above the peak of the Laffer curve. A   Ramsey plan can be implemented even in that case by adding instantaneous debt to
\citeauthor{LucasStokey1983}'s    contractible subspace and requiring that  continuation  governments preserve the   purchasing power of instantaneous debt.  Using the purchasing-power-adjusted instantaneous debt balance as the state variable, we formulate a Ramsey plan with a Bellman equation and use it to study deterministic settings that feature right continuous with left limit  (C\`adl\`ag) initial term debt structures and government spending sequencees.
%We  formulate a Ramsey plan with  a Bellman equation and use it to study settings  with several initial  term debt  structures and government spending sequencees. 
We extract  implications about  tax smoothing and effects of fiscal policies on bond prices.

\section*{to recycle}

\NW{A key idea underpinning our implementation is that there is no value of using term debt in the original Lucas-Stokey implementation. However, the government is stuck with the initial debt structure $b_{-1, i}; i\geq 0$. This is why the optimal thing to do is to keep $b_{-1, i}; i\geq 0$ as it is and then use short-term debt, supplementing necessary commitment so that the government can actually issue one-period debt at the desired Ramsey plan implied interest rate to implement the Ramsey plan. } 
%\WJ{In settings with incomplete markets, there will be room to use term debt for hedging/risk management purposes. But this is not the case in LS. Longer term is the source of the problem not the cure.} 
\NW{Long term debt issuers have incentives to dilute existing debt holders. But they ultimately bear the costs themselves. This is why shortening debt maturity (using one-period debt) and making sure the government can issue the debt at its Ramsey term implements the Ramsey plan.} 



%Below we sketch out the key as doing so maximizes the household's welfare from its perspective 
By committing to a one-period ahead tax policy as prescribed above, a period $i-1$ government  
voluntarily chooses to implement the Ramsey plan by actively using one-period debt to pay all debts due and finance the remaining accumulated liability. 



\TS{Neng and Wei: at this point we can some literary paragraphs that elaborate on how the implementation shuts down incentives (e.g., ``debt dilution'' motives) to deviate from the Ramsey plan.}

%\NW{Tom and Wei: Let's discuss where to discuss these observations after discussing the next two sections...} 

\NW{Tom: We will work on some literary paragraphs along the debt dilution/Coase-Stokey durable goods insights this weekend (hopefully.)} 

\WJ{ The following is copied from Slides on durable good monopoly interpretation}
 Connection to durable goods monopoly (Coase, Stokey, Bulow): Term debt motivates continuation governments to act
strategically (i.e., to dilute term debt holders).
Short-term debt  attenuates durable-goods monopoly problem when supplemented with a new timing protocol.


\newpage 


\newpage 


%\paragraph{Relation with partial commitment in \cite{chari2020optimal}}
%
%In \cite{chari2020optimal}, they study Ramsey equilibrium for a Chamley-Judd \TS{Tom: let's not call it a Chamley-Judd model; I can explain.} model and show that adding a partial commitment to the value of wealth makes the Ramsey plan time consistent. Our local commitment condition relates to \cite{chari2020optimal}'s partial commitment but have the following differences.
%
%First, because of the complete term debt structure in the Lucas-Stokey economy, our local commitment condition must be applied with a unique short-term debt only strategy. In contrast, in \cite{chari2020optimal}, the use of partial commitment is independent of the debt management policies.\footnote{In a discussion note, \cite{chari2020optimal} define value of wealth for a case of term debt with multiple maturities. For this case, the partial commitment to value of wealth in \cite{chari2020optimal} is strong because it requires the commitment of value for all term debt. In contrast, our local commitment condition requires committing to the value of one-period debt.}
%
%Interestingly, in an economy with only one-period debt, the partial commitment to value of the one-period debt is effectively a special case of our model that considers of all maturity of term debt. When there is only one-period debt, as an underlying assumption, the current government must leave only a balance of one-period debt to future government, which is effectively imposing the short-term debt only strategy described in our paper. By committing to the value of one-period debt, the future government's optimization problem is effectively not distorted by long term debt. As a result, the future government will optimally confirm the continuation of Ramsey plan. 
%
%%its Ramsey compare the Ramsey equilibrium under full commitment and the Markov equilibrium under partial commitment to the value of wealth. 
%%partial commitment to asset returns is different from partial commitment to instrument (one-period ahead policy). In contrast, our local commitment condition shows that partial commitment to the value in unit of utility coincides with the commitment to policy one-period ahead.
%
%
%
%Second, as pointed out by \cite{chari2020optimal}, the partial commitment to value of wealth is different from partial commitment to one-period ahead policy. In contrast, in our model, local commitment condition suggests that partial commitment to the asset return is equivalent to the partial commitment of one-period ahead policy.
%
%%unlike the Chamley model in \cite{chari2020optimal} where their analysis mainly focus on the one-period debt, ours is the Lucas-Stokey environment with a complete set of term debt with all maturities. As a consequence, our local commitment condition must be imposed together with a unique short-term debt only strategy. In contrast, \cite{chari2020optimal} does not have this requirement.
%
%
%
%%Third, in an economy with only one-period debt, the partial commitment to asset price or one-period ahead is effectively a special case of our model with all maturity of term debt. In fact, when there is only one-period debt, the current must leave only a balance of one-period debt to the government in the next period. Effectively, the future government's optimization won't be distorted by the unexistence of long term debt. As a result, the future government will optimally follow the continuation of Ramsey plan.



\end{document}

\section{Optimal Debt Management}

%\NW{What we do in this section is to infer the dynamics of cumulative deficits absent debt restructuring.} 

%\NW{Emphasize this is new.}  

%{Cumulative primary deficits without  debt restructuring  turns out to play  an important role in optimally  managing government debt .  }
With no restructurings of initial term debt $\{b_{-1,j}\}_{j=0}^\infty$, the time $i$ value of government liabilities accumulated  during the time interval $[t_0, t_{i-1}]$ would  evidently be
\begin{equation} 
\Pi_{i}^* = \sum_{j=0}^{i-1} \, \frac{q_{0, j}^*}{q_{0, i}^*} \left( b_{-1, j} - S_{0, j}^*  \right)\Delta\,, \quad \forall i\geq 1\,,
\end{equation}
where  ${q_{0, j}^*}/{q_{0, i}^*}$ converts a time $j$  price to a time $i$ price and  the  increment   to  the government's  accumulated liabilities  over  time interval of length $\Delta$  at time $j$  is $\left(b_{-1, j} - S_{0, j}^*\right) \Delta$. 
The stock  $\Pi_{i}^* $  is thus determined at time ${i-1}$ but valued in units of time $i$ consumption goods. It plays a crucial role in our   implementation of a Ramsey plan. 
%
%\NW{
%For $i=1$, 
%\begin{equation}
%\Pi_{1}^* =  \frac{q_{0, 0}^*}{q_{0, 1}^*} \left( b_{-1, 0} - S_{0, 0}^*  \right)\Delta\,,
%\end{equation}
%} 


Equation (\ref{Pistar}) implies that under a Ramsey plan the following difference equation connects  $\Pi^*_{i+1}$ to $\Pi^*_{i}$:
\begin{eqnarray}
\Pi^*_{i+1}
% & = & \Pi^*_{{i}}\frac{q_{0, {i}}^*}{q_{0, i+1}^*} + \left( b_{-1, {i}} -S_{0, {i}}^*\right)\Delta \cdot\frac{q_{0, {i}}^*}{q_{0, i+1}^*}  \\
%& = & \Pi^*_{{i}}e^{r_{0, {i}}^*\Delta} +  \left( b_{-1, {i}} -S_{0, {i}}^*\right)\Delta \cdot e^{r_{0, {i}}^*\Delta} \label{Pi_dynamics}
%\\ 
=  \big[ \Pi^*_{{i}} +  \left( b_{-1, {i}} -S_{0, {i}}^*\right)\Delta \big] \cdot e^{r_{0, {i}}^*\Delta}\, ,  \label{Pi_dynamics}
\end{eqnarray}
where $r_{0, {i}}^*$ is the short rate over the time interval  $[t_{i}, t_{i+1})$ under the Ramsey plan and 
\[ q_{0, i}^* = e^{-\sum_{j=0}^{i-1}\, r_{0, j}^*\Delta}, \quad \forall i\geq 1;\,\, q_{0, 0}^*=1\,.\]
\TS{Equation \eqref{Pi_dynamics} asserts that   the government's accumulated   liability  due  in period $i$ is the sum of a stock   $\Pi^\ast_{i}$ inherited  from the past and  a new   payment  $\left(  b_{-1, {i}} -S_{0, {i}}^*\right)\Delta$ due  in period $i.$   To finance this obligation at time $i$, the government issues one-period debt in an amount that yields 
proceeds just sufficient to cover the sum of these two components. This results in a time $t_{i+1}$ value of cumulative deficit $\Pi_{i+1}^*$  described by  equation \eqref{Pi_dynamics}. }%, which is the sum of  %Consequently, the total cumulative liability 
%$\left(  b_{-1, {i}} -S_{0, {i}}^*\right)\Delta.$ finances the sum of these two terms by issuing


%Equation (\ref{Pi_dynamics}) states that $\Pi^\ast_{i+1}$ due in period $i+1$ is simply given by  $\Pi^\ast_{i+1}$ due in period $i$ 

\TS{\textbf rewrite or delete the following paragraph.}

From time to time, a government must manage its debt to satisfy the dynamics budget constraint.
\WJ{Let $b_{i-1, j}\Delta$ for $j\geq i$ denote the debt purchases at $t_{i-1}$ and promised to pay off at time $t_j$. The dynamic budget constraint requires the time $t_i$ government to restructure its inherited debt $\{ b_{i-1, j}\Delta; j\geq i+1 \}$ to $\{ b_{i, j}\Delta; j\geq i+1 \}$ for paying off both debt that due at time $t_i$: $b_{i-1, i}\Delta$  and time $t_i$ primary deficit: $S^*_{0, i}\Delta$. As we will show later, managing dynamic budget constraint is effectively about managing the cumulative deficits $\Pi_{i+1}^*$}



%\paragraph{Short-term Debt Balance $B_{t_i}$.}
%The stock of cumulative liability $\{ \Pi_{t_{i}}^*\}$ must be financed from time to time such that the dynamic budget constraint can be satisfied. To highlight the use of short-term debt to finance $\{ \Pi_{t_{i}}^*\}$, we let $B_{t_i}$ denote the stock of short-term debt issued at $t_{i-1}$ and due at time $t_i$ to finance this $\Pi_{t_i}^*$. Note $B_{t_i}$ is different from the one-period coupon rate $b_{t_{i-1}, t_i}$ over $[t_{i-1}, t_i]$. 

%\NW{We shall frequently remind readers how tractable and clean our analysis is, compared to the literature.} 

\paragraph{Managing $\Pi_{i+1}^*$ with Short-term (One-period) Debt Only.}
\TS{We describe a debt management policy that keeps all term debt swith times to maturity  weakly larger than $2\Delta$ equal to ones   implied by the initial term debt GGHH and adjusts  one-$\Delta$ maturity  debt to finance accumulated  liability $\Pi^*_{i+1}$.}

\begin{lemma}
At any time $i\geq 1$ a time $i$, the following continuation debt structure satisfies a time $i$ government's dynamic budget constraint under the Ramsey plan:
\begin{eqnarray} 
b_{i, i+1}\Delta &=& \Pi_{i+1}^*+b_{-1, i+1}\Delta; \label{finance_d10} \\ 
\, b_{i, j}\Delta &=& b_{-1, j}\Delta, \quad j\geq i+2\,. \label{finance_d20}
\end{eqnarray}
\end{lemma}
\noindent Equations \eqref{finance_d}-\eqref{finance_d2} state that one-period debt due at time ${i+1}$ consists of two parts: the accumulated liability  $\Pi_{i+1}^*$ under the Ramsey plan and the initial term  debt $b_{-1, i+1}\Delta$ due at time $i+1$. 
%Equation (\ref{finance_d2}) states that t
The time $i$ leaves unaltered all continuation term debts with maturity equal or longer than $2\Delta$. 	 %by induction. 
%The (indirect) equilibrium effect of longer term debt is captured by $\Pi_{i+1}^*$. 

\begin{proof} %the time $t_0$
We can use (\ref{finance_d})-(\ref{finance_d2})  and mathematical  induction to prove that  all governments $i \geq 0$ satisfy  their dynamic budget constraints.
The $i=0$ government  issues one-period debt due at time $i=1$ in   amount  $b_{0, 1}\Delta=\Pi^*_1 + b_{-1, 1}\Delta$. Using definition \eqref{Pistar}, $ b_{0, 1}\Delta = \left(b_{-1, 0} -S_{0, 0}^* \right)\Delta \frac{q_{0, 0}^*}{q_{0, 1}^*} +b_{-1, 1}\Delta $ consists of two parts: $ \left(b_{-1, 0} -S_{0, 0}^* \right)\Delta \frac{q_{0, 0}^*}{q_{0, 1}^*}$  collects proceeds of $\left(b_{-1, 0} -S_{0, 0}^* \right)\Delta$ that allow the government to to service the  time $0$ primary deficit, while indicates that it reissues  initial term  debt $b_{-1, 1}\Delta$. The government also  reissues all initial debts with times to maturity equal or longer than $2\Delta$: $b_{0, j}\Delta=b_{-1, j}\Delta; j\geq 2$.  That  generates total proceeds: 
\begin{equation}
\sum_{j=1}^\infty \frac{q_{0, 1}^*}{q_{0, 0}^*} \left(b_{0, j}\Delta - b_{-1, j}\Delta \right) = \left(b_{-1, 0} -S_{0, 0}^* \right)\Delta \,.
\end{equation}
Consequently,  the time $0$ dynamic budget constraint holds:
\begin{equation*}
G_0\Delta +b_{-1, 0}\Delta = \tau_{0, 0}^*N_{0, 0}^* +\sum_{j=1}^\infty \frac{q_{0, 1}^*}{q_{0, 0}^*} \left(b_{0, j}\Delta - b_{-1, j}\Delta \right)
\end{equation*} 
We next can show that if  \eqref{finance_d}-\eqref{finance_d2} satisfy dynamic budget constraint at time $i-1$ for $i \geq 1$,   then it also satisfies dynamic budget constraint at time $i$.  At time $i$, the government inherits one-period debt $b_{i-1, i}\Delta=\Pi_{{i}}^*+b_{-1, i}\Delta$
as well as  initial term debt  $\{ b_{i-1, j}\Delta=b_{-1, j}\Delta; j>i\}$. If the time $i$ issues  one-period debt equal to  $b_{i, i+1}\Delta=\Pi_{{i+1}}^*+b_{-1, i+1}\Delta$ as well as longer term  debt from the initial debt structure $\{ b_{i, j}\Delta=b_{-1, j}\Delta; j>i+1\}$,  it raises  \begin{equation}
\sum_{j=i+1}^\infty \frac{q_{0, j}^*}{q_{0, i}^*} \left(b_{i, j}\Delta - b_{i-1, j}\Delta \right) = \frac{q_{0, i+1}^*}{q_{0, i}^*} \left(b_{i, i+1}\Delta - b_{-1, i+1}\Delta \right)=\frac{q_{0, i+1}^*}{q_{0, i}^*} \Pi_{i+1}^*\,.
\end{equation}
Using \eqref{Pi_dynamics}, it follows that the time $i$ dynamic budget constraint holds:
\begin{equation}
G_{i}\Delta + b_{i-1, i}\Delta = \tau_{0, i}^*N_{0, i}^* + \sum_{j=i+1}^\infty \frac{q_{0, j}^*}{q_{0, i}^*} \left(b_{i, j}\Delta - b_{i-1, j}\Delta \right)\,.
\end{equation}
Therefore, financing rules   \eqref{finance_d}-\eqref{finance_d2} satisfy  dynamic budget constraints for all time $i \geq 0$
governments. 
\end{proof}



The time $i$ government's issues short-term debt consisting of three parts: 1) it  rolls over  the short-term debt it has inherited at the beginning of period $I$ to time $i+1$ so that  $b_{i-1, i}\Delta e^{r_{0, i-1}^*\Delta}  = \left(\Pi_{{i}}^*  + b_{-1, i}\Delta\right)e^{r_{0, i-1}^*\Delta} $; 2) it issues enough additional  short-term debt to service the time  $i$ primary deficit: $\left( G_{{i}}- \tau_{0, {i}}^*N_{0, i}^*\right)\Delta e^{r_{0, i-1}^*\Delta} $; 3) it reissues  longer term debts that will due at time $t_{i+1}$: $b_{i-1, i+1}\Delta=b_{-1, i+1}\Delta$. Consequently  short-term debt $b_{i, i+1}$ due at time $i+1$ follows
\begin{equation} \label{key}
b_{i, {i+1}}\Delta = b_{i-1, {i}}\Delta e^{r_{0, i}^*\Delta} +  \left( G_{{i}}- \tau_{0, {i}}^*N_{0, {i}^*}\right)\Delta e^{r_{0, {i}}^*\Delta} + b_{-1, i+1}\Delta = \Pi_{{i+1}}^* + b_{-1, i+1}\Delta \,.
\end{equation} 
%Using \eqref{Pi_dynamics}, we show that the short-term balance  perfectly tracks the value of cumulative deficits, i.e., for all $i\geq 0$, %netting the due initial debt
%\begin{equation} 
%	b_{i, i+1}\Delta = \Pi_{{i+1}}^* + b_{-1, i+1}\Delta \,.
%\end{equation}
Note that the short-term debt balance netting of initial term due at the same time perfectly tracks the cumulative deficits.
\TS{Neng and Wei and Tom: the preceding sentence needs rewriting/clarification.}


%\NW{The intuition for (\ref{key}) is as follows.} 

%
%\NW{
%For $i=0$, 
%\begin{equation}
%\Pi_{1}^* = b_{0, 1}\Delta- b_{-1, 1}\Delta % \frac{q_{0, 0}^*}{q_{0, 1}^*} \left( b_{-1, 0} - S_{0, 0}^*  \right)\Delta\,,
%\end{equation}
%} 
%
%\NW{
%For $i=0$, 
%\begin{equation}
%\Pi_{1}^* =  \frac{q_{0, 0}^*}{q_{0, 1}^*} \left( b_{-1, 0} - S_{0, 0}^*  \right)\Delta\,,
%\end{equation}
%} 
%

%Equation is key as it shows how we can roll over the cumulative financing deficits using one-period debt only. 


\paragraph{Managing $\Pi_{t_i}^*$ with All Term Debt}  \citeauthor{LucasStokey1983}'s continuation Ramsey planner restructures term debts of  all maturities. %, including one-period debt. %Let ${b}_{{i}, j}\Delta$ denote the government debt issued at time $t_{i}$, which is promised to repay one unit of consumption at $t_j>t_i$. 
%Let $\{{b}_{{i-1}, j}\Delta, j\geq i\}$ denote the term-debt structure purchased at time $t_{i-1}$ that leaves for the time $t_i$ government. 
For any time $i\geq 0$,  \citeauthor{LucasStokey1983}'s debt rescheduling policies $\{{b}_{{i}, j}\Delta, j\geq i+1\}$  must satisfy the dynamic budget constraint
\begin{equation}
G_i\Delta + b_{i-1, i}\Delta = \tau_{0, i}^*N_{0, i}^* + \sum_{j=i+1}^\infty \frac{q_{0, j}^*}{q_{0, i}^*} \left(b_{i, j}\Delta - b_{i-1, j}\Delta \right) \,.
\end{equation}
Equation \eqref{Pi_dynamics} implies that 
\begin{equation}
\sum_{j=i+1}^\infty \left( {b}_{{i}, j} - {b}_{{i-1}, j} \right)\Delta\frac{q_{0, j}^*}{q_{0, i}^*} =	\Pi_{{i+1}}^*e^{-r_{0, {i}}^*\Delta} -\Pi_{{i}}^* +{b}_{{i-1}, {i}}\Delta   \,.
\end{equation}
%Evidently,  unlike our earlier policy that does not restructure longer than one-period debts,  \citeauthor{LucasStokey1983}'s  implementation typically  modifies  term debt balances from time to time. %, i.e., $ b_{i-1, j}\Delta \neq b_{-1, j}\Delta$ for most $j\geq i$. 

%and promises to repay $\Pi_{t_{i-1}}^*e^{r_{0, t_{i-1}}^*\Delta}$ at time $t_{i}$. Meanwhile, the time$-t_{i}$ government has to issue additional amount of short-term debt
%$\left( b_{0, t_i} +G_{t_i}- \tau_{0, t_i}N_{0, t_i}\right)\Delta$. As a result, the total balance of short-term debt for time$-t_i$ government is $B_{t_{i}} = \Pi_{t_i}^*$.





\end{document}

\subsection{Consequence of lack of local commitment}
In our strategy, on the one hand, using short-term debt allows us to leave term debt structure untouched, which effectively preserving the distortions exerted by the original term debt structure for the continuation Ramsey planner. On the other hand, our strategy relies on the local commitment condition to constraint the future government to deviate from the Ramsey plan. What if the continuation Ramsey planner lacks local commitment on the purchasing power of the short-term debt? Leveraging the economy of \cite{debortoli2021optimal} as example, we show that our strategy cannot implement Ramsey plan without the local commitment condition. 



\begin{example}
In an economy that households follow utility function $U(C, N)=\log C-\eta N$ and government faces initial debt structure $\{b_{0, t_0}=b, b_{0, t_j}=0; j\geq 1 \}$ and constant spendings $\{G_{0, t_j}=G; j\geq 0\}$,  the Ramsey plan cannot be implemented by leaving $\{b_{0, t_j}; j\geq 1\}$ untouched and adjusting instantaneous debt according to \eqref{implement}.\end{example}

Let's focus on the time$-t_1$ government's continuation Ramsey problem. According to our strategy, the time$-t_1$ government confronts short term balance $\Pi_{t_1}^*$ and term debt structure $\{b_{0, t_j}=0; j\geq 1 \}$, where
\[\Pi_{t_1}^*= \frac{1}{q_{0, t_1}^*} \left( b+G-\tau_{0, t_0}N_{0, t_0} \right)\Delta= e^{\rho \Delta} C_{0, t_1}^*/C_{0, t_0}^*\left( b+G-\tau_{0, t_0}N_{0, t_0} \right)\Delta \,.\]
The time$-t_1$ government's Lagrangian is given by
\begin{equation}
\sum_{j=1}^{\infty} \, e^{-(t_j-t_1) \Delta} \left[ \log C_{t_1, t_j} -\eta ( C_{t_1, t_j}+G) + \Phi_1  \left (1-\eta(C_{t_1, t_j}+G) \right) \right]\Delta - \Phi_1\Pi_{t_{1}}^*/C_{t_1, t_1}.
\end{equation}

Without local commitment condition, the FOCs for the optimal consumption are given by
\begin{align}
\left[1/C_{t_1, t_1} -\eta(1+ \Phi_1)\right]\Delta -\Phi_1\Pi_{t_1}^*/C_{t_1, t_1}^2 = 0 \label{FOC_d1} \\
1/C_{t_1, t_j} -\eta(1+ \Phi_1)  =0,\, \text{for $j\geq 2$}. \label{FOC_d2}
%C_{t_1, t_1}  &=  \argmax \left\{ \bigg[ U(C_{t_1, t_1}, N_{t_1, t_1})  +\Phi_i \left (1 - (C_{t_1, t_1}+G) -b_{0, t_1}/C_{0, t_j}\right) \bigg]\Delta -\Phi_1\Pi_{t_{1}}^*/C_{0, t_1} \right\}, \\
%C_{t_i, t_j}  & =  \argmax \left\{ U(C_{t_i, t_j}, N_{t_i, t_j})  +\Phi_i \left (C_{t_i, t_j}U_{C, t_j} + (C_{t_i, t_j}+G_{t_j})U_{N, t_j} -b_{0, t_j}U_{C, t_j}\right)  \right\},\, \text{for $j\geq 2$}.
\end{align}
We then obtain $C_{t_1, t_j}=C_{t_1, t_2}$ for $j\geq 2$, and the associated implementability condition satisfies
\begin{align} \label{IC_d}
& \left(1-\eta ( C_{t_1, t_1}+G)  \right) \Delta +\frac{e^{-\rho\Delta}}{(1-e^{-\rho\Delta})}\left(1-\eta( C_{t_1, t_2}+G)  \right) \Delta \notag \\
& \quad\quad\quad\quad\quad+ \frac{C_{0, t_0}^*}{C_{t_1, t_1}} \left( 1- \frac{b}{C_{0, t_0}^*}-\eta(C_{0, t_0^*}+G) \right)e^{\rho \Delta}=0\,.
\end{align}
We can easily verify that $C_{t_1, t_1}=C_{t_1, t_2}=C_{0, t_1}^*$ satisfies this implementability condition. However, $\Phi_0^*$  and $C_{t_1, t_1}=C_{t_1, t_2}=C_{0, t_1}^*$ do not satisfy the FOCs \eqref{FOC_d1}-\eqref{FOC_d2}. As a result, the continuation of the Ramsey plan is not optimal for the continuation Ramsey planner. 


\subsection{From discrete-time to continuous-time}

Even when the future government lacks commitment, we show that our strategy can implement the Ramsey plan whenever the Lucas-Stokey method works if we let $\Delta\to 0$ and assume the consumption sequence subject to the following {C\`adl\`ag property: $\lim_{s\downarrow t} C_{0, s} = C_{0, t}$ for all $t\geq 0$.
%\begin{assumption}\label{assump2}
%	Exogenous flows of government expenditures  $\overrightarrow {G_0} = \{G_t;t\geq 0\}$ and 
%	debt-service payouts $\overrightarrow {b_0} = \{b_{0,\,t}; t\geq 0 \}$ are {C\`adl\`ag sequencees, i.e., $\lim_{s\downarrow t} b_{0, s} = b_{0, t}$ and $\lim_{s\downarrow t} G_{s} = G_{t}$ for all $t\geq 0$;  and $\lim_{s\uparrow t} b_{0, s}$ and $\lim_{s\uparrow t} G_{s}$ exist for all $t>0$.}
%\end{assumption}   

To illustrate, let's use the settings in Example 1 and fix $t_1=1$ and initial debt from 0 to 1 is $b$ and that for time $1$ beyond is zero. The cumulative liability valued at $t_1$ is then given by
\[\Pi_{t_1}^*= \int_{0}^{t_1}\frac{q_{0, s}^*}{q_{0, t_1}^*} \left( b+G-\tau_{0, t_0}N_{0, t_0} \right)ds= \frac{e^{\rho}-1}{\rho} \frac{C_{0, t_1}^*}{C_{0, 0}^*}\left( b-C_{0, 0}^*(1-\eta(C_{0, 0}^*+G)) \right).\]
When $\Delta\to 0$, the FOC \eqref{FOC_d1} disappears. Instead, the time $t_1$ consumption $C_{t_1, t_1}$ also satisfies \eqref{FOC_d2} because of the {C\`adl\`ag property: $C_{t_1, s}=1/(\eta(1+\Phi_1))$ for all $s\geq t_1$. Applying $\Delta\to 0$ to the implementability condition \eqref{IC_d}, we obtain the following equation for $\Phi_1$:
\begin{equation}\label{eqPhi_1}
\frac{1}{\rho} \left( 1-\eta G-\frac{1}{1+\Phi_1}\right) -\eta \Pi_{t_1}^* (1+\Phi_1) =0\,,
\end{equation}
which can be written as
\begin{equation}
\rho \eta \Pi_{t_1}^*\Phi_1^2+\left(2 \rho \eta \Pi_{t_1}^* +\eta G -1 \right) \Phi_1 +\rho\eta \Pi_{t_1}^*+\eta G =0\,.
\end{equation}
As one can see, the equation \eqref{eqPhi_1} is a quadratic function of $\Phi_1$ and $\Phi_1=\Phi_0^*$ is one of the root of this equation. Moreover, when $2 \rho \eta \Pi_{t_1}^* +\eta G -1<0$ or  $\Pi_{t_1}^* < \frac{1-\eta G}{2\rho \eta}$, this quadratic equation has two positive roots. We then can verify that when if and only if the initial debt $b$ is lower than the threshold $b^*$ defined in \cite{debortoli2021optimal}, the time$-t_1$ government will optimally choose $\Phi_0^*$, rather than another positive root of \eqref{eqPhi_1}, making the Ramsey plan implementable.

%Note that the value $\frac{1-\eta G}{2\rho \eta}$ is exactly the cumulative liabilities at time $t_1$ under a Ramsey plan that $C_{0, s}^*=C^{\text{Laffer}}$ for $s\geq t_1$. Let $b^*$ denote the initial debt level such that the Ramsey plan satisfies $C_{s}=C^{\text{Laffer}}$ for $s\geq t_1$. Then when initial debt is too high, $b>b^*$, such that the government has to distort the consumption substantially, the cumulative liability $\Pi_{t_1}^*$ then satisfies  $\Pi_{t_1}^* < \frac{1-\eta G}{2\rho \eta}$, and the time$-t_1$ government optimally choose another positive root of \eqref{eqPhi_1}, in which Ramsey plan cannot be implemented.


%Actually, with the entire forward curve $f_{0,t}=r_{0,t}$ in our notation, we then have the entire yield curve. In effect, with the entire forward curve (or equivalently spot curve), $y_{0,t}$ is the long rate for borrowing over (0,t): 
%\[ e^{\int_0^t r_{0,s} d s} = e^{\int_0^t y_{0,t} d s } = e^{y_{0,t} t}  
%\] 
%
%Mapping back to discrete time, to capture our short-term debt, we shall have 
%a pair: 
%using continuous-time compounding convention for interest rate. Note that this is a forward contract (for short-term debt); 


%key issues posed in LS: 83
%\begin{itemize} 
%	\item  the problem/challenge in LS: 83 comes from the existing ... strcture, which can be aribitrary -- so instead of dreaming up a solution by restrcuturing the term debt, it is more natural to go to the other end by NOT touching and creating more problems with restrcuturing the term debt, but instead trying to track / implement the Rarmsey plan  by using isntantaneous debt to track XXx. 
%	\begin{itemize} 
%		\item 		\item with a restrcutred term debt, a new problem arises -- how are you ogoing to do this every period and knowing that it always works -- there is no recursive structure anywhere in this approach.... 
%	\end{itemize} 
%	\item Chari;s (proposed) soln doesnt work b/c doing a PV of $b$ destroys information embedded in term dbt $XXX$. in contrast, our impelemtation stategy sllows us to keep all the info from b0t. and not creating further problem. 
%	\begin{itemize} 
%		\item 	with the PV variable fancy W in Chari, compressing $b$ via PV loses info in $b$ that should impact cosnumption for HH going forward... 
%		\item with a restrcutred term debt, a new problem arises -- how are you ogoing to do this every period and knowing that it always works -- there is no recursive structure anywhere in this approach.... 
%		\item LS had it for shortterm ddebt 
%	\end{itemize} 
%\end{itemize} 
%
%Summary and comparison: 
%\begin{itemize} 
%	\item  Chari;'s original discrete time frmulation: 
%	\item our cont time formulation of Chari:
%	\item our moidel: 
%	
%	%\end{itemize} 
%\end{itemize} 
%
%some observations to include in the future 
%\begin{itemize} 
%	\item LS lit is about characterizing and honor Ramsey forward curve/term structure: note that $r_{0,t}$ is time-0 the (instatntaous) forward rate starting at $t$. by a (instatntaous) forward rate, we mean the forward contract  maturely insrtantly. 
%	\item in a way, Ramsey problem is about to get cont. Ramsey planner to confirm the continuation of Ramsey forward curve.  
%	\item  in a broad and loose snse (at least for now), this relates to instantaneous forward curve approach of term structure: HJM (Ho-Lee Hull-White)
%	\item regarding the proposal of using a combo of term debt and instantaneous debt can sometimes work but
%	does not expand the set of implementable ramsey plan. In other words, if istant debt + local commitment cannot do the trick, then any comb of term debt and instant. debt cannot either. tighten up language later. this is response to less negative Referee 2/ 
%\end{itemize} 







\end{document}

%
%\section{One period term debt in the discrete-time model versus instantaneous debt in the continuous-timer model}
%
%To show the difference of one period term debt in the discrete-time model and the instantaneous debt in the continuous-timer model, we solve the Ramsey problem for the discrete-time model with and letting the length of time period goes to zero. 
%
%Let the time stamps be $\{0, \Delta, 2\Delta, \cdots, i\Delta, \cdots\}$.
%Using the primary approach, we obtain the first order necessary conditions for the households are
%\begin{eqnarray}
%1-\tau_{0, i\Delta} &=& -\frac{U_{N, i\Delta}}{U_{C, i\Delta}} \\
%q_{0, i\Delta} & = & e^{-\rho (i\Delta)} \frac{U_{C, i\Delta}}{U_{C, 0}}\,.
%\end{eqnarray}
%
%Assume the government faces only an one period debt, i.e., $b_{0, 0}>0$ and $b_{0, i\Delta}=0$ for $i\geq 1$.
%The Ramsey planner's problem is to mamximize the household's objective
%\begin{equation}
%\sum_{i=0}^{\infty} e^{-\rho i\Delta } U(C_{0, i\Delta}, N_{0, i\Delta}) 
%\end{equation}
%subject to the implementability condition
%\begin{equation}
%\sum_{i=0}^\infty \, e^{-\rho i\Delta } \left( U_{C,  i\Delta} C_{0, i\Delta}   + U_{N,  i\Delta}   N_{0, i\Delta}     \right)\geq  \sum_{i=0}^\infty \, e^{-\rho i\Delta } U_{C,  i\Delta}   b_{0, i\Delta} = U_{C, 0}   b_{0, 0}\,.
%\end{equation}
%
%Formulating the Lagrangian and solve the optimality conditions, we obtain that for $i\geq 1$
%\begin{align*}
%&&(1+\Lambda_0)\left(U_{C, 0} + U_{N, 0}\right) +  \Lambda_0\left( \left(C_{0, 0}-b_{0, 0} \right) \left( U_{CC, 0}+U_{CN, 0} \right) + U_{NC, 0}+U_{NN, 0} \right) = 0 \\
%&& (1+\Lambda_0)\left(U_{C, i\Delta} + U_{N, i\Delta}\right) +  \Lambda_0\left( C_{0, i\Delta}\left( U_{CC, i\Delta}+U_{CN, i\Delta} \right) + U_{NC, i\Delta}+U_{NN, i\Delta} \right) = 0  
%\end{align*}
%
%Taking \citeauthor{debortoli2021optimal}'s utility function for instance, letting $\eta=\gamma=1$, we have
%\begin{align*}
%&&-1-\Lambda_0 +\frac{1}{C_{0, 0}} + \Lambda_0\frac{b_{0. 0}}{C^2_{0, 0}} =0 \\
%&& -1-\Lambda_0 +\frac{1}{C_{0, i\Delta}} =0\,, \quad i\geq 1 \,.
%\end{align*}
%Substituting into the implementability condition, we have
%\begin{equation*}
%1-C_{0, 0}-G + \frac{e^{-\rho\Delta}}{1-e^{-\rho\Delta}}\left(1-C_{0, \Delta}-G \right) =\frac{b_{0, 0}}{C_{0, 0}}\,.
%\end{equation*}
%As $\Delta\to 0$, $\frac{e^{-\rho\Delta}}{1-e^{-\rho\Delta}}\to\infty$, the Ramsey planner optimally choose $C_{0, 0} \neq \lim_{\Delta\to 0} C_{0, \Delta}$ to satisfy the implementability condtion and meanwhile to maximize the households' objective. 
%
%This limiting results of one-period term debt is quite different from that of the instantaneous debt. Facing instantaneous debt, the government has no incentive to choose $C_{0, 0} \neq C_{0, \Delta}$ for any $\Delta>0$ as its ability to inflate the purchasing power of instantaneous debt is constrained.
%
%
%\begin{figure}
%	\centering
%	\begin{minipage}{\textwidth}
%		\includegraphics[width=0.45\linewidth, height=0.26\textheight]{figure_202403/Delta_A.eps}
%		\includegraphics[width=0.45\linewidth, height=0.26\textheight]{figure_202403/Delta_B.eps}
%	\end{minipage}
%	\caption{Ramsey plan as $\Delta\to 0$. $b_{0, 0}=0.5, \eta=\gamma=1, \rho=0.6931$.}
%\end{figure}
%
%


\section{To recycle}

Let  $\{C_{0, t}^*, N_{0, t}^*; t\geq 0\}$ be the Ramsey plan derived in section \ref{sec:environ} and 
note that  the associated product $X_t^\ast$ sequence  follows dynamics
\begin{eqnarray}\label{dyast}  
d X_t^\ast  =  \left[\rho X_{t}^\ast-C^\ast_{0, t}U_{C, t} -N^\ast_{0, t}U_{N, t}  +b_{0, t}U_{C, t}\right]dt \, . % \\
%&  & +\, .
\end{eqnarray}
where $N^\ast_{0, t} = C^\ast_{0, t}+G_t$. Intergrating this differential equation forward gives
%Then the ``forward looking" of $X_t^\ast$ implied by the implementability constraint suggests that
\begin{eqnarray} \label{Xastforward}
X_t^\ast = \int_{t}^{\infty} e^{-\rho(s-t)}\left[ C^\ast_{0, s}U_{C^\ast, s} +N^\ast_{0, s}U_{N, s}  -b_{0, s}U_{C, s} \right] ds
\end{eqnarray}

For  $t\in[0, \infty)$,
define  a continuation value at state $(X^\ast_t, t)$  as
\begin{equation}
V_t^\ast = %\max_{\{C_{0, s\geq t} \}}\; 
\int_t^{\infty} e^{-\rho(s-t)} U({C}_{0, s}^\ast, {N}_{0, s}^\ast) d s \,.
\end{equation}
Evidently,  $V^\ast$ is a function of $X_t^\ast$ and $t$. Thus, use (\ref{Xastforward}) to deduce
\begin{eqnarray*}
V_t^* %&= & \int_t^{\infty} e^{-\rho(s-t)} U({C}_{0, s}^\ast, {N}_{0, s}^\ast) d s \,\\
& = & \int_t^{\infty} e^{-\rho(s-t)} U({C}_{0, s}^\ast, {N}_{0, s}^\ast) d s  \\
& & \quad  +  \Lambda_{0}^*\left\{-X_t^\ast +\int_{t}^{\infty} e^{-\rho(s-t)}\left[ C^\ast_{0, s}U_{C, s} +N^\ast_{0, s}U_{N, s}  -b_{0, s}U_{C, s} \right] ds    \right\} \\
& = & -\Lambda_{0}^\ast X_t^\ast + \int_t^{\infty} e^{-\rho(s-t)} \bigg( U({C}_{0, s}^\ast, {N}_{0, s}^\ast) + \Lambda_0^\ast \left[C^\ast_{0, s}U_{C, s} +N^\ast_{0, s}U_{N, s}  -b_{0, s}U_{C, s} \right]  \bigg)dt  \\
& = & -\Lambda_{0}^\ast X_t^\ast+ H^*(t; \Lambda_{0}^\ast)  \,.
\end{eqnarray*}
where $H^*(t; \Lambda_{0}^\ast) = \int_t^{\infty} e^{-\rho(s-t)} \big( U({C}_{0, s}^\ast, {N}_{0, s}^\ast) + \Lambda_0^\ast \left[C^\ast_{0, s}U_{C, s} +N^\ast_{0, s}U_{N, s}  -b_{0, s}U_{C, s} \right]  \big)dt $ satisfies the  ordinary differential equation:
\begin{align}\label{HJBHast} %\nonumber
\rho H^\ast 
& =   \frac{d H^\ast}{dt} +\; U(C^\ast_{0, t}, N^\ast_{0, t})  + \Phi_0\left(C^\ast_{0, t} U_{C, t}+N^\ast_{0, t}U_{N, t} -b_{0, t}U_{C, t} \right)\,    % \\
\end{align}
so that $V_t^\ast$ is a function of $X_t^\ast$ and $t$: 
\begin{equation} \label{VHast}
V^\ast_t = V(X_t^\ast, t) = -\Lambda_{0}^\ast X_t^\ast + H^\ast(t; \Lambda_{0}^\ast); \, \;t\geq 0 \, .
\end{equation}
We can  verify that for all $t\geq 0$  value function $V^\ast(X_t,t)$   satisfies differential  equation: % \TS{I propose to replace $0+$ with $0$ throughtout.}
\begin{align} \label{Vepsilonast}
\rho V^\ast  & =   \frac{\partial V^\ast }{\partial t}+ %\max_{C_{0, t}} \;\; 
U(C_{0, t}^\ast, C_{0, t}^\ast+G_t) + \frac{\partial V^\ast }{\partial X_t} \left(\rho X_t^\ast-C_{0, t}^\ast, U_{C, t}-(C_{0, t}^\ast+G_t)U_{N, t} +b_{0, t}U_{C, t} \right)\,.
\end{align}
Recall from section \ref{sec:environ} that  %$C^\ast_{0, t}= C(\Lambda_0^\ast, G_t, b_{0, t})$ is the maximand of the optimization problem:
\begin{equation} \label{opt-Cast}
C^\ast_{0, t} = \argmax_{C_{0, t}} \;\; U(C_{0, t}, C_{0, t}+G_t)  + \Lambda_0^\ast\left(C_{0, t} U_{C, t}+(C_{0, t}+G_t)U_{N, t} -b_{0, t}U_{C, t} \right).
\end{equation}
and that the associated  labor supply is 
\begin{equation}\label{opt-Nast}
N^\ast_{0, t} = C^\ast(\Lambda_0^\ast, G_t, b_{0, t})+G_t. 
\end{equation}








\end{document}



\newpage
\section{Revisit DNY}

As our model could be extended to piece-wise continuous debt maturity distribution, we are able to revisit the findings in DNY. To be precise, 
The initial given debt structure is %In addition, the time$-0$ government faces the following special initial debt
\begin{equation} \label{DNYDebt}
b_{0, t}=b_0\geq 0, t\in[0, T);\;  b_{0, t}=0, t\in[T, \infty)\, ; B_{0} = 0\, ,
\end{equation}
where $T$ is fixed. 
Note that here  $T$ can be any finite value rather than one. 
The  government inherits from the previous government  a constant stream of debt services at the rate of $b$ per unit over the time interval $[0, T)$  that it has to honor. To keep close to the DNY, we assume there is no initial short-term debt: $B_{0} = 0$.

\paragraph{Optimal commitment policy} Given the initial debt structure, we follow the our proposed dynamic programing method to derive the optimal commitment policy for the government as follows.

\begin{corollary}
A Ramsey plan for this problem, when exists, is given by
\begin{eqnarray*}
& C_{0, t}^{*}  =  C(\Phi_{0}^{*}, b_{0}, G)1_{\{t\in[0, T]\}} +  C(\Phi_{0}^{*}, 0, G)1_{\{t\in(T, \infty)\}};
N_{0, t}^{*}  =  C_{0, t}^{*}+G ;  \\
& 	\tau_{0, t}^{*}   =   1+\frac{U_{N, t}}{U_{C, t}};
q^*_{0,t}  = 
e^{-\rho t}\frac{U_{C, t}}{U_{C, 0}}
\end{eqnarray*}
where 
$\Phi_{0}^*= \argmax_{\{\Phi_{0}\geq 0: \Psi_0(\Phi_{0})=0\} } \;  \int_{0}^{\infty} e^{-\rho t}U(C(\Phi_0, b_{0, t}, G), C(\Phi_0, b_{0, t}, G)+G)dt ,$
in which $\Psi_0(\Phi_{0})$ is given by
\begin{align}
\Psi_0(\Phi_{0}) = \frac{1-e^{-\rho T}}{\rho}\left( 1-\eta(C(\Phi_0, b_0, G)+G)^\gamma-\frac{b_0}{C(\Phi_0, b_0, G)} \right)+\frac{e^{-\rho T}}{\rho}\left(1-\eta(C(\Phi_0, 0, G)+G)^\gamma\right)\,,
\end{align}
and
$C(\Phi_{0}^{}, b_{0}, G)$ and $C(\Phi_{0}^{}, 0, G)$
are the unique solution of 
\begin{equation*}
\frac{1}{C}+ \Phi_0^{}\frac{b_{0}}{C^2} -\left(\eta+\Phi_0^{}\eta\gamma^2\right)(C_{}+G)^{\gamma-1} = 0\,,
\end{equation*}
and 
\begin{equation*}
\frac{1}{C} -\left(\eta+\Phi_0^{}\eta\gamma^2\right)(C_{}+G)^{\gamma-1} = 0\,,
\end{equation*}
respectively, for $\Phi_0>0$.
\end{corollary}



\paragraph{Time-consistency and the sign of shadow price}

Because the initial debt $\{b_{0, t}\}$ is constant zero over time $t\geq T$, we can apply our criterion of time-consistency in Proposition \ref{prop:9} to decide whether the Ramsey plan is time consistent for time $T$ onward. 

\begin{corollary}
The time$-0$ Ramsey plan is time consistent if and only if the that time$-0$ shadow price of servicing debt $\Phi_0^*$ is the lowest among the two positive roots of $\widetilde{\Psi}_T(\cdot)=0$, i.e., 
\begin{equation}
\Phi_0^* = \min\{\Phi>0: \widetilde{\Psi}_T(\Phi)=0\}	
\end{equation} 
where $\widetilde{\Psi}_T(\Phi)=\frac{1}{\rho}\left(1 - (C(\Phi, 0, G)+G)^\gamma \right)C(\Phi, 0, G) - \Pi_T$ and $\Pi_T$ is the short-term debt the time $T$ government inherited from the previous government.
\end{corollary}

From corollary (\ref{}), we know that $C(\Phi, 0, G)$ is a decreasing function of $\Phi$. Therefore, the equation $\widetilde{\Psi}_T(\cdot)=0$ usually has two positive roots.

Introducing short-term debt entirely changes DNY's argument about the sign of Lagrange multiplier. Another interpretation of DNY if only term-debt restructuring policy is adopted is given below. When $\lambda=0$, following equation (\ref{phiu0}), the associate $\Phi_T(0)$ satisifes
\begin{equation}
\Phi_T(0)   =  \Phi_0^*   \left(1 - \frac{\Pi_T U_{C^*, T}}{\int_{T}^{\infty} e^{-\rho (t-T)}  A_{C^*, t}U_{C^*, t} dt} \right)^{-1} \,.
%\frac{\int_{u}^{\infty} e^{-\rho (t-u)} \left(b_{0, t} -A_{u, t} \right)U_C dt }{\int_{u}^{\infty} e^{-\rho (t-u)} \left(b_{0, t} -A_{u, t} \right)U_C dt+B_u U_C(C_{0, u} ,C_{0, u}+G_u)} \,, \\
\end{equation}
When $b_0$ is higher, the MRS adjusted accumulative primary deficit $\Pi_TU_{C^*, T}$ is higher, and the $\Phi_T(0)$ could be negative, such that the optimality of this potential debt restructuring policy is not achievable.





\begin{figure}
\centering
\includegraphics[width=0.7\linewidth]{figure_d/Yared_Laffer_2}
\caption{When initial debt $b_0$ is too high, time $0$ government chooses (1) a primary surplus level for $T$ onward and (2) a smaller consumption level that supports this primary surplus, to generate higher MRS adjusted primary surplus. This smaller consumption level is associated with a negative marginal MRS adjusted funding gap (i.e., $\frac{\partial}{\partial C}\left(\frac{S(C)-b_1}{C_1}\right)<0$). Without short-term debt, time $T$ government facing restructured debt $\widehat{b}_1$ left by time $T-$ government, optimally chooses the consumption level which supports the same primary surplus level chosen by time $0$ government for $T$ onward but has negative marginal MRS adjusted gap.  }
\label{fig:yaredlaffer2}
\end{figure}

\begin{table}[!h]
\caption{Key values of Figure 2 ($b_0=50, b_1=0$). $S(C)=C(1-\eta(C+G))$. $C_1$ and $\widehat{C}_1$ are two roots of $S(C)=\widehat{b}_1=\rho \widehat{B}_{T-}$, and they satisfy $C_1\widehat{C}_1=\widehat{b}_1$. In addition $\frac{S(C_1)-b_1}{C_1}=\frac{\widehat{b}_1}{C_1}=\widehat{C}_1$. In our model with short-term debt, one can use short-term debt to roll over government's funding gap, i.e., $S(C_{0, t})-b_{0, t}$, and restructure the short-term debt and leave it to future government. In this case, time $T-$ government accumulates short-term debt $B_{T-}=91.9356$ and restructures it to be $\widehat{B}_{T-}=1.5748$ and leave it to time $T$ government. Facing $\widehat{B}_{T-}=1.5748$ and $b_1=0$, time $T$ government solve the optimization problem again, and find there two positive Lagrangian multiplies $\Lambda_{T}^*=7.7703, \quad 0.4597$ satisfies the FOCs of the optimization problem. Lagrange multiplier represents the shadow price of \NW{marginal value of MU adj funding gap}. In contrast, in DNY(2021), there is no short-term debt, time $T-$ government accumulates funding gap $B_{T-}=91.9356$ and restructures it to be long-term debt $\widehat{b}_{1}=0.0787$ and leave it to time $T$ government. The future government facing only $\widehat{b}_{1}=0.0787$ solves the optimization problem, and find there is a unique positive Lagrange multiplier satisfying FOCs. For both models, time $T$ government will renege and choose a higher consumption level after $T$ which supports the same primary surplus level with the consumption level of commitment policy for time $T$ onward.}
\begin{tabular}{c|l|l}
\hline 
Laffer consumption&  \multicolumn{2}{c}{$C^{laffer} = 0.4 $}	  \\
First-best consumption &  \multicolumn{2}{c}{$C^{FB} = 1  $ }  \\
Commitment policy	&  \multicolumn{2}{c}{$C_0= 6.71   ;\quad  C_1 = 0.1149 $
} \\ 
&\multicolumn{2}{c}{$N_0= 6.91    ;\quad N_1 = 0.1349 $
} \\
&\multicolumn{2}{c}{$\tau_0= -5.71    ;\quad \tau_1 = 0.8851 $
} \\
Value with $C_0, C_1$ at $0$ & \multicolumn{2}{c}{ -52.0320}   \\
Two consumption levels&\multicolumn{2}{c}{$C_1= 0.1149    ;\quad \widehat{C}_1 = 0.6851 $
} \\
supporting primary surplus & \multicolumn{2}{c}{} \\

\hline
&   Ours (with short-term debt)  &  Yared (no short-term debt) \\\hline
Accumulated funding gap &  $B_{T-} = 91.9356 $ & $B_{T-} = 91.9356 $\\
Debt restructuring  & $\widehat{B}_{T-} = \frac{B_{T-}}{C_0}C_1= 1.5748$  & $\widehat{b}_1 = \rho \frac{B_{T-}}{C_0}C_1 = 0.0787$ \\
&& \\
MU adj funding gap at $C_1$ & $\frac{S(C_1)-b_1}{C_1} =0.6851 $ &  $\frac{S(C_1)-\widehat{b}_1}{C_1} = 0$ \\
& $\frac{\partial}{\partial C_1}\frac{S(C_1)-b_1}{C_1} = -1 $ &  $ \frac{\partial}{\partial C_1}\frac{S(C_1)-\widehat{b}_1}{C_1} = 4.9573$ \\
Lagrange multiplier at $C_1$  & $\Lambda_{T}^* = 7.7703 $    & $\Lambda_{T}^* = -1.5524 $  \\
Value with $C_1$ at $T$ & -49.5660 & -49.5660   \\

&& \\
MU adj funding gap at $\widehat{C}_1$ & $\frac{S(\widehat{C}_1)-b_1}{\widehat{C}_1} = 0.1149$ &  $\frac{S(\widehat{C}_1)-\widehat{b}_1}{\widehat{C}_1} = 0$  \\
& $\frac{\partial}{\partial \widehat{C}_1}\frac{S(\widehat{C}_1)-b_1}{\widehat{C}_1} = -1$ &  $\frac{\partial}{\partial \widehat{C}_1}\frac{S(\widehat{C}_1)-\widehat{b}_1}{\widehat{C}_1} =-0.8323 $  \\
Lagrange multiplier at $\widehat{C}_1$  & $\Lambda_{T}^* = 0.4597$   & $\Lambda_{T}^* = 0.5524$  \\
Value with $\widehat{C}_1$ at $T$ & -25.2662 & -25.2662   \\
%Value with $C_0, \widehat{C}_1$ at $0$ & -28.9172 & -28.9172   \\
\hline 
\end{tabular} 
\end{table}

\clearpage
\begin{figure}
\centering
\includegraphics[width=0.7\linewidth]{figure_d/Yared_Laffer_3}
\caption{When initial debt $b_0$ is not too high, time $0$ government optimally chooses (1) a primary surplus level for $T$ onward and (2) a higher consumption level that supports this primary surplus, to generate enough MRS adjusted primary surplus. Without short-term debt, time $T$ government facing restructured debt $\widehat{b}_1$ left by time $T-$ government, optimally chooses the same consumption level as the commitment policy. The negative marginal MRS adjusted gap for the consumption level ensures its optimality for time $T$ government. }
\label{fig:yaredlaffer3}
\end{figure}

\begin{table}[!h]
\caption{Key values of Figure 3 ($b_0=10, b_1=0$). $S(C)=C(1-\eta(C+G))$. $C_1$ and $\widehat{C}_1$ are two roots of $S(C)=\widehat{b}_1=\rho \widehat{B}_{T-}$, and they satisfy $C_1\widehat{C}_1=\widehat{b}_1$. In addition $\frac{S(C_1)-b_1}{C_1}=\frac{\widehat{b}_1}{C_1}=\widehat{C}_1$.  In this case, time $T-$ government accumulates short-term debt $B_{T-}=14.5539$ and restructures it to be $\widehat{B}_{T-}=2.9694$ and leave it to time $T$ government. Facing $\widehat{B}_{T-}=2.9694$ and $b_1=0$, time $T$ government solve the optimization problem again, and find there two positive Lagrangian multiplies $\Lambda_{T}^*=0.9709, \quad 2.4173$ satisfies the FOCs of the optimization problem. Lagrange multiplier represents the shadow price of \NW{marginal value of MU adj funding gap}. In contrast, in DNY(2021), there is no short-term debt, time $T-$ government accumulates funding gap $B_{T-}=14.5593$ and restructures it to be long-term debt $\widehat{b}_{1}=0.1485$ and leave it to time $T$ government. The future government facing only $\widehat{b}_{1}=0.1485$ solves the optimization problem, and find there is a unique positive Lagrange multiplier satisfying FOCs. For both models, time $T$ government will not renege and choose the same consumption level with that of commitment policy for time $T$ onward.   }
\begin{tabular}{c|l|l}
\hline 
Laffer consumption&  \multicolumn{2}{c}{$C^{laffer} = 0.4 $}	  \\
First-best consumption &  \multicolumn{2}{c}{$C^{FB} = 1  $ }  \\
Commitment policy	&  \multicolumn{2}{c}{$C_0= 2.4877   ;\quad  C_1 = 0.5074 $
} \\ 
&\multicolumn{2}{c}{$N_0= 2.6877    ;\quad N_1 = 0.5274 $
} \\
&\multicolumn{2}{c}{$\tau_0= -1.4877    ;\quad \tau_1 = 0.4926 $
} \\
Value with $C_0, C_1$ at $0$ & \multicolumn{2}{c}{-28.0985}   \\
Two consumption levels &\multicolumn{2}{c}{$C_1= 0.5074   ;\quad \widehat{C}_1 = 0.2926 $
} \\
supporting primary surplus & & \\
\hline
&   Ours (with short-term debt)  &  Yared (no short-term debt) \\\hline
Accumulated funding gap &  $B_{T-} = 14.5593 $ & $B_{T-} = 14.5593 $\\
Debt restructuring  & $\widehat{B}_{T-} = \frac{B_{T-}}{C_0}C_1= 2.9694$  & $\widehat{b}_1 = \rho \frac{B_{T-}}{C_0}C_1 = 0.1485$ \\
&& \\
MU adj funding gap at $C_1$ & $\frac{S(C_1)-b_1}{C_1} =0.2926 $ &  $\frac{S(C_1)-\widehat{b}_1}{C_1} = 0$ \\
& $\frac{\partial}{\partial C_1}\frac{S(C_1)-b_1}{C_1} = -1 $ &  $ \frac{\partial}{\partial C_1}\frac{S(C_1)-\widehat{b}_1}{C_1} = -0.4231$ \\
Lagrange multiplier at $C_1$  & $\Lambda_{T}^* = 0.9709 $    & $\Lambda_{T}^* = 2.2932 $  \\
Value with $C_1$ at $T$ & -27.7169 & -27.7169   \\

&& \\
MU adj funding gap at $\widehat{C}_1$ & $\frac{S(\widehat{C}_1)-b_1}{\widehat{C}_1} = 0.5074$ &  $\frac{S(\widehat{C}_1)-\widehat{b}_1}{\widehat{C}_1} = 0$  \\
& $\frac{\partial}{\partial \widehat{C}_1}\frac{S(\widehat{C}_1)-b_1}{\widehat{C}_1} = -1$ &  $\frac{\partial}{\partial \widehat{C}_1}\frac{S(\widehat{C}_1)-\widehat{b}_1}{\widehat{C}_1} =0.7345 $  \\
Lagrange multiplier at $\widehat{C}_1$  & $\Lambda_{T}^* = 2.4173$   & $\Lambda_{T}^* = -3.2938$  \\
Value with $\widehat{C}_1$ at $T$ & -34.4296 & -34.4296  \\
%Value with $C_0, \widehat{C}_1$ at $0$ & -34.4831 & -34.4831  \\
\hline 
\end{tabular} 
\end{table}





\paragraph{Time-consistency and consumption over the Laffer curve}
Another important interpretation of DNY is that they think a Ramsey plan is time-inconsistent because the allocation of time$-0$ Ramsey plan is above the peak of Laffer curve for time $T$ onward.

However, as we show in the following example, when government's initial debt satisfies 
\begin{equation}
b_{0, t}=b_0= 0, t\in[0, T);\;  b_{0, t}=b_1>0, t\in[T, \infty)\, ; B_{0} = 0\, ,
\end{equation}
where $b_1$ is high enough, its time$-0$ Ramsey plan has a consumption above the peak of Laffer curve for time $T$ onward. This is because, the government optimally save in the period $[0, T)$ and then pass the saving to the time$-T$ government such that $\Pi_T<0$. 

We can show that this time$-0$ Ramsey plan is time consistent when $\eta=1, \gamma=1$. Therefore, even the consumption is above the peak of Laffer curve, it is still time consistent.

\begin{figure}[h!]
\centering
\includegraphics[width=\linewidth, height=0.6\textheight]{figure_yared/Counter_example_1}
\caption{Optimal commitment policy with varying $\eta$. Initial long-term debt $b_{0-, t}=b_01_{\{t< T\}}+b_11_{\{t\geq T\}}$ satisfies $b_0=0, b_1=0.3$. Both governments have the same initial short-term debt $B_{0-}=0$ and the same government spending $G_t=0.2$, but have different labor costs $\eta=2$ v.s. $\eta=1$ Other parameter values are $\rho=0.05, \eta=1, T=10$.  \label{fig:commitpolicyeta}}

\end{figure}




%
%
%%\begin{proposition}
%%	For utility $U(C, N)=C-\eta\frac{N^{1+\gamma}}{1+\gamma}$ and an initial debt $\{b_{0-, t}, B_{0-}, G_t, t\geq 0\}$, if the commitment policy exists, it must be time-consistent.
%%\end{proposition}
%
%%\begin{proposition}
%%	If $\{b_{0-, t}=0; t\geq 0\}$, and $\{G_t=g; t\geq 0\}$, then the commitment policy is time-consistent.
%%\end{proposition}
%
%%\WJ{Conjecture: For Markov-Chain model, if $\{b_{0-, t}\}$ is state-dependent but not time-dependent, then the optimal commitment policy is still time-consistent.}
%%
%%\WJ{For the case $\{b_{0-, t}=0; t\geq 0\}$, there are two ways to implement the commitment policy which is time-consistent:
%%	\begin{itemize}
%%		\item Method 1: keep using short-term debt $B_t$ to roll over the debt, and the government at each period pay out $B_0\rho$ which is exactly the primary surplus at that period.
%%		\item Method 2: restructure the short-term debt at time $0$ $B_0$ to be a term structure debt $b_{0, t}=B_0\rho$. So the government will use the primary surplus to pay off its due debt at each period in future. 
%%		\item Implication: Short-term debt is equivalent to a constant term structure debt.
%%\end{itemize}}
%
%
%
%
%
%%
%%\subsection{A Different Timing Protocol}
%%
%%In this part, we consider a new timing protocol for the government. Let $B_{0-}$ denote the lump-sum quantity due at $0$ (short-term debt). In this case, 
%%Time $0$ government takes debt structure $B_{0-}$  and $\{ b_{0-, t}; \, t\geq 0\}$ and government spending $G_t$ as given. Then the government chooses $\{C_{0, t}, N_{0, t}; t\geq 0\}$, a restructured term debt $\{b_{0, t}; t\geq 0\}$ and short-term debt $B_0$ simultaneously to solve
%%\begin{eqnarray} \label{obj1}
%%&& \max_{\{C_{0, t}, b_{0, t}; t\geq 0\}} \, \int_{0}^{\infty} U(C_{0, t}, C_{0, t}+G_t) dt \\
%%&s.t. &  \quad  \int_{0}^{\infty} e^{-\rho t} ( C_{0, t}U_C+(C_{0, t}+G_t)U_N )  dt  \geq   \int_{0}^{\infty} e^{-\rho t}b_{0, t} U_C(C_{0, t}, C_{0, t}+G_t) dt + B_{0} U_C(C_{0, 0}, C_{0, 0}+G_0) , \label{imconstraint} \\
%%&& \quad  \int_{0}^{\infty} e^{-\rho t}b_{0, t}^{} U_C(C_{0, t}, C_{0, t}+G_t) dt+ + B_{0} U_C(C_{0, 0}, C_{0, 0}+G_0)  \nonumber \\
%%&& \quad  = \int_{0}^{\infty} e^{-\rho t}b_{0-, t} U_C(C_{0, t}, C_{0, t}+G_t) dt  + B_{0-} U_C(C_{0, 0}, C_{0, 0}+G_0) . \label{restrdebt}  
%%\end{eqnarray}
%%
%%Actually, we can show that this problem with a new timing protocol is equivalent to Lucas-Stokey problem.
%%\begin{proposition}
%%Under commitment, Lucas-Stokey problem (\ref{obj})-(\ref{impletconstraint}) is equivalent to the new timing protocol problem (\ref{obj1})-(\ref{restrdebt}). 
%%\end{proposition} 
%%
%%\begin{proof}
%%	We first show Lucas-Stokey's optimal policy is admissible to the new timing protocol problem. 
%%	Assume Lucas-stokey's optimal policy is $\{\widehat{C}_{0, t}, \widehat{N}_{0, t}, \widehat{\tau}_{0, t}; t\geq 0\}$. For the timing protocol problem, we just let $b_{0, t}=b_{0-, t}$ and $B_0= B_{0-}$. 
%%	
%%	Substituting $\{\widehat{C}_{0, t}, \widehat{N}_{0, t}, \widehat{\tau}_{0, t}, b_{0, t}; t\geq 0\}$ and $B_0$ into problem (\ref{obj1})-(\ref{restrdebt}), we can easily verify that constraints (\ref{imconstraint}) and (\ref{restrdebt}) are well satisfied. Hence, this is an admissible strategy to the new timing protocol problem.
%%		
%%	Next, we show the new timing protocol problem's optimal policy is admissible to the Lucas-Stokey problem. Assume new timing protocol problem's optimal policy is $\{\widetilde{C}_{0, t}, \widetilde{N}_{0, t}, \widetilde{\tau}_{0, t}; t\geq 0\}$. Then simpliy substitue this strategy into Lucas-Stokey problem (\ref{obj})-(\ref{impletconstraint}). Since equations (\ref{imconstraint})-(\ref{restrdebt}) hold, which implies that
%%	\begin{equation*}
%%	\int_{0}^{\infty} e^{-\rho t} ( \widetilde{C}_{0, t}U_C+(\widetilde{C}_{0, t}+G_t)U_N )  dt  \geq   \int_{0}^{\infty} e^{-\rho t}b_{0-, t} U_C(\widetilde{C}_{0, t}, C_{0, t}+G_t) dt + B_{0-} U_C(\widetilde{C}_{0, 0}, \widetilde{C}_{0, 0}+G_0) . 
%%	\end{equation*}
%%	Therefore $\{\widetilde{C}_{0, t}, \widetilde{N}_{0, t}, \widetilde{\tau}_{0, t}; t\geq 0\}$ satisfy the constraint (\ref{impletconstraint}) and is admissible to the Lucas-Stokey problem.
%%	
%%    In summary, Lucas-Stokey problem (\ref{obj})-(\ref{impletconstraint}) is equivalent to the new timing protocol problem (\ref{obj1})-(\ref{restrdebt}).
%%\end{proof}
%
%
%%\section{Time-Consistency of Optimal Policy}
%%
%%Time$-u$ government inherited debt from previous government: term debt $b_{0-, t}, t\geq u$ and short-term debt $B_{u}=Y_u/U_C(C_{0, u}^*, C_{0, u}^*+G_u)$.
%%
%%If the time$-u$ government's can reevaluate it's policy given the inherited debt, whether it will follow the continuation of time$-0$ government's policy?  
%%
%%The time$-u$ government chooses $\{C_{u, t}, N_{u, t}, \tau_{u, t}; t\geq u\}$ to solve the following problem
%%
%%\begin{eqnarray} \label{obj}
%%&& \max_{\{C_{u, t}; t\geq 0\}} \, \int_{0}^{\infty}  e^{-\rho t} U(C_{u, t}, C_{u, t}+G_t) dt \\
%%&s.t. &  \int_{0}^{\infty} e^{-\rho t} ( C_{0, t}U_C+(C_{0, t}+G_t)U_N )  dt  \geq   \int_{0}^{\infty} e^{-\rho t}b_{0-, t} U_C(C_{0, t}, C_{0, t}+G_t) dt+B_0U_{C},  \label{impletconstraint} %\\
%%\end{eqnarray}   
%
%
%
%\section{Application to DNY} 
%\WJ{ Plan in this section.
%\begin{itemize}
%	\item For a simple debt structure used in DNY, we provide a necessary and sufficient condition under which the optimal commitment policy is time consistent.
%	\item Show numerically the government's optimal commitment policy and the associated shadow prices for different debt restructuring strategies.
%\end{itemize}
%}
%
%
%
%\newpage
%
%In this section, we apply our solution method to the continuous-time formulation of the economy  analyzed by DNY: a utility and initial debt structure, and a constant government spending in DNY (2021). We  first characterize the optimal commitment policy and then discuss the time-consistency property of the commitment solution as DNY (2021). 
%
%
%Our continuous-time formulation of the DNY (2021) economy is as follows. First, the agent has the following utility function: %consider an economy with isoelastic preferences over consumption $c$ and labor $n$, where 
%\begin{equation} \label{DNYUility}
%U(C, N) = \log C - \eta \frac{N^\gamma}{\gamma}.
%\end{equation}
%where $\eta>0$ and $\gamma\geq 1$. %which corresponds to a commonly used utility function for the evaluation of optimal fiscal policy. 
%Second,  government spending is exogenous and constant over time, i.e.,
%\begin{equation} \label{DNYG}
%G_t =G; \; t\in[0, \infty).
%\end{equation}
%Third, the given debt structure is %In addition, the time$-0$ government faces the following special initial debt
%\begin{equation} \label{DNYDebt}
%B_{0-} = 0;\;  b_{0-, t}=b_0\geq 0, t\in[0, T);\;  b_{0-, t}=b_1\geq 0, t\in[T, \infty)\, , 
%\end{equation}
%where $T$ is fixed. 
%Note that here  $T$ can be any finite value rather than one. 
%The  government inherits from the previous government  a constant stream of debt services at the rate of $b$ per unit over the time interval $[0, T)$  that it has to honor But there is no  short-term debt due at $0$: $B_{0-} = 0$. 
%
%%\WJ{A general assumption of Utility:
%%\begin{enumerate}
%%	\item $U(C, C+G)$ is concave in $C$ 
%%	\item Primary surplus $S(C)=\tau N-G =C+\frac{U_N}{U_C}(C+G)$ is concave in $C$ 
%%	\item MRS adjusted gap $\frac{S(C)-b}{C}$ is concave in $C$
%%\end{enumerate}
%%}
%
%\subsection{Laffer Curve Subject to Agent's Static Optimality} 
%
%
%\WJ{
%What is Laffer curve and when does Laffer curve critique matter? 
%We show government's maximization of tax revenue via choosing of tax rate is consistent with household's utility maximization. Therefore, Lucas-Stokey time-consistent requires the government always chose the low tax rate and stay in the right side of the Laffer curve given the amount to tax revenue the government wants to raise. We examine DNY's result in a model with short-term debt. We show that, unlike the change of sign of the Lagrangian multiplier for future government in DNY, the Lagrangian multiplier are always positive from period to period but the future government may renege because it may face non-unique positive Lagrangian multipliers.
%}
%
%\NW{Let $S_t$ denote the government's primary surplus at $t$: 
%	\begin{equation}
%	S_t  =  \tau_t N_t -G_t \, . 
%	\end{equation} 
%} 
%
%DNY (2021) define the following static Laffer curve problem maximizing the primary surplus $S$:
%\begin{equation}
%\max  \quad S %\tau N -G 
%\end{equation} 
%subject to the resource constraint 
%$C+ G = N$ and the agent's FOC: $1-\tau = -U_N/U_C$. 
%Using (\ref{DNYUility}), it is straightforward to show 
%\begin{equation} \label{laffer} 
%\tau N -G =(1-\eta C(C+G)^{\gamma-1})(C+G)-G = C(1-\eta(C+G)^\gamma)\, . 
%\end{equation} 
%
%\WJ{Define $C^{L}$ as the level of consumption that maximizes the primary surplus:
%	\begin{equation}
%	C^{L} = \arg \max_{C} \; C(1-\eta(C+G)^\gamma).
%	\end{equation} 
%	Therefore, $C^{L}$ is the level of consumption associated with the maximal tax revenue at the peak of the Laffer curve under tax rate $\tau^L$. }
%
%In DNY setup, the household's utility is higher when government chooses a consumption higher than $C^L$ for a given primary surplus level lying between 0 and $C^L(1-\eta(C^L+G)^\gamma)$. More specifically, for a given primary surplus level $S^* \in [0, C^L(1-\eta(C^L+G)^\gamma))$, there usually exists two consumption levels $C_1<C^L<C_2$ that can generate surplus $S^*$. Under the specific utility in DNY, we can show that $U(C_2, C_2+G)>U(C_1, C_1+G)$, implying the consumption level higher than $C^L$ usually produces higher utility for household. 
%
%\begin{figure}
%	\centering
%	\includegraphics[width=0.7\linewidth]{figure_d/Yared_Laffer}
%	\caption{The maximum of $U(C, N)=\log C-\eta \frac{N^\gamma}{\gamma}$ subject to resource constraint $C+G = N$ is achieved at $C^{FB} $ (blue dashed line) which satisfies $C^{FB}(C^{FB}+G)^{\gamma-1}=1/\eta$. On the other hand, $S^*=0$ is least primary surplus level such that there are two consumption levels, i.e., $C_1=0, (C_2+G)^\gamma = 1/\eta$ (red dashed line), that can generate $S^*=0$. For $C\leq C^{FB}$, $U(C, N)$ is an increasing function of $C$. If $C_2 \leq C^{FB}$, then we have 
%		$U(C_2, C_2+G)>U(C_1, C_1+G)$ for any $S^* \in [0, C^L(1-\eta(C^L+G)^\gamma))$.}
%	\label{fig:yaredlaffer}
%\end{figure}
%
%For general utility $U(C, N)$, given resource constraint $C+G=N$ and the agent's FOC $1-\tau = -U_N/U_C$, the primary surplus, as a function of consumption, is given
%\begin{equation}\label{psurplus}
%S(C) = \tau N -G = (1+\frac{U_N}{U_C})(C+G)-G = C + \frac{U_N}{U_C} (C+G)
%\end{equation}
%
%On the other hand, household's first-best policy solves the following
%\begin{eqnarray}
%\max_{C, N}  \; U(C, N) \\
%s.t.\quad  C+G = N
%\end{eqnarray}
%which suggests the first-best consumption $C^{FB}$ satisfies
%\begin{equation}\label{focfb}
%U_C+U_N = 0.
%\end{equation}
%
%Substituting (\ref{focfb}) into (\ref{psurplus}), we have for the first-best $C^{FB}$
%\begin{equation}
%S(C^{FB}) =(1+\frac{U_N}{U_C}\Big|_{C=C^{FB}})(C^{FB}+G)-G=-G<0.
%\end{equation} 
%Therefore, if both $U(C, C+G)$ and $S(C)$ are both concave functions of $C$, for $C_1<C_2$ that generate the same non-negative primary surplus $S(C_1)=S(C_2)\geq 0$, we have  $C_1<C_2<C^{FB}$ and  $U(C_1, C_1+G)<U(C_2, C_2+G)$.
%
%
%\NW{Time$-0$ government and time$-T$ government faces the same primary surplus for time $T$ onward. However, they may choose different level of consumption to support the primary surplus. }
%
%
%
%\subsection{Optimal commitment policy}
%\begin{proposition}
%	Given the utility (\ref{DNYUility}) and the initial debt (\ref{DNYDebt}),  the time$-0$ government's optimal commitment policy exists and is unique: 
%	
%	\begin{enumerate}
%		\item $C_{0, t}^* = C_0 1_{\{t\in[0, T)\}}+C_1 1_{\{t\in[T, \infty)\}} $,
%		where $C_0, C_1$ are the unique solutions to the following system of equations for some $\Lambda_0^*>0$
%		\begin{eqnarray}
%		&	\frac{1}{C_0}- \eta(C_0+G)^{\gamma-1}+\Lambda_0^*\left( \frac{b_0}{C_0^2}-\eta\gamma(C_0+G)^{\gamma-1} \right) = 0, \\
%		&	\frac{1}{C_1}- \eta(C_1+G)^{\gamma-1}+\Lambda_0^*\left(\frac{b_1}{C_1^2} -\eta\gamma(C_1+G)^{\gamma-1} \right) = 0, \\
%		&	\left(1-\eta(C_0+G)^{\gamma}\right)(1-e^{-\rho T})+\left(1-\eta(C_1+G)^{\gamma}\right)e^{-\rho T}=\frac{b_0}{C_0}(1-e^{-\rho T})+\frac{b_1}{C_1}e^{-\rho T}.	
%		\end{eqnarray}
%		\item Optimal labor supply and tax policy are given
%		\begin{eqnarray}
%		N_{0, t}^* & = & (C_0+G) 1_{\{t\in[0, T)\}}+(C_1+G) 1_{\{t\in[T, \infty)\}}\\
%		\tau_{0, t}^* &= & \left(1-\eta(C_0+G)^{\gamma-1}C_0\right) 1_{\{t\in[0, T)\}}+\left(1-\eta(C_1+G)^{\gamma-1}C_1\right)1_{\{t\in[T, \infty)\}}.
%		\end{eqnarray}
%		
%		\item The optimal value function for the time$-0$ government is
%		\begin{equation*}
%		L_0^* = \left(\log C_0 - \eta\frac{(C_0+G)^\gamma}{\gamma}\right)(1-e^{-\rho T})+\left(\log C_1 - \eta\frac{(C_1+G)^\gamma}{\gamma}\right)e^{-\rho T}.
%		\end{equation*}
%	\end{enumerate}	
%\end{proposition}
%
%\begin{proof}
%	\begin{itemize}
%		\item The implementability holds equality
%		\item Lagrangian is concave (both objective and constraint are concave). Therefore the solution is unique with a positive Lagrangian multiplier.
%		\item The nonlinear equation for the Lagrangian multiplier (\ref{eqPhi}) has a unique real root
%	\end{itemize}
%\end{proof}
%
%
%\begin{corollary}
%	Given the utility (\ref{DNYUility}) and the initial debt (\ref{DNYDebt}), if $\gamma=1$, then the time$-0$ government's optimal commitment policy is uniquely given as
%	\begin{eqnarray}
%	C_{0, t}^* & = & \frac{2\Lambda_0^*b}{-1+\sqrt{1+4b\eta\Lambda_0^*(1+\Lambda_0^*)}}1_{\{t\in[0, T)\}}+ \frac{1}{\eta(1+\Lambda_0^*)}1_{\{t\in[T, \infty)\}}; \\
%	N_{0, t}^* & = & C_{0, t}^*+G;  \\
%	\tau_{0, t}^* & = & 1-\eta C_{0, t}^*.
%	\end{eqnarray}
%	where $\Lambda_0^*>0$ is the unique root of the equation
%	\begin{eqnarray*}
%		\frac{-1+\sqrt{1+4b\eta\Lambda_0(1+\Lambda_0)}}{2\Lambda_0}(1-e^{-\rho T}) & =&  \left[ 1-\eta\left(\frac{2\Lambda_0 b}{-1+\sqrt{1+4b\eta\Lambda_0(1+\Lambda_0)}}+G \right) \right](1-e^{-\rho T}) \\
%		&& \; + \left[ 1-\eta\left(\frac{1}{\eta(1+\Lambda_0)}+G \right) \right]e^{-\rho T}.
%	\end{eqnarray*}
%\end{corollary}
%
%Define the the following debt level
%\begin{equation}
%\bar{b} = \max_{\widetilde{C}} \; \widetilde{C}\Big\{ [1-\eta(\widetilde{C}+G)^\gamma](1-e^{-\rho T})+(1-\eta G^\gamma)e^{-\rho T}  \Big\}.
%\end{equation}
%Then $\bar{b}$ represents the highest value of $b$ that is feasible.
%\WJ{Show that $\bar{b}$ is a decreasing function of $T$.}
%
%\begin{proposition}
%	There exists $b^* \in (0, \bar{b})$ such that the solution admits $C_1>C^{L}$ if $b<b^*$ and $C_1>C^{L}$ if $b>b^*$. Moreover, $b^*$ is an decreasing function of $T$.
%\end{proposition}
%
%\NW{All discrete-time models are with simple debt structure. Continuous-time models have more flexible debt structure.} 
%
%\subsection{Time Consistency in the sense of Lucas-Stokey}
%
%Next, we show that using short-term debt to roll over gap between the government's primary surplus and its due debt from time to time, we are able to make the optimal commitment policy time-consistent by restructuring short-term debt from $T-$ to $T$.
%
%\begin{proposition}
%	The short-term debt for time $T-$ government satisfies
%	\begin{equation}
%	B_{T-} = -\frac{1}{\rho}(e^{\rho T}-1)[S({C}_0)-b_0]+e^{\rho T}B_0
%	\end{equation} 
%	where $S({C}_0) = C_0(1-\eta(C_0+G)^\gamma)$ is the primary surplus.
%\end{proposition}
%
%Following Lucas definition of time-consistency, we try to show that the time $T-$ government can restructure its short-term debt from $B_{T-}$ to be $\widehat{B}_{T-}$ and leave it to the future government, then the time $T$ government will still follow time $0$ government's commitment policy.
%
%\begin{proposition}
%	Time $T-$ government restructure its short-term debt to be $\widehat{B}_{T-} = \frac{B_{T-}}{C_0}C_1$. If $b_0<b^*$, the time $T$ government's optimal consumption is $\widehat{C}_{T, t}=C_1, t\geq T$. Therefore, the commitment policy is time-consistent. 
%\end{proposition}
%
%\begin{proof}
%	After time $T-$ government's debt restructuring, time $T$ government faces debt $\widehat{B}_{T-}=\frac{B_{T-}}{C_0}C_1$ and long-term debt $b_{T-, t}=b_{0-, t}=b_1, t\geq T$. 
%	
%	Time $T-$ chooses a policy $\{ \widehat{C}_{T, t}, \widehat{N}_{T, t}; t\in[T, \infty)\} $ to solve the following optimization problem
%	\begin{eqnarray*}
%		&\max \int_{T}^{\infty} e^{-\rho (t-T)}U(\widehat{C}_{T, t}, \widehat{N}_{T, t}) dt.  \\
%		&\text{s.t.} \quad \int_{T}^{\infty} e^{-\rho (t-T)}(\widehat{C}_{T, t}U_C+\widehat{N}_{T, t}U_N)\geq \int_{T}^{\infty} e^{-\rho (t-T)} b_{0-, t}U_C dt+ \widehat{B}_{T-}U_C(\widehat{C}_{T, T}, \widehat{N}_{T, T}).
%	\end{eqnarray*}
%	
%	The optimal solution $\widehat{C}^*_{T, t}=\widehat{C}_1$ satisfies
%	\begin{eqnarray}
%	& \frac{1}{\widehat{C}_1}- \eta(\widehat{C}_1+G)^{\gamma-1}+\Lambda_T^*\left(\frac{b_1}{\widehat{C}_1^2} -\eta\gamma(\widehat{C}_1+G)^{\gamma-1} \right) = 0, \\
%	& \frac{1}{\rho}\left(1-\eta(\widehat{C}_1+G)^{\gamma}\right)=\frac{1}{\rho}\frac{b_1}{\widehat{C}_1}+\frac{\widehat{B}_{T-}}{\widehat{C}_1}.	
%	\end{eqnarray}
%	where $\Lambda_T^*$ is positive.
%	
%	Actually, we can show that equations system (32)-(33) admits a unique solution satisfying $\Lambda_T^*>0$. In fact, $S(\widehat{C}_1)$ and $U(\widehat{C}_1, \widehat{C}_1+G)$ both are concave functions of $\widehat{C}_1$, The first-best of $U(\widehat{C}_1, \widehat{C}_1+G)$ is higher that any consumption with non-negative primary surplus (i.e., Figure 1). For any given primary surplus level $b_1+\rho \widehat{B}_{T-}$, there exists two consumption levels that can support this given primary surplus. For these two consumption level, both has positive marginal utility and both has negative marginal MRS adjusted funding gap (i.e., $\frac{\partial}{\partial C}\left(\frac{S(C)-b_1}{C}\right)<0$). Optimality implies that the higher consumption that supports primary surplus $b_1+\rho \widehat{B}_{T-}$ is the optimal solution. 
%	
%	On the other hand, using time $0$ government's commitment policy, we can verify that $\widehat{C}_1 = C_1$ is the solution of (32)-(33) which admits a positive $\Lambda_T^*$ satisfying $\Lambda_T^*=\Lambda_0^*$.
%	
%\end{proof}
%
%
%
%
%
%
%\subsubsection{The Argument in DNY (2021)} 
%
%In this section, we show that If short-term is not traded in the market, then Yared's argument works.
%If short-term debt doesn't exists in the market, time $T-$ government needs to restructure its accumulated gap ${B}_{T-}$ to a long-term debt $\widehat{b}_{T-, t}, t\geq T$ for the time $T$ government. The time $T$ government chooses a policy $\{ \widehat{C}_{T, t}, \widehat{N}_{T, t}; t\in[T, \infty)\} $ to solve the following optimization problem
%\begin{eqnarray*}
%	&\max \int_{T}^{\infty} e^{-\rho (t-T)}U(\widehat{C}_{T, t}, \widehat{N}_{T, t}) dt.  \\
%	&\text{s.t.} \quad \int_{T}^{\infty} e^{-\rho (t-T)}(\widehat{C}_{T, t}U_C+\widehat{N}_{T, t}U_N)\geq \int_{T}^{\infty} e^{-\rho (t-T)} \widehat{b}_{T-, t}U_C dt+ \widetilde{B}_{T-}U_C(\widehat{C}_{T, T}, \widehat{N}_{T, T}).
%\end{eqnarray*}
%
%\begin{proposition}
%	If Lucas time-consistency restructuring method exists, then this debt restructuring method is unique, i.e.,
%	\begin{equation}
%	\widehat{b}_{T-, t} =\widehat{b}_1 = b_1+\rho \frac{B_{T-}}{C_0}C_1
%	\end{equation}
%	where $B_{T-}=-\frac{1}{\rho}(e^{\rho T}-1)[S({C}_0)-b_0]+e^{\rho T}B_0$ is the accumulated gap for time $T-$ government.
%\end{proposition}
%
%\begin{proposition}
%	If $b_0>b^*$, the time $T$ government faces $\widehat{b}_{T-, t}, t\geq T$ will renege. If $b_0<b^*$, the optimal commitment policy is time-consistent.
%\end{proposition}
%
%\begin{proof}
%	Time $T$ government faces $\widehat{b}_{T-, t}, t\geq T$ and its optimal policy satisfies
%	
%	\begin{eqnarray}
%	& \frac{1}{\widehat{C}_1}- \eta(\widehat{C}_1+G)^{\gamma-1}+\Lambda_T^*\left(\frac{\widehat{b}_1}{\widehat{C}_1^2} -\eta\gamma(\widehat{C}_1+G)^{\gamma-1} \right) = 0, \\
%	& \frac{1}{\rho}\left(1-\eta(\widehat{C}_1+G)^{\gamma}\right)=\frac{1}{\rho}\frac{b_1}{\widehat{C}_1}+\frac{\frac{B_{T-}}{C_0}C_1}{\widehat{C}_1}.	\\
%	& \widehat{C}_1\left(1-\eta(\widehat{C}_1+G)^{\gamma}\right)=\widehat{b}_1
%	\end{eqnarray}
%	
%	Note that equation (36) and equation (33) are the same. However, equation (35) is different from equation (32). 
%	
%	Actually, $b^*$ is the level of $b_0$ such that the commitment policy satisfies $C_1= C^{L}$ where 
%	\begin{equation}
%	C^{L} = \arg \max_{C} S(C)
%	\end{equation}
%	
%	
%\end{proof}
%
%
%\begin{figure}
%	\centering
%	\includegraphics[width=0.7\linewidth]{figure_d/Yared_Laffer_2}
%	\caption{When initial debt $b_0$ is too high, time $0$ government chooses (1) a primary surplus level for $T$ onward and (2) a smaller consumption level that supports this primary surplus, to generate higher MRS adjusted primary surplus. This smaller consumption level is associated with a negative marginal MRS adjusted funding gap (i.e., $\frac{\partial}{\partial C}\left(\frac{S(C)-b_1}{C_1}\right)<0$). Without short-term debt, time $T$ government facing restructured debt $\widehat{b}_1$ left by time $T-$ government, optimally chooses the consumption level which supports the same primary surplus level chosen by time $0$ government for $T$ onward but has negative marginal MRS adjusted gap.  }
%	\label{fig:yaredlaffer2}
%\end{figure}
%
%\begin{table}[!h]
%	\caption{Key values of Figure 2 ($b_0=50, b_1=0$). $S(C)=C(1-\eta(C+G))$. $C_1$ and $\widehat{C}_1$ are two roots of $S(C)=\widehat{b}_1=\rho \widehat{B}_{T-}$, and they satisfy $C_1\widehat{C}_1=\widehat{b}_1$. In addition $\frac{S(C_1)-b_1}{C_1}=\frac{\widehat{b}_1}{C_1}=\widehat{C}_1$. In our model with short-term debt, one can use short-term debt to roll over government's funding gap, i.e., $S(C_{0, t})-b_{0, t}$, and restructure the short-term debt and leave it to future government. In this case, time $T-$ government accumulates short-term debt $B_{T-}=91.9356$ and restructures it to be $\widehat{B}_{T-}=1.5748$ and leave it to time $T$ government. Facing $\widehat{B}_{T-}=1.5748$ and $b_1=0$, time $T$ government solve the optimization problem again, and find there two positive Lagrangian multiplies $\Lambda_{T}^*=7.7703, \quad 0.4597$ satisfies the FOCs of the optimization problem. Lagrangian multiplier represents the shadow price of \NW{marginal value of MU adj funding gap}. In contrast, in DNY(2021), there is no short-term debt, time $T-$ government accumulates funding gap $B_{T-}=91.9356$ and restructures it to be long-term debt $\widehat{b}_{1}=0.0787$ and leave it to time $T$ government. The future government facing only $\widehat{b}_{1}=0.0787$ solves the optimization problem, and find there is a unique positive Lagrangian multiplier satisfying FOCs. For both models, time $T$ government will renege and choose a higher consumption level after $T$ which supports the same primary surplus level with the consumption level of commitment policy for time $T$ onward.}
%	\begin{tabular}{c|l|l}
%		\hline 
%		Laffer consumption&  \multicolumn{2}{c}{$C^{laffer} = 0.4 $}	  \\
%		First-best consumption &  \multicolumn{2}{c}{$C^{FB} = 1  $ }  \\
%		Commitment policy	&  \multicolumn{2}{c}{$C_0= 6.71   ;\quad  C_1 = 0.1149 $
%		} \\ 
%		&\multicolumn{2}{c}{$N_0= 6.91    ;\quad N_1 = 0.1349 $
%		} \\
%		&\multicolumn{2}{c}{$\tau_0= -5.71    ;\quad \tau_1 = 0.8851 $
%		} \\
%		Value with $C_0, C_1$ at $0$ & \multicolumn{2}{c}{ -52.0320}   \\
%		Two consumption levels&\multicolumn{2}{c}{$C_1= 0.1149    ;\quad \widehat{C}_1 = 0.6851 $
%		} \\
%		supporting primary surplus & \multicolumn{2}{c}{} \\
%		
%		\hline
%		&   Ours (with short-term debt)  &  Yared (no short-term debt) \\\hline
%		Accumulated funding gap &  $B_{T-} = 91.9356 $ & $B_{T-} = 91.9356 $\\
%		Debt restructuring  & $\widehat{B}_{T-} = \frac{B_{T-}}{C_0}C_1= 1.5748$  & $\widehat{b}_1 = \rho \frac{B_{T-}}{C_0}C_1 = 0.0787$ \\
%		&& \\
%		MU adj funding gap at $C_1$ & $\frac{S(C_1)-b_1}{C_1} =0.6851 $ &  $\frac{S(C_1)-\widehat{b}_1}{C_1} = 0$ \\
%		& $\frac{\partial}{\partial C_1}\frac{S(C_1)-b_1}{C_1} = -1 $ &  $ \frac{\partial}{\partial C_1}\frac{S(C_1)-\widehat{b}_1}{C_1} = 4.9573$ \\
%		Lagrangian multiplier at $C_1$  & $\Lambda_{T}^* = 7.7703 $    & $\Lambda_{T}^* = -1.5524 $  \\
%		Value with $C_1$ at $T$ & -49.5660 & -49.5660   \\
%		
%		&& \\
%		MU adj funding gap at $\widehat{C}_1$ & $\frac{S(\widehat{C}_1)-b_1}{\widehat{C}_1} = 0.1149$ &  $\frac{S(\widehat{C}_1)-\widehat{b}_1}{\widehat{C}_1} = 0$  \\
%		& $\frac{\partial}{\partial \widehat{C}_1}\frac{S(\widehat{C}_1)-b_1}{\widehat{C}_1} = -1$ &  $\frac{\partial}{\partial \widehat{C}_1}\frac{S(\widehat{C}_1)-\widehat{b}_1}{\widehat{C}_1} =-0.8323 $  \\
%		Lagrangian multiplier at $\widehat{C}_1$  & $\Lambda_{T}^* = 0.4597$   & $\Lambda_{T}^* = 0.5524$  \\
%		Value with $\widehat{C}_1$ at $T$ & -25.2662 & -25.2662   \\
%		%Value with $C_0, \widehat{C}_1$ at $0$ & -28.9172 & -28.9172   \\
%		\hline 
%	\end{tabular} 
%\end{table}
%
%\clearpage
%\begin{figure}
%	\centering
%	\includegraphics[width=0.7\linewidth]{figure_d/Yared_Laffer_3}
%	\caption{When initial debt $b_0$ is not too high, time $0$ government optimally chooses (1) a primary surplus level for $T$ onward and (2) a higher consumption level that supports this primary surplus, to generate enough MRS adjusted primary surplus. Without short-term debt, time $T$ government facing restructured debt $\widehat{b}_1$ left by time $T-$ government, optimally chooses the same consumption level as the commitment policy. The negative marginal MRS adjusted gap for the consumption level ensures its optimality for time $T$ government. }
%	\label{fig:yaredlaffer3}
%\end{figure}
%
%\begin{table}[!h]
%	\caption{Key values of Figure 3 ($b_0=10, b_1=0$). $S(C)=C(1-\eta(C+G))$. $C_1$ and $\widehat{C}_1$ are two roots of $S(C)=\widehat{b}_1=\rho \widehat{B}_{T-}$, and they satisfy $C_1\widehat{C}_1=\widehat{b}_1$. In addition $\frac{S(C_1)-b_1}{C_1}=\frac{\widehat{b}_1}{C_1}=\widehat{C}_1$.  In this case, time $T-$ government accumulates short-term debt $B_{T-}=14.5539$ and restructures it to be $\widehat{B}_{T-}=2.9694$ and leave it to time $T$ government. Facing $\widehat{B}_{T-}=2.9694$ and $b_1=0$, time $T$ government solve the optimization problem again, and find there two positive Lagrangian multiplies $\Lambda_{T}^*=0.9709, \quad 2.4173$ satisfies the FOCs of the optimization problem. Lagrangian multiplier represents the shadow price of \NW{marginal value of MU adj funding gap}. In contrast, in DNY(2021), there is no short-term debt, time $T-$ government accumulates funding gap $B_{T-}=14.5593$ and restructures it to be long-term debt $\widehat{b}_{1}=0.1485$ and leave it to time $T$ government. The future government facing only $\widehat{b}_{1}=0.1485$ solves the optimization problem, and find there is a unique positive Lagrangian multiplier satisfying FOCs. For both models, time $T$ government will not renege and choose the same consumption level with that of commitment policy for time $T$ onward.   }
%	\begin{tabular}{c|l|l}
%		\hline 
%		Laffer consumption&  \multicolumn{2}{c}{$C^{laffer} = 0.4 $}	  \\
%		First-best consumption &  \multicolumn{2}{c}{$C^{FB} = 1  $ }  \\
%		Commitment policy	&  \multicolumn{2}{c}{$C_0= 2.4877   ;\quad  C_1 = 0.5074 $
%		} \\ 
%		&\multicolumn{2}{c}{$N_0= 2.6877    ;\quad N_1 = 0.5274 $
%		} \\
%		&\multicolumn{2}{c}{$\tau_0= -1.4877    ;\quad \tau_1 = 0.4926 $
%		} \\
%		Value with $C_0, C_1$ at $0$ & \multicolumn{2}{c}{-28.0985}   \\
%		Two consumption levels &\multicolumn{2}{c}{$C_1= 0.5074   ;\quad \widehat{C}_1 = 0.2926 $
%		} \\
%		supporting primary surplus & & \\
%		\hline
%		&   Ours (with short-term debt)  &  Yared (no short-term debt) \\\hline
%		Accumulated funding gap &  $B_{T-} = 14.5593 $ & $B_{T-} = 14.5593 $\\
%		Debt restructuring  & $\widehat{B}_{T-} = \frac{B_{T-}}{C_0}C_1= 2.9694$  & $\widehat{b}_1 = \rho \frac{B_{T-}}{C_0}C_1 = 0.1485$ \\
%		&& \\
%		MU adj funding gap at $C_1$ & $\frac{S(C_1)-b_1}{C_1} =0.2926 $ &  $\frac{S(C_1)-\widehat{b}_1}{C_1} = 0$ \\
%		& $\frac{\partial}{\partial C_1}\frac{S(C_1)-b_1}{C_1} = -1 $ &  $ \frac{\partial}{\partial C_1}\frac{S(C_1)-\widehat{b}_1}{C_1} = -0.4231$ \\
%		Lagrangian multiplier at $C_1$  & $\Lambda_{T}^* = 0.9709 $    & $\Lambda_{T}^* = 2.2932 $  \\
%		Value with $C_1$ at $T$ & -27.7169 & -27.7169   \\
%		
%		&& \\
%		MU adj funding gap at $\widehat{C}_1$ & $\frac{S(\widehat{C}_1)-b_1}{\widehat{C}_1} = 0.5074$ &  $\frac{S(\widehat{C}_1)-\widehat{b}_1}{\widehat{C}_1} = 0$  \\
%		& $\frac{\partial}{\partial \widehat{C}_1}\frac{S(\widehat{C}_1)-b_1}{\widehat{C}_1} = -1$ &  $\frac{\partial}{\partial \widehat{C}_1}\frac{S(\widehat{C}_1)-\widehat{b}_1}{\widehat{C}_1} =0.7345 $  \\
%		Lagrangian multiplier at $\widehat{C}_1$  & $\Lambda_{T}^* = 2.4173$   & $\Lambda_{T}^* = -3.2938$  \\
%		Value with $\widehat{C}_1$ at $T$ & -34.4296 & -34.4296  \\
%		%Value with $C_0, \widehat{C}_1$ at $0$ & -34.4831 & -34.4831  \\
%		\hline 
%	\end{tabular} 
%\end{table}
%
%
%\clearpage
%\section{Numerical results}
%
%
%\begin{figure}[h!]
%	\centering
%	\includegraphics[width=\linewidth, height=0.6\textheight]{figure_d/commit_policy}
%	\caption{Optimal commitment policy. Let $b_{0-, t}=b_01_{\{t< T\}}+b_11_{\{t\geq T\}}$. One inherits debt $b_0=50, b_1=0$ and the other inherits debt $b_0=10, b_1=0$. Both governments have the same initial short-term debt $B_{0-}=0$ and the same government spending $G_t=0.2$. Other parameter values are $\rho=0.05, \eta=1, T=1$. For the government with $b_0=50$, its consumption after $T=1$ is lower than the Laffer consumption, which implies Lucas time-consistency is not possible. \label{fig:commitpolicy}}
%	
%\end{figure}
%
%
%
%\begin{figure}[h!]
%	\centering
%	\includegraphics[width=\linewidth, height=0.6\textheight]{figure_d/commit_policy_eta}
%	\caption{Optimal commitment policy with varying $\eta$. Initial long-term debt $b_{0-, t}=b_01_{\{t< T\}}+b_11_{\{t\geq T\}}$ satisfies $b_0=10, b_1=0$. Both governments have the same initial short-term debt $B_{0-}=0$ and the same government spending $G_t=0.2$, but have different labor costs $\eta=2$ v.s. $\eta=1$ Other parameter values are $\rho=0.05, T=1$. The household with higher $\eta$, supplies lower labor and consumes less. Government facing higher $\eta$ may be more likely to renege after $T$ since its tax rate after $T$ is above the peak of Laffer curve.  \label{fig:commitpolicyeta}}
%	
%\end{figure}
%
%\begin{figure}[h!]
%	\centering
%	\includegraphics[width=\linewidth, height=0.6\textheight]{figure_d/commit_policy_G}
%	\caption{Optimal commitment policy with varying $G$. Initial long-term debt $b_{0-, t}=b_01_{\{t< T\}}+b_11_{\{t\geq T\}}$ satisfies $b_0=10, b_1=0$. Both governments have the same initial short-term debt $B_{0-}=0$, but have different government spending $G=0.6$ v.s. $G=0.2$. Other parameter values are $\rho=0.05, \eta=1, T=1$. The government facing higher spending tax less before $T$ and tax more after $T$, which motivates household to consume more before $T$ and consume less after $T$. The high spending government's after $T$ tax rate is above the peak of Laffer curve, implying Lucas-Stokey time-consistent may not be feasible.  \label{fig:commitpolicyG}}
%	
%\end{figure}
%
%\clearpage
%\section{Conclusions}


\section{Piece-wise Continuous Debt Maturity Distribution}

In this section, we extend the continuouse debt maturity distribution to piece-wise continous distribution. Precisely, we assume the following initial debt and government spending structure.

\begin{assumption}\label{assump2}
Both the initial term debt $\{b_{0, t}; t\geq0 \}$ and government spending $\{G_t; t\geq0\}$ are C\`adl\`ag sequencees, i.e., $\lim_{s\downarrow t} b_{0, s} = b_{0, t}$ and $\lim_{s\downarrow t} G_{s} = G_{t}$ for all $t\geq 0$, and $\lim_{s\uparrow t} b_{0, s}$ and $\lim_{s\uparrow t} G_{s}$ exists for all $t>0$. 
\end{assumption} 

Given assumption \ref{assump2}, we focus on the allocations that both consumptions and labor supplies are piece-wise continuous on time. Specifically, we impose the following condition

\begin{condition}\label{cond2}
Under assumption \ref{assump2},	we focus on allocations $\{C_{0, t}, N_{0, t}; t\geq 0\}$ satisfying that, for all $t\geq 0$,
\begin{equation} \label{smooth2}
\lim_{s\downarrow t} C_{0, s} = C_{0, t}\,.
\end{equation}
\end{condition}

As we show in the Appendix XX, under the condition \ref{cond2}, if there is no jump in consumption at time $t\geq 0$, then the government commits to its short-term debt, i.e., $\lim_{s\downarrow t}B_sq_{0, s}=B_tq_{0, t}$. However, if there is jump in consumptions at time $t$, i.e., $C_{t-}\neq C_{t}$. Then commitment to the short-term debt condition requires a time $t$ government to adjust it's inherited short-term debt $B_{t}$ such that short-term debt holder's purchasing power has not been manipulated, i.e.,
\begin{equation}\label{cond2_add}
B_{t}q_{0, t}=B_{t-}q_{0, t-}\,.
\end{equation}


Given condition \ref{cond2} and the additional commitment condition \ref{cond2_add}, and assuming all government deficits are financed by short-term debt, we can show that the MRS adjusted government short-term debt, i.e., $X_t$ is continuous given by equation (\ref{dy}). Moreover,because of the condition \ref{cond1} and that $\{b_{0, t}, t\geq 0\}$ is C\`adl\`ag sequence, the consumption is smooth at time $0$, i.e.,
$C_{0, 0} = \lim_{s\downarrow 0} C_{0, s}$. Proposition 5 for the optimal commitment policy still holds for the piece-wise continuous debt maturity distribution. 


Moreover, to implement time-consistency for this extension, a government only needs to follow two rules
\begin{itemize}
\item If there is no jump of allocation at time $t$, the government simply follow the rescheduling policies described in section 5 to the future government.
\item If there is jump of allocation at time $t$, to ensure the time $t$ government to hornor the purchasing power of accumulative primary deficit up to time $t-$, $\Pi_{t-}$, the time $t-$ government needs to ensure $\Pi_{t-}q_{0, t-}=\Pi_{t}q_{0, t}$. If the government $t-$ chooses only short-term debt to meet its budget constraint, i.e., $\Pi_{t-}=B_{t-}$, it needs to leave $B_t$ that satisfies equation \ref{cond2_add} to the time $t$ government. 
\end{itemize}
We leave details of the derivation to the Appendix XX. 




\newpage 



\section{Debt Rescheduling Irrelevance}

For a commitment government,
once a feasible allocation (i.e., Ramsey plan) is chosen at time$-0$, then its commitment requires the future government to comply with this feasible allocation at any future time $t$ ($t\geq 0$) by rescheduling debt structures to satisfy the dynamic budget constraint at time $t$.  
In this section, we propose the use of short-term debt to reschedule government debt and show the irrelevance of rescheduling policy. 

\paragraph{Complete debt market} We assume the government faces a complete debt market such that the government can issue debt with all kinds of maturities, including both short-term debt and term debt,\footnote{Lucas-Stokey and many others only consider term debt. AMSS consider short-term debt, however, it neglects the term debt.} to finance its primary deficits. Let $B_t$ denote the short-term debt balance sequence for the government. The stock  $B_t$  is a ``demand deposit''  that earns instantaneous interest at the ``force of interest'' $r_t$.
%The function $\{r_t\}_{t \geq 0}$ is determined as part of a competitive equilibrium.   
Over a small   time interval $dt$,
payments  $r_{t} B_{t} dt $ are due on  instantaneous  debt $B_t$.
Let  $\{ b_{u, t}; u\leq t\}$ denote a flow of promised payouts that   the government issues at  time $u$. % due at time $t$. XXXX
%implemented with all future government at date $u<\nu\leq t$ choosing $b_{\nu, t} = b_{u, t}$.   
{We follow Lucas and Stokey (1983) and   Angeletos (2002) by 
assuming that  the government continuously  buying back all claims on that flow  and then issuing fresh claims. At instant $u > 0$, the government issues a new coupon flow
$ \{b_{u,t}\}_{t \geq u}.$ }At time $t$,  the government issues net new  coupons  $\{\partial_t b_{t, s}, s\geq t\}$ due at time $s$.  Thus,  the time$-t$ market value of net new coupon flows  issued at time $t$ is $ \left(\int_{s\geq t} q_{t, s}\partial_tb_{t, s} ds \right) dt$.  % 


%The stock  $B_t$  is a ``demand deposit''  that earns instantaneous interest at the ``force of interest'' $r_t$.
%The function $\{r_t\}_{t \geq 0}$ is determined as part of a competitive equilibrium.   
%Over a small   time interval $dt$,
%payments  $r_{t} B_{t} dt $ are due on  instantaneous  debt $B_t$. The time-$t$ price of one unit of the consumption good at  time $s\geq t$ is %. In equilibrium:  


\paragraph{Commitment on short-term debt} The avalibility of shor-term debt in debt market such that the government can use it to finance its primary deficit requires that the government commits to its short-term debt. That is, the government shall ensure the purchasing power of short-term debt will not be manipulated. Therefore, we impose the following condition which ensures such commitment.

\begin{condition}\label{cond1}
Under assumption \ref{assump1}, we focus on allocations $\{C_{0, t}, N_{0, t}; t\geq 0\}$ satisfying that, for all $t\geq 0$,
\begin{equation} \label{smooth}
\lim_{s\to t} C_{0, s} = C_{0, t}\,.
\end{equation}
\end{condition}

To illustrate, this condition, together with the first-order condition (\ref{foch2}), suggests that the debt price is a smooth function of time:
\begin{equation*}
\lim_{s\to t} q_{0, s} = q_{0, t} \,.
\end{equation*}
It ensures that, the purchasing power of holding $B$ units of short-term debt for an instantaneous time period from $t$ to $t+dt$ doesn't change because their time$-0$ values satisfies $\lim_{dt\to 0} Bq_{0, t+dt}=Bq_{0, t}$.



\paragraph{Dynamic budget constraint} Over the period $[t, t+dt]$, the government uses both short-term debt and term debt to finance its primary deficit such that the following dynamic budget constraint holds. \WJ{We assume that the government  continuously repurchases all outstanding flow coupons and then issues new ones.}
%A coupon flow rescheduling policy is a choice of $\{\{\partial_t b_{t, s}\}_{s\geq t}\}_{t=0}^\infty$. 
\begin{eqnarray} \label{dynamicbudgetconstraint1}
d\Pi_t-r_t\Pi_tdt \leq \left( dB_t-r_tB_tdt \right)  + \left( \left(\int_{s\geq t} q_{t, s}\partial_tb_{t, s} ds \right) dt - \left(\int_{0}^t \partial_u b_{u, t} du\right)dt \right)
\end{eqnarray}
The left hand side is the primary deficit over this $dt$ period. The first term of the right hand side is the net issuance of short-term debt over this period, and the second term is the financing over this $dt$ period from additional issuance of term debt netting repayment to the due debt that additionally issued in the past .


% For example, if the government at time $u$ issues debt due at date $t>u$ of size $b_{u, t}$, which then it holds to maturity without issuing additional debt, then this can equivalently implemented with all future government at date $u<\nu\leq t$ choosing $b_{\nu, t} = b_{u, t}$.  
According to equation (\ref{eqPi}), we rewrite the government's time $t$ flow budget  constraint as follows %pay off the inherited debt from time before $t\leq v$ and satisfies
{\begin{eqnarray} \label{dynamicbudgetconstraint2}
\tau_tN_{0, t} dt +   \left(\int_{s\geq t} q_{t, s}\partial_tb_{t, s} ds \right) dt + dB_t \geq G_t dt + r_tB_t dt+  b_{t, t}dt \, . 
\end{eqnarray} 
}
where $b_{t, t}$ is post-rescheduling  time $t$ coupon flow
\begin{equation} \label{btt}
b_{t, t} = b_{0, t} + \int_{0}^t \partial_u b_{u, t} du\, .
\end{equation}


\begin{example}{(Rescheduling policy in Lucas Stokey) }
In Lucas-Stokey, the debt rescheduling policy uses only term debt. That is, a coupon flow rescheduling policy of $\{\partial_t b_{t, s};s\geq t, t\geq 0 \}$ and $\{B_t=0; t\geq 0\}$ is used to satisfy the dynamic budget constraint, i.e.,
{\begin{eqnarray} \label{dynamicbudgetconstraint3}
&	d\Pi_t-r_t\Pi_tdt \leq \left( \left(\int_{s\geq t} q_{t, s}\partial_tb_{t, s} ds \right) dt - \left(\int_{0}^t \partial_u b_{u, t} du\right)dt \right)\\
&	\tau_tN_{0, t} dt +   \left(\int_{s\geq t} q_{t, s}\partial_tb_{t, s} ds \right) dt  \geq G_t dt +  b_{t, t}dt \, . 
\end{eqnarray} 
}
Dynamic budget constraint requires that such pure term debt rescheduling policy to persistently adjust its term debt for all maturities.
\end{example}

\begin{example}(Rescheduling policy with only short-term debt)
Our complete debt market assumption suggests a simple debt rescheduling policy---the use of only short-term debt. That is, the government sets  net new claims  
$b_{u, t} = b_{0, t}$ for all $0<u< t$ , so that  $\{\partial_t b_{t, s}; s\geq t, t\geq 0\}$ is  identically zero for all $t \geq 0$ and chooses $B_t$ such that
{\begin{eqnarray} \label{dynamicbudgetconstraint4}
& d\Pi_t-r_t\Pi_tdt \leq \left( dB_t-r_tB_tdt \right)  \\
&	\tau_tN_{0, t} dt  + dB_t \geq G_t dt + r_tB_t dt+  b_{0, t}dt \, . 
\end{eqnarray} 
}
The benefit of such pure short-term debt rescheduling policy at least threefold. First, it does not touch the term debt balances that passes to the future government. Second, as we will show later, use of only short-term debt allows us to recursify this Ramsey problem and solve the optimal commitment policy via a dynamic programming method. Third, we will further show that such rescheduling policy suggests a simple way implement time consistency in the sense of Lucas-Stokey.

\end{example}


%\WJ{We assume that the government  continuously repurchases all outstanding flow coupons and then issues new ones.}
%A coupon flow rescheduling policy is a choice of 
%$\{\{\partial_t b_{t, s}\}_{s\geq t}\}_{t=0}^\infty$. 
%The 
%To implement  a  no-reschedule policy,  the government at date $u<\nu\leq t$  sets   net new claims  
%$b_{\nu, t} = b_{u, t}$, so that  $\{\partial_t b_{t, s}, s\geq t\}$ is  identically zero for all $t \geq 0$.

Next, we show that the rescheduling policy is independent of the implementability condition (\ref{impletconstraint}).
We assume the government runs no bubble such that the following transversality condition is satisfied 
\begin{equation}\label{condB}
\lim_{t\to \infty} q_{0, t} B_t =0. 
\end{equation}

\begin{proposition}\label{propo::tom1}
The coupon flow rescheduling  policy does not affect the implementability constraint.
\end{proposition}

\begin{proof}

%Define
%\begin{equation}\label{sdf0}
%\frac{M_{0, t}}{M_{0, 0}} = q_{0, t} = e^{-\rho t} \frac{U_{C, t}}{U_{C, 0}}\, .
%\end{equation}
Substitute (\ref{btt}) into the dynamic budget constraint (\ref{dynamicbudgetconstraint2}), multiply both sides by $q_{0, t},$  
and integrate (\ref{dynamicbudgetconstraint2})  to get
{
\begin{eqnarray} \label{imc0}
& & \int_{0}^{\infty} \left(\tau_tN_t-G_t \right)q_{0, t}dt +  \int_{0}^{\infty}  \int_{s\geq t} q_{0, s}\partial_tb_{t, s} ds dt + \int_{0}^{\infty} q_{0, t} dB_t \notag \\
& \geq &  \int_{0}^{\infty} r_tB_t q_{0, t} dt+ \int_{0}^{\infty} q_{0, t}\left(b_{0, t} + \int_{0}^t \partial_u b_{u, t} du \right)dt \, .
\end{eqnarray}
}
Apply Fubini's theorem to obtain
\begin{equation} \label{fubini}
\int_{0}^{\infty}  \int_{s\geq t} q_{0, s}\partial_tb_{t, s} ds dt = \int_{0}^{\infty} q_{0, t} \int_{0}^t \partial_u b_{u, t} du dt \,.
\end{equation}
Use  (\ref{foch2}) and (\ref{qprice}) to deduce
\begin{eqnarray} \label{integralB}
\int_{0}^{\infty} q_{0, t} dB_t - \int_{0}^{\infty} r_tB_t q_{0, t} dt & = &  \int_{0}^{\infty} e^{-\rho t}\frac{U_{C, t}}{U_{C, 0}} dB_t - \int_{0}^{\infty} r_tB_t e^{-\rho t}\frac{U_{C, t}}{U_{C, 0}} dt \nonumber \\
& = & \frac{1}{U_{C, 0}}\int_{0}^{\infty} d\left(e^{-\rho t}B_t U_{C, t}  \right)  \nonumber\\
& = &   \lim_{t\to \infty} e^{-\rho t}\frac{U_{C, t}}{U_{C, 0}} B_t - B_0    \nonumber\\
& = & 0  \, , 
\end{eqnarray}
where the last equality imposes condition (\ref{btt}) and $B_0=0$ because there is no initial short-term debt. 

Substitute (\ref{fubini}) and (\ref{integralB}) into (\ref{imc0}) to obtain the implementability constraint that competitive equilibrium 
allocations must obey:
\begin{equation}
%\int_{0}^{\infty} \left(C_{0, t} U_C + (C_{0, t}+G_t)U_N \right) dt = \int_{0}^{\infty} b_{0, t}U_C dt \,.
\int_{0}^{\infty} e^{-\rho t} ( C_{0, t}U_{C, t}+(C_{0, t}+G_t)U_{N, t} )  dt  \geq   \int_{0}^{\infty} e^{-\rho t}b_{0, t} U_{C, t} dt \, .  \label{impletconstraint1}
\end{equation}
Rescheduling  policy $\{\partial_t b_{t, s};s\geq t, t\geq 0 \}$ does not appear in 
%the implementability constraint
\eqref{impletconstraint}.
\end{proof}



\newpage 


\subsection{Dynamics of $B_t$ and $X_t$}
%Using (\ref{dB}), 
%Let  a state variable $X_t$ which is the MRS adjusted short-term debt, i.e. 
%\begin{equation}
%X_t = B_t U_{C}(C_{0, t}, N_{0, t}).
%\end{equation}


Proposition \ref{propo::tom1} implies that  in posing a Ramsey problem we can use only short-term debt to rescheduling debt to meet the dynamic budget constraint. That is, we are free to set coupon-flow management policy
\begin{equation*}
\partial_tb_{t, s}   = 0\,, s\geq t, t\geq 0 ,
\end{equation*}
and 
\begin{equation}\label{rollover}
b_{t, t}  = b_{0, t}\,, t\geq 0 \,.
\end{equation}
The government's dynamic budget constraint then becomes \footnote{Precisely, the dynamic budget constraint is an inequality. However, optimality requires that the government issues minimum debt such that dynamic budget constraint holds equality and no resource is wasted.}
% by using the roll-over of $B_t$ as follows: %under the pure short-term debt rolling strategy is given by
\begin{equation}\label{eqn:Bmotion} 
dB_t =r_tB_t dt+ \left(G_t + b_{0,t} - \tau_tN_{0, t}  \right)dt  . %+ b_{0, t}dt . 
\end{equation}
Equation \eqref{eqn:Bmotion}  describes how   the government adjust  its ``deposit account'' (short-term debt balance)
$B_t$ to cover its cash-flow deficit $\left(G_t + b_{0,t} - \tau_tN_{0, t}  \right)dt$ under  coupon-flow management policy \eqref{rollover}.  
The deposit balance $B_t$ is both a record of past government extravagance (or frugality) and of future government
frugality (or extravagance).  Thus,  integrating \eqref{eqn:Bmotion}  from time $0$ to $t$, we deduce 
\begin{eqnarray}\label{eqB}
B_t 
= \int_{0}^{t} e^{\int_{s}^t r_{u}dt}\left[b_{0, s} - {(\tau_{0, s}N_{0, s} - G_s )} \right] ds     
\end{eqnarray} 
so that $B_t$  equals the government's accumulated primary deficit  evaluated at time $t$ prices, i.e., $B_t=\Pi_t $.
Integrating  \eqref{eqn:Bmotion}  from time $t$ to infinity, we obtain 	\begin{equation}\label{Bforward}
B_t =% \int_{t}^{\infty} \frac{M_{0,s}}{M_{0,t}} \left(\tau_{0, s}N_{0, s}-G_s - b_{0, s} \right) d t = 
\int_{t}^{\infty} e^{-\int_{s}^t r_{u}dt}  \left[ {(\tau_{0, s}N_{0, s} - G_s )} - b_{0, s} \right] ds  ,
\end{equation}
so that  $B_t$ also equals the present value of the government's prospective  net cash flows after servicing its promised coupon flow. 






\section{Time Consistency of Ramsey Plan}


In this section, we examine the time consistency of optimal commitment policy under the definition of Lucas-Stokey. We establish, for our economy with complete debt market, that the Ramsey problem can achieve time consistent via debt restructuring methods using both term debt and short-term debt that have not seen in Lucas-Stokey. We emphasize a unique short-term debt only restructuring method and provide sufficient conditions under which the time consistent is attainable. 



%In this section, we will first discuss all potential debt restructuring policies that could make the optimal commitment policy chosen the time$-0$ government time-consistent in the sense of Lucas and Stokey. Unlike the Lucas and Stokey (1983) where such potential debt restructuring policy is unique (if exists), we will show that when short-term debt is included, there exists a series of potential debt restructuring policies. Then we will provide a sufficient condition such that a debt restructuring that rolls over only short-term debt could make optimal commitment policy time-consistent in the sense of Lucas and Stokey.  

\NW{The complication all comes from the fact that $X_t$ depends on $U_{C, t}$ and $B_t$, but the key is $C_t$ may be different for time-0 planner and time-$t$ planner. That is, $C_{0,t}$ may differ from $C_{t,t}$. In other words, the planner has incentives to manipulate  the interest rate. In standard HJB, state variable is simply  inherited and non changeable.} 



%Time$-u$ government inherited debt from previous government: long-term debt $b_{0, t}, t\geq u$ and short-term debt $B_{u}=Y_u^*/U_C(C_{0, u}^*, C_{0, u}^*+G_u)$.

%If the time$-u$ government's can reevaluate it's policy given the inherited debt, whether it will follow the continuation of time$-0$ government's policy?  

%\subsection{Lagrangian Method}
%The time$-u$ government chooses $\{C_{u, t}, N_{u, t}, \tau_{u, t}; t\geq u\}$ to solve the following problem
%
%\begin{eqnarray} \label{obj1}
%	&& \max_{\{C_{u, t}; t\geq 0\}} \, \int_{u}^{\infty}  e^{-\rho t} U(C_{u, t}, C_{u, t}+G_t) dt \\
%	&s.t. &  \int_{u}^{\infty} e^{-\rho t} ( C_{u, t}U_C+(C_{u, t}+G_t)U_N )  dt  \geq   \int_{u}^{\infty} e^{-\rho t}b_{0-, t} U_C(C_{u, t}, C_{u, t}+G_t) dt+B_uU_{C},  \label{impletconstraint1} %\\
%\end{eqnarray}   
%
%Similarly, we can use Lagrangian method to solve this problem for time$-u$ government. We formulate the following Lagrangian
%\begin{eqnarray}
%	{\mathcal L}^{}_u &=&  \;   \int_{u}^{\infty}  e^{-\rho t}\left[ U(C_{u, t}, C_{u, t}+G_t)   \right] dt \notag \\
%	&&\; - \Lambda_{u}^{} \left[ \int_{u}^{\infty} e^{-\rho t} ( C_{u, t}U_C+(C_{u, t}+G_t)U_N )  dt -   \int_{u}^{\infty} e^{-\rho t}b_{0-, t}^{} U_C(C_{u, t}, C_{u, t}+G_t) dt -B_{u}U_{C} \right]
%	%	 \notag \\
%	%    &&\; - \Psi_0^{} \left[ \int_{0}^{\infty} e^{-\rho t}b_{0, t}^{} U_C(C_{0, t}, C_{0, t}+G_t) dt -\int_{0}^{\infty} e^{-\rho t}b_{0-, t} U_C(C_{0, t}, C_{0, t}+G_t) dt \right] 
%	\notag \\
%	& = & \int_{u}^{\infty}  e^{-\rho t}\left[ U(C_{u, t}, C_{u, t}+G_t)- \Lambda_{u}^{}\left( C_{u, t}U_C+(C_{u, t}+G_t)U_N-b_{0-, t}^{} U_C\right) %- \Psi_0(b_{0, t}^{} U_C- b_{0-, t} U_C)
%	\right]dt+\Lambda_{u}B_{u}U_{C}
%\end{eqnarray}
%
%The optimal policy under commitment is characterizes as follows: For each time $t\geq u$, consumption is chosen to solve for the period $[t, t+dt]$
%\begin{equation} \label{optp-L}
%	\max_{C_{u, t}}  \left[ U(C_{u, t}, C_{u, t}+G_t)- \Lambda_{u}^{}\left( C_{u, t}U_C+(C_{u, t}+G_t)U_N-b_{0-, t} U_C \right) \right] dt
%\end{equation}
%Note that $C(\Lambda_{u}, b_{0-, t}, G_t)$ denote optimal solution for this optimization problem. The optimal labor supply and tax policy are $N(\Lambda_{u}, b_{0-, t}, G_t)$ and $\tau(\Lambda_{u}, b_{0-, t}, G_t)$, respectively.
%
%The Lagrangian multiplier $\Lambda_{u}$ can be obtained by solve the implementability condition:
%\begin{eqnarray}\label{eqPhiu}
%	&& \int_{u}^{\infty} e^{-\rho (t-u)} ( C(\Lambda_{u}, b_{0-, t}, G_t)U_C+(C(\Lambda_{u}, b_{0-, t}, G_t)+G_t)U_N )  dt \nonumber \\
%	&&  =  \int_{u}^{\infty} e^{-\rho (t-u)}b_{0-, t} U_C dt+B_uU_{C}( C_{u, u} ,C_{u, u}+G_u)
%\end{eqnarray}
%where $U_C=U_C(C_{u, t}, C_{u, t}+G_u)= U_C(C(\Lambda_{u}, b_{0-, t}, G_t), C(\Lambda_{u}, b_{0-, t}, G_t)+G_t)$ and $U_N= U_N(C_{u, t}, C_{u, t}+G_u)= U_N(C(\Lambda_{u}, b_{0-, t}, G_t), C(\Lambda_{u}, b_{0-, t}, G_t)+G_t)$.
%
%
%Note that $B_u$ satisfies
%\begin{eqnarray} \nonumber
%	B_{u}& = & \frac{Y_u^*}{U_C(C_{0, u}^*, C_{0, u}^*+G_u)} \\ \nonumber
%	&=& \frac{e^{-\rho u}B_{0-}U_C^* +\int_{0}^{u}  e^{-\rho (s-u)}\left[-C_{0, s}^*U_{C}^*-N_{0, s}^*U_{N}^* +b_{0-, s}U_{C}^*\right]ds}{U_C(C_{0, u}^*, C_{0, u}^*+G_u)} \\
%	& = & \frac{\int_{u}^{\infty}  e^{-\rho (s-u)}\left[C_{0, s}^*U_{C}^*+N_{0, s}^*U_{N}^* -b_{0-, s}U_{C}^*\right]ds}{U_C(C_{0, u}^*, C_{0, u}^*+G_u)}.\label{Bu}
%\end{eqnarray}
%where $U_C^*= U_C(C_{0, u}^*, C_{0, u}^*+G_u)= U_C(C(\Lambda_{0}^*, b_{0-, t}, G_t), C(\Lambda_{0}^*, b_{0-, t}, G_t)+G_t)$ and $U_N^*= U_N(C_{0, u}^*, C_{0, u}^*+G_u)= U_N(C(\Lambda_{0}^*, b_{0-, t}, G_t), C(\Lambda_{0}^*, b_{0-, t}, G_t)+G_t)$
%
%Substituting (\ref{Bu}) into (\ref{eqPhiu}), we have
%\begin{eqnarray} \label{eqPhiu2}
%	\frac{\int_{u}^{\infty} e^{-\rho (s-u)} ( C(\Lambda_{u}, b_{0-, s}, G_t)U_C+(C(\Lambda_{u}, b_{0-, s}, G_s)+G_s)U_N -b_{0-, s} U_C)  dt}{U_C(C_{u, u}, C_{u, u}+G_u)}  \nonumber \\
%	=\frac{\int_{u}^{\infty}  e^{-\rho (s-u)}\left[C_{0, s}^*U_{C}+N_{0, s}^*U_{N} -b_{0-, s}U_{C}\right]ds}{U_C(C_{0, u}^*, C_{0, u}^*+G_u)}.
%\end{eqnarray}
%
%It can be rewritten as
%\begin{eqnarray}
%\frac{	\int_{u}^{\infty} e^{-\rho (s-u)} ( C(\Lambda_{u}, b_{0-, s}, G_t)U_C+(C(\Lambda_{u}, b_{0-, s}, G_s)+G_s)U_N -b_{0-, s} U_C)  ds }{U_C(C(\Lambda_{u}, b_{0-, u}, G_u), C(\Lambda_{u}, b_{0-, u}, G_u)+G_u)} \nonumber \\
%	= \frac{\int_{u}^{\infty} e^{-\rho (s-u)} ( C(\Lambda_{0}^*, b_{0-, s}, G_t)U_C^*+(C(\Lambda_{0}^*, b_{0-, s}, G_s)+G_s)U_N^* -b_{0-, s} U_C^*)ds}{U_C(C(\Lambda_{0}^*, b_{0-, u}, G_u), C(\Lambda_{0}^*, b_{0-, u}, G_u)+G_u)}
%\end{eqnarray}
%
%%To make the MRS we use to calculate $B_u$ and the one we use to solve time$-u$ government's problem consistent, we impose the following
%%$$C_{u, u}=C_{0, u}^*.$$
%
%Define the following function
%\begin{eqnarray} \nonumber
%\Psi_{u}(\Lambda_{u}) & =& \frac{	\int_{u}^{\infty} e^{-\rho (s-u)} ( C(\Lambda_{u}, b_{0-, s}, G_t)U_C+(C(\Lambda_{u}, b_{0-, s}, G_s)+G_s)U_N -b_{0-, s} U_C)  ds }{U_C(C(\Lambda_{u}, b_{0-, u}, G_u), C(\Lambda_{u}, b_{0-, u}, G_u)+G_u)}  \\
%&&- \frac{\int_{u}^{\infty} e^{-\rho (s-u)} ( C(\Lambda_{0}^*, b_{0-, s}, G_t)U_C^*+(C(\Lambda_{0}^*, b_{0-, s}, G_s)+G_s)U_N^* -b_{0-, s} U_C^*)ds}{U_C(C(\Lambda_{0}^*, b_{0-, u}, G_u), C(\Lambda_{0}^*, b_{0-, u}, G_u)+G_u)}
%\end{eqnarray}
%
%Implementability condition for time$-u$ government implies that the Lagrangian multiplier $\Lambda_{u}$ satisfies
%\begin{equation}
%\Psi_u(\Lambda_{u}) = 0.
%\end{equation}
%
%We can easily verify that $\Psi_u(\Lambda_{0}^*)=0$. Note that equation (\ref{eqPhiu2}) is nonlinear and may have multiple roots. The key to time-consistent policy is whether $\Lambda_{u}=\Lambda_{0}^*$ is the root of (\ref{eqPhiu2}) that maximizes time$-u$ government's objective $\mathcal{L}_u$.



\subsection{Potential debt restructuring policies}

Lucas-Stokey's definition of time consistency requires a government leaves a speacial debt structure to the future government such that, if the future government has an option to reoptimize, it still optimally choose the same plan with the one chosen by the time$-0$ government. In Lucas-Stokey, when debt market consists only term debt, there exists a unique potential debt restructuring method such that time consistency is achievable. In this part,
we characterize all potential debt restructuring policies that can pass to the time$-u$ government such that the time$-0$ optimal commitment policy satisfies the FOCs for the time$-u$ government, and we will show that potential debt restructuring policies are inifite. Note that a potential debt restructuring policy may not achieve Lucas-Stokey's time consistency. We will provide sufficient condition for a potential debt restructuring to achieve time consistency later.


\paragraph{Time-$u$ debt value invariant equation}
Under time$-0$ Ramsey plan, the time$-u$ value of debt that faces the time$-u$ government consists of two parts: the part accumulated from time $0$ to $u$, measured by the accumulative primary deficit $\Pi_u$, and the part from initial debt that due at time beyond $u$. Given the time$-0$ optimal commitment policy $\{C_{0, t}^*, t\geq0\}$,%\footnote{For the ease of notation, we omit the $^*$ in the optimal commitment policy} 
the first part $\Pi_{u}$ satisfies
\begin{eqnarray}
\Pi_u & = & \int_{0}^u\, \frac{M_{0, t}}{M_{0, u}} \left(b_{0, t} +G_t-\tau_s N_{0, t} \right)dt  \nonumber\\
& = & \int_{0}^u \frac{e^{\rho(u-t)}}{U_{C, u}} \left[ b_{0, t}U_{C, t}-C_{0, t}^*U_{C, t}-(C_{0, t}^*+G_t) U_{N, t}   \right]dt\,. 
\end{eqnarray}
For the second part, its time$-u$ value is measured by
\begin{eqnarray}
& &\int_{u}^\infty\, \frac{M_{0, t}}{M_{0,u}} \left(b_{0, t} +G_s-\tau_t N_{0, t} \right)dt \nonumber \\
& = & \int_u^{\infty} \frac{e^{-\rho(t-u)}}{U_{C, u}} \left[ b_{0, t}U_{C, t}-C_{0, t}^*U_{C, t}-(C_{0, t}^*+G_t)U_{N, t}  \right]dt\,.
\end{eqnarray}

First, a potential debt restructuring policy must satisfies that the time$-u$ value of debt remains unchanged under the time$-0$ Ramsey plan. 
%The debt restructuring is all about choosing some fraction of short-term debt and some term debt to cover the funding gap $B_u$, and then passing them to the future government. 
Let $\lambda$ denote the fraction of time$-u$ accumulative primary deficit restructured via short-term. The short-term debt that passed to the time$-u$ government is then $\widehat{B}_u = \lambda \Pi_u$. The term debt passed to the future government, denoted by $\{\widehat{b}_{u, t}; t\geq u\}$, satisfies
\begin{eqnarray}\label{termdebt_structure}
%&&B_u U_C + \int_u^{\infty} e^{-\rho(t-u)} \left[ -C_{0, t}U_C-(C_{0, t}+G_t)U_N +b_{0, t}U_C \right]dt   \nonumber \\
&  &  \int_u^{\infty} \frac{e^{-\rho(t-u)}}{U_{C, u}} \left[\widehat{b}_{0, t}U_{C, t} -C_{0, t}^*U_{C, t}-(C_{0, t}^*+G_t)U_{N, t} \right]dt \nonumber\\
& = & (1-\lambda)\Pi_u+ \int_u^{\infty} \frac{e^{-\rho(t-u)}}{U_{C, u}} \left[ b_{0, t}U_{C, t}-C_{0, t}^*U_{C, t}-(C_{0, t}^*+G_t)U_{N, t}  \right]dt\,.
\end{eqnarray}

Using implementability condition (\ref{impletconstraint}), we obtain the following {\it time$-u$ debt value invariant equation}
\begin{align}
\int_{u}^{\infty} e^{-\rho (t-u)} ( C_{0, t}^*U_{C, t}+(C_{0, t}^*+G_t)U_{N, t} )  dt
=  \int_{u}^{\infty} e^{-\rho (t-u)}\widehat{b}_{u, t} U_{C, t} dt+\widehat{B}_uU_{C, u}( C_{0, u}^* ,C_{0, u}^*+G_u) \label{time_u_value_invariant}
\end{align}

Specifically, when $\lambda=1$, all accumulative primary deficit passed to the future government is restructured purely via the short-term debt, i.e., $\widehat{B}_u=\Pi_u$, and the term debt passed to the future government is the same with that inherited directly from the time$-0$ government, i.e., $\{\widehat{b}_{u, t} = b_{0, t}; t\geq u\}$. This is a nature debt restructuring policy that has not been studied in Lucas and Stokey (1983) and DNY (2019). Note that there are many $\{\widehat{b}_{u, t}; t\geq u\}$ that may satisfy the equation (\ref{time_u_value_invariant}). %However, there exists only one $\{\widehat{b}_{u, t}; t\geq u\}$ which also satisfy the time$-u$ government's FOCs, as we will show next.

However, we show next that, for a given $\lambda$, there exists only one potential debt restructuring policy $\{\widehat{B}_u, \widehat{b}_{u, t}; t\geq u\}$ such that %Lucas and Stokey time-consistency may hold. 
%For a time$-u$ government, we will show that there is a unique debt structure $\{\widehat{B}_u, \widehat{b}_{u, t}, t\geq u\}$ which satisfies 
the necessary FOCs that ensures Lucas-Stokey time consistency hold. %such that the time$-u$ government will not deviate from the time$-0$ Ramsey plan.

\paragraph{Necessary FOCs}
%\subsection{Dynamic programming method}
Facing an initial debt $\{\widehat{B}_u, \widehat{b}_{u, t}; t\geq u\}$ and given the reoptimization option, the time$-u$ government's optimization problem can be decomposed to be a {\it continuation Ramsey problem for time $t\geq u+$ onward} in which the government chooses allocation $\{C_{u, t}, t\geq u+\}$ and a {\it time$-u$ Ramsey problem} in which the government chooses $\{C_{u, u}\}$ and $\Phi_u$ at time $u$. Then the MRS adjusted accumulative primary deficit for the time$-u$ government, denoted by $\widehat{X}_t, t\geq u$, evolves as
\begin{eqnarray}\label{dyhat}  
d \widehat{X}_t & = & \left[\rho \widehat{X}_{t}-C_{u, t}U_{C, t} -N_{u, t}U_{N, t}  +\widehat{b}_{0, t}U_{C, t}\right]dt\,. % \\
%&  & +\, .
\end{eqnarray} 
%with $\widehat{X}_u = \widehat{B}_uU_{C, t}$. %Here $U_{C, t}, U_{N, t}$ are short notation for $U_{C, t}(C_{u, t}, N_{u, t}), U_{N, t}(C_{u, t}, N_{u, t})$, respectively. 
Note the starting value of $\widehat{X}_t$ is $\widehat{X}_u=\widehat{B}_uU_{C, t}$, which is not trivial because $\widehat{B}_u$ does not necessary equal to zero, contrasting with $X_0=0$ for the time$-0$ government. In addition, the time$-u$ government faces a time$-u$ implementability condition
\begin{equation}\label{imcbu}
\int_{u}^{\infty} e^{-\rho (t-u)} ( C_{u, t}U_{C, t}+(C_{u, t}+G_t)U_{N, t} )  dt
=  \int_{u}^{\infty} e^{-\rho (t-u)}\widehat{b}_{u, t} U_{C, t} dt+\widehat{B}_uU_{C, u}( C_{u, u} ,C_{u, u}+G_u)
\end{equation}


We consider time$-u$ government's continuation Ramsey problem for $t\geq u+$ onward.
Let $\widehat{V}(\widehat{X}, t)$ be the value of a continuation Ramsey planner when $\widehat{X}_t=\widehat{X}$ at $t\geq u+$. The Bellman equation for the time $u$ continuation Ramsey planner is 
%	\begin{equation}
%		V(Y, t; 0) = \max_{\{C_{0, s}, s\geq t \}}\int_{t}^{\infty} e^{-\rho (s-t)} U(C_{0,s}, N_{0,s}) ds 
%	\end{equation} where $(Y, t)$ are the state variables. 
%	The HJB equation is given as
\begin{eqnarray}\label{HJBW} \nonumber
\rho \widehat{V} (\widehat{X}, t) & = &  \frac{\partial \widehat{V}}{\partial t}+\max_{C_{u, t}} \;\; U(C_{u, t}, C_{u, t}+G_t) \\
& & + \frac{\partial \widehat{V}(\widehat{X}, t) }{\partial X} \left(\rho \widehat{X}-C_{u, t} U_{C, t}-(C_{u, t}+G_t)U_{N, t} +\widehat{b}_{u, t}U_{C, t} \right);\, t\geq u+   %\\
%&& \quad   \;\;\; %\WJ{- \frac{\partial V}{\partial X} \left(\sum_{k\neq s}\lambda_{s k}\beta(s, k) \right) }
%+\WJ{\sum_{k\neq s}\lambda_{sk}\left[ V(X+\beta{(s, k)}X, k)-V(X, s) - \frac{\partial V(X, s)}{\partial X} \beta{(s, k)}X   \right] }.
\end{eqnarray}

Solving the equation we obtain
\begin{equation}\label{Wvalue}
\widehat{V}(\widehat{X}, t; \Phi_u) = -\Phi_u \widehat{X}+\widehat{H}(t; \Phi_u); t\geq u+
\end{equation} 
where $\Phi_u\geq 0$ will be determined by the following time$-u$ Ramsey problem and $\widehat{H}(t; \Phi_u)$ satisfies
\begin{eqnarray}\label{HJBHu} %\nonumber
\rho \widehat{H}(t; \Phi_u) 
& = &  \frac{d \widehat{H}}{d t}+ \max_{C_{u, t}} \;\; U(C_{u, t}, C_{u, t}+G_t) \notag \\
& & + \Phi_u\left(C_{u, t} U_{C, t}+(C_{u, t}+G_t)U_{N, t} -\widehat{b}_{u, t}U_{C, t} \right); t\geq u   % \\
%&& +\sum_{k\neq s}\lambda_{sk}\left[ \WJ{H(k; \Phi_1)}-H(s; \Phi_1) \right].
\end{eqnarray}

The optimal allocations $\{C_{u, t}; t\geq u+\}$ then must satisfy the FOCs for all $t\geq u+$
\begin{equation} \label{FOCbu}
(1+\Phi_u^{})(U_{C, t}+U_{N, t})+ \Phi_u \left[ (C_{u, t}-\widehat{b}_{u, t})(U_{CC, t}+U_{CN, t})+(C_{u, t}+G_t)(U_{NC, t}+U_{NN,t }) \right] = 0\,.
\end{equation}
We require $\Phi_u\geq 0$ to ensure the optimization problem in (\ref{HJBHu}) admits interior solution.


\paragraph{Potential debt restructuring policies} We first give the definition and then show the potential debt restructuring policy is unique. Note that we focus on continuous flow of debt-payout, i.e., $\widehat{b}_{u, s} = \lim_{t\downarrow s} b_{u, t}, s\geq u$.

\begin{definition}
A potential debt restructuring policy is a debt structure $\{\widehat{B}_u, \widehat{b}_{u, t}; t\geq u\}$ that faces the time$-u$ government such that (\ref{time_u_value_invariant}) holds and (\ref{FOCbu}) holds for some $\Phi_u$ for the time$-0$ Ramsey plan $\{C_{0, t}^*, t\geq u\}$.
\end{definition}

\begin{lemma}
Let $\widehat{B}_u$ be fixed for a $\lambda$, i.e., $\widehat{B}_u=\lambda\Pi_u$. There is a unique sequence $\{\widehat{b}_{u, t}; t\geq u\}$ such that $\{\widehat{B}_u, \widehat{b}_{u, t}; t\geq u\}$ is a potential debt restructuring policy, i.e.,
\begin{align}\label{eqhatbut}
\widehat{b}_{u, t} & =  \frac{\Phi_0^*}{\Phi_u} b_{0,t} + \left(1-  \frac{\Phi_0^*}{\Phi_u} \right) A_{u, t}\,.
\end{align}
where $\Phi_0^*$ is given in Proposition 5, $\Phi_u$ is uniquely determined by
\begin{align}
\Phi_u  & =  \Phi_0^*\frac{\int_{u}^{\infty} e^{-\rho (t-u)} \left(b_{0, t} -A_{u, t} \right)U_{C, t} dt }{\int_{u}^{\infty} e^{-\rho (t-u)} (  C_{0, t}^*  U_{C, t}+(C_{0, t}^*+G_t)U_{N, t} - A_{u, t}  U_{C, t})  dt -\lambda \Pi_u U_{C, u}(C_{0, u}^* ,C_{0, u}^*+G_u)} \,, 
\end{align}
and 
\begin{equation}
A_{u, t} = \frac{U_{C, t}+U_{N, t} +C_{0, t}^*(U_{CC, t}+U_{CN, t})+(C_{0, t}^*+G_t)(U_{CC, t}+U_{CN, t})}{U_{CC, t}+U_{CN, t}} \,,
\end{equation} 
\end{lemma}

%For a given $\lambda$ such that $\widehat{B}_u=\lambda\Pi_u$, we solve (\ref{hatbut}) for $\widehat{b}_{u, t}$ and substitute into (\ref{imc0t}).
\begin{proof}

A necessary condition for the time$-u$ government to be time-consistent is that the debt restructuring policy $\{\widehat{B}_u, \widehat{b}_{u, t}, t\geq u\}$ is chosen such that the time$-0$ Ramsey plan $\{C_{0, t}^*, t\geq u\}$ satisfies the necessary FOCs (\ref{FOCbu}) and the time$-u$ debt value invariant condition (\ref{time_u_value_invariant}). That is, there exists some $\Phi_u$, such that $\{C_{0, t}^*, t\geq u\}$ satisfies the following equations for $t\geq u+$:
\begin{align}
& (1+\Phi_u^{})(U_{C, t}+U_{N, t})+ \Phi_u \left[ (C_{0, t}^*-\widehat{b}_{u, t})(U_{CC, t}+U_{CN, t})+(C_{0, t}^*+G_t)(U_{NC, t}+U_{NN, t}) \right] = 0;    \label{focc0t}\\
& \int_{u}^{\infty} e^{-\rho (t-u)} ( C_{0, t}^*U_{C, t}+(C_{0, t}^*+G_t)U_{N, t} )  dt
=  \int_{u}^{\infty} e^{-\rho (t-u)}\widehat{b}_{u, t} U_{C, t} dt+\widehat{B}_uU_{C, u}( C_{0, u}^* ,C_{0, u}^*+G_u) \label{imc0t}
\end{align}
Note that the debt restructuring equation (\ref{termdebt_structure}) and time$-0$ implementability condition (\ref{impletconstraint}) (which holds equality under optimality) ensures equation (\ref{imc0t}) holds. 

On the other hand, $\{C_{0, t}^*, t\geq u\}$ also satisfy the FOCs (for $t\geq u+$) and implementability condition for time$-0$ government's Ramsey problem:
\begin{align}
& (1+\Phi_0^{*})(U_{C, t}+U_{N, t})+ \Phi_0^* \left[ (C_{0, t}^*-{b}_{0, t})(U_{CC, t}+U_{CN,t })+(C_{0, t}^*+G)(U_{NC, t}+U_{NN, t}) \right] = 0;   \label{focc0t1}\\
& \int_{u}^{\infty} e^{-\rho (t-u)} ( C_{0, t}^*U_{C, t}+(C_{0, t}^*+G_t)U_{N, t} )  dt
=  \int_{u}^{\infty} e^{-\rho (t-u)}{b}_{0, t} U_{C, t} dt+{\Pi}_uU_{C, u}( C_{0, u}^* ,C_{0, u}^*+G_u) \label{imc0t1}
\end{align}
where equation (\ref{imc0t1}) is derived from time$-0$ implementability condition (\ref{impletconstraint}) (which holds equality under optimality). 

Combining (\ref{focc0t})-(\ref{focc0t1}), we obtain
\begin{eqnarray} \label{hatbut}
\Phi_u \widehat{b}_{u, t} = \Phi_0^* b_{0, t} + (\Phi_u-\Phi_0^*) A_{u, t}, \, t\geq u+
\end{eqnarray} 
where the coefficient $A_{u, t}$ satisfies
\begin{equation*}
A_{u, t} = \frac{U_{C, t}+U_{N, t} +C_{0, t}^*(U_{CC, t}+U_{CN, t})+(C_{0, t}^*+G_t)(U_{CC, t}+U_{CN, t})}{U_{CC, t}+U_{CN, t}} \,.
\end{equation*}
Because $\widehat{b}_{u, t}, t\geq u$ and $\{C^*_{0, t}, t\geq u\}$ both are continuous in $t$, we further obtain that, at time $t=u$,
\begin{eqnarray} \label{hatbuu}
\Phi_u \widehat{b}_{u, u} = \Phi_0^* b_{0, u} + (\Phi_u-\Phi_0^*) A_{u, u},.
\end{eqnarray} 
Thus equation (\ref{eqhatbut}) follows for all $t\geq u$. 
\end{proof}

%where $\lambda B_u = \widehat{B}_u$. 

Knowing that $\Phi_u$ and $\widehat{b}_{u, t}, t\geq u$ are functions of $\lambda$, i.e., $\Phi_u(\lambda), \widehat{b}_{u, t}(\lambda), t\geq u$. When $\lambda=1$, it refers to a debt restructuring strategy that uses only short-term debt to restructure debt and pass it to the time$-u$ government. Then we can easily verify that
\begin{eqnarray}
\Phi_u(1) & = & \Phi_0^* \\
\widehat{b}_{u, t}(1) &=& b_{0, t}\,.
\end{eqnarray}

On the other hand, when $\lambda=0$, it refers to a debt restructuring strategy that uses only term debt to restructure debt and pass it to the time$-u$ government. This particular strategy is the same that have been studied in Lucas and Stokey (1983). In this case, we have
\begin{equation}\label{phiu0}
\Phi_u(0)   =  \Phi_0^*   \left(1 + \frac{\Pi_u U_{C, u}(C_{0, u}^*, C_{0, u}^*+G_u)}{\int_{u}^{\infty} e^{-\rho (t-u)} \left(b_{0, t} -A_{u, t} \right)U_{C, t} dt} \right)^{-1} \,.
%\frac{\int_{u}^{\infty} e^{-\rho (t-u)} \left(b_{0, t} -A_{u, t} \right)U_C dt }{\int_{u}^{\infty} e^{-\rho (t-u)} \left(b_{0, t} -A_{u, t} \right)U_C dt+B_u U_C(C_{0, u} ,C_{0, u}+G_u)} \,, \\
\end{equation}

\begin{proposition}
For a potential debt resturcturing policy that leaves short-term debt of a balance $\widehat{B}_u=\lambda\Pi_u$ to the time$-u$ government, then $\Phi_u(\lambda)$ satisfies 
\begin{equation} \label{Phiulambda}
\Phi_u(\lambda) = \frac{\Phi_u(1)\Phi_u(0)}{(1-\lambda)\Phi_u(1) +\lambda \Phi_u(0)}\,.
\end{equation}
\end{proposition}

A potential debt restructuring policy may not necessary ensure Lucas-Stokey time-consistency because it may not necessary solve the time$-u$ Ramsey problem. In the following, we provide a sufficient condition such that a potential debt restructuring policy achieves time consistency. 






\subsection{Necessary and sufficient conditions for time-consistency}

An important element to verify whether a potential debt structuring policy achieves time consistency is to verify whether it solve the following time$-u$ Ramsey problem. 

\paragraph{Time$-u$ Ramsey problem}
Facing the continuation value function $\widehat{V}(\widehat{X}_t, t); t\geq u$, the time$-u$ Ramsey planner chooses $\widehat{X}_u$ at time $u$. Continuity of $\{C_{u, t}\}$ at $t=u$ from condition 2 requires that
\begin{equation} \label{cuucontinuous}
C_{u, u} = \lim_{s\downarrow u} C_{u, s} = C(\Phi_u, \widehat{b}_{u, u}, G_u)\,,
\end{equation}
which implies that $\Phi_u$ is a function of $C_{u, u}$. The time$-u$ Ramsey planer takes $\{\widehat{B}_u, \widehat{b}_{u, t}; t\geq u\}$ as given and solves
\begin{equation}
\widehat{W}(\{\widehat{B}_u, \widehat{b}_{u, t}; t\geq u\}) = \max_{C_{u, u}} \widehat{V} (\widehat{X}_u, u)\,,
\end{equation}
where $\widehat{X}_u= \widehat{B}_uU_{C, u}$ and the optimization is subject to the implementability condition (\ref{imcbu}).


%The time$-u$ government's value function $V(X, t; \Phi_u)$, for a given $\Phi_u$, also satisfies (\ref{Wvalue}) and (\ref{HJBHu}) for $t=u$. This suggests that the consumption policy at time $t=u$ solves 
%\begin{equation*}
%\max_{C_{u, u}} \;\; U(C_{u, u}, C_{u, u}+G_u) 
%+ \Phi_u\left(C_{u, u} U_{C}+(C_{u, u}+G_t)U_N -\widehat{b}_{u, u}U_{C} \right)
%\end{equation*}
%
%
%Then, the time$-u$ government chooses $\Phi_u$ to solve the following optimization problem
%\begin{equation}
%J(\widehat{B}_{u}, \{\widehat{b}_{u, t}, t\geq u\}) = \max_{\Phi_u} W(Y_u, u; \Phi_u)
%\end{equation}
%where $W(Y_u, u; \Phi_u)$ satisfies (\ref{HJBHu}) when $t=u$
%such that $\Phi_u$ satisfies the implementability condition
%\begin{eqnarray}\label{Phiueq}
%&& \int_{u}^{\infty} e^{-\rho (t-u)} ( C(\Phi_{u}, \widehat{b}_{u, t}, G_t)U_C+(C(\Phi_{u}, \widehat{b}_{u, t}, G_t)+G_t)U_N )  dt \nonumber \\
%&&  =  \int_{u}^{\infty} e^{-\rho (t-u)}\widehat{b}_{u, t} U_C dt+\widehat{B}_uU_{C}( C(\Phi_{u}, \widehat{b}_{u, u}, G_u) ,C(\Phi_{u}, \widehat{b}_{u, u}, G_u)+G_u)
%\end{eqnarray}
%
%In this section, we provide a sufficient condition for the optimal commitment policy to be time-consistent in the sense of Lucas and Stokey. In particular, this sufficient condition applies to the pure short-term debt restructuring strategy: time$-u$ government inherits $\{B_u, \{b_{0, t}, t\geq u\}\}$.




Define the following function
\begin{eqnarray} \nonumber
\Psi_{u}(\Phi) & =& \frac{	\int_{u}^{\infty} e^{-\rho (t-u)} ( C(\Phi, \widehat{b}_{u, t}, G_t)U_{C, t}+(C(\Phi, \widehat{b}_{u, t}, G_t)+G_t)U_{N, t} -\widehat{b}_{u, t} U_{C, t})  dt }{U_{C, t}(C(\Phi, \widehat{b}_{u, u}, G_u), C(\Phi, \widehat{b}_{u, u}, G_u)+G_u)}  \\ \nonumber
&&- \lambda\, \frac{\int_{u}^{\infty} e^{-\rho (t-u)} ( C(\Phi_{0}^*, b_{0, t}, G_t)U_C+(C(\Phi_{0}^*, b_{0, t}, G_t)+G_s)U_N -b_{0, t} U_C)dt}{U_C(C(\Phi_{0}^*, b_{0, u}, G_u), C(\Phi_{0}^*, b_{0, u}, G_u)+G_u)}\,,
\end{eqnarray}
where $\widehat{b}_{u, t}, t\geq u$ are given by (\ref{eqhatbut}). 
Note that the implementability condition requires the unique $\Phi$ must satisfy %from a potential debt restructuring policy 
\begin{equation} \label{imcPhiu}
\Psi_u(\Phi) = 0 \,.
\end{equation}

Using (\ref{cuucontinuous}), the time$-u$ Ramsey planner's problem is equivalent to
\begin{align*}
\widehat{W}(\{\widehat{B}_u, \widehat{b}_{u, t}; t\geq u\}) & =  \max_{\{\Phi\geq 0: \Psi_u(\Phi)=0\} } \widehat{V} (\widehat{B}_uU_{C, u}, u)  \\
& =   \max_{\{ \Phi\geq 0: \Psi_u(\Phi)=0 \}} \int_{u}^\infty e^{-\rho (t-u)} U(C(\Phi, \widehat{b}_{u, t}, G_t), C(\Phi, \widehat{b}_{u, t}, G_t)+G_t) dt\,.
\end{align*}


\subsubsection{A General Necessary and Sufficient Condition}

We then provide an necessary and sufficient condition for a Ramsey plan to be time-consistency for arbitrary initial debt and government spendings.

\begin{proposition}
A time$-0$ Ramsey plan is time consistency at time $-u$ if and only if there exists a $\lambda$ such that the uniquely determined $\widehat{b}_{u, t}, t\geq u$ and $\Phi_u(\lambda)$ in Lemma 1 solves 
\begin{eqnarray}\label{Phiuoptimal}
\Phi_u(\lambda) = \arg \max _{\{\Phi\geq 0: \Psi_u(\Phi)=0\}} \int_{u}^\infty e^{-\rho (t-u)} U(C(\Phi, \widehat{b}_{u, t}, G_t), C(\Phi, \widehat{b}_{u, t}, G_t)+G_t) dt\,.
\end{eqnarray}
\end{proposition}


{\bf Remark:} This general condition is difficult to verify as there are infinitely many potential debt restructuring policies. When the debt market is incomplete in the sense that no short- term debt exists, $\lambda$ must be zero and the associated $(0, \widehat{b}_{u, t})$ changes for different $u$. As a result, the optimization object changes substantially for different $u$.


{\bf Remark:} DNY provide an example in which, when $\lambda=0$, the condition (\ref{Phiuoptimal}) is violated because $\Phi_u(0)<0$, and thus time-consistency under the definition of Lucas-Stokey may not always hold. 

{\bf Remark:} DNY's argument to relate the sign of shadow price to the time-consistency may fail when there is short-term debt in the market (therefore $\lambda=0$). However, when there is short-term debt such that $\lambda$ could be non-zero, we argue that there could be the case in which $\Phi_u(\lambda)\geq 0$ still hold but time consistency still fails. Because $\Psi_u(\cdot)=0$ is a nonlinear equation, there might exists other positive root of $\Psi_u(\cdot)=0$ that maximizes (\ref{Phiuoptimal}), and thus time consistency still fails.

%Note that $\Psi_u(\Phi)=0$ is a non-linear equation of $\Phi$ and may have multiple roots. We can easily verify that $\Psi_u(\Phi_0)=0$. The time-consistency for the optimal commitment policy depends on whether $\Phi_0$ is the optimal root for the non-linear equation $\Psi_u(\Phi)=0$ for all $u>0$. 

\subsubsection{A Special Necessary and Sufficient Condition}

Because our model allows that a government uses the potential debt restructuring policy that relies on only short-term debt, i.e., $\lambda=1$, we are particularly interest in how such special potential debt restructuring policy attain time consistency. Knowing that when $\lambda=1$, $\widehat{B}_u = \Pi_u, \widehat{b}_{u, t}=b_{0, t}, t\geq u$, and $\Phi_u(1)=\Phi_0^*$. The function $\widetilde{\Psi}_u(\cdot)$ is then given by 
\begin{eqnarray} \nonumber
\widetilde{\Psi}_{u}(\Phi) & =& \frac{	\int_{u}^{\infty} e^{-\rho (t-u)} ( C(\Phi, {b}_{0, t}, G_t)U_{C, t}+(C(\Phi, {b}_{0, t}, G_t)+G_t)U_{N, t} -{b}_{0, t} U_{C, t})  dt }{U_{C, t}(C(\Phi, {b}_{0, u}, G_u), C(\Phi, {b}_{0, u}, G_u)+G_u)}  \\ \nonumber
&&-  \frac{\int_{u}^{\infty} e^{-\rho (t-u)} ( C(\Phi_{0}^*, b_{0, t}, G_t)U_C+(C(\Phi_{0}^*, b_{0, t}, G_t)+G_s)U_N -b_{0, t} U_C)dt}{U_C(C(\Phi_{0}^*, b_{0, u}, G_u), C(\Phi_{0}^*, b_{0, u}, G_u)+G_u)}\,.
\end{eqnarray}
We can easily verify that $\Phi_0^*$ is automatically a root of $\widetilde{\Psi}_u(\cdot)=0$. 

The question is that whether checking the pure short-term potential debt restructuring policy is neccesary to verify the time-consistency of the Ramsey plan.

The following proposition shows that, on the one hand, if $\Phi_0^*$, as a root of $\widetilde{\Psi}_u(\cdot)=0$, maximizes the government's objective among other roots of $\widetilde{\Psi}_u(\cdot)=0$, then time-consistency is achievable.
In addition, we can furhter show that, if the initial debt $\{b_{0, t}; t\geq 0\}$ and government spending $\{G_t; t\geq 0\}$ both are constants for all time $t\geq u$, it is also a necessary condition for a government's time$-0$ Ramsey plan to be time-consistent in the sense of Lucas and Stokey. 

\begin{proposition} \label{prop:9} Assume the initial debt $\{b_{0, t}; t\geq 0\}$ and government spending $\{G_t; t\geq 0\}$ both are constants over time $t\geq u$. A time$-0$ Ramsey plan is time-consistent over time $t\geq u$ if and only if the following holds:
\begin{equation}\label{tc_condition2}
\Phi_0^* = \argmax_{\{\Phi_{}\geq 0: \widetilde{\Psi}_u(\Phi)=0\} } \; \int_{u}^\infty e^{-\rho (t-u)} U(C(\Phi, {b}_{0, t}, G_t), C(\Phi, {b}_{0, t}, G_t)+G_t) dt\,.
\end{equation}
\end{proposition}

{\bf Remark:} This condition is easier to verify as it only requires to check the positive roots of a relative simple equation $\widetilde{\Psi}_u(\cdot)=0$. How many of positive roots of $\widetilde{\Psi}_u(\cdot)=0$ depends on the initial debt structure $\{b_{0, t}; t\geq 0\}$. 

To prove this, we first provide a lemma that if a potential debt restructuring policy with only short-term is not time-consistentcy at time $u$, then all potential debt restructuring policies are not time-consistency at time $u$. 
\begin{lemma}\label{lemma2}
Assume the initial debt $\{b_{0, t}; t\geq 0\}$ and government spending $\{G_t; t\geq 0\}$ both are constants over time $t\geq u$. If a potential debt restructuring policy with only short-term debt $(\Pi_u, {b}_{0, t})$ cannot achieve time-consistency at time $u$, then the anyother potential debt restructuring policy with $\lambda\neq 1$, i.e., $(\lambda\Pi_u, \widehat{b}_{u, t})$ also cannot achieve time-consistency at time $u$.
\end{lemma}

\begin{proof}
Assume that potential debt restructuring policy $(\Pi_u, {b}_{0, t})$ cannot achieve time-consistency at time $u$. It means that for the time-$u$ government that inherits $(\Pi_u, {b}_{0, t})$ from time $u-$ government, there must exists another $\breve{\Phi}_u \neq \Phi_0^*$ such that the it solves the optimization problem (87), i.e., 
\begin{eqnarray*} %\label{Phiuoptimal}
\breve{\Phi}_u = \arg \max _{\{\Phi\geq 0: \Psi_u(\Phi)=0\}} \int_{u}^\infty e^{-\rho (t-u)} U(C(\Phi, {b}_{0, t}, G_t), C(\Phi, {b}_{0, t}, G_t)+G_t) dt\,.
\end{eqnarray*}		

The corresponding new allocation $\{\breve{C}_{u, t} = C(\breve{\Phi}_u, {b}_{0, t}, G_t); \breve{N}_{u, t} = C(\breve{\Phi}_u, {b}_{0, t}, G_t)+G_t; t\geq u\}$ is flat for $t\geq u$ and solves the FOCs and implementability condition given below:
\begin{align}
& (1+\breve{\Phi}_u^{})(U_{\breve{C}, t}+U_{\breve{N}, t})+ \breve{\Phi}_u \left[ (\breve{C}_{u, t}-{b}_{0, t})(U_{\breve{C}\breve{C}, t}+U_{\breve{C}\breve{N}, t})+(\breve{C}_{u, t}+G_t)(U_{\breve{N}\breve{C}, t}+U_{\breve{N}\breve{N}, t}) \right] = 0;    \label{focc0t2}\\
& \int_{u}^{\infty} e^{-\rho (t-u)} (\breve{C}_{u, t}U_{\breve{C}, t}+(\breve{C}_{u, t}+G_t)U_{\breve{N}, t} )  dt
=  \int_{u}^{\infty} e^{-\rho (t-u)}{b}_{0, t} U_{\breve{C}, t} dt+\Pi_uU_{\breve{C}, u} \label{imc0t2}
\end{align}
Note that the equation (\ref{focc0t2}) is the same for all $t\geq u$ because $b_{0, t}$ is constant over time for $t\geq u$. Most importantly, the new allocation surpasses the time$-0$ Ramsey plan allocation $\{C_{0, t}^*, N_{0, t}^*; t\geq u\}$, i.e., 
\begin{equation*}\label{eq_surpass}
\int_{u}^\infty e^{-\rho (t-u)} U(\breve{C}_{u, t}, \breve{C}_{u, t}+G_t) dt > \int_{u}^\infty e^{-\rho (t-u)} U(C_{0, t}^*, C_{0, t}^*+G_t) dt
\end{equation*}

We then show there exists a $\widetilde{\Phi}\neq \Phi_u$ such that there exists a better allocation for $(\lambda\Pi_u, \widehat{b}_{u, t})$. 

By construction, we let 
\begin{equation}\label{tildePhi_u}
\widetilde{\Phi}_u = \frac{\Phi_u Q_{\breve{C}, t}}{Q_{C^*, t}+(A_{C^*, t}-A_{\breve{C}, t})\Phi_u}
\end{equation}
where $Q_{C, t} = \frac{U_{C, t}+U_{N, t}}{U_{CC, t}+U_{CN, t}}$. Knowing that $b_{0, t}$ is constant over time $t\geq u$, $Q_{C^*, t}, Q_{\breve{C}, t}, A_{C^*, t}, A_{\breve{C}, t}$ are all constants over time $t\geq u$. Therefore, $\widetilde{\Phi}_u$ is a constant over time.

We then verify that $\widetilde{\Phi} $ satisfies the following necessary conditions such that $\{\breve{C}_{u, t}; \breve{N}_{u, t}\}$ is a feasible allocation:
\begin{align}
& (1+\widetilde{\Phi}_u^{})(U_{\breve{C}, t}+U_{\breve{N}, t})+ \widetilde{\Phi}_u \left[ (\breve{C}_{u, t}-{\widehat{b}}_{u, t})(U_{\breve{C}\breve{C}, t}+U_{\breve{C}\breve{N}, t})+(\breve{C}_{u, t}+G_t)(U_{\breve{N}\breve{C}, t}+U_{\breve{N}\breve{N}, t}) \right] = 0;    \label{focc0t3}\\
& \int_{u}^{\infty} e^{-\rho (t-u)} (\breve{C}_{u, t}U_{\breve{C}, t}+(\breve{C}_{u, t}+G_t)U_{\breve{N}, t} )  dt
=  \int_{u}^{\infty} e^{-\rho (t-u)}{\widehat{b}}_{u, t} U_{\breve{C}, t} dt+\lambda\Pi_uU_{\breve{C}, u} \label{imc0t3}
\end{align}

We rewrite FOCs (\ref{focc0t}), (\ref{focc0t1}), and (\ref{focc0t2}) for different debt restructuring policies as following
\begin{eqnarray*}
Q_{C^*, t} +\Phi_0^* A_{C^*, t} & = & \Phi_0^* b_{0, t} \\
Q_{C^*, t} +\Phi_u A_{C^*, t} & = & \Phi_u \widehat{b}_{0, t} \\	
Q_{\breve{C}, t} + \breve{\Phi}_u A_{\breve{C}, t} & =& \breve{\Phi}_u b_{0, t} 
\end{eqnarray*}

Combining with equation (\ref{tildePhi_u}), we have
\begin{eqnarray*}
\widetilde{\Phi}_u(\widehat{b}_{u, t} - A_{\breve{C}, t}) = \widetilde{\Phi}_u( Q_{C^*, t}/\Phi_u +A_{C^*, t} - A_{\breve{C}, t}) = Q_{\breve{C}, t},
\end{eqnarray*}
which suggests the equation (\ref{focc0t3}) holds. This result also implies that $\breve{C}_{u, t} = C(\widetilde{\Phi}, \widehat{b}_{u, t}, G_t)$. 

We then verify the equation (\ref{imc0t3}). We rewrite IMCs (\ref{imc0t}), (\ref{imc0t1}), and (\ref{imc0t2}) for different debt restructuring policies as following
\begin{eqnarray*}
C^{*}_{u, t}U_{C^*, t}+(C^{*}_{u, t}+G_t)U_{N^*, t} = b_{0, t} U_{C^*, t}+\rho \Pi_u U_{C^*, t}\\
C^{*}_{u, t}U_{C^*, t}+(C^{*}_{u, t}+G_t)U_{N^*, t} = \widehat{b}_{u, t} U_{C^*, t}+\rho \lambda \Pi_u U_{C^*, t} \\
\breve{C}_{u, t}U_{\breve{C}, t}+(\breve{C}_{u, t}+G_t)U_{\breve{N}, t} = {b}_{0, t} U_{\breve{C}, t}+\rho  \Pi_u U_{\breve{C}, t} 
\end{eqnarray*}
The first two conditions above suggests that $b_{0, t}+\rho\Pi_u = \widehat{b}_{u, t}+\rho \lambda\Pi_u$. Substituting this into the right hand side of the third condition above, we obtain that
\begin{equation*}
\breve{C}_{u, t}U_{\breve{C}, t}+(\breve{C}_{u, t}+G_t)U_{\breve{N}, t} = \widehat{b}_{u, t} U_{\breve{C}, t}+\rho \lambda \Pi_u U_{\breve{C}, t} 
\end{equation*}
which suggests the implementability condition (\ref{imc0t3}) is satisfied. Overall, we've veried that the allocation $\{\breve{C}_{u, t}; \breve{N}_{u, t}\}$ is feasible for the potential debt restructuring policy $\{\lambda\Pi_u, \widehat{b}_{u, t}\}$. Because of the condition (\ref{eq_surpass}), the potential debt restructuring policy $\{\lambda\Pi_u, \widehat{b}_{u, t}\}$ is better than that from time-0 Ramsey plan even a government adopts the potential debt restructuring policy $\{\lambda\Pi_u, \widehat{b}_{u, t}\}$. 

\end{proof}







\begin{proof}[Proof of Porposition 9]
The sufficieny is obvious. The necessity of condition (\ref{tc_condition2}) is obtained by Lemma \ref{lemma2}.
\end{proof}



As application, we can show for some special initial debt, the positive roots of $\widetilde{\Psi}(\cdot)=0$ is unique and thus time consistency is attained.

\begin{corollary}
If $\{b_{0, t}=b; t\geq 0\}$, $\{G_t=g; t\geq 0\}$, if the Ramsey plan exists, it must be time-consistent for any time $t$.
\end{corollary}

%\begin{proposition}
%	If $\{b_{0, t}; t\geq 0\}$ and $\{G_t; t\geq 0\}$ both are periodic functions, if the commitment policy exists, it must be time-consistent.
%\end{proposition}




\end{document}


\pagebreak 

\section{To Do}


\WJ{
\begin{itemize}
\item DNY relates the failure of time-consistency to the sign of $\Phi_u(0)$. When $\Phi_u(0)<0$ for time$-u$ government, the optimality of the time$-0$ government's optimal commitment policy is violated because of the negative shadow price. 
\item However, in our framework, this argument may fail because there exists some debt restructuring strategy such that the shadow price for time$-u$ government is always positive. For example, for the pure short-term debt restructuring strategy with $\lambda=1$, we the shadow price for the time$-u$ government is always positive, i.e., $\Phi_u(1)=\Phi_0>0$.
\item Moreover, we can provide an example (See Table 1.) that even when $\Phi_u(1)=\Phi_0>0$, the Lucas and Stokey time-consistency could still fail. This is because of the existence of the multiple positive shadow prices. 
\end{itemize}
}



\WJ{
\begin{itemize}
\item DNY implies that the time-inconsistency in the sense of Lucas and Stokey may associate with conditions that a government in future taxes above the peak of the Laffer curve. 
\item Our sufficient condition is different from DNY's Laffer curve argument.
\item We can further provide an example such that even if the government taxes above the peak of the Laffer curve, the optimal commitment policy is still time consistent in the sense of Lucas and Stokey; see Figure 1.
\end{itemize}
}






{\bf Regarding time-consistency of optimal commitment policy}
\begin{itemize}
\item A complete characterization of potential debt restructuring strategies, which uses both short-term debt and term debt, to implement time-consistent in the sense of Lucas Stokey.
\item A pure short-term debt restructuring strategy that coincides the commitment government's pure debt roll-over strategy is provided and sufficient condition for Lucas-Stokey time-consistency is provided.
\item DNY arguments on Lucas-Stokey time-consistency is studied in our general framework
\begin{itemize}
\item The sign of Lagrange multiplier that remains positive for future government may be not necessary to maintain Lucas-Stokey time consistency.
\item The ``a government in the future will not necessarily tax above the peak of the Laffer curve" argument may fail.
\end{itemize}
\end{itemize}




Substituting (\ref{foch1}) and (\ref{foch2}) into the household's budget constraint (\ref{budgetc}), we obtain
\begin{equation}\label{imc}
\int_{0}^{\infty} e^{-\rho t} ( C_{0, t}U_C+N_{0, t}U_N )  dt  =  \int_{0}^{\infty} e^{-\rho t}b_{0, t} U_C(C_{0, t}, N_{0, t}) dt%+B_{0}U_{C}(C_{0, 0}, N_{0, 0}) .
\end{equation}
This is the implementability constraint that enters into the government's optimization problem. 



\newpage 

For example, if from time $0$ to time $t$, the government issues no additional debt that matures at time $t$, then we have 
\begin{equation*}
\partial_u b_{u, t} =0, \, u\leq t\,,
\end{equation*} 
which suggests that
\begin{equation*}
b_{t, t} = b_{0, t} \,.
\end{equation*}

\newpage 


Let $B_t$ denote the government's short-term debt sequence and let $r_t$ denote the equilibrium interest rate sequence to be determined later. 
At any time $t$, taking $r_t$ and pre-determined long-term debt holdings $b_{0,t}$ as given, the government chooses consumption, labor supply, and short-term asset holdings subject to  the following  law of motion:\footnote{We impose a  transversalty condition: $\lim_{t\rightarrow \infty} e^{-\int_t^{\infty} r_s d s} B_s = 0.$ 
} %and its decision problem is in essence a static one. %which can be replicated via dynamically trading a few long-lived assets as we show in Section \ref{}. %\footnote{See Arrow (19XX), Black and Scholes (1973), Harrison and Kreps (1979), Duffie and Huang (1985) for classic contributions on how dynamic trading makes markets dynamically complete.}  
%The following present-value budget constraint for the household holds at $t=0$: 
%\begin{equation}\label{imc}
%\mathbb{E}_0\int_{0}^{\infty} \left[ (1-\tau_t)N_tM_t -C_tM_t  \right]dt+B_0 = 0
%\end{equation}
\begin{eqnarray}\label{budgetB}
d B_t = \left(r_{t} B_{t}+ G_{t}-\tau_{0, t} {N_{t}} +b_{0, t}\right) dt\, . 
\end{eqnarray}
Note that the government can roll over the balance of $B_t$ at the then prevailing interest rate $r_t$. 




Let $W_t$ denote the household's holdings of the government's short-term bond at time $t$. 
At any time $t$, taking $r_t$ and endowed  long-term) asset holdings $b_{0,t}$ as given, the household chooses consumption, labor supply, and short-term asset holdings subject to  the following  law of motion: %and its decision problem is in essence a static one. %which can be replicated via dynamically trading a few long-lived assets as we show in Section \ref{}. %\footnote{See Arrow (19XX), Black and Scholes (1973), Harrison and Kreps (1979), Duffie and Huang (1985) for classic contributions on how dynamic trading makes markets dynamically complete.}  
%The following present-value budget constraint for the household holds at $t=0$: 
%\begin{equation}\label{imc}
%\mathbb{E}_0\int_{0}^{\infty} \left[ (1-\tau_t)N_tM_t -C_tM_t  \right]dt+B_0 = 0
%\end{equation}
\begin{eqnarray}\label{budgetW}
d W_t  = \left[r_{t} W_{t}-C_{t}+ (1-\tau_{0, t}) {N_{t}} +b_{0, t}\right]dt\, . 
\end{eqnarray}
We impose the corresponding transversalty condition: $\lim_{t\rightarrow \infty} e^{-\int_t^{\infty} r_s d s} W_s = 0.$ 

In equilibrium, short-term debt supply is equal to demand at all $t$: 
\begin{equation}
W_t = B_t \, . 
\end{equation} 
The long-term debt market also clears.

\newpage 

Note that this implementability condition that time$-0$ government must comply to is independent of debt roll over policy for the future government. Because of this, in our model, we only focus on the debt roll over policy that manages only short-term debt to satisfy the period by period dynamic budget constraint. In this case, we have
\begin{equation*}
\left(\int_{s\geq t} q_{t, s}\partial_tb_{t, s} ds \right)  = 0\, t\geq 0,
\end{equation*}
and thus 
\begin{equation*}
b_{t, t}  = b_{0, t}\,, t\geq 0.
\end{equation*}
Therefore, the household's dynamic budget constraint can be simplified as
\begin{eqnarray} \label{dynamicbudgetconstraintsimple}
\left(\tau_tN_t-G_t \right)dt  + dB_t =r_tB_t dt+  b_{0, t}dt \, , 
\end{eqnarray}
which is exactly the budget constraint (\ref{budgetB}). 


\newpage  


\section{First-best (FB) benchmark.} 
The FB problem boils down to a static optimization for each period and the FB solution is attained by setting $U_C(C,N) = -U_N(C,N)$, as the marginal product of labor is one unit of consumption and hence the marginal utility of consumption is equal to the marginal disutility of labor. Finally, government spending at $t$ is financed by a lump-sum non-distortionary tax that is equal to $G_t$.  
In the FB economy, the Richardian equivalence holds and hence government debt does not matter for resource allocation.  

\NW{In effect, both government and households commitment and taxes are nondistortionally.} 
\subsubsection{The materials below will be relegated to appendices.} 
Rewriting equation (\ref{dy}) to be 
\begin{equation}
d (e^{-\rho t} X_t) = e^{-\rho t}\left[-C_{0, t}U_{C}-N_{0, t}U_{N} +b_{0, t}U_{C}\right]dt,
\end{equation}
Integrating both sides from $0$ to $t$, we have
\begin{equation}
e^{-\rho t} X_t -X_0 = \int_{0}^{t}  e^{-\rho s}\left[-C_{0, s}U_{C}-N_{0, s}U_{N} +b_{0, s}U_{C}\right]ds \,.
\end{equation}	
%We then obtain the following  equation for $B_t$ 
%\begin{eqnarray}
% B_t  &=& e^{\rho t}B_{0} + \int_{0}^{t}  e^{\rho (t-s)}\frac{1}{U_{C}(C_0,N_0)} \left[-C_{0, s}U_{C}-N_{0, s}U_{N} +b_{0, s}{U_{C}} \right]ds \nonumber\\ 
% &=& e^{\rho t} \left[ B_{0} + \int_{0}^{t}\frac{ M_{0,s}}{M_{0,0}} \left[G_{0, s}+ \tau_sN_{0, s}+b_{0, s} \right]ds\right] 
%\end{eqnarray}	


Under the transaversality condition $\lim_{t\to \infty}e^{-\rho t} X_t = 0$, the dynamics of $X_t$ implies the following forward-looking equation: , 
\begin{eqnarray}
-X_0 = \int_{0}^{\infty}  e^{-\rho s}\left[-C_{0, s}U_{C}-N_{0, s}U_{N} +b_{0-, s}U_{C}\right]ds \\  
\int_{0}^{\infty}  e^{-\rho s}\left[C_{0, s}U_{C}+N_{0, s}U_{N} \right]ds= \int_{0}^{\infty}  e^{-\rho s}\left[b_{0, s}U_{C}\right]ds+ B_{0}U_C(C_{0, t}, N_{0, t})
\end{eqnarray}
which is exactly the implementability condition (\ref{impletconstraint}). Therefore, we have verified that $X_t$ is a qualified state variable in this problem.



\section{Cyclical Ramsey plans}

In this section, we discuss optimal commitment policies given that the term debt and government spending are periodic.

Consider an economy with isoelastic preferences over consumption $C$ and labor supply $N$, i.e.,
\begin{equation}
U(C, N) = \log C - \eta \frac{N^\gamma}{\gamma} \,.
\end{equation}

The period by period optimization problem that faces the government is
\begin{equation}
\max_{C_{0, t}} \quad \log C_{0, t}- \eta \frac{(C_{0, t}+G_t)^\gamma}{\gamma} + \Phi_0 \left (1-\eta\gamma (C_{0, t}+G_t)^\gamma -\frac{b_{0, t}}{C_{0, t}}  \right)\,, t\geq 0
\end{equation}
The FOCs implies that
\begin{eqnarray}\label{FOCspecialU}
\frac{1}{C_{0, t}}+ \Phi_0\frac{b_{0, t}}{C_{0, t}^2} -\left(\eta+\Phi_0\eta\gamma^2\right)(C_{0, t}+G_t)^{\gamma-1} = 0\,, t\geq 0 \,.
\end{eqnarray}

The existence of interior solution requires that $\Phi_0>0$. When $\Phi_0>0$, the FOC (\ref{FOCspecialU}) admits a unique solution. 

In particular, when $\gamma=1$, we have
\begin{equation*}
C_{0, t}^* = C(\Phi_0, b_{0, t}, G_t) = \frac{2 b_{0, t}\Phi_0}{\sqrt{1+4\eta b_{0, t}\Phi_0(1+\Phi_0)}-1} \,. %1_{\{b_{0,t}>0\}} + \frac{1}{\eta(1+\Phi_0)}1_{\{b_{0, t}=0\}}\,.
\end{equation*}

The implementability condition requires that $\Phi_0$ solves the following equation
\begin{eqnarray*}
\Psi_0(\Phi_{}) = \int_{0}^{\infty} e^{-\rho t} \left( 1-\eta (C(\Phi_{}, b_{0, t}, G_t)+G_t)-\frac{b_{0, t}}{C(\Phi_{}, b_{0, t}, G_t)} \right) dt-\frac{B_{0}}{ C(\Phi_{}, b_{0, 0}, G_0)} =0
\end{eqnarray*}

Numerical results show that this equation admits a unique positive root, i.e., $\Psi_0(\Phi_0^*)=0$\,.

\begin{figure}[h!]
\centering
\includegraphics[width=\linewidth, height=0.6\textheight]{figure_202305/fig_b_sin_g_c.eps}
\caption{Optimal commitment policy. $b_t = b_0+\theta_b \sin(\pi t), b_0=0.1, \theta_b=0.1; g=0.2, \rho=0.05, \eta=1$. $\Phi_0^*=0.6998$. }

\end{figure}


\begin{figure}[h!]
\centering
\includegraphics[width=\linewidth, height=0.6\textheight]{figure_202305/fig_b_c_g_sin.eps}
\caption{Optimal commitment policy. $g_t = g_0+\theta_g \sin(\pi t), g_0=0.2, \theta_g=0.1; b_t=0.1, \rho=0.05, \eta=1$. $\Phi_0^*=0.7329$. }

\end{figure}


\begin{figure}[h!]
\centering
\includegraphics[width=\linewidth, height=0.6\textheight]{figure_202305/fig_b_sin_g_sin.eps}
\caption{Optimal commitment policy. $b_t = b_0+\theta_b \sin(\pi t), b_0=0.1, \theta_b=0.1; g_t = g_0+\theta_g \sin(\pi t), g_0=0.2, \theta_g=0.1; \rho=0.05, \eta=1$. $\Phi_0^*=0.7071$.  }

\end{figure}

\begin{figure}[h!]
\centering
\includegraphics[width=\linewidth, height=0.6\textheight]{figure_202305/fig_b_sin_g_cos.eps}
\caption{Optimal commitment policy. $b_t = b_0+\theta_b \sin(\pi t), b_0=0.1, \theta_b=0.1; g_t = g_0+\theta_g \cos(\pi t), g_0=0.2, \theta_g=0.1; \rho=0.05, \eta=1$. $\Phi_0^*=0.6999$. }

\end{figure}


\begin{figure}[h!]
\centering
\includegraphics[width=\linewidth, height=0.6\textheight]{figure_202305/fig_b_sin_g_sin_2.eps}
\caption{Optimal commitment policy. $b_t = b_0+\theta_b \sin(\pi t), b_0=0.1, \theta_b=0.1; g_t = g_0+\theta_g \sin(\pi t/2 ), g_0=0.2, \theta_g=0.1; \rho=0.05, \eta=1$. $\Phi_0^*=0.7145$. }

\end{figure}

\begin{figure}[h!]
\centering
\includegraphics[width=\linewidth, height=0.6\textheight]{figure_202305/fig_b_sw_g_c.eps}
\caption{Optimal commitment policy. $b_t = b_0+\theta_b \sum_{k=1}^{5}\sin((2k-1)\pi t )/(2k-1), b_0=0.1, \theta_b=0.1; g_t = g_0+\theta_g \sin(\pi t), g_0=0.2, \theta_g=0; \rho=0.05, \eta=1$. $\Phi_0^*=0.6955$. }

\end{figure}


\begin{example}
Let $b_{0, t} = b_{0, 0}+\theta_b\sin(\pi*t)$, which is a $2-$periodic function, then consumption $C(\Phi_0, b_{0, t}, G_{t})$ is also a $2-$periodic function. Moreover, we have
\begin{equation}
C(\Phi_0, b_{0, t}, G_{t}) \sim a_0 +\sum_{n=1}^{\infty} \left(a_n \cos\frac{n\pi}{2} t + b_n \sin\frac{n\pi}{2}t \right)
\end{equation}  
where
\begin{eqnarray*}
a_n & = & \frac{1}{2} \int_{0}^{2}  C(\Phi_0, b_{0, t}, G_{t}) \cos\frac{n\pi}{2} t dt  \\
b_n & = & \frac{1}{2} \int_{0}^{2}  C(\Phi_0, b_{0, t}, G_{t}) \sin\frac{n\pi}{2} t dt\,, n\leq 1 \\
a_0 & = & \frac{1}{2} \int_{0}^{2}  C(\Phi_0, b_{0, t}, G_{t}) dt\,.
\end{eqnarray*}
\end{example}


%\begin{proposition}
%	If $b_{0, t}$ is a Lipschitz continuous, periodic function with respect to time $t$, then we can find a Fourier series $b^{(n)}_{0, t}$, which are continuous in time $t$, such that the optimal commitment policy for initial debt $\{b_{0, t}\}$, i.e., $C(\Phi_0^*, b_{0, t}, G_{0, t})$, can be approximated by the optimal commitment policy for initial debt $\{b^{(n)}_{0, t}\}$, i.e.,  $C(\Phi_0^*, b^{(n)}_{0, t}, G_{0, t})$. That is
%	\begin{equation}
%	\lim _{n\to \infty}  \, C(\Phi_0^*, b^{(n)}_{0, t}, G_{0, t}) = C(\Phi_0^*, b_{0, t}, G_{0, t})\,.
%	\end{equation}
%\end{proposition}



\TS{I more or less stopped editing here but have some off-line comments and questions on the subsequent
fascinating material.}
\begin{figure}[h!]
\centering
\includegraphics[width=\linewidth, height=0.6\textheight]{figure_202305/multiplier_1.eps}
\caption{Optimal commitment policy. $B_0=2.5, b_t = b_0, b_0=0.01; g=0.2, \rho=0.05, \eta=1$. $\Phi_0^*=0.8180$. $C^*_t = 0.5581$; }

\end{figure}

\begin{figure}[h!]
\centering
\includegraphics[width=\linewidth, height=0.6\textheight]{figure_202305/multiplier_2.eps}
\caption{Optimal commitment policy. $B_0=2.5, b_t = b_0, b_0=0.01; g=0.2, \rho=0.05, \eta=1$. $\Phi_0^*=3.7803$. $C^*_t = 0.2419$.}

\end{figure}


\NW{
\begin{itemize}
\item Lucas Stokey led us a wrong way to implement optimal commitment policy
\item To implement optimal commitment consumption choosing at $t=0$, we need to twist both Lagrange Multiplier and $b$ structure which is crazy
\item Short-term debt $B$ substantially simplifies the implementation strategy
\item In essence, LS is not complete market model, because it  
\end{itemize}
}


\clearpage




\pagebreak
\section*{Appendices}
%\appendix{\bf \centering \LARGE \noindent Appendices}
\appendix{}
%\setcounter{page}{1} 
%\setcounter{Theorem}{0} 
%\renewcommand{\thepage}{I-\arabic{page}}
%\setcounter{theorem}{0}
%\renewcommand{\thetheorem}{\Alph{section}.\arabic{theorem}}
\setcounter{equation}{0}
\renewcommand{\theequation}{I-\arabic{equation}}
\setcounter{section}{0}
\renewcommand{\thesection}{\Alph{section}}
\renewcommand{\thesubsection}{\Alph{subsection}}
%\setcounter{figure}{0}
%\renewcommand{\thefigure}{I-\arabic{figure}}
%\setcounter{footnote}{0}
\bigskip%


\subsection{Digression on Debt Management: Proposed Modified Debt Stream to include ''Dirac debt''}

\noindent
Let $T_D = t_0, t_1, \ldots, t_T$ be a possibly countably infinite sequence of nonnegative dates.
We'll set $t_0 = 0$ and require  $T$ to be finite.
Let $\{\hat B_{0,t}\}_{t=t_0}^{t_T}$ be a sequence of finite   real numbers.
$\hat B_{0,t}$ is a discrete payment due at time $t \in [0, + \infty)$, so $\hat B_{0,t}$ is a  ``stock'', not a flow.
We let $b_{0,t}$ be a continuous (or almost continuous with a countable number of jumps)  coupon steam, where coupon payments $b_{0,t} dt $ are due in interval $dt = t, t+ \epsilon$.
Let $q_{0,t}$ be a time $0$ (Arrow-Debreu) pricing kernel.
At time $0$, the government owes a payout stream of 
$$ 
p_{0,t} = b_{0,t} + \sum_{k=0}^\infty \delta (t-t_k) \hat B_{0,t_k}, \quad t \in [0, \infty)
$$
where $\delta(t)$ is a Dirac generalized function that verifies the property

$$
\int_{0}^\infty  \delta (t-t_k) \hat B_{0,t_k} dt 
= \hat B_{0,t_k} .
$$
At time $0$, the present value of the government's promised payment stream is
$$
PV(0) = \int_0^\infty q_{0,t} b_{0,t} dt + \sum_{k=0}^\infty q_{0,t_k} \hat B_{0,t_k} 
$$


\textbf{Why do this}

\begin{itemize}
\item It allows sharpening presentation  of differences between what we do and what is done in discrete-time papers in the Lucas-Stokey (1983) tradition.  
\item It quietly makes a case that what we do is more general than Lucas-Stokey's in the sense of expanding choice sets and so on.
\item It provides a way of emphasizing the point that the Ramsey planner (who chooses everything at time $0$ and then walks away) moves $B_{t,t}$ over time but leaves $B_{t,t_i} = B_{0, t_i}$ for all $i > 1$.  This sharpens the message that $B_{t,t} U_{c,t}$ is the key state variable. (Here we might have to alter my  notation to make things clearer -- I probably obscured what I want  to say.)
\end{itemize}

At a minimum, all of this could be done in a footnote or two if we don't want to "complicate" things too much at first.  



\subsection{Optimal Commitment Policy and Time Consistency for Piece-wise Continuous Debt Maturity Distribution}



\section*{Tom's January 23 suggestions}

\begin{enumerate}
\item Define ``covenant'' precisely early in paper.
\item Say exactly what the covenant is in the one-period debt contract. Discuss attitudes of borrower and lender toward that covenant.
\item Improve English, consolidate and avoid repetition.  Increase impact by saying things concisely once.
\item Tighten the English text about how things work.  
\item Tighten description of timing and how ``flows'' depict integrals over $\Delta$ versus point in time objects.  Link to econometric work on aggregation over time.

\end{enumerate}


\section{Introduction}



% that opened new possibilities for  macroeconomics generally.
%Before  \citet{LucasStokey1983}, 
% \citet{lucas1972econometric,lucas1972expectations,lucas1975equilibrium}  modeled 
%  government policies as 
%exogenous  stochastic processes.\footnote{ \citet{lucas1978asset}  had   taken a representative agent's consumption process and the associated  stochastic discount factor as exogenous. A key aspect of the  \citet{LucasStokey1983} analysis  was to design a tax plan to manipulate a stochastic discount factor process.}    %The time  inconsistency of Ramsey plans demonstrated by 
% \citet{kydland1977rules} and  \citet{calvo1978time} argued 
%% as it had \citet{kydland1977rules} and  \citet{calvo1978time},
% that Ramsey plans are implausible because  they are not implementable  under  the  sequential timing protocols that 
% actually characterize  government decision-making processes.\footnote{\citet{kydland1980dynamic}   offered  a concise  way of representing   Ramsey plans  recursively, but 
%%those plans were still not implementable.}  But  
% \citet{LucasStokey1983} showed how  to implement a Ramsey plan
% when  successive government administrations choose continuation tax plans  sequentially. 
%That opened the door  to analyzing  optimal government policy in settings with sequential  government decision making.   
% As their laboratory, % for studying implementability of a Ramsey plan, 
%\citeauthor{LucasStokey1983}
%turned to \citet{barro1979determination}. 
%, a principal conclusion of which   is  that  intertemporal properties of an optimal  tax rate are 
%disconnected from intertemporal properties of the exogenous government expenditure process that must  be
%financed.
%  \citeauthor{LucasStokey1983}  reexamined %understand sources of  
%Barro's findings in a general equilibrium context. They  %that  striking tax-smoothing outcome as well as other features of Barro model,  

%\citet{ramsey1927}

By building on  \citeauthor{ramsey1927}'s (1927) static  framework with distorting flat-rate taxes, \citet{LucasStokey1983} studied optimal monetary and fiscal policies in  a closed, dynamic, stochastic, complete-markets economy.\footnote{While their  question overlapped with   \citet{barro1979determination},   their approach was  different.} 
%  with distorting flat-rate taxes
%
% To study optimal monetary and fiscal policies,  \citet{LucasStokey1983} formulated a Ramsey plan within a  complete-markets,  dynamic closed economy  with  distorting flat-rate taxes.\footnote{While their  question overlapped with   \citet{barro1979determination},   their approach was  different.} 
They showed how  to implement a Ramsey plan
when  successive government administrations choose continuation tax plans  sequentially. They described policies for managing   financial obligations that  induce a sequence of governments to confirm a Ramsey plan for tax rates.\footnote{\citet{kydland1977rules} and  \citet{calvo1978time} argued 
% as it had \citet{kydland1977rules} and  \citet{calvo1978time},
that Ramsey plans are implausible because  they are not implementable  under  the  sequential timing protocols that 
actually characterize  government decision-making processes. \citet{kydland1980dynamic}   offered  a concise  way of representing   Ramsey plans  recursively, but 
those plans were still not implementable.  A by product of   \citeauthor{LucasStokey1983}'s  analysis is a determinate optimal term structure of government debt. In addition, 
their closed-economy, complete-markets analysis overturned  \citet*{barro1979determination}'s finding that intertemporal properties of an optimal tax rate are disconnected from the intertemporal properties of  government expenditures.}  
%{ What initially convinced Lucas to set aside his reservations about the implausible timing protocols assumed in posing a Ramsey problem was his desire to %\citetauthor{LucasStokey1983} set out  
% formalize and therefore to understand  the striking tax-smoothing outcome of   
%\citet{barro1979determination}, whose principal conclusion  was that the intertemporal properties of an optimal  tax rate are 
%completed disconnected from the intertemporal properties of that exogenous government expenditure process that is to  be
%financed. 
%Under    particular   assumptions about the promises that  continuation governments are  bound to respect, 
% \citeauthor{LucasStokey1983} showed how to implement a   Ramsey plan.  
%  responded to \citet*{kydland1977rules}'s and
%citet*{calvo1978time}'s doubts about the plausibility of  Ramsey plans.  
% \citet*{LucasStokey1983} obtained strikingly different  conclusions outcomes than Barro.

%\WJ{Instead, in \citeauthor{LucasStokey1983}'s model,  intertemporal properties of the tax rate typically  mirror those  of government 
%expenditures.\footnote{Subsection \ref{subsec:continuous2} of this paper describes a special case of \citeauthor{LucasStokey1983}'s
%model in which Barro's disconnection result still  prevails.} } \NW{Tom: do we need the last sentence as this ``mirroring'' result is not general and actually does not  hold in DNY's example. Maybe we can also remove footnote 4?} 
%\NW{Tom: shall we promote this footnote to the main body so that we discuss  this important intellectual history and relate it to our work?} \NW{BTW, shall we use disconnect or disconnection? I know both can be used as a noun.} 
%
%\citet{lucas1975equilibrium}, \citet{barro1979determination}, \citet{lucas1972expectations}, \citet{kydland1980dynamic},
%\citet{kydland1977rules}, \citet{calvo1978time}


%\citet*{debortoli2021optimal}, \citet{LucasStokey1983}
This paper resolves a difficulty   that \citet*{debortoli2021optimal} detected with \citeauthor{LucasStokey1983}'s (1983, [Sec.~3])  prescription for managing the  
term structure of government debt. In  \citeauthor{LucasStokey1983}'s model, each of 
a sequence of governments   finances an exogenous  joint stochastic process for government expenditures $ \{G_{i}\Delta\}_{i=0}^{\infty}$ and debt service coupons $ \{b_{-1, i}\Delta \}_{i=0}^{\infty}$. % for time $t \geq 0$. 
At time $0$,  a Ramsey planner chooses a process of distorting flat-rate taxes and possibly a restructured debt service coupon process  $\{{b}_{0, i}\Delta \}_{i=1}^{\infty}$. In period $ 1$, continuation Ramsey planners  are free to  redesign the continuation of  the  flat-rate tax process and also  to reschedule government debt  from $t $ onward; but they  must  honor  continuation debt service coupon process that they inherited. 
%\WJ{{\bf How about we comment out the following sentence in red?} \citeauthor{LucasStokey1983} provided examples  in which  appropriate   continuations $ \{\hat b_{t, s}\}_{s=t}^{\infty}$ of {term structures of government debt} induce  continuation  planners to choose to continue  the original Ramsey tax rate process.  }  % However, 
\citet*{debortoli2021optimal} constructed examples in which initial debt is so high that the Ramsey plan sets the tax rate above the peak of the Laffer curve and  \citeauthor{LucasStokey1983}'s way of restructuring government debt fails to motivate  continuation planners   to continue the Ramsey tax plan. % confirm the Ramsey plan.  
%call for this tax policy.} %it is ex ante optimal to do so.} 

To set  the stage for  our expansion of \citeauthor{LucasStokey1983}'s    contractible subspace, it is useful to read how
\citet{aguiar2019take}   contrasted  \citeauthor{LucasStokey1983}'s (1983) model with theirs: 
{\it 
\begin{quotation}
\citet{LucasStokey1983} studied optimal fiscal policy with complete markets and discussed
at length how maturity choice is a useful tool to provide incentives to a government
that lacks commitment to taxes and debt issuance, but cannot default. The government
has an incentive to manipulate the risk-free real interest rate, by changing taxes which affects
investors’ marginal utility, to alter the value of outstanding long-term bonds, something
ruled out by our small open-economy framework with risk-neutral investors. Their
main result is that the maturity of debt should be spread out, resembling the issuance of
consols. Our model instead emphasizes default risk, something absent from their work.
Our main result is also the reverse, providing a force for the exclusive use of short-term
debt.
\end{quotation}
}


Short-term debt also  plays an essential role in our model. But unlike \citeauthor{aguiar2019take},
we retain almost all other parts  of   \citeauthor{LucasStokey1983}'s and \citeauthor{debortoli2021optimal}'s structure, including complete markets,
a closed economy setting, the presence of 
incentives that  continuation Ramsey planners have  to use flat-rate taxes to manipulate interest rates, 
and  obligations of  governments to honor all government debt that they inherit.  
%  to finance its cumulative deficits via instantaneous debt and  
In contrast to   \citeauthor{LucasStokey1983}'s  and \citeauthor{debortoli2021optimal}'s  implementation that relies on debt restructurings 
$\{ {b}_{0, i} \Delta \}_{i=1}^\infty \neq \{b_{-1, i} \Delta \}_{i=1}^{\infty}$, 
we  require only %add instantaneous debt in the form of an 
one-period debt balance $b_{0, 1}\Delta $ that  period-$1$ continuation Ramsey planners must service, together with  a 
covenant for one-period-ahead tax policy. 
By managing  % the {\em stock} 
$b_{0, 1}\Delta$  appropriately,   \citeauthor{LucasStokey1983}'s Ramsey plan can be implemented without restructuring 
$\{b_{-1, i}\}_{i=1}^{\infty}$, i.e., by setting  $\{ \hat{b}_{0, i}\}_{i=1}^\infty = \{b_{-1, i}\}_{i=1}^{\infty}$. %  at all $t$ 
(A plan that  can be implemented  is often  called  ``time consistent.'')\footnote{ %\citet[ch.~21]{Ljungqvist_Sargent_2025} 
\citet[ch.~20]{sargent2000recursive}
assume that the government has access to only one-period debt in a discrete time version of  \citeauthor{LucasStokey1983}'s model. It is  a counterpart to instantaneous debt in our continuous time model.  It also leads to Bellman equations that can be used to express a discrete-time counterpart to the local commitment constraint that we impose in this paper.}
{Table \ref{tab1} compares}   the contractible subspace in our model  with  \citeauthor{LucasStokey1983}'s and \citeauthor{debortoli2021optimal}'s.   


%\begin{table}[h!] 
%	%\centering
%	\caption{Comparison with existing methods \label{tab1}}	
%	\resizebox{.95\textwidth}{!}{
%		{\def\arraystretch{2}\tabcolsep=8pt
%			\begin{tabular}{c|c|c}
%				\hline
%				& Contractible space & Implementable?  \\ 
%				\hline 
%				\cite{LucasStokey1983}	& $\{ {b}_{0, i}\}_{i=1}^\infty \neq \{b_{-1, i}\}_{i=1}^{\infty}$ & Not always  \\   %\rule{0pt}{5ex}
%				\hline 
%				\citet*{debortoli2021optimal}	& $\{ {b}_{0, i}\}_{i=1}^\infty \neq \{b_{-1, i}\}_{i=1}^{\infty}$ & Not always  \\   %\rule{0pt}{5ex}
%				\hline 
%				\multirow{2}{*}{This paper} 	& ${b}_{0, 1}\Delta ,\, \{ {b}_{0, i}\}_{i=1}^\infty = \{b_{-1, i}\}_{i=1}^{\infty}$
%				&\multirow{2}{*}{Always} \\  %\rule{0pt}{5ex}
%				& Covenant for one-period-ahead tax policy  &      \\
%				\hline 
%			\end{tabular} 
%		}
%	}
%\end{table}



\begin{table}[h!]
\caption{Comparison with existing methods \label{tab1}}	
%	\resizebox{.95\textwidth}{!}{
%\caption{Comparison with existing methods \label{tab1}}
\centering
{\def\arraystretch{1.5}\tabcolsep=3pt
\begin{tabular}{c|c|c}
\hline
& Contractible space & Implements Ramsey Plan?  \\
\hline
%Lucas and Stokey (1983)	
LS and DNY & $\{ {b}_{0, i}\}_{i=1}^\infty \neq \{b_{-1, i}\}_{i=1}^{\infty}$ & Not always  \\   %\rule{0pt}{5ex}
% \hline
% %Debortoli, Nunes, and Yared (2021)	
% DNY (2021) & $\{ {b}_{0, i}\}_{i=1}^\infty \neq \{b_{-1, i}\}_{i=1}^{\infty}$ & Not always  \\   %\rule{0pt}{5ex}
\hline
\multirow{2}{*}{This paper} 	& $\TS{\big( b_{0, 1},  \tau_{0,1}\big)}$
&\multirow{2}{*}{\bf Always} \\  %\rule{0pt}{5ex}
&  $\NW{\{ {b}_{0, i}\}_{i=2}^\infty = \{b_{-1, i}\}_{i=2}^{\infty}}$  &      \\
\hline  % (Local commitment)
\end{tabular}
%		}
}
\end{table}	



%Why do we propose and how we specify the local commitment condition Lucas-Stokey

%First, we  point out that a Ramsey plan uniquely pins down  the cumulative government liability at time $t$: $
% absent any term debt restructuring up to $t$

%To explain  how we  %use $\widehat B_t$ to 
%implement   \citeauthor{LucasStokey1983}'s  Ramsey plan, 
A bond price process $\{q_{0,j}^\ast; j\geq 0\}$ and a  primary surplus process $\{S_{0,j}^\ast\Delta; j\geq 0\}$ are affiliated with \citeauthor{LucasStokey1983}'s Ramsey plan.  Together with the initial term debt structure $\{b_{-1,j}\Delta; j\geq 0\}$, these two objects uniquely determine a process: $\Pi_i^\ast= \sum_{j=0}^{i-1} \frac{q_{0,j}^\ast}{q_{0,i}^\ast} \big(b_{-1,j}- S_{0,j}^\ast\big)\Delta$ that   accumulates the   government's unpaid  liabilities from period $0$ to $i-1$, conditional on the government not having rescheduled  the initial  term debt process  $\{b_{-1, j}\Delta\}_{j=0}^{\infty}$ before period  $i$.  Time $j$ flow liability $\big(b_{-1,j}- S_{0,j}^\ast\big)\Delta $ times   price  $\frac{q_{0,j}^\ast}{q_{0,i}^\ast} $  contributes to  a period $i$ value of government liabilities. Our implementation of a Ramsey plan refrains from  restructuring  an initial  debt term structure $\{b_{-1, j}\Delta\}_{j=0}^{\infty}$ and instead  instructs continuation governments to manage  only  short-term debt. % (i.e., zero in the limit). 
%to   track  the government's cumulative liability over time. 
This arrangement takes into account all  relevant    quantity-and-equilibrium-price  information about the  initial term debt structure  that concerns the Ramsey planner.


The Ramsey planner  thus implements its plan by  leaving the initial  debt term structure with period-to-maturity weakly larger than two untouched and  paying for all  government purchases  and debt servicing by issuing short-term debt.  
%Absent  term debt structuring, the instantaneous debt balance is the representative household's (additional) net worth.  
Continuation governments are obligated  to commit to a covenant condition of one-period-ahead tax policy that relates to the interest rate under Ramsey plan used in issuing the short-term debt. This covenant condition effectively preserve the  ``purchasing power'' of cumulative  government liabilities  $\Pi_i^\ast$, defined as the product of  $\Pi_i^\ast$ and the representative household's marginal utility of consumption $U_{C,i}^\ast$. %As we use instantaneous debt to track government's cumulative deficits in 
The Ramsey planner leaves each future government an short-term debt balance $b_{i-1, i}\Delta$ that equals $\Pi_i^\ast+b_{-1, i}\Delta$, the interest rate of which  continuation governments must commit. %of instantaneous debt. } 
% the product of instantaneous debt and the representative household's marginal utility of consumption. %should be preserved. 
%to define a state variable that allows us to reformulates the Ramsey plan  in that 
%the cumulative deficits under the Ramsey plan 
%
%
%\NW{  Among  debts of various maturities,  instantaneous debt has the shortest  maturity and is least  vulnerable to manipulation by future governments.  Moreover, the {instantaneous debt} balance  includes information about both  the equilibrium bond price and the  initial term debt ({${b}_{0,t}$}) of the  government's liability under the Ramsey plan. Therefore, the {instantaneous debt} balance is a sufficient statistic for government finance provided that continuation Ramsey planner is willing to confirm the Ramsey plan. 
%  }
Continuation governments  are tempted to deviate from  the  Ramsey  tax plan in order to manipulate that interest rate. % of government liabilities.  % absent commitment. %However, even instantaneous debt is still potentially subject to government's manipul. 
To prevent them from doing that and thus to   implement the Ramsey plan  we augment \citeauthor{LucasStokey1983}'s assumptions about commitments assigned to continuation governments by adding   a  covenant condition that requires that  governments  must commit to one-period ahead tax policy. Sections \ref{sec:im} and \ref{sec:dny}  show that this restriction  is sufficient to induce them to confirm the Ramsey plan.
% chosen by the time 0 government.} %Because instantaneous debt has infinitesimal maturity, }   device prevent  from diluting the  ``purchasing power'' of  instantaneous debt, inducing continuation Ramsey planners
%When the Ramsey planner has access to instantaneous debt, 
%}


%  : 	%\begin{equation*}
%		%\hat{B}_{1} =
%$	\widehat{B}_{1}=	\int_1^\infty e^{-\rho(t-1)} {\hat{b}_1} d t = 	 { {\hat{b}_1}}/{\rho} .$ %=\Pi_1^\ast %= \frac{S(C_1^*)}{\rho} 
%		%= \frac{\left(e^{\rho}-1\right)}{\rho} \left( \frac{b_{0}}{C_0^*} -1 + (C_0^*+G)   \right) C_1^*\, . 
%	%	\end{equation*}
%%presented to continuation planner. 
%The {(local) commitment} condition (\ref{key2}) prohibits  time-$1$ continuation planner  from  diluting  the household's ``purchasing power'' of $\widehat{B}_1$, which is  something that the continuation Ramsey planner  wants to do when it
%confronts the  C\`adl\`ag  $\{\hat b_{0,\,t}; t\geq 0 \}$ debt term structure without instantaneous debt  in \citeauthor{debortoli2021optimal}'s counterexample.


%Because instantaneous debt has the shortest maturity among all debt contracts, its value is least vulnerable to manipulation by future governments. Moreover, 
%
%\NW{Under complete markets, nstantaneous debt has the shortest maturity among all debt contracts, its value is least vulnerable to manipulation by future governments..} 

%The presence of instantaneous debt and the local commitment is what allows  our restructuring policy \ref{restruct1}  to implement the Ramsey plan. 
% Note how  the {instantaneous debt} balance includes information about both  the price  and the  quantity ({$\hat{b}_1$}) of the  government's liability: 	%\begin{equation*}
%\hat{B}_{1} =
%$	\widehat{B}_{1}=	\int_1^\infty e^{-\rho(t-1)} {\hat{b}_1} d t = 	 { {\hat{b}_1}}/{\rho} .$ %=\Pi_1^\ast %= \frac{S(C_1^*)}{\rho} 
%= \frac{\left(e^{\rho}-1\right)}{\rho} \left( \frac{b_{0}}{C_0^*} -1 + (C_0^*+G)   \right) C_1^*\, . 
%	\end{equation*}
%presented to continuation planner. 
%	Implementing Ramsey plan with 
%		The last equality has to hold because 



%This is why 
%In our  implementation, %since there is no adjustment to the initial term debt and the entire  cumulative liability $\Pi^\ast$ is booked with instantaneous debt balance: $B_t=\Pi_t^\ast$. it is  always feasible (though not necessarily optimal) for the continuation Ramsey planner to confirm the Ramsey plan. In order to induce the continuation Ramsey planner to confirm the Ramsey plan, we need a local commitment condition that preserves the purchasing power of instantaneous debt, removing the continuation  planner's incentive to manipulate debt price always from the one implied by the Ramsey plan. In sum,  
%
%\NW{In sum, our implementation proceeds in two steps. First, because the instantaneous debt  books the government's  accumulated liability without touching the initial term debt $b_{0,t}$ under the Ramsey plan,  it is always feasible (though not necessarily optimal) for the continuation Ramsey planner to choose the Ramsey plan; Second, we incentivize the continuation Ramsey planner to  confirm the Ramsey plan by using the local (instantaneous) commitment condition discussed above. 
%} 

%\TS{TOM: We moved this paragraph here.} 


%\NW{ That is, no matter whether the optimal tax rate is below or above  the peak of  the Laffer curve in the original Ramsey plan, a continuation Ramsey planner will choose to confirm the continuation of the original  plan. Recall that in \cite{debortoli2021optimal}'s implementation, a  continuation Ramsey planner does not want to choose a tax rate above  the peak of  the Laffer curve even when this tax policy is optimal under 
%%confirm the continuation of the  Ramsey plan if 
%the original  plan. %even when  the (initial) Ramsey plan calls for such a tax policy. That is,  in our implementation,  a  tax rate above  the peak of  the Laffer curve if  the (initial) Ramsey plan calls for such a tax policy. 
%A key insight of \citet*{debortoli2021optimal} underpinning their  counterexamples when initial debt is high so that the Ramsey planner has incentives to reduce the value of its debt. } 

%in  in which


% . 


%How can our way of  implementing   a Ramsey plan %\footnote{As in \citet*{debortoli2021optimal}, we specify conditions under which a Ramsey plan exists.} %When initial debt is too high, Ramsey plans may not exist.  
%we can implement the Ramsey plan. %They showed that in their counterexamples when initial debt servicing is high enough the Ramsey tax rate is   above  the peak of  the Laffer curve, 
%work despite the failure of \cite{debortoli2021optimal}'s  implementation  when  the Ramsey planner sets the tax rate above the peak of  the Laffer curve? 
\cite{debortoli2021optimal} show that when initial debt $b$ is high, %\NW{initial  debt service path $ \{b_{0, t}\}_{t=0}^{\infty}$ is sufficiently large},  
the Ramsey planner sets the   tax rate  above  the peak of  the Laffer curve.
% in order to manipulate the equilibrium interest rate in order to devalue government debt. reduce its  value  by .  
By doing that, it manipulates  the competitive equilibrium bond price process to  reduce the value of the initial  debt $b$.
%  to  reduce  consumption and hence bond price. 
But in \citeauthor{debortoli2021optimal}'s implementation,  %facing a restructured debt $\{\hat{b}_{0,t}\}$, 
a   continuation Ramsey planner  wants to lower that above-the-peak-of-the-Laffer-curve  tax rate. %is  no longer required to confirm %This is confirming the plan is simply too costly.  there is no commitment device that constrains
In contrast,  our implementation uses only short-term debt to satisfy the dynamic budget constraints.
We %, and \NW{moreover our implementation} 
require continuation Ramsey planners committing to a covenant condition of one-period-ahead tax policy. %accumulated government  liabilities. % under the original Ramsey plan. 
%We impose this requirement by adding to the contractible subspace  instantaneous debt that tracks the government's liability implied by the Ramsey plan and also  a  commitment to preserve the purchasing power of that debt  over the next instant. 
With these additional opportunities and constraints,  our    continuation Ramsey planner voluntarily chooses the Ramsey tax rate even when it is above the peak of the Laffer curve. 
%even when this tax policy is optimal under  honor a (local) commitment condition that
%confirm the continuation of the  Ramsey plan if the original  plan.


% the Ramsey planner has no additional instrument to constraint 

%initial debt servicing is so high  that  

%Consider \cite{debortoli2021optimal}'s  counterexample where initial debt servicing costs is high but a Ramsey plan still exists. While they showed that the original Ramsey plan cannot be confirmed by a government in the future, }  

%Sometimes   our implementation and \citeauthor{debortoli2021optimal}'s  implementation both  work, for example when initial debt servicing costs are not too high.  

\WJ{Need to add discussions on the Laffer curve argument, Lagrangian perspective, and necessity of covenant condition}


When the initial debt $b$  is not too large, the  Ramsey planner sets the  tax rate   below  the peak of  the Laffer curve. In this case,  \citeauthor{debortoli2021optimal}'s implementation and ours both work; actually,  we can dispense with  our  local commitment requirement because a tax rate below the peak of the Laffer curve does not tempt a continuation planner to deviate from the interest rate under Ramsey plan for the short-term debt. 
In situations in which  \citeauthor{debortoli2021optimal}'s implementation and ours both work,
they   are associated with  different  processes for the government's ``gross of interest surplus'', i.e.,
the  primary surplus plus debt service payments. In the context of \citeauthor{debortoli2021optimal}'s example, subsection \ref{sub52} shows that    \citeauthor{LucasStokey1983}'s implementation sets  the government's gross of interest surplus   equal to  zero always, while in our implementation  the government  runs a (positive) gross of interest deficit  when term debt  service $\{b_{0,t}\}$ is  initially large , then runs   a zero gross of interest deficit after  term debt service   payments $b_{0,t}$ drop to zero. This is because the Ramsey planner restructures the government's term debt in \citeauthor{LucasStokey1983}'s and \citeauthor{debortoli2021optimal}'s implementation, but not in ours. %\footnote{Our  implementation connects   to  Barro's tax smoothing recommendation and is reminiscent of Milton  Friedman's permanent-income hypothesis. \TS{I think that this footnote needs more elaboration or it should be dropped.} \WJ{Agreed and let's drop this one.}} 
%
%\NW{In sum, our first contribution  is to show how and why using instantaneous debt and imposing a local commitment condition allow us to implement Ramsey plans for \citeauthor{debortoli2021optimal}'s settings (see Sections \ref{sec:counterexample} and \ref{sec:expandcontract}).} 
%But if you take their implementation seriously, there is really no notion of using fiscal deficits to smooth tax distortions.
%In a way, this is also why the profession thinks that term debt management gets us close to LS Ramsey, but as we show it’s exactly the opposite!}

%\NW{One thing that Wei and I think we’d highlight even more (implicit in our current version) is that the "implied fiscal deficits"
%\textbf{Neng: please see below where I say that we want a definition of "fiscal deficits" that differentiates it sharply from "primary" or "net of interest deficit".} in a LS implementation of Ramsey (if feasible, i.e., for the b0=0.3 case in DNY) is zero, but of course it is not zero in our implementation. This has a lot to say about plausibility of our implementation compared to LS implementation. } 
%
%To sum up  in the first half of our paper (Sections 4 to 5), we 
%resolve a difficulty   that \citet*{debortoli2021optimal} detected with the prescription for an optimal 
% term structure of government debt proposed by \citet[Sec.~3]{LucasStokey1983} by expanding the contractible subspace of \cite{LucasStokey1983} to include instantaneous debt and adding a local (instantaneous) commitment condition to preserve the purchasing power of the instantaneous debt. 
% 
% 

%Sections \ref{sec:bellmanramsey} and \ref{sec:quantitative}  %we generalize our previous findings and achieve the following. First, we 
%describe   a recursive formulation of the Ramsey problem for  general right continuous with left limit (C\`adl\`ag) $\{b_{0,t}, G_t\}$ processes. A key  state variable $X_t$  equals the product of the government's cumulative liability $\Pi_t$ and   contemporaneous marginal utility $U_{C,t}$. The value function is additively separable in $X_t$ and time $t$ and so can be represented as  $-\Phi_0 ^\ast X_t + H(t;\Phi_0 ^\ast)$, where %. We show that
%$\Phi_0 ^\ast$ is the  multiplier  on the implementability constraint in a  Lagrangian formulation. %$-\Phi_0 ^\ast X_t + H(t;\Phi_0 ^\ast)$. We show that $\Phi_0 ^\ast$ corresponds to the  multiplier for the implementability constraint in the standard Lagrangian solution method.} %primal approach. } 
%
%
%{Sections \ref{sec:instantaneous} and \ref{sec:implement}  use our recursive formulations of the Ramsey problem  to   
%	implement  Ramsey plan for general C\`adl\`ag $\{b_{0,t}, G_t\}$ processes.  
%	%using a  local commitment condition similar to the one introduced in our analysis of the \citeauthor{debortoli2021optimal}'s counterexample. 
%	%Third, 
%	In particular, when   $ \{b_{0,t}, G_t\}$  are continuous  functions of time, the  local commitment condition that facilitates  implementing  the Ramsey plan simply requires   that the consumption rate is continuous in time. Furthermore,  
%	%\footnote{The requirement (for our implementation) that  the consumption rate is  continuous is  more general and flexible than those in discrete-time models. This is because  consumption  over any finite  time interval in our implementation is allowed to be time varying, while by construction consumption is  constant within a single time period in  discrete-time models.} %, construct an associated discrete-time model by forming a discrete time consumption process  $\{C_{0, t} \Delta\}_{t=0}^\infty$ and so on.  The discrete time consumption process  is ``discontinuous in $t$''  even though the consumption rate $C_{0, t}$ is continuous.
%	Ramsey plans  for settings with general C\`adl\`ag $\{b_{0,t}, G_t\}$ processes  are well approximated by  limits of Ramsey plans with appropriate continuous $ \{b_{0,t}, G_t\}$ processes. %These results allow us to conveniently apply our  methodology to characterize and implement (i.e., deliver ``time-consistent'') Ramsey plans.  
%}
%%






\end{document}


